%============================================================
%  CHAPTER 2 — DEPTH-FIRST SEARCH
%  Entry file: chapter opener + \input's only.
%
%  Folder layout
%     lessons/               Part A: the call stack, timestamps, the DFS tree, cycles, the DFS toolbox
%     p1_cycles/             theory, 2 solved problems, exercises
%     p2_subtrees/           theory, 3 solved problems, exercises
%     p3_timestamps/         theory, 3 solved problems, exercises
%     p4_backtracking/       theory, 2 solved problems, exercises
%     p5_reverse/            theory, 2 solved problems, exercises
%     summary.tex, problem_set.tex, problem_set_solutions.tex
%============================================================
\chapter{Depth-First Search}
\label{ch:dfs}

\begin{tcolorbox}[enhanced, colback=white, colframe=secondaryblue, boxrule=1pt, arc=2mm,
                  left=10pt, right=10pt]
\textbf{\large\textcolor{primaryblue}{One recursive idea --- go as deep as possible, then back up --- and everything it reveals: cycles, subtrees, timestamps and the DFS tree. Learn to read a DFS as a story of entering and leaving vertices, and most of the structure of a graph falls out of that story for free.}}
\end{tcolorbox}

\chapterfacts{1000}{1900}{5}{12}{30}

\subsection*{How This Chapter Is Organised}

\begin{diagrambox}
\begin{tikzpicture}[node distance=3mm and 3mm,
    th/.style={rectangle, rounded corners=3pt, draw=theoryborder, fill=theorybg, line width=1pt,
               font=\small\bfseries, align=center, minimum height=10mm, text width=13.2cm},
    pt/.style={rectangle, rounded corners=3pt, draw=patternborder, fill=patternbg, line width=1pt,
               font=\scriptsize, align=center, minimum height=12mm, text width=2.35cm},
    pr/.style={rectangle, rounded corners=2pt, draw=primaryblue, fill=white,
               font=\scriptsize, align=center, text width=2.35cm, minimum height=6mm}]
    \node[th] (T) {Part A --- How DFS Works and What It Reveals (\S\ref{sec:ch2-lessons})\\[-1pt]
        {\scriptsize\mdseries the call stack $\cdot$ colours \& timestamps $\cdot$ the DFS tree \& edge types $\cdot$ cycles $\cdot$ the DFS toolbox}};
    \node[pt, below=6mm of T.south west, anchor=north west, xshift=1mm] (P1) {\textbf{P1 Cycle}\\\textbf{Detection}\\find it, print it};
    \node[pt, right=of P1] (P2) {\textbf{P2 Subtree}\\\textbf{Aggregation}\\up and down};
    \node[pt, right=of P2] (P3) {\textbf{P3 Entry \&}\\\textbf{Exit Times}\\intervals \& orders};
    \node[pt, right=of P3] (P4) {\textbf{P4 Back-}\\\textbf{tracking}\\all simple paths};
    \node[pt, right=of P4] (P5) {\textbf{P5 Reverse}\\\textbf{Search}\\start at the target};
    \foreach \p/\a/\b in {P1/Round Trip/Round Trip II,
                          P2/White-Black Balanced Subtrees/Count Nodes With the Highest Score,
                          P3/Military Problem/Count Descendants,
                          P4/All Paths From Source to Target/Path with Maximum Gold,
                          P5/All Ancestors in a DAG/Solve The Maze}{
        \node[pr, below=of \p] (\p a) {\a};
        \node[pr, below=2mm of \p a] (\p b) {\b};
    }
    \node[pr, below=2mm of P2b] (P2c) {Ki};
    \node[pr, below=2mm of P3b] (P3c) {Maze};
    \foreach \p in {P1,P2,P3,P4,P5} \draw[flow] (T.south -| \p.north) -- (\p.north);
    \node[th, anchor=north, at={([yshift=-5mm]P3b.south |- P3c.south)}, draw=successgreen, fill=solutionbg] (S) {Summary, Cheat Sheet, Practice Set \& Solution Sketches};
    \node[font=\small\bfseries\color{patternborder}, rotate=90, anchor=south] at ([xshift=-3mm]P1a.west |- P1a.south) {Part B};
\end{tikzpicture}
\end{diagrambox}

\subsection*{Before You Start}

You need Chapter~1: adjacency lists, the \texttt{visited} array, and the recursive and iterative DFS of Lesson~\ref{ssec:ch1-first-traversal}. This chapter starts from exactly those two programs and asks what else they can tell us. Recursion in C++ (what a call stack is, what happens when a function returns) is assumed; Lesson~\ref{ssec:ch2-call-stack} recalls the parts that matter for graphs.

\subsection*{After This Chapter You Will Be Able To}
\begin{enumerate}[leftmargin=*]
    \item \textbf{Write DFS three ways} --- recursive, iterative with a per-vertex iterator, and ``order array'' style --- and choose by the recursion depth the input allows.
    \item \textbf{Timestamp} a DFS with colours and entry/exit times, and test ``is $u$ an ancestor of $v$?'' in \bigO{1} with the parenthesis theorem.
    \item \textbf{Classify} every edge as tree, back, forward or cross, and use the classification: a directed graph has a cycle iff DFS finds a back edge.
    \item \textbf{Find and print} a cycle in undirected and directed graphs, including the multi-edge and self-loop traps.
    \item \textbf{Compute} a value for every subtree (sizes, counts, sums) bottom-up, and push values down from the root top-down.
    \item \textbf{Enumerate} all simple paths with backtracking when the graph is tiny, and \textbf{reverse} a search so that it starts from the target.
\end{enumerate}

\chaptercontents

%============================================================
\section{Part A --- How DFS Works and What It Reveals}
\label{sec:ch2-lessons}
%============================================================
%============================================================
% CH2 · LESSON 1 — RECURSION AND THE CALL STACK
%============================================================
\subsection{Recursion and the Call Stack}
\label{ssec:ch2-call-stack}

Chapter~1 used DFS only to answer ``which vertices can I reach?''. For that question the \emph{order} of the visit did not matter --- BFS gave the same set. This chapter is about the order. DFS has a very particular rhythm: it \textbf{enters} a vertex, disappears into the first unexplored neighbour, comes back, disappears into the next one, and finally \textbf{leaves} the vertex for good. Almost everything in this chapter --- timestamps, the DFS tree, cycle detection, subtree sizes --- is read off that rhythm of entering and leaving.

% local drawing of the running example D: a pic whose vertex / edge styles are overridden per snapshot
\tikzset{
    g2-v1/.style=gv, g2-v2/.style=gv, g2-v3/.style=gv, g2-v4/.style=gv,
    g2-v5/.style=gv, g2-v6/.style=gv, g2-v7/.style=gv, g2-v8/.style=gv,
    g2-e12/.style=gde, g2-e15/.style=gde, g2-e16/.style=gde, g2-e23/.style=gde, g2-e24/.style=gde,
    g2-e31/.style=gde, g2-e43/.style=gde, g2-e54/.style=gde, g2-e56/.style=gde, g2-e67/.style=gde,
    g2-e76/.style=gde, g2-e87/.style=gde,
    g2-D/.pic={
        \node[g2-v1] (d1) at (0,0) {1};
        \node[g2-v2] (d2) at (1.6,0.9) {2};
        \node[g2-v3] (d3) at (3.2,1.5) {3};
        \node[g2-v4] (d4) at (3.2,0.2) {4};
        \node[g2-v5] (d5) at (1.6,-1.0) {5};
        \node[g2-v6] (d6) at (3.2,-1.0) {6};
        \node[g2-v7] (d7) at (4.8,-1.0) {7};
        \node[g2-v8] (d8) at (4.8,0.2) {8};
        \draw[g2-e12] (d1) -- (d2); \draw[g2-e15] (d1) -- (d5);
        \draw[g2-e16] (d1) to[bend right=48] (d6);
        \draw[g2-e23] (d2) -- (d3); \draw[g2-e24] (d2) -- (d4);
        \draw[g2-e31] (d3) to[bend right=40] (d1);
        \draw[g2-e43] (d4) -- (d3); \draw[g2-e54] (d5) -- (d4); \draw[g2-e56] (d5) -- (d6);
        \draw[g2-e67] (d6) to[bend left=22] (d7); \draw[g2-e76] (d7) to[bend left=22] (d6);
        \draw[g2-e87] (d8) -- (d7);
    },
    % G3 of Chapter 1 (undirected), redrawn here so that this chapter compiles on its own
    g2-u1/.style=gv, g2-u2/.style=gv, g2-u3/.style=gv, g2-u4/.style=gv, g2-u5/.style=gv, g2-u6/.style=gv, g2-u7/.style=gv,
    g2-f12/.style=ge, g2-f14/.style=ge, g2-f23/.style=ge, g2-f25/.style=ge, g2-f45/.style=ge, g2-f56/.style=ge,
    g2-G3/.pic={
        \node[g2-u1] (w1) at (0,0) {1};
        \node[g2-u2] (w2) at (1.4,0.9) {2};
        \node[g2-u4] (w4) at (1.4,-0.9) {4};
        \node[g2-u3] (w3) at (2.9,0.9) {3};
        \node[g2-u5] (w5) at (2.9,-0.9) {5};
        \node[g2-u6] (w6) at (4.3,-0.9) {6};
        \node[g2-u7] (w7) at (4.3,0.9) {7};
        \draw[g2-f12] (w1) -- (w2); \draw[g2-f14] (w1) -- (w4); \draw[g2-f23] (w2) -- (w3);
        \draw[g2-f25] (w2) -- (w5); \draw[g2-f45] (w4) -- (w5); \draw[g2-f56] (w5) -- (w6);
    }
}

\subsubsection*{The running example $D$}

Part~A uses one small directed graph throughout. It is drawn so that it contains everything a DFS can meet: two cycles, a vertex that is reached twice, an edge that ``jumps ahead'' and a vertex (8) that cannot be reached from 1.

\begin{diagrambox}
\begin{tikzpicture}[x=1.15cm, y=1cm]
    \pic {g2-D};
    \node[draw=darkgray, fill=lightgray!60, rounded corners=2pt, font=\ttfamily\scriptsize, align=left, inner sep=3pt, anchor=west] (inp) at (6.6,0.1)
        {8 12\\1 2\\1 5\\1 6\\2 3\\2 4\\3 1\\4 3\\5 4\\5 6\\6 7\\7 6\\8 7};
    \node[font=\scriptsize\color{darkgray}, above=1pt of inp] {input (directed)};
    \node[note, anchor=west, text width=3.4cm, align=left] at (8.3,0.1) {adj[1] = 2 5 6\\adj[2] = 3 4\\adj[3] = 1\\adj[4] = 3\\adj[5] = 4 6\\adj[6] = 7\\adj[7] = 6\\adj[8] = 7};
\end{tikzpicture}
\figcap{The directed graph $D$: $n = 8$, $m = 12$. Neighbours are tried in the order of the adjacency lists, i.e.\ in input order; this fixes every number in Part~A. Cycles: $1 \to 2 \to 3 \to 1$ and $6 \rightleftarrows 7$.}
\end{diagrambox}

%------------------------------------------------------------
\subsubsection*{What the call stack holds}

Here is the recursive DFS of Chapter~1, with one change: it now does something both when it \emph{enters} a vertex and when it \emph{leaves} it.

\begin{lstlisting}[style=cpp]
#include <bits/stdc++.h>
using namespace std;

// Recursive DFS with the two "hooks": code before the loop runs when u is ENTERED,
// code after the loop runs when u is LEFT (all its neighbours are done).
vector<vector<int>> adj;
vector<bool> visited;
vector<int> pre, post;                 // entry order and exit order

void dfs(int u) {
    visited[u] = true;
    pre.push_back(u);                  // hook 1: pre-order (entering u)
    for (int v : adj[u])
        if (!visited[v])
            dfs(v);                    // while dfs(v) runs, u waits on the call stack
    post.push_back(u);                 // hook 2: post-order (leaving u)
}

int main() {
    int n, m;
    cin >> n >> m;
    adj.assign(n + 1, vector<int>());
    for (int i = 0; i < m; i++) {
        int a, b;
        cin >> a >> b;
        adj[a].push_back(b);           // directed edge a -> b
    }
    visited.assign(n + 1, false);
    for (int s = 1; s <= n; s++)       // restart from every vertex not reached yet
        if (!visited[s]) dfs(s);
    cout << "pre :";
    for (int v : pre) cout << " " << v;
    cout << "\npost:";
    for (int v : post) cout << " " << v;
    cout << "\n";
}
\end{lstlisting}

\begin{lstlisting}[style=trace]
input : the input of D above
output: pre : 1 2 3 4 5 6 7 8
        post: 3 4 2 7 6 5 1 8
\end{lstlisting}

Every call \texttt{dfs(u)} owns a \textbf{stack frame} that stores two things: the vertex $u$, and \emph{how far the loop over \texttt{adj[u]} has got}. When \texttt{dfs(v)} returns, the frame of $u$ resumes exactly where it stopped. The frames that are open at a given moment form the \textbf{call stack}, and the call stack always spells a path in the graph: from the start vertex, along the edges through which DFS went deeper, down to the vertex currently being explored.

\begin{diagrambox}
\begin{tikzpicture}[x=0.6cm, y=0.6cm,
    fr/.style={draw=recursionred, fill=recursionred!14, minimum width=1.25cm, minimum height=0.42cm,
               font=\scriptsize\ttfamily, inner sep=1pt},
    frtop/.style={fr, fill=currentfill, draw=accentorange, line width=0.9pt}]
    \foreach \k/\lbl/\st in {0/{enter 1}/{1}, 1/{enter 2}/{1,2}, 2/{enter 3}/{1,2,3}, 3/{leave 3, enter 4}/{1,2,4},
                             4/{leave 4, 2; enter 5}/{1,5}, 5/{enter 6}/{1,5,6}, 6/{enter 7}/{1,5,6,7}, 7/{leave 7, 6, 5, 1}/{},
                             8/{restart: enter 8}/{8}}{
        \begin{scope}[xshift=\k*1.55cm]
            \foreach \v [count=\h from 0] in \st {
                \node[fr] at (0,\h*0.75) {dfs(\v)};
            }
            \draw[darkgray, thick] (-1.15,-0.45) -- (1.15,-0.45);
            \node[font=\tiny, text width=1.45cm, align=center, anchor=north] at (0,-0.6) {\lbl};
        \end{scope}
    }
    \node[font=\scriptsize\bfseries, text=recursionred, anchor=west] at (-1.2,3.6) {call stack (top = vertex being explored)};
\end{tikzpicture}
\figcap{The call stack of the DFS on $D$ after each step. The deepest moment is \texttt{1 5 6 7}: four open frames. After vertex 1 is left the stack is empty, and the outer loop restarts the search from 8, the only vertex not reached from 1.}
\end{diagrambox}

\begin{insightbox}[The call stack is the path from the root]
At every moment the open frames, read from bottom to top, are a path $s = x_0 \to x_1 \to \dots \to x_k = u$ in the graph, where each step $x_{i} \to x_{i+1}$ is the edge through which DFS entered $x_{i+1}$. A vertex is \emph{on the stack} exactly between the moment it is entered and the moment it is left. Lessons~\ref{ssec:ch2-timestamps}--\ref{ssec:ch2-cycles} turn this one sentence into the parenthesis theorem, the edge types and cycle detection.
\end{insightbox}

\subsubsection*{Two hooks, two orders}

\begin{itemize}[leftmargin=*]
    \item \textbf{Pre-order} (code before the loop): a vertex is processed \emph{before} everything below it. Use it to pass information \emph{down}: depths, the path from the root, a value accumulated from the root (Problem~\ref{ssec:ch2-ki}).
    \item \textbf{Post-order} (code after the loop): a vertex is processed \emph{after} everything below it. Use it to collect information \emph{up}: subtree sizes, counts, sums (Pattern~2), and the finishing order behind topological sorting (Chapter~5).
\end{itemize}
On $D$ the pre-order is $1,2,3,4,5,6,7,8$ and the post-order is $3,4,2,7,6,5,1,8$. Vertex 1 is first in pre-order and almost last in post-order: everything reachable from it is entered after it and left before it. That observation, made precise, is the parenthesis theorem of the next lesson.

%------------------------------------------------------------
\subsubsection*{Iterative DFS that also knows when a vertex is left}

The iterative DFS of Chapter~1 replaced the call stack by a vector \texttt{st} and remembered, for every vertex, the index \texttt{it[u]} of the next neighbour to try --- that index \emph{is} the ``how far the loop has got'' part of a stack frame. The vertex is left when \texttt{it[u]} reaches the end of its list. So both hooks are available, and the program below prints the same two orders as the recursive one.

\begin{lstlisting}[style=cpp]
#include <bits/stdc++.h>
using namespace std;

// Iterative DFS with entry AND exit events. st[] plays the call stack, it[u] the loop position.
// Same order as the recursive version, no recursion depth limit.
int main() {
    int n, m;
    cin >> n >> m;
    vector<vector<int>> adj(n + 1);
    for (int i = 0; i < m; i++) {
        int a, b;
        cin >> a >> b;
        adj[a].push_back(b);
    }
    vector<bool> visited(n + 1, false);
    vector<size_t> it(n + 1, 0);
    vector<int> st, pre, post;
    for (int s = 1; s <= n; s++) {
        if (visited[s]) continue;
        visited[s] = true; pre.push_back(s); st.push_back(s);      // enter s
        while (!st.empty()) {
            int u = st.back();
            if (it[u] < adj[u].size()) {
                int v = adj[u][it[u]++];
                if (!visited[v]) {
                    visited[v] = true; pre.push_back(v);              // enter v ("call" dfs(v))
                    st.push_back(v);
                }
            } else {
                post.push_back(u);                                    // leave u ("return")
                st.pop_back();
            }
        }
    }
    cout << "pre :";
    for (int v : pre) cout << " " << v;
    cout << "\npost:";
    for (int v : post) cout << " " << v;
    cout << "\n";
}
\end{lstlisting}

\begin{lstlisting}[style=trace]
input : the input of D above
output: pre : 1 2 3 4 5 6 7 8          (identical to the recursive version)
        post: 3 4 2 7 6 5 1 8
\end{lstlisting}

This is the version kept in the book's library (\texttt{code/lib/dfs\_iterative.hpp}, Appendix~A). It costs \VE{} time like the recursive one, and \texttt{st} never holds more than one copy of a vertex, so its memory is \bigO{n}.

\subsubsection*{Two tempting shortcuts that are \emph{not} DFS}

Many tutorials write ``iterative DFS'' by pushing all neighbours at once. Both variants below are fine for plain reachability, but neither gives a real DFS.

\begin{diagrambox}
\begin{tikzpicture}[x=0.95cm, y=0.95cm]
    \begin{scope}
        \tikzset{g2-e16/.style={gtree, -{Stealth[length=2.4mm]}}, g2-e12/.style={gtree, -{Stealth[length=2.4mm]}},
                 g2-e15/.style={gtree, -{Stealth[length=2.4mm]}}, g2-e67/.style={gtree, -{Stealth[length=2.4mm]}},
                 g2-e54/.style={gtree, -{Stealth[length=2.4mm]}}, g2-e43/.style={gtree, -{Stealth[length=2.4mm]}},
                 g2-v6/.style=gvB}
        \pic {g2-D};
        \node[font=\small\bfseries, anchor=west] at (-0.4,2.45) {(a) mark when pushed};
        \node[font=\scriptsize, text width=6.3cm, align=left, anchor=north west] at (-0.4,-2.0)
            {Pop $u$, then push every unmarked neighbour and mark it at once. Order $1,2,5,6,7,4,3$ (8 not reached from 1). Vertex 6 gets parent 1 (red) although a real DFS enters it from 5: the ``tree'' is wrong.};
    \end{scope}
    \begin{scope}[xshift=7.6cm]
        \tikzset{g2-v6/.style=gvB}
        \pic {g2-D};
        \node[font=\small\bfseries, anchor=west] at (-0.4,2.45) {(b) mark when popped};
        \node[font=\scriptsize, text width=6.3cm, align=left, anchor=north west] at (-0.4,-2.0)
            {Pop $u$; skip it if marked; else mark it and push all unmarked neighbours (in reverse, so the first neighbour is on top). Order $1,2,3,4,5,6,7$ --- correct here. But vertex 6 sits on the stack \textbf{twice} (pushed by 1 and by 5), the stack is not a path, and nothing tells you when a vertex is \emph{left}.};
    \end{scope}
\end{tikzpicture}
\figcap{The two shortcuts on $D$ (orders produced by running them). Neither has an exit event, so neither can compute post-order values, exit times or cycles.}
\end{diagrambox}

\begin{warningbox}[Use the iterator version whenever the order matters]
Variant (a) is not even a DFS order. Variant (b) visits vertices in a DFS order, but its stack can hold up to $m$ entries and it cannot tell when a vertex is finished. Whenever a solution needs exit times, post-order values or ``is $v$ on the current path?'', write the version with \texttt{it[u]} (or recursion).
\end{warningbox}

%------------------------------------------------------------
\subsubsection*{How deep can recursion go?}

Each open frame occupies stack memory, and the stack is small: the operating system reserves a fixed amount for it, independent of the memory limit of the problem. When the recursion goes deeper than the stack allows, the program dies with a \stackoverflow{} --- usually reported as ``runtime error'' with no message at all.

To get real numbers we ran the recursive \texttt{dfs} above on a path $0 - 1 - 2 - \dots - (n-1)$ (depth $= n$) on a Windows machine with MinGW g++ 6.3 and its default 2\,MB stack:

\begin{center}
\small
\begin{tabular}{@{}lccc@{}}
\toprule
\textbf{compiled with} & \textbf{largest $n$ that worked} & \textbf{bytes per frame} & \textbf{$n = 10^6$} \\
\midrule
\texttt{g++ -O2} & about $65\,000$ & $\approx 32$ & crash \\
\texttt{g++ -O0} (no optimisation) & about $43\,000$ & $\approx 48$ & crash \\
\texttt{g++ -O2 -Wl,--stack=268435456} (256\,MB) & --- & $\approx 32$ & works \\
\bottomrule
\end{tabular}
\end{center}

\begin{diagrambox}
\begin{tikzpicture}[x=1cm, y=1cm]
    \draw[->, thick, darkgray] (0,0) -- (12.2,0) node[right, font=\scriptsize] {depth};
    \foreach \x/\l in {0/0, 2/$10^4$, 4/$5{\cdot}10^4$, 6/$10^5$, 8/$2{\cdot}10^5$, 10/$10^6$}{
        \draw[darkgray] (\x,0.08) -- (\x,-0.08); \node[font=\scriptsize, below] at (\x,-0.1) {\l};}
    \fill[successgreen!25] (0,0.25) rectangle (4.6,0.75);
    \node[font=\scriptsize] at (2.3,0.5) {safe everywhere};
    \fill[accentorange!25] (4.6,0.25) rectangle (8.7,0.75);
    \node[font=\scriptsize] at (6.65,0.5) {depends on the judge};
    \fill[warningred!20] (8.7,0.25) rectangle (12,0.75);
    \node[font=\scriptsize] at (10.35,0.5) {only with a huge stack};
    \draw[warningred, thick, dashed] (4.6,0.1) -- (4.6,1.2) node[above, font=\tiny, text=warningred] {2\,MB default: crash near $6.5\cdot10^4$};
\end{tikzpicture}
\figcap{Rule of thumb. Codeforces runs solutions with a large stack (256\,MB), so depth $10^6$ is usually fine there; at home, on many other judges and inside LeetCode's harness the default is a few megabytes.}
\end{diagrambox}

\textbf{What to do.} Before writing a recursive DFS, ask how deep it can go: the worst case is the longest path the DFS can walk, which is $n$ on a path graph. If that can exceed a few $10^4$ and you do not know the judge's stack, use the iterative version. A second, often simpler, way out is the \emph{order array} trick below.

%------------------------------------------------------------
\subsubsection*{Post-order work without recursion: the order array}

Many problems need post-order work only (``compute something for every vertex from its children''). Then a full iterative DFS is not necessary:
\begin{enumerate}[leftmargin=*]
    \item run any iterative search from the root that records \texttt{order} (vertices in the order they are reached) and \texttt{parent};
    \item loop over \texttt{order} \textbf{backwards}: every vertex appears after its parent in \texttt{order}, so going backwards every vertex is processed after all its children, and can add its value to its parent.
\end{enumerate}
\begin{lstlisting}[style=snippet]
for (int i = n - 1; i >= 1; i--) {      // order[0] is the root, it has no parent
    int v = order[i];
    sz[parent[v]] += sz[v];              // all children of v were handled before v
}
\end{lstlisting}
On the DFS forest of $D$ (tree edges $1{\to}2, 2{\to}3, 2{\to}4, 1{\to}5, 5{\to}6, 6{\to}7$; vertex 8 alone) this gives the subtree sizes $\mathrm{sz} = (7,3,1,1,3,2,1,1)$ for vertices $1..8$. The trick works with \emph{any} search order in which parents come before children --- even BFS order --- because it only uses the tree, not the DFS rhythm. It is the workhorse of Pattern~2 when the input is a tree with up to $2\cdot10^5$ vertices.

%------------------------------------------------------------
\subsubsection*{Check your understanding}

\begin{enumerate}[leftmargin=*]
    \item In $D$, change \texttt{adj[1]} to \texttt{6 5 2}. What are the new pre-order and post-order?
    \item While the DFS on $D$ is inside \texttt{dfs(4)}, which frames are open, and which vertices are already left?
    \item Why can variant (b) above hold the same vertex twice on its stack, but the iterator version never?
    \item A tree with $2\cdot 10^5$ vertices is given as a parent array with $p_i < i$. Give an \bigO{n} way to compute all subtree sizes with no search at all.
    \item Your recursive DFS passes the samples and gets ``runtime error'' on test 7, where $n = 2\cdot10^5$. Name the most likely cause and two fixes.
\end{enumerate}

\begin{solutionbox}[Answers]
\begin{enumerate}[leftmargin=*]
    \item DFS from 1 now enters 6, 7, then 5 (6 is done, 5 tries 4), 4, 3, and finally 2 (both its neighbours are done), then 8. Running the program on the changed input prints \textbf{pre: 1 6 7 5 4 3 2 8} and \textbf{post: 7 6 3 4 5 2 1 8}.
    \item Open frames (bottom to top): \texttt{dfs(1)}, \texttt{dfs(2)}, \texttt{dfs(4)}; the path $1 \to 2 \to 4$. Already left: only 3.
    \item Variant (b) pushes every unmarked neighbour, and a vertex can be unmarked when two different vertices push it (6 is pushed by 1 and by 5 in $D$). The iterator version pushes $v$ only at the moment it marks $v$, so each vertex is pushed at most once.
    \item Since $p_i < i$, every child has a larger label than its parent. Loop $i = n, n-1, \dots, 2$ and do \texttt{sz[p[i]] += sz[i]}: the labels themselves are an order with parents before children (the order-array trick with $\texttt{order} = 1,2,\dots,n$).
    \item A stack overflow on a deep tree (a path, or a long ``caterpillar''). Fix 1: the iterative DFS with \texttt{it[u]}. Fix 2: if only post-order work is needed, the order-array trick. (On Codeforces the large stack often hides the problem; on other judges it does not.)
\end{enumerate}
\end{solutionbox}

%============================================================
% CH2 · LESSON 2 — COLOURS, ENTRY AND EXIT TIMES
%============================================================
\subsection{Colours, Entry and Exit Times}
\label{ssec:ch2-timestamps}

The \texttt{visited} array answers one question: \emph{has DFS been here?} For everything else in this chapter we need a finer answer: \emph{has DFS been here, and is it still here?} The call stack of Lesson~\ref{ssec:ch2-call-stack} suggests three states, and a clock that ticks whenever DFS enters or leaves a vertex turns them into numbers.

\subsubsection*{Three colours}

\begin{definitionbox}[White, gray, black]
During a DFS every vertex has one of three colours:
\begin{center}
\legendchip{gv}{\textbf{white} --- not entered yet} \qquad \legendchip{gvseen}{\textbf{gray} --- entered, not left (on the stack)} \qquad \legendchip{gvdone}{\textbf{black} --- left (finished)}
\end{center}
A vertex turns gray when \texttt{dfs(u)} starts and black when it returns. It never changes colour again. The gray vertices, at any moment, are exactly the vertices of the call stack --- a path from the root to the current vertex.
\end{definitionbox}

In code the colours are one \texttt{int} per vertex, \texttt{0/1/2}; the old \texttt{visited[u]} is simply ``colour $\ne$ white''. The difference between gray and black is what Lesson~\ref{ssec:ch2-cycles} needs to detect cycles in directed graphs.

\subsubsection*{Entry and exit times}

Keep a global counter \texttt{timer}. When DFS enters $u$, record $\tin{u} = \texttt{timer++}$; when it leaves $u$, record $\tout{u} = \texttt{timer++}$. With one clock for both events, a DFS of $n$ vertices uses the times $0, 1, \dots, 2n-1$, each exactly once.

\begin{lstlisting}[style=snippet]
void dfs(int u) {
    color[u] = 1; tin[u] = timer++;          // u turns gray
    for (int v : adj[u])
        if (color[v] == 0) { parent[v] = u; dfs(v); }
    color[u] = 2; tout[u] = timer++;         // u turns black
}
\end{lstlisting}

\begin{diagrambox}
\begin{tikzpicture}[x=0.62cm, y=0.62cm]
    \tikzset{gv/.append style={minimum size=6mm, font=\scriptsize\bfseries}}
    \begin{scope}[g2-v1/.style=gvseen, g2-v2/.style=gvseen, g2-v3/.style=gvdone, g2-v4/.style=gvcur,
                  g2-e12/.style={gtree, -{Stealth[length=2.2mm]}}, g2-e23/.style={gtree, -{Stealth[length=2.2mm]}},
                  g2-e24/.style={gtree, -{Stealth[length=2.2mm]}}]
        \pic {g2-D};
        \node[font=\scriptsize, anchor=west] at (-0.5,-2.6) {\textbf{time 4:} enter 4 \quad stack \texttt{1 2 4}};
    \end{scope}
    \begin{scope}[xshift=4.6cm, g2-v1/.style=gvseen, g2-v2/.style=gvdone, g2-v3/.style=gvdone, g2-v4/.style=gvdone,
                  g2-v5/.style=gvseen, g2-v6/.style=gvseen, g2-v7/.style=gvcur,
                  g2-e12/.style={gtree, -{Stealth[length=2.2mm]}}, g2-e23/.style={gtree, -{Stealth[length=2.2mm]}},
                  g2-e24/.style={gtree, -{Stealth[length=2.2mm]}}, g2-e15/.style={gtree, -{Stealth[length=2.2mm]}},
                  g2-e56/.style={gtree, -{Stealth[length=2.2mm]}}, g2-e67/.style={gtree, -{Stealth[length=2.2mm]}}]
        \pic {g2-D};
        \node[font=\scriptsize, anchor=west] at (-0.5,-2.6) {\textbf{time 9:} enter 7 \quad stack \texttt{1 5 6 7}};
    \end{scope}
    \begin{scope}[xshift=9.2cm, g2-v1/.style=gvdone, g2-v2/.style=gvdone, g2-v3/.style=gvdone, g2-v4/.style=gvdone,
                  g2-v5/.style=gvdone, g2-v6/.style=gvdone, g2-v7/.style=gvdone, g2-v8/.style=gvcur,
                  g2-e12/.style={gtree, -{Stealth[length=2.2mm]}}, g2-e23/.style={gtree, -{Stealth[length=2.2mm]}},
                  g2-e24/.style={gtree, -{Stealth[length=2.2mm]}}, g2-e15/.style={gtree, -{Stealth[length=2.2mm]}},
                  g2-e56/.style={gtree, -{Stealth[length=2.2mm]}}, g2-e67/.style={gtree, -{Stealth[length=2.2mm]}}]
        \pic {g2-D};
        \node[font=\scriptsize, anchor=west] at (-0.5,-2.6) {\textbf{time 14:} enter 8 \quad stack \texttt{8}};
    \end{scope}
\end{tikzpicture}
\figcap{Three moments of the DFS on $D$. Orange $=$ current vertex, light orange $=$ gray (on the stack), blue $=$ black. Thick blue arrows are the edges along which DFS entered a new vertex.}
\end{diagrambox}

Running the DFS on $D$ (outer loop over $s = 1..8$) gives:
\begin{center}
\small
\begin{tabular}{@{}l*{8}{c}@{}}
\toprule
vertex $u$ & 1 & 2 & 3 & 4 & 5 & 6 & 7 & 8 \\
\midrule
$\tin{u}$ & 0 & 1 & 2 & 4 & 7 & 8 & 9 & 14 \\
$\tout{u}$ & 13 & 6 & 3 & 5 & 12 & 11 & 10 & 15 \\
parent & --- & 1 & 2 & 2 & 1 & 5 & 6 & --- \\
\bottomrule
\end{tabular}
\end{center}

\begin{diagrambox}
\begin{tikzpicture}[x=0.78cm, y=0.5cm]
    \draw[->, darkgray] (-0.3,0) -- (16.2,0) node[right, font=\scriptsize] {time};
    \foreach \t in {0,...,15} {\draw[darkgray] (\t,0.12) -- (\t,-0.12); \node[font=\tiny, below] at (\t,-0.12) {\t};}
    % intervals: vertex/tin/tout/level/colour
    \foreach \v/\a/\b/\l/\c in {1/0/13/5/primaryblue, 2/1/6/4/successgreen, 3/2/3/3/accentorange, 4/4/5/3/accentorange,
                                5/7/12/4/successgreen, 6/8/11/3/accentorange, 7/9/10/2/softpurple, 8/14/15/5/primaryblue}{
        \draw[\c, line width=2.2pt] (\a,\l*0.62) -- (\b,\l*0.62);
        \draw[\c, line width=1pt] (\a,\l*0.62-0.25) -- (\a,\l*0.62+0.25);
        \draw[\c, line width=1pt] (\b,\l*0.62-0.25) -- (\b,\l*0.62+0.25);
        \node[font=\scriptsize\bfseries, fill=white, inner sep=1pt, text=\c] at ({(\a+\b)/2},\l*0.62+0.42) {\v};
    }
    \node[font=\ttfamily\small, anchor=west] at (-0.3,-1.5) {( 1 ( 2 ( 3 3 ) ( 4 4 ) 2 ) ( 5 ( 6 ( 7 7 ) 6 ) 5 ) 1 ) ( 8 8 )};
\end{tikzpicture}
\figcap{The intervals $[\tin{u}, \tout{u}]$ on the time line. Writing ``\texttt{(}$u$'' at $\tin{u}$ and ``$u$\texttt{)}'' at $\tout{u}$ gives a correctly nested bracket sequence --- hence the name of the next theorem.}
\end{diagrambox}

\subsubsection*{The parenthesis theorem}

\begin{theorembox}[Parenthesis theorem]
For any two vertices $u \ne v$ reached by a DFS, exactly one of the following holds:
\begin{enumerate}[leftmargin=*]
    \item the intervals $[\tin{u}, \tout{u}]$ and $[\tin{v}, \tout{v}]$ are \textbf{disjoint}, and neither vertex is a descendant of the other in the DFS forest;
    \item $[\tin{v}, \tout{v}] \subset [\tin{u}, \tout{u}]$, and $v$ is a \textbf{descendant} of $u$;
    \item $[\tin{u}, \tout{u}] \subset [\tin{v}, \tout{v}]$, and $u$ is a descendant of $v$.
\end{enumerate}
Two intervals never overlap partially.
\end{theorembox}

\textbf{Proof.} Say $\tin{u} < \tin{v}$ (otherwise swap the names). Look at the moment $\tin{v}$ when $v$ is entered.
\emph{Case $u$ is still gray then}, i.e.\ $\tin{v} < \tout{u}$. Gray vertices form the current stack, which is a path from the root through $u$ down to the vertex that is entering $v$ --- so $v$ is entered while \texttt{dfs(u)} is still open and becomes a descendant of $u$. A call returns only after all calls it made have returned, so \texttt{dfs(v)} returns before \texttt{dfs(u)}: $\tout{v} < \tout{u}$. That is case 2.
\emph{Case $u$ is already black}, i.e.\ $\tout{u} < \tin{v}$: the intervals are disjoint. $v$ cannot be a descendant of $u$ (descendants are entered while $u$ is gray), and $u$ cannot be a descendant of $v$ (it was entered first). That is case 1. \hfill$\blacksquare$

\begin{insightbox}[The \bigO{1} ancestor test]
\[
    u \text{ is an ancestor of } v \text{ (or } u = v) \iff \tin{u} \le \tin{v} \;\text{ and }\; \tout{v} \le \tout{u}.
\]
One DFS (\VE) and then every query ``is $u$ above $v$ in the DFS tree?'' costs two comparisons. On $D$: $1$ is an ancestor of $7$ ($0 \le 9$, $10 \le 13$); $2$ is not an ancestor of $6$ ($\tout{6} = 11 > 6$).
\end{insightbox}

\begin{warningbox}[Ancestor in the DFS tree is not the same as ``can reach'']
In $D$ there is an edge $5 \to 4$, so 5 can reach 4 --- but 4 is \emph{not} a descendant of 5: its interval $[4,5]$ lies before $[7,12]$. The DFS reached 4 earlier, from 2. The parenthesis theorem is about the tree that this particular DFS built, not about reachability. (For a rooted \emph{tree} given as input, the two notions coincide, which is why the test is so useful in Patterns~2 and~3.)
\end{warningbox}

\subsubsection*{The white-path theorem}

Which vertices become descendants of $u$? Not ``all vertices reachable from $u$'' (vertex 4 above), but almost:

\begin{theorembox}[White-path theorem]
$v$ becomes a descendant of $u$ in the DFS forest if and only if, at time $\tin{u}$, there is a path from $u$ to $v$ that consists of \textbf{white} vertices only.
\end{theorembox}

\textbf{Proof sketch.} ($\Rightarrow$) The tree path from $u$ to $v$ consists of vertices entered after $u$, so they were white at $\tin{u}$. ($\Leftarrow$) Suppose some vertex of the white path is not a descendant of $u$; let $w$ be the first such vertex and $x$ its predecessor on the path ($x = u$ or a descendant of $u$). DFS examines the edge $x \to w$ before $x$ is left; at that moment $w$ is either still white --- then it becomes a child of $x$ --- or it was entered after $\tin{u}$ and before $\tout{x} \le \tout{u}$, which by the parenthesis theorem makes it a descendant of $u$. Either way a contradiction. \hfill$\blacksquare$

On $D$ at time $\tin{5} = 7$, vertex 4 is already black, so the path $5 \to 4$ is not white and 4 is not a descendant of 5; the path $5 \to 6 \to 7$ is white, and indeed 6 and 7 end up below 5. The white-path theorem is the key step in the correctness proofs of cycle detection (Lesson~\ref{ssec:ch2-cycles}) and, much later, of strongly connected components (Chapter~19).

\subsubsection*{A second convention: pre-order numbers and subtree sizes}

The ``one clock'' times above are perfect for comparisons. For \emph{storing} subtree data in an array, a second convention is more convenient: number the vertices $0, 1, \dots, n-1$ in pre-order (call it $\mathrm{pos}[u]$) and compute the subtree sizes $\mathrm{sz}[u]$. Then the subtree of $u$ occupies exactly the block
\[
    \mathrm{pos}[u],\ \mathrm{pos}[u] + 1,\ \dots,\ \mathrm{pos}[u] + \mathrm{sz}[u] - 1
\]
of the pre-order, and $u$ is an ancestor of $v$ iff $\mathrm{pos}[u] \le \mathrm{pos}[v] < \mathrm{pos}[u] + \mathrm{sz}[u]$.

\begin{diagrambox}
\begin{tikzpicture}[x=1cm, y=1cm]
    \foreach \v/\i in {1/0, 2/1, 3/2, 4/3, 5/4, 6/5, 7/6, 8/7}{
        \node[gcell, minimum size=8mm, font=\small\bfseries] (c\i) at (\i*0.9,0) {\v};
        \node[font=\tiny, text=darkgray] at (\i*0.9,-0.6) {pos \i};
    }
    \draw[primaryblue, line width=1.4pt, decorate, decoration={brace, amplitude=5pt}] (-0.4,0.5) -- (5.8,0.5)
        node[midway, above=5pt, font=\scriptsize, text=primaryblue] {subtree of 1: pos $0..6$ (sz 7)};
    \draw[successgreen, line width=1.2pt, decorate, decoration={brace, amplitude=4pt, mirror}] (0.5,-0.85) -- (3.1,-0.85)
        node[midway, below=4pt, font=\scriptsize, text=successgreen] {subtree of 2: $1..3$};
    \draw[successgreen, line width=1.2pt, decorate, decoration={brace, amplitude=4pt, mirror}] (3.2,-0.85) -- (5.8,-0.85)
        node[midway, below=4pt, font=\scriptsize, text=successgreen] {subtree of 5: $4..6$};
\end{tikzpicture}
\figcap{The pre-order of $D$ is $1,2,3,4,5,6,7,8$ and the sizes are $(7,3,1,1,3,2,1,1)$. Every subtree is a contiguous block. Turning subtree questions into range questions on this array (``flattening'') is the subject of Chapter~13 (tree fundamentals); answering range queries on it fast is Chapter~15 (subtree and path queries). Here we only need the ancestor test and the block itself (Problem~\ref{ssec:ch2-military}).}
\end{diagrambox}

\subsubsection*{Code: timestamps and ancestor queries}

\begin{lstlisting}[style=cpp]
#include <bits/stdc++.h>
using namespace std;

// Colours, entry/exit times (one clock) and O(1) ancestor queries.
// Input: n m, m directed edges, then q queries "u v": is u an ancestor of v in the DFS forest?
vector<vector<int>> adj;
vector<int> color, tin, tout, par;
int timer = 0;

void dfs(int u) {
    color[u] = 1; tin[u] = timer++;              // white -> gray
    for (int v : adj[u])
        if (color[v] == 0) { par[v] = u; dfs(v); }
    color[u] = 2; tout[u] = timer++;             // gray -> black
}

bool isAncestor(int u, int v) {                  // u == v counts as an ancestor
    return tin[u] <= tin[v] && tout[v] <= tout[u];
}

int main() {
    int n, m;
    cin >> n >> m;
    adj.assign(n + 1, vector<int>());
    for (int i = 0; i < m; i++) {
        int a, b;
        cin >> a >> b;
        adj[a].push_back(b);
    }
    color.assign(n + 1, 0); tin.assign(n + 1, -1); tout.assign(n + 1, -1); par.assign(n + 1, 0);
    for (int s = 1; s <= n; s++)
        if (color[s] == 0) dfs(s);
    int q;
    cin >> q;
    while (q--) {
        int u, v;
        cin >> u >> v;
        cout << u << " above " << v << "? " << (isAncestor(u, v) ? "yes" : "no") << "\n";
    }
}
\end{lstlisting}

\begin{lstlisting}[style=trace]
input : D, then 4 queries   1 7   2 6   5 4   8 7
output: 1 above 7? yes
        2 above 6? no
        5 above 4? no        (5 -> 4 is an edge, but 4 was entered from 2)
        8 above 7? no        (8 -> 7 is an edge, but 7 was finished long before)
\end{lstlisting}

\subsubsection*{Check your understanding}

\begin{enumerate}[leftmargin=*]
    \item How many different values does \texttt{timer} take during a DFS of $n$ vertices? What are $\tin{r}$ and $\tout{r}$ for the root $r$ of a DFS from a single start that reaches all $n$ vertices?
    \item For a leaf $u$ of the DFS tree, what is $\tout{u} - \tin{u}$? For a general vertex, express $\tout{u} - \tin{u}$ through the subtree size.
    \item Can $\tin{u} < \tin{v} < \tout{u} < \tout{v}$ ever happen? Why not?
    \item In $D$, list every vertex that is gray at time 10.
    \item With the pre-order convention, which pre-order positions does the subtree of 6 occupy in $D$?
\end{enumerate}

\begin{solutionbox}[Answers]
\begin{enumerate}[leftmargin=*]
    \item $2n$ values, $0..2n-1$: every vertex is entered once and left once. The root is entered first and left last: $\tin{r} = 0$, $\tout{r} = 2n-1$.
    \item A leaf is left right after it is entered: $\tout{u} - \tin{u} = 1$. In general all $\mathrm{sz}[u]$ vertices of the subtree are entered and left inside $u$'s interval, which uses $2\,\mathrm{sz}[u]$ ticks: $\tout{u} - \tin{u} = 2\,\mathrm{sz}[u] - 1$. (On $D$: vertex 5 has $12 - 7 = 5 = 2 \cdot 3 - 1$.)
    \item No --- that would be a partial overlap, excluded by the parenthesis theorem: if $v$ is entered while $u$ is gray, $v$ is left before $u$.
    \item At time 10 vertex 7 is being left; just before that tick the stack was \texttt{1 5 6 7}. After it, gray are \textbf{1, 5, 6} (7 has just turned black).
    \item $\mathrm{pos}[6] = 5$ and $\mathrm{sz}[6] = 2$: positions \textbf{5 and 6} (vertices 6, 7).
\end{enumerate}
\end{solutionbox}

%============================================================
% CH2 · LESSON 3 — THE DFS TREE AND THE FOUR EDGE TYPES
%============================================================
\subsection{The DFS Tree and the Four Edge Types}
\label{ssec:ch2-dfs-tree}

Every vertex except the start vertices is entered through exactly one edge: the edge $u \to v$ that DFS was examining when it found $v$ white. These edges form the \textbf{DFS tree} (one tree per start vertex of the outer loop --- together a \textbf{DFS forest}). The other edges are not wasted: each of them was examined once and found a vertex that was gray or black, and \emph{which} colour it found says a lot about the graph.

\begin{definitionbox}[The four edge types (directed graphs)]
When DFS, sitting at $u$, examines the edge $u \to v$:
\begin{center}
\small
\renewcommand{\arraystretch}{1.3}
\begin{tabular}{@{}llll@{}}
\toprule
\textbf{colour of $v$} & \textbf{extra test} & \textbf{type} & \textbf{meaning in the DFS forest} \\
\midrule
white & --- & \textbf{tree} edge & $v$ becomes a child of $u$ \\
gray & --- & \textbf{back} edge & $v$ is an ancestor of $u$ (or $v = u$, a self-loop) \\
black & $\tin{u} < \tin{v}$ & \textbf{forward} edge & $v$ is a descendant of $u$, but not its child \\
black & $\tin{v} < \tin{u}$ & \textbf{cross} edge & neither is an ancestor of the other \\
\bottomrule
\end{tabular}
\end{center}
\end{definitionbox}

Why these meanings? A gray $v$ is on the stack, i.e.\ on the path from the root to $u$: an ancestor. A black $v$ with $\tin{v} > \tin{u}$ was entered while $u$ was gray, so it is a descendant (parenthesis theorem), and it is not a child since $u \to v$ is not its tree edge. A black $v$ with $\tin{v} < \tin{u}$ was entered before $u$ and is already finished, so its interval ends before $u$'s begins: the two are unrelated.

\begin{diagrambox}
\begin{tikzpicture}[x=1.15cm, y=1cm]
    % the DFS forest of D, drawn as rooted trees
    \node[gvsrc] (t1) at (2,0) {1};
    \node[gv] (t2) at (0.8,-1.2) {2};
    \node[gv] (t5) at (3.2,-1.2) {5};
    \node[gv] (t3) at (0.1,-2.4) {3};
    \node[gv] (t4) at (1.5,-2.4) {4};
    \node[gv] (t6) at (3.2,-2.4) {6};
    \node[gv] (t7) at (3.2,-3.6) {7};
    \node[gvsrc] (t8) at (5.2,0) {8};
    \foreach \a/\b in {1/2, 1/5, 2/3, 2/4, 5/6, 6/7} \draw[gtree, -{Stealth[length=2.4mm]}] (t\a) -- (t\b);
    \draw[gback] (t3) to[bend left=40] (t1);
    \draw[gback] (t7) to[bend right=50] (t6);
    \draw[gfwd] (t1.east) .. controls (4.5,-0.3) and (4.4,-2.0) .. (t6.east);
    \draw[gcross] (t4) -- (t3);
    \draw[gcross] (t5) to[bend left=15] (t4);
    \draw[gcross] (t8) to[bend left=30] (t7);
    \foreach \v/\a/\b/\p in {1/0/13/left, 2/1/6/left, 3/2/3/below, 4/4/5/below, 5/7/12/left, 6/8/11/left, 7/9/10/left, 8/14/15/right}
        \node[gdist, text=darkgray, \p=0.5mm of t\v] {\tiny[\a,\b]};
    % legend
    \node[anchor=west, font=\scriptsize, align=left] at (6.8,-0.6) {%
        \tikz[baseline=-0.5ex]\draw[gtree, -{Stealth[length=2mm]}] (0,0) -- (0.7,0);\ tree: $1{\to}2$, $2{\to}3$, $2{\to}4$, $1{\to}5$, $5{\to}6$, $6{\to}7$\\[3pt]
        \tikz[baseline=-0.5ex]\draw[gback] (0,0) -- (0.7,0);\ back: $3{\to}1$, $7{\to}6$\\[3pt]
        \tikz[baseline=-0.5ex]\draw[gfwd] (0,0) -- (0.7,0);\ forward: $1{\to}6$\\[3pt]
        \tikz[baseline=-0.5ex]\draw[gcross] (0,0) -- (0.7,0);\ cross: $4{\to}3$, $5{\to}4$, $8{\to}7$};
    \node[anchor=west, font=\scriptsize\itshape, text=darkgray, text width=6.3cm] at (6.8,-2.7)
        {Children are drawn left to right in the order DFS entered them. Every cross edge points to the \textbf{left} (to an earlier, finished part of the forest) --- never to the right.};
\end{tikzpicture}
\figcap{The DFS forest of $D$ with every one of its 12 edges classified (the program below prints exactly this). Green border $=$ root of a DFS tree; small numbers $= [\tin{}, \tout{}]$.}
\end{diagrambox}

\begin{insightbox}[Cross edges only point back in time]
An edge $u \to v$ with $\tin{u} < \tin{v}$ can never be a cross edge: when DFS examines it, $v$ is either white (tree edge), or it was entered after $u$ and hence is $u$'s descendant (forward edge). So a cross edge always goes from a later vertex to an earlier, already finished one --- in pictures, from right to left. This asymmetry is used in Chapter~5 (topological order) and Chapter~19 (strongly connected components).
\end{insightbox}

\subsubsection*{Undirected graphs: only tree and back edges}

In an undirected graph each edge $\{u, v\}$ is stored twice and DFS examines it from both ends. Classify an edge by the \emph{first} time it is examined.

\begin{theorembox}[No forward or cross edges in undirected graphs]
In a DFS of an undirected graph every edge is a tree edge or a back edge (it joins a vertex to one of its ancestors).
\end{theorembox}

\textbf{Proof.} Let $\{u, v\}$ be an edge with $\tin{u} < \tin{v}$. While $u$ is gray, DFS scans all of $u$'s neighbours, including $v$; if $v$ is still white then, $\{u,v\}$ becomes a tree edge. Otherwise $v$ was entered while $u$ was gray, so $v$ is a descendant of $u$. Then $v$ is entered and scanned before $u$ is left, and the edge is first examined from $v$'s side, where it leads to the gray ancestor $u$: a back edge. A cross edge would need two vertices with disjoint intervals joined by an edge --- but the earlier one would have pulled the later one into its subtree. \hfill$\blacksquare$

\begin{diagrambox}
\begin{tikzpicture}[x=1.15cm, y=1cm]
    \begin{scope}
        \pic {g2-G3};
        \node[font=\scriptsize\bfseries] at (2.15,-1.8) {$G_3$ (Chapter 1)};
    \end{scope}
    \begin{scope}[xshift=7.2cm, yshift=1.0cm]
        \node[gvsrc] (u1) at (0,0) {1};
        \node[gv] (u2) at (0,-1) {2};
        \node[gv] (u3) at (-0.9,-2) {3};
        \node[gv] (u5) at (0.9,-2) {5};
        \node[gv] (u4) at (0.3,-3) {4};
        \node[gv] (u6) at (1.5,-3) {6};
        \node[gvsrc] (u7) at (2.6,0) {7};
        \foreach \a/\b in {1/2, 2/3, 2/5, 5/4, 5/6} \draw[gtree] (u\a) -- (u\b);
        \draw[gback, -] (u4) to[bend left=55] (u1);
        \node[anchor=west, font=\scriptsize, text width=4.4cm] at (2.3,-1.7) {DFS from 1 (then from 7): five tree edges and one back edge $4$--$1$, which joins 4 to its ancestor 1. No edge joins two different branches.};
    \end{scope}
\end{tikzpicture}
\figcap{The undirected graph $G_3$ of Lesson~\ref{ssec:ch1-first-traversal} and its DFS forest. The edge $4$--$5$ is a tree edge (5 entered 4); the edge $1$--$4$ is first examined from 4, where 1 is gray: a back edge.}
\end{diagrambox}

\begin{warningbox}[The edge back to the parent is not a back edge]
From $v$'s side, the tree edge $\{p, v\}$ to the parent $p$ leads to a gray vertex too. It must be \textbf{skipped}, otherwise every tree edge looks like a cycle of length 2. Skip it by remembering the \emph{edge id} through which $v$ was entered, not the parent \emph{vertex}: with two parallel edges between $p$ and $v$, the second one really is a back edge (a cycle of length 2), and ``skip the parent vertex'' would miss it. Lesson~\ref{ssec:ch2-cycles} builds cycle detection on exactly this rule.
\end{warningbox}

\subsubsection*{Code: classify every edge}

\begin{lstlisting}[style=cpp]
#include <bits/stdc++.h>
using namespace std;

// Classify every edge of a DIRECTED graph as tree / back / forward / cross.
vector<vector<int>> adj;
vector<int> color, tin;
int timer = 0;
map<string, int> cnt;

void dfs(int u) {
    color[u] = 1; tin[u] = timer++;
    for (int v : adj[u]) {
        string type;
        if (color[v] == 0) type = "tree";
        else if (color[v] == 1) type = "back";           // v is on the stack
        else type = (tin[u] < tin[v]) ? "forward" : "cross";
        cout << u << "->" << v << " " << type << "\n";
        cnt[type]++;
        if (color[v] == 0) dfs(v);
    }
    color[u] = 2; timer++;                               // the exit tick (tout) is not needed here
}

int main() {
    int n, m;
    cin >> n >> m;
    adj.assign(n + 1, vector<int>());
    for (int i = 0; i < m; i++) {
        int a, b;
        cin >> a >> b;
        adj[a].push_back(b);
    }
    color.assign(n + 1, 0); tin.assign(n + 1, -1);
    for (int s = 1; s <= n; s++)
        if (color[s] == 0) dfs(s);
    for (auto& p : cnt) cout << p.first << ": " << p.second << "\n";
}
\end{lstlisting}

\begin{lstlisting}[style=trace]
input : D
output: 1->2 tree      2->3 tree      3->1 back      2->4 tree
        4->3 cross     1->5 tree      5->4 cross     5->6 tree
        6->7 tree      7->6 back      1->6 forward   8->7 cross
        back: 2   cross: 3   forward: 1   tree: 6        (printed one per line)
\end{lstlisting}

For undirected graphs the same program needs edge ids, to skip the parent edge and to report each edge once:

\begin{lstlisting}[style=cpp]
#include <bits/stdc++.h>
using namespace std;

// Undirected graph: every edge is a tree edge or a back edge. adj stores (neighbour, edge id).
vector<vector<pair<int, int>>> adj;
vector<int> color, parentEdge;
vector<bool> done;                                       // edge already classified?

void dfs(int u) {
    color[u] = 1;
    for (auto e : adj[u]) {
        int v = e.first, id = e.second;
        if (id == parentEdge[u] || done[id]) continue;   // skip the edge we came in by
        done[id] = true;
        if (color[v] == 0) {
            cout << u << "-" << v << " tree\n";
            parentEdge[v] = id;
            dfs(v);
        } else {
            cout << u << "-" << v << " back\n";          // v is gray: an ancestor of u
        }
    }
    color[u] = 2;
}

int main() {
    int n, m;
    cin >> n >> m;
    adj.assign(n + 1, {});
    for (int i = 0; i < m; i++) {
        int a, b;
        cin >> a >> b;
        adj[a].push_back({b, i});
        adj[b].push_back({a, i});
    }
    color.assign(n + 1, 0); parentEdge.assign(n + 1, -1); done.assign(m, false);
    for (int s = 1; s <= n; s++)
        if (color[s] == 0) dfs(s);
}
\end{lstlisting}

\begin{lstlisting}[style=trace]
input : G3 (7 6 / 1 2 / 1 4 / 2 3 / 2 5 / 4 5 / 5 6)
output: 1-2 tree   2-3 tree   2-5 tree   5-4 tree   4-1 back   5-6 tree
\end{lstlisting}

\begin{comparisonbox}[What each edge type is used for]
\small
\begin{tabularx}{\linewidth}{@{}>{\bfseries}l X@{}}
tree & the DFS tree itself: parents, subtrees, depths (Patterns~2 and~3) \\
back & \textbf{cycles}: a back edge $u \to v$ plus the tree path $v \leadsto u$ is a cycle (Lesson~\ref{ssec:ch2-cycles}, Pattern~1); in undirected graphs, low-links, bridges and cut vertices (Chapter~18) \\
forward & harmless for cycles (the tree path already goes there); matters only in algorithms that care about multiple paths \\
cross & harmless for cycles; important for strongly connected components (Chapter~19) \\
\end{tabularx}
\end{comparisonbox}

\subsubsection*{Check your understanding}

\begin{enumerate}[leftmargin=*]
    \item A DFS of a connected undirected graph with $n$ vertices and $m$ edges: how many tree edges and how many back edges does it find?
    \item Can a forward edge appear in a DFS of an undirected graph? Can a self-loop $u \to u$ in a directed graph be anything other than a back edge?
    \item In $D$, the edge $8 \to 7$ is a cross edge. Restart the outer loop at $s = 8$ first (then $1, 2, \dots$). What is the type of $8 \to 7$ now, and of $6 \to 7$?
    \item Why does the classification of a directed edge need $\tin{}$ only for black targets?
    \item A directed graph has a DFS with no back edges. What can you say about it?
\end{enumerate}

\begin{solutionbox}[Answers]
\begin{enumerate}[leftmargin=*]
    \item $n - 1$ tree edges (one per vertex except the root) and $m - n + 1$ back edges --- every other edge.
    \item No (theorem above). A self-loop finds $u$ itself, which is gray: always a back edge, i.e.\ a cycle of length 1.
    \item DFS from 8 enters 7 via $8 \to 7$ (\textbf{tree}) and 6 via $7 \to 6$ (tree); from 6 the edge $6 \to 7$ finds 7 gray: a \textbf{back} edge. The later search from 1 finds 6 black and entered earlier through both $5 \to 6$ and $1 \to 6$: two cross edges (so $1 \to 6$ changes from forward to cross). The cycle $6 \rightleftarrows 7$ is still detected, through a different edge. (Output of the trace program with the start order $8, 1, 2, \dots$.)
    \item White means tree and gray means back, whatever the times. Only for a black target are two answers possible, and the order of entry decides between descendant (forward) and unrelated (cross).
    \item It is acyclic (a DAG): the next lesson proves that a directed graph has a cycle iff some DFS finds a back edge --- and then \emph{every} DFS finds one.
\end{enumerate}
\end{solutionbox}

%============================================================
% CH2 · LESSON 4 — CYCLE DETECTION AND RECOVERY
%============================================================
\subsection{Cycle Detection and Recovering the Cycle}
\label{ssec:ch2-cycles}

``Does this graph have a cycle, and if so, show me one'' is the most common question DFS answers in contests: dependency graphs must be acyclic, a round trip must exist, a process must not loop forever. Lesson~\ref{ssec:ch2-dfs-tree} already did most of the work: back edges are the cycles.

\subsubsection*{The cycle theorem}

\begin{theorembox}[Back edges and cycles]
\begin{enumerate}[leftmargin=*]
    \item \textbf{Directed:} a directed graph has a cycle if and only if a DFS of it (from every vertex, as in the outer loop) finds a back edge --- an edge $u \to v$ to a \emph{gray} vertex $v$.
    \item \textbf{Undirected:} an undirected graph has a cycle if and only if a DFS finds an edge, other than the one it came in by, to an already visited vertex.
\end{enumerate}
In both cases the cycle is the back edge plus the tree path from $v$ down to $u$.
\end{theorembox}

\textbf{Proof (directed).} ($\Leftarrow$) A back edge $u \to v$ goes to a gray vertex, i.e.\ an ancestor of $u$; the tree path $v \leadsto u$ plus the edge $u \to v$ is a directed cycle.
($\Rightarrow$) Let $C$ be a directed cycle and $v$ the vertex of $C$ that DFS enters first. At time $\tin{v}$ all other vertices of $C$ are white, so the rest of the cycle is a white path from $v$ to its predecessor $u$ on $C$. By the white-path theorem $u$ becomes a descendant of $v$, and when DFS examines the edge $u \to v$, $v$ is still gray (it is left only after its descendants): a back edge. \hfill$\blacksquare$

The undirected case is the same argument with the extra care of Lesson~\ref{ssec:ch2-dfs-tree}: the edge to the parent is skipped, so the non-tree edges are exactly the back edges.

\begin{diagrambox}
\begin{tikzpicture}[x=0.95cm, y=0.95cm]
    \begin{scope}
        \tikzset{g2-v1/.style=gvsel, g2-v2/.style=gvsel, g2-v3/.style=gvcur,
                 g2-e12/.style={gtree, -{Stealth[length=2.4mm]}}, g2-e23/.style={gtree, -{Stealth[length=2.4mm]}},
                 g2-e31/.style=gback}
        \pic {g2-D};
        \node[font=\scriptsize, text width=6cm, align=left, anchor=north west] at (-0.4,-2.0)
            {The first back edge DFS meets on $D$ is $3 \to 1$ (at time 2, stack \texttt{1 2 3}). Walk the parents from 3 up to 1: $3, 2, 1$; reverse it and close the loop: \textbf{cycle $1 \to 2 \to 3 \to 1$}.};
    \end{scope}
    \begin{scope}[xshift=7.6cm]
        \tikzset{g2-u1/.style=gvsel, g2-u2/.style=gvsel, g2-u5/.style=gvsel, g2-u4/.style=gvcur,
                 g2-f12/.style=gtree, g2-f25/.style=gtree, g2-f45/.style=gtree, g2-f14/.style={gback, -}}
        \pic {g2-G3};
        \node[font=\scriptsize, text width=6cm, align=left, anchor=north west] at (-0.4,-2.0)
            {On $G_3$ the first non-parent edge to a visited vertex is $4$--$1$ (stack \texttt{1 2 5 4}). Parents from 4 up to 1: $4, 5, 2, 1$: \textbf{cycle $1, 2, 5, 4, 1$}.};
    \end{scope}
\end{tikzpicture}
\figcap{Recovering a cycle. Orange $=$ the vertex where the back edge (red, dashed) is found; purple $=$ the rest of the cycle, read off the parent pointers.}
\end{diagrambox}

\subsubsection*{Recovering the cycle}

When DFS sits at $u$ and finds the back edge $u \to v$, the cycle is on the call stack: it is the part of the stack from $v$ to $u$. Two ways to read it:
\begin{itemize}[leftmargin=*]
    \item \textbf{parent pointers:} $x = u$; while $x \ne v$: output $x$, $x = \texttt{parent}[x]$; output $v$. This lists the cycle backwards ($u, \dots, v$); reverse it for a directed cycle.
    \item \textbf{explicit path:} keep the stack as a vector \texttt{path}; the cycle is \texttt{path} from the position of $v$ to the end. (Problem~\ref{ssec:ch2-round-trip-ii} uses parents, Exercise E1.1 uses the path.)
\end{itemize}
Both cost \bigO{\text{cycle length}} once, after an \VE{} search.

\begin{lstlisting}[style=cpp]
#include <bits/stdc++.h>
using namespace std;

// Find ONE directed cycle with colours 0/1/2 and print it, or report that the graph is acyclic.
vector<vector<int>> adj;
vector<int> color, par;
int cycStart = -1, cycEnd = -1;                     // the back edge cycEnd -> cycStart

bool dfs(int u) {
    color[u] = 1;
    for (int v : adj[u]) {
        if (color[v] == 0) {
            par[v] = u;
            if (dfs(v)) return true;                // stop as soon as a cycle is known
        } else if (color[v] == 1) {                 // gray: a back edge u -> v
            cycStart = v; cycEnd = u;
            return true;
        }                                           // black: forward or cross edge, harmless
    }
    color[u] = 2;
    return false;
}

int main() {
    int n, m;
    cin >> n >> m;
    adj.assign(n + 1, vector<int>());
    for (int i = 0; i < m; i++) {
        int a, b;
        cin >> a >> b;
        adj[a].push_back(b);
    }
    color.assign(n + 1, 0); par.assign(n + 1, 0);
    for (int s = 1; s <= n && cycStart == -1; s++)
        if (color[s] == 0) dfs(s);
    if (cycStart == -1) { cout << "acyclic\n"; return 0; }
    vector<int> cyc;
    for (int x = cycEnd; x != cycStart; x = par[x]) cyc.push_back(x);
    cyc.push_back(cycStart);
    reverse(cyc.begin(), cyc.end());                // cycStart -> ... -> cycEnd
    cyc.push_back(cycStart);                        // close the loop
    for (int x : cyc) cout << x << " ";
    cout << "\n";
}
\end{lstlisting}

\begin{lstlisting}[style=trace]
input : D                                   output: 1 2 3 1
input : 3 3 / 1 2 / 1 3 / 2 3   (a DAG)     output: acyclic
input : 2 1 / 2 2               (self-loop) output: 2 2
\end{lstlisting}

\subsubsection*{Three traps}

\begin{diagrambox}
\begin{tikzpicture}[x=1cm, y=1cm]
    % trap 1: visited is not enough
    \begin{scope}
        \node[font=\small\bfseries] at (1,1.55) {(1) a black target};
        \node[gvdone] (a1) at (0,0.6) {1};
        \node[gvdone] (a2) at (2,0.6) {2};
        \node[gvcur] (a3) at (1,-0.6) {3};
        \draw[gtree, -{Stealth[length=2.4mm]}] (a1) -- (a2);
        \draw[gtree, -{Stealth[length=2.4mm]}] (a2) -- (a3);
        \draw[gfwd] (a1) -- (a3);
        \node[font=\scriptsize, text width=3.6cm, align=center] at (1,-1.6) {DAG $1{\to}2{\to}3$, $1{\to}3$: when 1 examines $1 \to 3$, vertex 3 is \textbf{black}. ``visited $\Rightarrow$ cycle'' would be wrong.};
    \end{scope}
    % trap 2: parallel edges
    \begin{scope}[xshift=5.2cm]
        \node[font=\small\bfseries] at (1,1.55) {(2) parallel edges};
        \node[gvseen] (b1) at (0,0) {1};
        \node[gvcur] (b2) at (2,0) {2};
        \draw[gtree] (b1) to[bend left=25] node[gw, above] {e0} (b2);
        \draw[gback, -] (b2) to[bend left=25] node[gw, below] {e1} (b1);
        \node[font=\scriptsize, text width=3.6cm, align=center] at (1,-1.6) {Undirected, two roads between 1 and 2: a cycle $1, 2, 1$ of length 2. ``Skip the parent \emph{vertex}'' misses it; skip the parent \emph{edge} (e0).};
    \end{scope}
    % trap 3: self loops
    \begin{scope}[xshift=10.4cm]
        \node[font=\small\bfseries] at (1,1.55) {(3) self-loops};
        \node[gvcur] (c1) at (1,0) {2};
        \draw[gback] (c1) to[out=60, in=120, looseness=7] (c1);
        \node[font=\scriptsize, text width=3.6cm, align=center] at (1,-1.6) {$2 \to 2$ is a back edge to itself: a cycle of length 1. Fine if the problem allows it --- many statements (Round Trip) forbid short cycles by guaranteeing simple graphs.};
    \end{scope}
\end{tikzpicture}
\figcap{The three classic situations a cycle detector must handle correctly.}
\end{diagrambox}

\begin{warningbox}[Directed graphs: gray, not visited]
The test must be ``$v$ is gray'' (on the current path). A black vertex was finished by an earlier part of the search; reaching it again only means two paths lead there (trap 1). With the test changed to ``$v$ is not white'', the program above declares a cycle at the edge $1 \to 3$ of the 3-vertex DAG in the trace --- and then the parent walk from 1 never meets 3, because 3 is not an ancestor of 1: our run looped until it died with \texttt{bad\_alloc}. A wrong cycle test often shows up as a crash in the \emph{recovery} code, not as a wrong answer.
\end{warningbox}

\begin{warningbox}[Undirected graphs: skip the edge, not the vertex]
In an undirected multigraph, two parallel edges $e_0, e_1$ between $p$ and $v$ form a cycle. When $v$ was entered through $e_0$, skip only $e_0$; then $e_1$ is a non-tree edge to the visited $p$ and the cycle $p, v, p$ is found. If the input guarantees a simple graph (no parallel edges), skipping the parent vertex is equivalent and slightly shorter.
\end{warningbox}

\subsubsection*{Undirected graphs: counting instead of searching}

For the \emph{yes/no} question in an undirected graph, a search is not even needed: a forest with $n$ vertices and $c$ components has exactly $n - c$ edges, so
\[
    \text{the graph has a cycle} \iff m > n - c.
\]
Counting the components $c$ still needs a traversal (or a DSU, Chapter~6), but this formula is handy for checks like ``is the graph a tree?'' ($m = n - 1$ and connected) or ``does every component have exactly one cycle?'' ($m = n$ per component, Exercise E1.2). To \emph{print} a cycle you need the DFS.

\subsubsection*{Iterative versions}

On $10^5$ vertices a long path can make the recursion deep (Lesson~\ref{ssec:ch2-call-stack}). Both detectors run unchanged on the iterator-based iterative DFS: colour 1 when a vertex is pushed, colour 2 when it is popped, and check the colour of every examined neighbour. The library file \texttt{code/lib/cycle.hpp} (Appendix~A) contains both iterative versions, the undirected one with edge ids; Problems~\ref{ssec:ch2-round-trip} and~\ref{ssec:ch2-round-trip-ii} show them on real inputs.

\subsubsection*{Check your understanding}

\begin{enumerate}[leftmargin=*]
    \item In a directed graph DFS finds the back edge $7 \to 6$ of $D$ (if it had not stopped at $3 \to 1$). Which cycle does the parent walk print?
    \item True or false: if a directed graph has a cycle, DFS from vertex 1 alone finds it.
    \item An undirected graph has $n = 10$ vertices, $m = 9$ edges and 2 components. Does it have a cycle?
    \item Why is it safe to stop the whole search at the first back edge, even though the colours of the unfinished vertices are then wrong?
    \item A directed graph has $10^5$ vertices and $2\cdot10^5$ edges and you must print a cycle. Which of the three traps can appear, and what is the time complexity?
\end{enumerate}

\begin{solutionbox}[Answers]
\begin{enumerate}[leftmargin=*]
    \item $\texttt{par}[7] = 6$, so the walk is $7, 6$; reversed and closed: $6 \to 7 \to 6$.
    \item False: the cycle may not be reachable from 1 (in $D$, a cycle through 8 would be invisible from 1). That is why the outer loop starts a DFS from every still-white vertex.
    \item A forest with 10 vertices and 2 components has $10 - 2 = 8$ edges; $9 > 8$, so \textbf{yes}.
    \item We need one cycle, and it is fully known at that moment (the back edge plus the parent pointers, which are final for all gray vertices). The colours are never used again.
    \item Trap 1 (black targets) and trap 3 (self-loops, if allowed). Trap 2 does not exist in directed graphs: $u \to v$ and $v \to u$ are two different edges and genuinely form a cycle of length 2. Time \bigO{n + m}.
\end{enumerate}
\end{solutionbox}

%============================================================
% CH2 · LESSON 5 — DFS AS A TOOLBOX
%============================================================
\subsection{DFS as a Toolbox}
\label{ssec:ch2-toolbox}

The four lessons so far looked at DFS from the inside. This lesson turns it around: here are the five things a contest solution most often asks DFS to do, each one a small change to the same skeleton. They are previews of the five patterns of Part~B.

\begin{diagrambox}
\begin{tikzpicture}[
    tool/.style={rectangle, rounded corners=3pt, draw=patternborder, fill=patternbg, line width=0.9pt,
                 font=\scriptsize, align=center, text width=2.45cm, minimum height=2.25cm},
    core/.style={rectangle, rounded corners=3pt, draw=deepblue, fill=theorybg, line width=1.1pt,
                 font=\scriptsize\ttfamily, align=left, text width=4.6cm}]
    \node[core] (c) at (0,0) {dfs(u):\\
\ \ \textcolor{successgreen}{\textbf{[enter u]}}\\
\ \ for v in adj[u]:\\
\ \ \ \ if v is white: dfs(v)\\
\ \ \ \ \textcolor{recursionred}{\textbf{[look at v's colour]}}\\
\ \ \textcolor{accentorange}{\textbf{[leave u]}}\\
\ \ \textcolor{softpurple}{\textbf{[un-mark u?]}}};
    \node[tool] (t1) at (-5.6,1.4) {\textbf{1. Cycles}\\(Pattern 1)\\[2pt]edge to a gray $v$ $\Rightarrow$ cycle; read it off the parents};
    \node[tool] (t2) at (-5.6,-1.4) {\textbf{2. Subtrees}\\(Pattern 2)\\[2pt]combine children at \emph{leave}; pass values down at \emph{enter}};
    \node[tool] (t3) at (5.6,1.4) {\textbf{3. Timestamps}\\(Pattern 3)\\[2pt]$\tin{}$ at enter, $\tout{}$ at leave; subtrees are intervals};
    \node[tool] (t4) at (5.6,-1.4) {\textbf{4. Backtracking}\\(Pattern 4)\\[2pt]un-mark at leave $\Rightarrow$ every simple path};
    \node[tool] (t5) at (0,-2.9) {\textbf{5. Reverse search}\\(Pattern 5)\\[2pt]run it on reversed edges, from the target};
    \draw[flow] (t1.east) -- (c.west |- t1.east);
    \draw[flow] (t2.east) -- (c.west |- t2.east);
    \draw[flow] (t3.west) -- (c.east |- t3.west);
    \draw[flow] (t4.west) -- (c.east |- t4.west);
    \draw[flow] (t5.north) -- (c.south);
\end{tikzpicture}
\figcap{One skeleton, five tools. Each tool fills in one of the coloured slots.}
\end{diagrambox}

%------------------------------------------------------------
\subsubsection*{Tool 1 and 2: information up and down a rooted tree}

Most ``subtree'' problems give a \textbf{tree} as $n-1$ undirected edges and fix a root. DFS from the root \emph{is} the rooted tree: every edge is a tree edge (there are no back edges in a tree), and the only neighbour that is not a child is the parent. So instead of a \texttt{visited} array one usually passes the parent:

\begin{lstlisting}[style=cpp]
#include <bits/stdc++.h>
using namespace std;

// A rooted tree from undirected edges: depth goes DOWN (enter), size goes UP (leave).
vector<vector<int>> adj;
vector<int> depth_, sz;

void dfs(int u, int p) {                 // p = parent of u (0 for the root)
    depth_[u] = (p == 0) ? 0 : depth_[p] + 1;   // pre-order: the parent is already done
    sz[u] = 1;
    for (int v : adj[u])
        if (v != p) {                    // in a tree, "not the parent" == "not visited"
            dfs(v, u);
            sz[u] += sz[v];              // post-order: the child is finished
        }
}

int main() {
    int n;
    cin >> n;
    adj.assign(n + 1, vector<int>());
    for (int i = 0; i < n - 1; i++) {
        int a, b;
        cin >> a >> b;
        adj[a].push_back(b);
        adj[b].push_back(a);
    }
    depth_.assign(n + 1, 0); sz.assign(n + 1, 0);
    dfs(1, 0);
    for (int v = 1; v <= n; v++) cout << v << ": depth " << depth_[v] << ", size " << sz[v] << "\n";
}
\end{lstlisting}

\begin{diagrambox}
\begin{tikzpicture}[x=1.1cm, y=1cm]
    \node[gvsrc] (r1) at (0,0) {1};
    \node[gv] (r2) at (-1.8,-1) {2}; \node[gv] (r3) at (0,-1) {3}; \node[gv] (r4) at (1.8,-1) {4};
    \node[gv] (r5) at (-2.5,-2) {5}; \node[gv] (r6) at (-1.1,-2) {6}; \node[gv] (r7) at (1.8,-2) {7};
    \node[gv] (r8) at (1.8,-3) {8};
    \foreach \a/\b in {1/2, 1/3, 1/4, 2/5, 2/6, 4/7, 7/8} \draw[gtree] (r\a) -- (r\b);
    \foreach \v/\d/\s in {1/0/8, 2/1/3, 3/1/1, 4/1/3, 5/2/1, 6/2/1, 7/2/2, 8/3/1}{
        \node[gdist, text=accentorange, above left=-1pt and -1pt of r\v] {\tiny d\d};
        \node[gdist, text=primaryblue, above right=-1pt and -1pt of r\v] {\tiny s\s};
    }
    \draw[->, thick, accentorange] (3.4,0) -- (3.4,-3) node[midway, right, font=\scriptsize, align=left] {depth flows\\down (enter)};
    \draw[->, thick, primaryblue] (5.3,-3) -- (5.3,0) node[midway, right, font=\scriptsize, align=left] {size flows\\up (leave)};
\end{tikzpicture}
\figcap{The tree $T$ with edges $1{-}2, 1{-}3, 1{-}4, 2{-}5, 2{-}6, 4{-}7, 7{-}8$, rooted at 1: orange $d$ = depth, blue $s$ = subtree size, both printed by the program.}
\end{diagrambox}

\begin{lstlisting}[style=trace]
input : 8 / 1 2 / 1 3 / 1 4 / 2 5 / 2 6 / 4 7 / 7 8
output: 1: depth 0, size 8     2: depth 1, size 3     3: depth 1, size 1     4: depth 1, size 3
        5: depth 2, size 1     6: depth 2, size 1     7: depth 2, size 2     8: depth 3, size 1
\end{lstlisting}

Every Pattern~2 problem is this program with a different quantity in place of the size (a count of colours, a sum, a product of sizes), or in place of the depth (a sum of values along the root path).

%------------------------------------------------------------
\subsubsection*{Tool 3: path reconstruction}

DFS leaves behind \texttt{parent[]}. Walking it from $t$ back to $s$ gives \emph{a} path from $s$ to $t$:
\begin{lstlisting}[style=snippet]
vector<int> path;
for (int x = t; x != s; x = parent[x]) path.push_back(x);
path.push_back(s);
reverse(path.begin(), path.end());     // s ... t
\end{lstlisting}
On $G_3$ from $s = 1$ the parents are $\texttt{parent}[2] = 1$, $\texttt{parent}[3] = 2$, $\texttt{parent}[5] = 2$, $\texttt{parent}[4] = 5$, $\texttt{parent}[6] = 5$, and the path to $4$ is $1, 2, 5, 4$ --- three edges, although the edge $1$--$4$ exists. DFS paths are paths, not shortest paths; for shortest paths use BFS (Chapter~3).

%------------------------------------------------------------
\subsubsection*{Tool 4: backtracking --- un-mark on the way out}

DFS marks a vertex once and never visits it again; that is what makes it linear. Remove the mark when the vertex is \emph{left} and the search changes completely: it now explores every \textbf{simple path} from the start --- every way of walking without repeating a vertex. The marks then mean ``on the current path'', exactly the gray vertices.

\begin{lstlisting}[style=cpp]
#include <bits/stdc++.h>
using namespace std;

// Backtracking: print EVERY simple path from s to t. onPath[v] = v is on the current path.
vector<vector<int>> adj;
vector<bool> onPath;
vector<int> path;
int t, found = 0;

void go(int u) {
    onPath[u] = true; path.push_back(u);
    if (u == t) {
        found++;
        for (int x : path) cout << x << " ";
        cout << "\n";
    } else {
        for (int v : adj[u])
            if (!onPath[v]) go(v);
    }
    onPath[u] = false; path.pop_back();      // un-mark: other paths may use u again
}

int main() {
    int n, m, s;
    cin >> n >> m;
    adj.assign(n + 1, vector<int>());
    for (int i = 0; i < m; i++) {
        int a, b;
        cin >> a >> b;
        adj[a].push_back(b);
        adj[b].push_back(a);
    }
    cin >> s >> t;
    onPath.assign(n + 1, false);
    go(s);
    cout << found << " simple paths\n";
}
\end{lstlisting}

\begin{lstlisting}[style=trace]
input : G3, then s = 1, t = 6
output: 1 2 5 6
        1 4 5 6
        2 simple paths
\end{lstlisting}

The price: the number of simple paths can be exponential (a $6 \times 6$ grid graph already has $1\,262\,816$ simple paths between opposite corners --- the backtracking program counted them), so backtracking is only for tiny graphs or heavily pruned searches. Pattern~4 is about recognising when it is affordable.

%------------------------------------------------------------
\subsubsection*{Tool 5: reverse search}

``From which vertices can $t$ be reached?'' would need one search per start vertex. Reverse every edge instead: $u$ can reach $t$ in $G$ iff $t$ can reach $u$ in the reversed graph $G^R$. One search from $t$ in $G^R$ answers the question for all $u$ at once.
\begin{lstlisting}[style=snippet]
for (auto& e : edges) radj[e.second].push_back(e.first);   // edge a -> b becomes b -> a
dfs_on(radj, t);                                            // marked = everything that can reach t
\end{lstlisting}
On $D$ the vertices that can reach 4 are $\{1, 2, 3, 4, 5\}$: 3 gets there via $3 \to 1 \to 2 \to 4$, while 6, 7 and 8 cannot (the only edges out of $\{6, 7\}$ stay inside it). Pattern~5 also reverses \emph{roles}: start from the exit of a maze, from the border of a grid, from the sinks of a DAG.

%------------------------------------------------------------
\subsubsection*{Which tool? A first decision table}

\begin{center}
\small
\renewcommand{\arraystretch}{1.3}
\begin{tabularx}{\linewidth}{@{}>{\raggedright\arraybackslash}X l l@{}}
\toprule
\textbf{The question is about \dots} & \textbf{tool} & \textbf{cost} \\
\midrule
``is there a cycle / a deadlock / a contradiction in the order?'' & 1: colours, back edges & \VE \\
``for every vertex, something about the vertices below it'' & 2: post-order aggregation & \VE \\
``for every vertex, something about the path from the root'' & 2: pre-order push-down & \VE \\
``is $u$ an ancestor of $v$?'', ``the $k$-th vertex of $u$'s subtree in DFS order'' & 3: timestamps & \VE{} $+\ \bigO{1}$ per query \\
``all paths'', ``the best path'' on $\le 15$--$25$ vertices or cells & 4: backtracking & exponential \\
``which vertices can reach the target / the border / the exit?'' & 5: reverse search & \VE \\
\bottomrule
\end{tabularx}
\end{center}

\subsubsection*{Check your understanding}

\begin{enumerate}[leftmargin=*]
    \item In the tree $T$, re-root at 4. What are the new depth and subtree size of vertex 1?
    \item Why does the tree DFS need no \texttt{visited} array, while a DFS on $G_3$ does?
    \item The backtracking program is run on $G_3$ with $s = 3$, $t = 6$. How many simple paths does it print?
    \item Which vertices of $D$ can reach vertex 7?
    \item Without backtracking, can the plain DFS (marks never removed) print \emph{all} simple paths from $s$ to $t$? Give an example on $G_3$.
\end{enumerate}

\begin{solutionbox}[Answers]
\begin{enumerate}[leftmargin=*]
    \item Rooted at 4, vertex 1 is a child of 4: depth 1. Its subtree is everything except the branch $\{4, 7, 8\}$: size $8 - 3 = 5$. (The ``$n - \mathrm{sz}$'' trick returns in Problem~\ref{ssec:ch2-highest-score}.)
    \item In a tree the only already-visited neighbour of $u$ is its parent; skipping the parent is enough. $G_3$ has a cycle $1, 2, 5, 4$, so a vertex can be reached a second time along a non-tree edge.
    \item Every path from 3 must start $3, 2$; from 2 it continues $2, 5, 6$ or $2, 1, 4, 5, 6$: \textbf{2} paths (the program prints \texttt{3 2 1 4 5 6}, then \texttt{3 2 5 6}, because \texttt{adj[2]} lists 1 before 5).
    \item Search from 7 in the reversed graph: $7 \leftarrow 6$, $7 \leftarrow 8$; $6 \leftarrow 5$, $6 \leftarrow 1$; $1 \leftarrow 3$; $3 \leftarrow 2$, $3 \leftarrow 4$. So \textbf{all 8 vertices} can reach 7 --- even 4, along $4 \to 3 \to 1 \to 6 \to 7$.
    \item No: once 5 is marked by the path $1, 2, 5, 6$, the path $1, 4, 5, 6$ can never use 5 again. Plain DFS finds one path per reachable vertex.
\end{enumerate}
\end{solutionbox}


%============================================================
\clearpage
\section{Pattern 1: Cycle Detection and Recovery}
\label{sec:ch2-p1}
%============================================================
%============================================================
% CH2 · PATTERN 1 — CYCLE DETECTION AND RECOVERY: THEORY
%============================================================
\subsection{Pattern Theory}
\label{ssec:ch2-p1-theory}

\begin{patternbox}{Pattern 1 at a Glance}
Run DFS from every white vertex. Keep the vertices of the current path \textbf{gray}. The first time DFS, sitting at $u$, examines an edge to a vertex $v$ that is \textbf{still on the path} (gray; in an undirected graph: and not through the edge it came in by), the path from $v$ down to $u$ plus that edge is a cycle. If the search ends without such an edge, there is no cycle.

\begin{center}
\begin{tikzpicture}[x=1.25cm]
    \foreach \v/\i/\s in {s/0/gvsrc, a/1/gvseen, v/2/gvsel, b/3/gvsel, c/4/gvsel, u/5/gvcur}
        \node[\s] (p\i) at (\i,0) {$\v$};
    \foreach \i [evaluate=\i as \j using int(\i+1)] in {0,...,4} \draw[gtree, -{Stealth[length=2.4mm]}] (p\i) -- (p\j);
    \draw[gback] (p5.north) to[bend right=40] node[lbl, above]{back edge $u \to v$} (p2.north);
    \node[gvdone] (x) at (1.5,-1.3) {$x$};
    \draw[gtree, -{Stealth[length=2.4mm]}] (p1) -- (x);
    \draw[gcross] (p4) to[bend left=12] node[lbl, below, pos=0.6]{to black: harmless} (x);
    \draw[decorate, decoration={brace, amplitude=5pt}, darkgray] (-0.35,0.65) -- (5.35,0.65)
        node[midway, above=5pt, font=\scriptsize]{the stack (gray vertices) $=$ the current path};
    \node[note, text width=3.6cm, align=left, anchor=west] at (6.0,0) {cycle $= v \to b \to c \to u \to v$\\(purple + orange), read off \texttt{parent[]} from $u$ up to $v$};
\end{tikzpicture}
\end{center}
\begin{recurrencebox}
\centering\small
colour $0 \to 1$ when entered, $1 \to 2$ when left;\quad examine $u \to v$:\\[3pt]
$\text{colour}[v] = 0$: tree edge, \texttt{parent}$[v] = u$ \qquad $\text{colour}[v] = 1$: \textbf{cycle} $v \leadsto u \to v$ \qquad $\text{colour}[v] = 2$: ignore\\[3pt]
undirected: skip the edge id \texttt{parentEdge}$[u]$ \qquad total \VE, cycle printed in \bigO{\text{its length}}
\end{recurrencebox}
\end{patternbox}

%------------------------------------------------------------
\subsubsection*{From the lessons to the pattern}

Lesson~\ref{ssec:ch2-cycles} proved that back edges and cycles are the same thing: a graph has a cycle iff a DFS meets a back edge, and the cycle is sitting on the call stack at that very moment. So the pattern is not ``find a clever algorithm'', it is ``run the right DFS and read the answer off the stack''. The difficulty in contest problems lies elsewhere:
\begin{itemize}[leftmargin=*]
    \item recognising that the question \emph{is} about a cycle --- ``a round trip'', ``a deadlock'', ``contradictory orders'', ``an infinite loop'', ``a process that never ends'';
    \item getting the details right for the kind of graph: directed or undirected, simple or with parallel edges and self-loops, connected or not;
    \item printing the cycle in the format the judge wants (often with the first vertex repeated at the end);
    \item and not overflowing the stack on $10^5$ vertices.
\end{itemize}

%------------------------------------------------------------
\subsubsection*{When to recognise it (and when not)}

\begin{itemize}[leftmargin=*]
    \item ``\textbf{find} a cycle / a round trip / a route that returns to its start'' --- print any one;
    \item ``is the order consistent?'', ``can all courses be finished?'', ``do the screenshots contradict each other?'' --- a directed cycle means \emph{no};
    \item ``is the graph a tree / a forest / a cactus / a single cycle?'' --- cycle tests per component (often by counting, $m$ vs.\ $n - c$);
    \item ``does the simulation run forever?'' when the next step is \emph{not} forced (if it is forced, it is a functional graph: Chapter~1, Pattern~5).
\end{itemize}
\textbf{Not this pattern:} the \emph{shortest} cycle needs BFS from every vertex (Chapter~3, and Chapter~22 for faster methods); an \emph{odd} cycle is a bipartiteness test (Chapter~4); a cycle of negative weight is Bellman--Ford (Chapter~8); a cycle through \emph{every} vertex or \emph{every} edge is Hamiltonian (Chapter~23) or Eulerian (Chapter~21); counting all cycles is exponential in general.

%------------------------------------------------------------
\subsubsection*{The algorithm and its invariant}

\begin{invariantbox}[Why reading the stack gives a cycle]
\textbf{Invariant.} At every moment, the gray vertices are exactly the vertices on the stack, in order $x_0, x_1, \dots, x_k$ with $\texttt{parent}[x_{i+1}] = x_i$, and every edge $x_i \to x_{i+1}$ is a real edge of the graph.

\emph{Detection is sound.} If the edge $u \to v$ with $u = x_k$ finds $v$ gray, then $v = x_j$ for some $j \le k$, and $x_j \to x_{j+1} \to \dots \to x_k \to x_j$ is a closed walk without repeated vertices --- a cycle. In an undirected graph, $j = k-1$ is impossible unless the edge is a \emph{second} edge between $x_{k-1}$ and $x_k$, because the tree edge itself is skipped; so in a simple graph the cycle has at least 3 vertices.

\emph{Detection is complete.} By the cycle theorem of Lesson~\ref{ssec:ch2-cycles}, if any cycle exists, the DFS meets a back edge before it finishes.

\emph{Cost.} Every vertex is pushed once and every adjacency entry examined once: \VE. The parent walk visits each cycle vertex once.
\end{invariantbox}

\begin{diagrambox}
\begin{tikzpicture}[x=1cm, y=1cm]
    % directed example
    \begin{scope}
        \node[font=\small\bfseries] at (1.6,2.2) {directed: gray or black?};
        \node[gvseen] (a) at (0,0.6) {1};
        \node[gvseen] (b) at (1.6,0.6) {2};
        \node[gvcur] (c) at (3.2,0.6) {3};
        \node[gvdone] (d) at (1.6,-0.8) {4};
        \draw[gtree, -{Stealth[length=2.4mm]}] (a) -- (b);
        \draw[gtree, -{Stealth[length=2.4mm]}] (b) -- (c);
        \draw[gtree, -{Stealth[length=2.4mm]}] (a) -- (d);
        \draw[gcross] (c) -- node[lbl, right, pos=0.4]{black} (d);
        \draw[gback] (c) to[bend right=50] node[lbl, above]{gray} (a);
        \node[font=\scriptsize, text width=4.4cm, align=center] at (1.6,-1.85) {$3 \to 4$ meets a finished vertex: no cycle through it. $3 \to 1$ meets the stack: cycle $1, 2, 3$.};
    \end{scope}
    % undirected example
    \begin{scope}[xshift=6.4cm]
        \node[font=\small\bfseries] at (1.6,2.2) {undirected: which edge to skip?};
        \node[gvseen] (p) at (0,0.6) {1};
        \node[gvseen] (q) at (1.6,0.6) {2};
        \node[gvcur] (r) at (3.2,0.6) {3};
        \draw[gtree] (p) -- node[gw, above]{e0} (q);
        \draw[gtree] (q) -- node[gw, above]{e1} (r);
        \draw[gfaded, dashed] (r) to[bend left=35] node[gw, below]{e1 again} (q);
        \draw[gback, -] (r) to[bend right=55] node[lbl, above]{e2} (p);
        \node[font=\scriptsize, text width=4.6cm, align=center] at (1.6,-1.85) {At 3, the edge e1 (grey) leads back to the parent: skipped. e2 leads to the gray 1: cycle $1, 2, 3$.};
    \end{scope}
\end{tikzpicture}
\figcap{The two decisions of the pattern. Directed: only gray targets close cycles. Undirected: skip exactly the edge used to enter the current vertex.}
\end{diagrambox}

%------------------------------------------------------------
\subsubsection*{Printing the cycle}

Judges usually want the cycle as a list of vertices in walking order, often with the start repeated at the end (``$k$, then $k$ cities, first $=$ last''). With parent pointers:

\begin{diagrambox}
\begin{tikzpicture}[x=1.15cm]
    \node[font=\scriptsize\bfseries, anchor=east] at (-0.4,0) {stack};
    \foreach \v/\i/\s in {1/0/gvseen, 4/1/gvseen, 7/2/gvsel, 2/3/gvsel, 9/4/gvsel, 5/5/gvcur}
        \node[\s] (s\i) at (\i,0) {\v};
    \foreach \i [evaluate=\i as \j using int(\i+1)] in {0,...,4} \draw[gtree, -{Stealth[length=2.4mm]}] (s\i) -- (s\j);
    \draw[gback] (s5.north) to[bend right=40] (s2.north);
    \node[font=\scriptsize, anchor=west] at (6.4,0.35) {back edge $5 \to 7$};
    \node[font=\scriptsize\ttfamily, align=left, anchor=west] at (0,-1.1) {walk parents from 5 to 7:\ \ 5 9 2 7};
    \node[font=\scriptsize\ttfamily, align=left, anchor=west] at (0,-1.6) {reverse:\ \ \ \ \ \ \ \ \ \ \ \ \ \ \ \ \ \ \ 7 2 9 5};
    \node[font=\scriptsize\ttfamily, align=left, anchor=west] at (0,-2.1) {append the start again:\ \ \ \ 7 2 9 5 7\ \ \ (k = 5)};
    \node[note, text width=4.6cm, anchor=west, align=left] at (6.4,-1.6) {Undirected graphs need no reversal: the cycle is the same read in either direction. Directed graphs do --- \texttt{5 9 2 7} walks the edges backwards.};
\end{tikzpicture}
\end{diagrambox}

%------------------------------------------------------------
\subsubsection*{Variations map}

\begin{diagrambox}
\begin{tikzpicture}[
    var/.style={rectangle, rounded corners=2pt, draw=patternborder, fill=patternbg,
                font=\scriptsize, align=center, text width=3.25cm, minimum height=1.3cm, inner sep=3pt}]
    \node[var, fill=answerfill, draw=accentorange, text width=5cm, font=\scriptsize\bfseries] (root) at (0,0)
        {DFS with colours; a back edge closes a cycle};
    \node[var] (v1) at (-5.4,-2) {\textbf{undirected, simple}\\skip the parent, $\ge 3$ vertices\\(Round Trip, 2.1.1)};
    \node[var] (v2) at (-1.8,-2) {\textbf{directed}\\gray, not visited\\(Round Trip II, 2.1.2)};
    \node[var] (v3) at (1.8,-2)  {\textbf{one cycle per component?}\\count: $m = n$, connected\\(E1.2)};
    \node[var] (v4) at (5.4,-2)  {\textbf{grid cycles}\\same colour, length $\ge 4$\\(E1.3)};
    \node[var] (v5) at (-5.4,-3.6) {\textbf{which components are acyclic?}\\trees among components\\(E1.4)};
    \node[var] (v6) at (-1.8,-3.6) {\textbf{consistency of orders}\\constraints $\to$ arrows, cycle $=$ no\\(E1.5)};
    \node[var] (v7) at (1.8,-3.6)  {\textbf{which edge to delete?}\\cycle + in-degree 2\\(E1.6)};
    \node[var] (v8) at (5.4,-3.6)  {\textbf{find the cycle, then use it}\\pseudotrees, entry vertex\\(Practice set)};
    \node[draw=patternborder, dashed, rounded corners=4pt, fit=(v1)(v4)(v5)(v8), inner sep=5pt] (grp) {};
    \draw[flow] (root.south) -- (grp.north);
    \node[note, below=2mm of grp] {shortest / odd / negative / Hamiltonian / Euler cycles are different problems: Chapters 3, 4, 8, 23, 21};
\end{tikzpicture}
\end{diagrambox}

\textbf{Twists and how to handle them}
\begin{itemize}[leftmargin=*]
    \item \textbf{Parallel edges in an undirected graph} are a cycle of length 2; skip the parent \emph{edge id}, not the parent vertex (template 1).
    \item \textbf{Self-loops} are cycles of length 1 in both kinds of graph. If the problem forbids them as answers, check the statement: many inputs simply exclude them.
    \item \textbf{Disconnected input:} the cycle may live in any component --- the outer loop over all start vertices is not optional.
    \item \textbf{``Cycle through vertex $x$'':} a directed cycle through $x$ exists iff some successor of $x$ reaches $x$ (one search in the reversed graph from $x$, Pattern~5); an undirected one iff $x$ has two neighbours that are connected without $x$.
    \item \textbf{Only yes/no, undirected:} count, no search order needed: a cycle exists iff $m > n - c$ (template 3, or a DSU).
    \item \textbf{Only yes/no, directed:} Kahn's peeling of in-degree-0 vertices (Chapter~5) leaves vertices behind iff there is a cycle; the leftovers also let you \emph{print} a cycle (Approach~3 of Problem~\ref{ssec:ch2-round-trip-ii}).
\end{itemize}

%------------------------------------------------------------
\subsubsection*{Code templates}

\begin{lstlisting}[style=cpp]
#include <bits/stdc++.h>
using namespace std;

// (1) UNDIRECTED, multi-edges and self-loops allowed: skip the EDGE we came in by.
//     Returns a cycle v0 v1 ... vk-1 (closing edge vk-1 - v0), or {} for a forest. Iterative.
vector<int> undirectedCycle(int n, const vector<pair<int, int>>& edges) {
    vector<vector<pair<int, int>>> adj(n);             // (neighbour, edge id)
    for (int i = 0; i < (int)edges.size(); i++) {
        adj[edges[i].first].push_back({edges[i].second, i});
        if (edges[i].first != edges[i].second) adj[edges[i].second].push_back({edges[i].first, i});
    }
    vector<int> color(n, 0), par(n, -1), parEdge(n, -1);
    vector<size_t> it(n, 0);
    for (int s = 0; s < n; s++) {
        if (color[s]) continue;
        color[s] = 1;
        vector<int> st(1, s);
        while (!st.empty()) {
            int u = st.back();
            if (it[u] == adj[u].size()) { color[u] = 2; st.pop_back(); continue; }
            int v = adj[u][it[u]].first, id = adj[u][it[u]].second;
            it[u]++;
            if (id == parEdge[u]) continue;            // the edge we came in by
            if (color[v] == 0) { color[v] = 1; par[v] = u; parEdge[v] = id; st.push_back(v); }
            else if (color[v] == 1) {                  // back edge to an ancestor (or a self-loop)
                vector<int> cyc;
                for (int x = u; x != v; x = par[x]) cyc.push_back(x);
                cyc.push_back(v);
                reverse(cyc.begin(), cyc.end());
                return cyc;                            // v ... u, closed by the edge u - v
            }
        }
    }
    return {};
}

// (2) DIRECTED: an edge to a GRAY vertex closes a cycle. Returns v0 -> ... -> vk-1 (-> v0) or {}.
vector<int> directedCycle(const vector<vector<int>>& adj) {
    int n = adj.size();
    vector<int> color(n, 0), par(n, -1);
    vector<size_t> it(n, 0);
    for (int s = 0; s < n; s++) {
        if (color[s]) continue;
        color[s] = 1;
        vector<int> st(1, s);
        while (!st.empty()) {
            int u = st.back();
            if (it[u] == adj[u].size()) { color[u] = 2; st.pop_back(); continue; }
            int v = adj[u][it[u]++];
            if (color[v] == 0) { color[v] = 1; par[v] = u; st.push_back(v); }
            else if (color[v] == 1) {
                vector<int> cyc;
                for (int x = u; x != v; x = par[x]) cyc.push_back(x);
                cyc.push_back(v);
                reverse(cyc.begin(), cyc.end());
                return cyc;
            }
        }
    }
    return {};
}

// (3) Yes/no for an undirected graph without a search order: cycle iff m > n - (#components).
bool undirectedHasCycle(int n, const vector<pair<int, int>>& edges) {
    vector<int> p(n);
    iota(p.begin(), p.end(), 0);
    auto f = [&](int x) {                              // iterative find (no deep recursion)
        while (p[x] != x) x = p[x] = p[p[x]];
        return x;
    };
    for (auto& e : edges) {
        int a = f(e.first), b = f(e.second);
        if (a == b) return true;                       // this edge closes a cycle
        p[a] = b;
    }
    return false;
}

int main() {
    // the graph D of Part A, 0-indexed (vertex k is k-1)
    vector<vector<int>> D = {{1, 4, 5}, {2, 3}, {0}, {2}, {3, 5}, {6}, {5}, {6}};
    for (int x : directedCycle(D)) cout << x + 1 << " ";       // 1 2 3
    cout << "\n";
    // G3 of Chapter 1, 0-indexed
    vector<pair<int, int>> G = {{0, 1}, {0, 3}, {1, 2}, {1, 4}, {3, 4}, {4, 5}};
    for (int x : undirectedCycle(7, G)) cout << x + 1 << " ";   // 1 2 5 4
    cout << "\n";
    // a path 1 - 2 - 3 and two parallel edges between 6 and 7
    vector<pair<int, int>> F = {{0, 1}, {1, 2}, {5, 6}, {6, 5}};
    for (int x : undirectedCycle(7, F)) cout << x + 1 << " ";   // 6 7
    cout << "\n" << undirectedHasCycle(7, G) << undirectedHasCycle(3, {{0, 1}, {1, 2}}) << "\n";  // 10
}
\end{lstlisting}

\begin{insightbox}[Which template when?]
\renewcommand{\arraystretch}{1.25}
\begin{tabularx}{\linewidth}{@{}>{\bfseries}l X@{}}
(1) \texttt{undirectedCycle} & print a cycle of an undirected graph; safe with parallel edges and self-loops \\
(2) \texttt{directedCycle} & print a directed cycle, or prove acyclicity (then the finishing order is a topological order, Chapter~5) \\
(3) \texttt{undirectedHasCycle} & yes/no only, or ``the first edge that closes a cycle'' when edges arrive one by one (Chapter~6) \\
\end{tabularx}
\end{insightbox}

%------------------------------------------------------------
\subsubsection*{Common mistakes}

\begin{warningbox}[Mistake 1 --- ``visited'' instead of ``gray'' in a directed graph]
\begin{lstlisting}[style=snippet]
else if (color[v] != 0) { /* cycle! */ }      // WRONG: a black v is not on the path
else if (color[v] == 1) { /* cycle! */ }      // RIGHT
\end{lstlisting}
Every DAG with a vertex of in-degree $\ge 2$ makes the wrong version report a cycle (and usually crash while printing it, Lesson~\ref{ssec:ch2-cycles}).
\end{warningbox}

\begin{warningbox}[Mistake 2 --- forgetting to skip the parent in an undirected graph]
Without the skip, every tree edge $p - v$ is seen again from $v$ and reported as the ``cycle'' $p, v$ of length 2. Round Trip requires at least 3 cities, so this is \wa{} on almost every test.
\end{warningbox}

\begin{warningbox}[Mistake 3 --- starting only from vertex 1]
The cycle can be in a component that does not contain vertex 1. Always loop over all start vertices.
\end{warningbox}

\begin{warningbox}[Mistake 4 --- recursion on $10^5$ vertices]
A path of $10^5$ cities followed by one closing road is a legal input; a recursive DFS goes $10^5$ levels deep. Use the iterative templates, or make sure the judge has a large stack (Lesson~\ref{ssec:ch2-call-stack}).
\end{warningbox}

\begin{warningbox}[Mistake 5 --- printing the directed cycle backwards]
The parent walk from $u$ to $v$ lists the cycle against the direction of the edges. Reverse it before printing (template 2); a judge that checks every consecutive pair rejects the unreversed list.
\end{warningbox}

\begin{checklistbox}[``I'm stuck'' checklist for Pattern 1]
\begin{itemize}[leftmargin=*]
    \item[$\square$] Directed or undirected? Simple, or are parallel edges and self-loops possible?
    \item[$\square$] Does a cycle mean ``yes'' (round trip exists) or ``no'' (orders contradict, deadlock)?
    \item[$\square$] Do I need the cycle itself, or only its existence (then counting or Kahn may be simpler)?
    \item[$\square$] Am I starting a search from \emph{every} white vertex?
    \item[$\square$] Is the minimum length of a valid cycle (3 cities? 4 cells?) guaranteed by my skip rule?
    \item[$\square$] Output format: number of vertices, first vertex repeated, direction of the walk?
\end{itemize}
\end{checklistbox}

%============================================================
% CH2 · P1 · PROBLEM 1 — ROUND TRIP (CSES 1669)
%============================================================
\subsection{Problem 2.1.1: Round Trip}
\problemlabel{2.1.1}{ssec:ch2-round-trip}

\problemcard{2.1.1}{Round Trip}{CSES 1669}{https://cses.fi/problemset/task/1669}{\exeasy\ (CSES Graph Algorithms; the classic first cycle problem)}{Undirected cycle detection: DFS, skip the parent, a visited neighbour closes a cycle; print it from the parent pointers}

\subsubsection*{1. Problem statement}

Byteland has $n$ cities and $m$ two-way roads. Design a \textbf{round trip} that starts in a city, visits two or more other cities and returns to the start; all intermediate cities must be distinct. Print the number $k$ of cities on the route (the start counted twice) and the route itself, or \texttt{IMPOSSIBLE}. Any valid route is accepted.

\textbf{Constraints:} $1 \le n \le 10^5$, $1 \le m \le 2 \cdot 10^5$; every road joins two different cities, and there is at most one road between two cities (a \emph{simple} graph).

\begin{table}[H]
\centering\small
\begin{tabular}{@{}lll@{}}
\toprule
\textbf{Input} & \textbf{Official output} & \textbf{Our output} \\
\midrule
\texttt{5 6 / 1 3 / 1 2 / 5 3 / 1 5 / 2 4 / 4 5} & \texttt{4 / 3 5 1 3} & \texttt{4 / 1 5 3 1} \\
\bottomrule
\end{tabular}
\end{table}

\begin{diagrambox}
\begin{tikzpicture}[x=1.15cm, y=1cm]
    \begin{scope}
        \node[gv] (a1) at (1.6,1.3) {1}; \node[gv] (a3) at (0,0.3) {3}; \node[gv] (a2) at (3.2,0.3) {2};
        \node[gv] (a5) at (1.6,-0.7) {5}; \node[gv] (a4) at (3.2,-1.2) {4};
        \foreach \a/\b in {1/3, 1/2, 3/5, 1/5, 2/4, 4/5} \draw[ge] (a\a) -- (a\b);
        \node[font=\scriptsize\bfseries] at (1.6,-1.9) {the sample: 5 cities, 6 roads};
    \end{scope}
    \begin{scope}[xshift=5.4cm]
        \node[gvsel] (b1) at (1.6,1.3) {1}; \node[gvsel] (b3) at (0,0.3) {3}; \node[gv] (b2) at (3.2,0.3) {2};
        \node[gvsel] (b5) at (1.6,-0.7) {5}; \node[gv] (b4) at (3.2,-1.2) {4};
        \foreach \a/\b in {1/2, 2/4, 4/5} \draw[ge] (b\a) -- (b\b);
        \foreach \a/\b in {1/3, 3/5, 1/5} \draw[gpath] (b\a) -- (b\b);
        \node[font=\scriptsize\bfseries] at (1.6,-1.9) {our answer: $1 \to 5 \to 3 \to 1$};
    \end{scope}
    \begin{scope}[xshift=10.8cm]
        \node[gvsel] (c1) at (1.6,1.3) {1}; \node[gv] (c3) at (0,0.3) {3}; \node[gvsel] (c2) at (3.2,0.3) {2};
        \node[gvsel] (c5) at (1.6,-0.7) {5}; \node[gvsel] (c4) at (3.2,-1.2) {4};
        \foreach \a/\b in {1/3, 3/5} \draw[ge] (c\a) -- (c\b);
        \foreach \a/\b in {1/2, 2/4, 4/5, 1/5} \draw[gpath] (c\a) -- (c\b);
        \node[font=\scriptsize\bfseries] at (1.6,-1.9) {also valid: $1, 2, 4, 5, 1$};
    \end{scope}
\end{tikzpicture}
\figcap{The sample has three cycles ($1,3,5$; $1,2,4,5$; $1,2,4,5,3$); the judge accepts any of them. The official answer $3,5,1,3$ is the first triangle read from 3.}
\end{diagrambox}

\subsubsection*{2. The graph model}

\begin{modelbox}[Graph Model --- Round Trip]
\renewcommand{\arraystretch}{1.35}
\begin{tabularx}{\linewidth}{@{}>{\bfseries}l X@{}}
Vertices & the $n$ cities \\
Edges & undirected, one per road; simple graph (no loops, no parallel edges); possibly disconnected \\
Weights & none \\
Question & find any cycle; a cycle in a simple graph automatically has $\ge 3$ vertices, which is exactly ``two or more other cities'' \\
Size & $|V| \le 10^5$, $|E| \le 2\cdot10^5$: we need \VE; a DFS path can be $10^5$ deep \\
\end{tabularx}
\end{modelbox}

\subsubsection*{3. Key observations}

\begin{itemize}[leftmargin=*]
    \item \textbf{``Round trip with distinct intermediate cities'' $=$ ``cycle''.} The route repeats only its start, so it is a cycle of the graph; conversely every cycle is a valid route.
    \item \textbf{Length $\ge 3$ is free.} A cycle of length 1 needs a loop, one of length 2 needs two parallel roads; the input has neither. So the only thing to avoid is the fake ``cycle'' $p, v, p$ made by walking a tree edge back --- skip the parent.
    \item \textbf{Existence is easy, printing needs the DFS.} A forest with $c$ components has $n - c$ edges, so a cycle exists iff $m > n - c$. To print one we need the moment a non-tree edge is found.
    \item \textbf{Depth.} A path $1 - 2 - \dots - 10^5$ plus the road $10^5 - 1$ is a valid input whose DFS goes $10^5$ deep: write the DFS iteratively.
\end{itemize}

\subsubsection*{4. Approaches}

\begin{bruteforcebox}
\textbf{Idea: an edge lies on a cycle iff its ends stay connected without it.} Try every road $(a,b)$: run a BFS from $a$ that is not allowed to use this road. If it reaches $b$, the BFS path $a \leadsto b$ plus the road is a cycle.

\begin{lstlisting}[style=cpp]
#include <bits/stdc++.h>
using namespace std;

// Brute force: an edge (a, b) lies on a cycle iff a and b stay connected without it.
// Try every edge; BFS from a to b avoiding that edge; the BFS path + the edge is the cycle.
int main() {
    int n, m;
    cin >> n >> m;
    vector<pair<int, int>> e(m);
    vector<vector<pair<int, int>>> adj(n + 1);          // (neighbour, edge id)
    for (int i = 0; i < m; i++) {
        cin >> e[i].first >> e[i].second;
        adj[e[i].first].push_back({e[i].second, i});
        adj[e[i].second].push_back({e[i].first, i});
    }
    for (int skip = 0; skip < m; skip++) {
        int a = e[skip].first, b = e[skip].second;
        vector<int> par(n + 1, -1);
        queue<int> q;
        q.push(a); par[a] = a;
        while (!q.empty()) {
            int u = q.front(); q.pop();
            for (auto& pr : adj[u])
                if (pr.second != skip && par[pr.first] == -1) { par[pr.first] = u; q.push(pr.first); }
        }
        if (par[b] == -1) continue;                     // edge skip is on no cycle
        vector<int> cyc;
        for (int x = b; x != a; x = par[x]) cyc.push_back(x);
        cyc.push_back(a);
        cyc.push_back(b);                               // close the loop with the skipped edge
        cout << cyc.size() << "\n";
        for (int x : cyc) cout << x << " ";
        cout << "\n";
        return 0;
    }
    cout << "IMPOSSIBLE\n";
}
\end{lstlisting}

On the sample the first road $(1,3)$ already works: the BFS from 1 avoiding it reaches 3 via 5, and the program prints \texttt{4 / 3 5 1 3} --- by chance exactly the official answer.

\begin{complexitybox}
\textbf{Time:} up to $m$ searches of \bigO{n + m} each, \bigO{m(n+m)}. A bad input: a tree on $10^5 - 3$ cities listed first, then a separate triangle --- every tree road costs a BFS over the whole tree before the triangle is reached, about $10^{10}$ steps. \textbf{Space:} \bigO{n+m}. \textbf{Verdict:} \tle{}, but a perfect reference for testing.
\end{complexitybox}
\end{bruteforcebox}

\begin{algorithmbox}[Approach 2: one DFS, the first non-tree edge closes the cycle]
\textbf{Idea.} Run DFS from every unvisited city. Sitting at $u$, skip the parent; any other neighbour $v$ that is already visited is an \textbf{ancestor} of $u$ (undirected DFS has only tree and back edges, Lesson~\ref{ssec:ch2-dfs-tree}), so $v, u, \texttt{par}[u], \dots, v$ is a cycle. Iterative, with \texttt{it[u]} as in Lesson~\ref{ssec:ch2-call-stack}.

\begin{lstlisting}[style=cpp]
#include <bits/stdc++.h>
using namespace std;

// Round Trip (CSES 1669): find a cycle of length >= 3 in a simple undirected graph.
// Iterative DFS; skipping the parent vertex is safe because the graph is simple.
int main() {
    ios::sync_with_stdio(false);
    cin.tie(nullptr);
    int n, m;
    cin >> n >> m;
    vector<vector<int>> adj(n + 1);
    for (int i = 0; i < m; i++) {
        int a, b;
        cin >> a >> b;
        adj[a].push_back(b);
        adj[b].push_back(a);
    }
    vector<int> par(n + 1, 0);
    vector<char> seen(n + 1, 0);
    vector<size_t> it(n + 1, 0);
    for (int s = 1; s <= n; s++) {
        if (seen[s]) continue;
        seen[s] = 1;
        vector<int> st(1, s);
        while (!st.empty()) {
            int u = st.back();
            if (it[u] == adj[u].size()) { st.pop_back(); continue; }   // u is finished
            int v = adj[u][it[u]++];
            if (v == par[u]) continue;                 // the tree edge we came in by
            if (!seen[v]) {                            // tree edge u - v
                seen[v] = 1; par[v] = u;
                st.push_back(v);
            } else {                                   // back edge u - v: v is an ancestor of u
                vector<int> cyc;
                cyc.push_back(v);
                for (int x = u; x != v; x = par[x]) cyc.push_back(x);
                cyc.push_back(v);                      // v, u, par[u], ..., v
                cout << cyc.size() << "\n";
                for (int x : cyc) cout << x << " ";
                cout << "\n";
                return 0;
            }
        }
    }
    cout << "IMPOSSIBLE\n";
}
\end{lstlisting}

\begin{complexitybox}
\textbf{Time:} \VE{} --- every city is pushed once, every road examined twice; the cycle is printed in \bigO{n}. \textbf{Space:} \bigO{n+m}, no recursion. \textbf{Verdict:} \ac{}; on a path of $10^5$ cities closed by one road it answers in well under a second (the printed cycle has all $10^5$ cities).
\end{complexitybox}
\end{algorithmbox}

\begin{invariantbox}[Why the visited neighbour is an ancestor]
\textbf{Invariant.} The stack holds the current path $x_0, \dots, x_k = u$ with $\texttt{par}[x_{i+1}] = x_i$; every finished vertex has had all its roads examined.

Suppose $u$ examines a road to a visited $v \ne \texttt{par}[u]$. If $v$ were finished and a descendant of $u$, then $v$ had examined the same road earlier, found $u$ visited and not its parent, and the program would already have stopped. A finished non-descendant is impossible in an undirected DFS: it would have entered $u$ itself while scanning that road. So $v$ is on the stack, an ancestor $x_j$ with $j \le k - 2$ ($j = k-1$ would be the parent, and there is only one road between two cities). The walk $v = x_j, x_k, x_{k-1}, \dots, x_{j+1}, x_j$ is a cycle of $k - j + 1 \ge 3$ cities.

If the search ends without stopping, every road was a tree edge, the graph is a forest, and \texttt{IMPOSSIBLE} is correct. \hfill$\blacksquare$
\end{invariantbox}

\begin{alternativebox}[Approach 3: DSU --- the first road that joins two connected cities]
Add the roads one by one to a disjoint-set union (Chapter~6). A road whose ends are already in the same set closes a cycle; the cycle is that road plus the unique path between its ends in the forest of the roads accepted so far, found by one BFS. Union by size keeps the sets shallow --- with the naive ``\texttt{dsu[find(a)] = find(b)}'' and a recursive \texttt{find}, the path input above builds a chain of $10^5$ parents, and our first version of this program crashed with a stack overflow inside \texttt{find}.

\begin{lstlisting}[style=cpp]
#include <bits/stdc++.h>
using namespace std;

// Alternative: add the roads one by one to a DSU. The first road (a, b) whose ends are already
// connected closes a cycle: the forest path from a to b plus the road itself.
vector<int> dsu, sz;
int findSet(int x) {                              // iterative, with path halving
    while (dsu[x] != x) x = dsu[x] = dsu[dsu[x]];
    return x;
}

int main() {
    ios::sync_with_stdio(false);
    cin.tie(nullptr);
    int n, m;
    cin >> n >> m;
    dsu.resize(n + 1); sz.assign(n + 1, 1);
    iota(dsu.begin(), dsu.end(), 0);
    vector<vector<int>> forest(n + 1);
    for (int i = 0; i < m; i++) {
        int a, b;
        cin >> a >> b;
        int ra = findSet(a), rb = findSet(b);
        if (ra != rb) {                                // a new forest edge: union by size
            if (sz[ra] < sz[rb]) swap(ra, rb);
            dsu[rb] = ra; sz[ra] += sz[rb];
            forest[a].push_back(b);
            forest[b].push_back(a);
            continue;
        }
        // a and b already connected: BFS in the forest for the unique path a ... b
        vector<int> par(n + 1, 0);
        par[a] = a;
        queue<int> q;
        q.push(a);
        while (!q.empty()) {
            int u = q.front(); q.pop();
            for (int v : forest[u])
                if (!par[v]) { par[v] = u; q.push(v); }
        }
        vector<int> cyc;
        for (int x = b; x != a; x = par[x]) cyc.push_back(x);
        cyc.push_back(a);
        cyc.push_back(b);
        cout << cyc.size() << "\n";
        for (int x : cyc) cout << x << " ";
        cout << "\n";
        return 0;
    }
    cout << "IMPOSSIBLE\n";
}
\end{lstlisting}
On the sample the roads $1{-}3$, $1{-}2$, $5{-}3$ are accepted, and $1{-}5$ is the first road inside one set: output \texttt{4 / 5 3 1 5}.
\begin{complexitybox}
\textbf{Time:} \bigO{m\,\alpha(n)} for the DSU plus one \bigO{n} BFS. \textbf{Space:} \bigO{n}. \textbf{Verdict:} \ac{}. The DSU version finds the cycle closed by the \emph{earliest} possible road in input order --- useful when the problem asks for that (``the first road that creates a loop'').
\end{complexitybox}
\end{alternativebox}

\begin{tricksbox}[Approach 4: Implementation Refinements]
\begin{itemize}[leftmargin=*]
    \item \textbf{Recursive version.} Ten lines shorter, but it recurses $10^5$ levels deep on the path input; on a 2\,MB stack (Lesson~\ref{ssec:ch2-call-stack}) that crashes. Unless you know the judge's stack is large, prefer the iterative form for $n \ge 5\cdot10^4$.
    \item \textbf{Parent vertex vs.\ parent edge.} Skipping \texttt{v == par[u]} is correct only because the graph is simple. For multigraphs store edge ids (template 1 of the pattern).
    \item \textbf{Fast output.} The cycle can have $10^5$ cities; print with one \texttt{"\textbackslash n"} at the end, never \texttt{endl} after every number.
    \item \textbf{Root's parent.} \texttt{par[s] = 0} works because cities are numbered from 1, so 0 is never a neighbour.
\end{itemize}
\end{tricksbox}

\subsubsection*{5. Trace}

\begin{tracebox}[{Sample: adjacency lists 1: 3 2 5,\ 2: 1 4,\ 3: 1 5,\ 4: 2 5,\ 5: 3 1 4}]
\begin{center}
\begin{tikzpicture}[x=0.85cm, y=0.85cm]
    \tikzset{gv/.append style={minimum size=6.5mm, font=\scriptsize\bfseries}}
    \foreach \k/\t/\sa/\sb/\sc/\sd/\se in {
        0/{1. enter 1; 1--3 is new}/gvcur/gv/gv/gv/gv,
        1/{2. at 3: skip 1; 3--5 is new}/gvseen/gv/gvcur/gv/gv,
        2/{3. at 5: skip 3; 5--1 is visited!}/gvsel/gv/gvsel/gv/gvcur}{
        \begin{scope}[xshift=\k*5.4cm]
            \node[\sa] (t1) at (1.6,1.3) {1}; \node[\sc] (t3) at (0,0.3) {3}; \node[\sb] (t2) at (3.2,0.3) {2};
            \node[\se] (t5) at (1.6,-0.7) {5}; \node[\sd] (t4) at (3.2,-1.2) {4};
            \foreach \a/\b in {1/2, 2/4, 4/5} \draw[ge] (t\a) -- (t\b);
            \ifnum\k=0 \draw[ge] (t1) -- (t3); \draw[ge] (t3) -- (t5); \draw[ge] (t1) -- (t5); \fi
            \ifnum\k=1 \draw[gtree] (t1) -- (t3); \draw[ge] (t3) -- (t5); \draw[ge] (t1) -- (t5); \fi
            \ifnum\k=2 \draw[gtree] (t1) -- (t3); \draw[gtree] (t3) -- (t5); \draw[gback, -] (t5) -- (t1); \fi
            \node[font=\scriptsize\bfseries] at (1.6,2.05) {\t};
        \end{scope}
    }
\end{tikzpicture}
\end{center}
\begin{center}\small
\begin{tabular}{@{}clll@{}}
\toprule
\textbf{step} & \textbf{at $u$} & \textbf{examine} & \textbf{action / stack after} \\
\midrule
1 & 1 & $1{-}3$: 3 unvisited & tree edge, \texttt{par[3] = 1}; stack \texttt{1 3} \\
2 & 3 & $3{-}1$: parent & skip \\
3 & 3 & $3{-}5$: 5 unvisited & tree edge, \texttt{par[5] = 3}; stack \texttt{1 3 5} \\
4 & 5 & $5{-}3$: parent & skip \\
5 & 5 & $5{-}1$: 1 visited, not the parent & \textbf{back edge}: cycle \texttt{1}, \texttt{5}, \texttt{par[5]=3}, \texttt{par[3]=1} \\
\bottomrule
\end{tabular}
\end{center}
Output \texttt{4} and \texttt{1 5 3 1}. Cities 2 and 4 are never visited: the search stops at the first cycle.
\end{tracebox}

\subsubsection*{6. Test suite}

\begin{testbox}
\centering\small
\begin{tabular}{@{}p{6.4cm}lp{5cm}@{}}
\toprule
\textbf{Input} & \textbf{Our output} & \textbf{Why} \\
\midrule
\texttt{2 1 / 1 2} & \texttt{IMPOSSIBLE} & one road: walking it back is not a round trip \\
\texttt{3 2 / 1 2 / 2 3} & \texttt{IMPOSSIBLE} & a path \\
\texttt{5 4 / 1 2 / 1 3 / 1 4 / 1 5} & \texttt{IMPOSSIBLE} & a star: many visited neighbours of 1, all children \\
\texttt{3 3 / 1 2 / 2 3 / 3 1} & \texttt{4 / 1 3 2 1} & smallest cycle \\
\texttt{6 5 / 1 2 / 2 3 / 4 5 / 5 6 / 6 4} & \texttt{4 / 4 6 5 4} & the cycle is not in vertex 1's component \\
\texttt{4 4 / 1 2 / 2 3 / 3 4 / 4 2} & \texttt{4 / 2 4 3 2} & the cycle does not contain the root \\
path $1..10^5$ + road $10^5{-}1$ & $k = 100001$ & deepest DFS, longest cycle \\
\bottomrule
\end{tabular}
\end{testbox}
All outputs come from running Approach~2. In addition, Approaches 1--3 were run on $3000$ random graphs with $n \le 9$; a checker verified every printed route (an edge between consecutive cities, no repeated city, at least 3 cities) and compared ``cycle or not'' with a DSU.

\subsubsection*{7. Common mistakes}

\begin{warningbox}[Mistake 1 --- no parent skip]
Without \texttt{if (v == par[u]) continue;}, the very first tree edge is reported: on the sample the program would print the ``round trip'' \texttt{1 3 1}, which visits only one other city.
\end{warningbox}

\begin{warningbox}[Mistake 2 --- searching only from city 1]
The fifth test above has its only cycle in the component $\{4,5,6\}$. A single DFS from 1 prints \texttt{IMPOSSIBLE}.
\end{warningbox}

\begin{warningbox}[Mistake 3 --- forgetting to close the route]
The judge wants $k$ cities \emph{including} the start twice: \texttt{4 / 1 5 3 1}, not \texttt{3 / 1 5 3}.
\end{warningbox}

\begin{warningbox}[Mistake 4 --- walking the parents past $v$]
The loop must stop at $v$ (\texttt{x != v}), not at the root. Walking to the root prints a ``cycle'' that leaves the cycle and comes back along a tree path, repeating cities.
\end{warningbox}

\subsubsection*{8. Takeaways}

\begin{itemize}[leftmargin=*]
    \item In an undirected DFS, a visited neighbour that is not the parent is an ancestor: the cycle is the tree path between them plus that edge.
    \item ``Two or more other cities'' is the statement's way of saying ``simple graph, ignore the parent edge''.
    \item Print by walking parents from $u$ up to $v$; close the loop for the output format.
    \item $n = 10^5$ means the DFS can be $10^5$ deep: iterative, or a large stack.
\end{itemize}

\begin{variantsbox}
\variant{One-way roads}{\varin{https://cses.fi/problemset/task/1678}}
  {every road is directed.}
  {the same search, but a visited vertex closes a cycle only if it is still on the stack.}
  {colours 0/1/2 and a gray target (Problem~\ref{ssec:ch2-round-trip-ii}).}
\variant{Parallel roads allowed}{\varnone}
  {two cities may be joined by several roads; a round trip may then use two of them.}
  {cycles of length 2 now exist.}
  {store edge ids and skip only the edge used to enter $u$ (template 1 of the pattern, which finds \texttt{6 7} on two parallel edges).}
\variant{Exactly one cycle, connected}{\varin{https://codeforces.com/problemset/problem/103/B}}
  {decide whether the graph is connected and has exactly one cycle.}
  {count instead of search: a connected graph has $n - 1$ edges plus one per independent cycle.}
  {connected and $m = n$ (Exercise E1.2).}
\variant{Components without a cycle}{\varin{https://codeforces.com/problemset/problem/659/E}}
  {count the components that are trees.}
  {per component, compare its edges and vertices.}
  {one DFS per component counting both; a tree has $e = v - 1$ (Exercise E1.4).}
\variant{A round trip through city $x$}{\varnone}
  {the trip must start and end in a given city $x$.}
  {$x$ needs two neighbours $a \ne b$ that stay connected when $x$ is removed.}
  {delete $x$, label the components of the rest, and look for two neighbours of $x$ with the same label; the cycle is $x, a \leadsto b, x$. \bigO{n+m}; checked against Approach~1 restricted to edges at $x$ on random graphs.}
\end{variantsbox}

%============================================================
% CH2 · P1 · PROBLEM 2 — ROUND TRIP II (CSES 1678)
%============================================================
\subsection{Problem 2.1.2: Round Trip II}
\problemlabel{2.1.2}{ssec:ch2-round-trip-ii}

\problemcard{2.1.2}{Round Trip II}{CSES 1678}{https://cses.fi/problemset/task/1678}{\exmedium\ (the directed twin of Problem~2.1.1 --- one word changes the algorithm)}{Directed cycle detection with three colours: only an edge to a \emph{gray} vertex closes a cycle; print it by reversing the parent walk}

\subsubsection*{1. Problem statement}

Byteland has $n$ cities and $m$ \textbf{one-way} flights. Design a round trip that starts in a city, visits one or more other cities and returns to the start, with all intermediate cities distinct. Print the number $k$ of cities on the route (start counted twice) and the route, or \texttt{IMPOSSIBLE}. Any valid route is accepted.

\textbf{Constraints:} $1 \le n \le 10^5$, $1 \le m \le 2\cdot10^5$; every flight goes from a city to a \emph{different} city. Two flights $a \to b$ and $b \to a$ may both exist.

\begin{table}[H]
\centering\small
\begin{tabular}{@{}lll@{}}
\toprule
\textbf{Input} & \textbf{Official output} & \textbf{Our output} \\
\midrule
\texttt{4 5 / 1 3 / 2 1 / 2 4 / 3 2 / 3 4} & \texttt{4 / 2 1 3 2} & \texttt{4 / 1 3 2 1} \\
\bottomrule
\end{tabular}
\end{table}

\begin{diagrambox}
\begin{tikzpicture}[x=1.2cm, y=1cm]
    \begin{scope}
        \node[gv] (a1) at (0,0) {1}; \node[gv] (a3) at (1.6,0.8) {3};
        \node[gv] (a2) at (1.6,-0.8) {2}; \node[gv] (a4) at (3.2,0) {4};
        \draw[gde] (a1) -- (a3); \draw[gde] (a2) -- (a1); \draw[gde] (a2) -- (a4);
        \draw[gde] (a3) -- (a2); \draw[gde] (a3) -- (a4);
        \node[font=\scriptsize\bfseries] at (1.6,-1.65) {the sample};
    \end{scope}
    \begin{scope}[xshift=5cm]
        \node[gvsel] (b1) at (0,0) {1}; \node[gvsel] (b3) at (1.6,0.8) {3};
        \node[gvsel] (b2) at (1.6,-0.8) {2}; \node[gv] (b4) at (3.2,0) {4};
        \draw[gpath, -{Stealth[length=3mm]}] (b1) -- (b3); \draw[gpath, -{Stealth[length=3mm]}] (b3) -- (b2);
        \draw[gpath, -{Stealth[length=3mm]}] (b2) -- (b1);
        \draw[gde] (b2) -- (b4); \draw[gde] (b3) -- (b4);
        \node[font=\scriptsize\bfseries] at (1.6,-1.65) {the only cycle: $1 \to 3 \to 2 \to 1$};
    \end{scope}
    \begin{scope}[xshift=10cm]
        \node[gv] (c1) at (0,0) {1}; \node[gv] (c3) at (1.6,0.8) {3};
        \node[gv] (c2) at (1.6,-0.8) {2}; \node[gvdst] (c4) at (3.2,0) {4};
        \draw[gde] (c3) -- (c4); \draw[gde] (c2) -- (c4);
        \node[note, text width=3.2cm] at (1.6,-1.55) {city 4 has two incoming flights but lies on no cycle};
    \end{scope}
\end{tikzpicture}
\figcap{The sample. Undirected, $\{2,3,4\}$ would be a triangle too --- but $2 \to 4 \leftarrow 3$ cannot be walked around.}
\end{diagrambox}

\subsubsection*{2. The graph model}

\begin{modelbox}[Graph Model --- Round Trip II]
\renewcommand{\arraystretch}{1.35}
\begin{tabularx}{\linewidth}{@{}>{\bfseries}l X@{}}
Vertices & the $n$ cities \\
Edges & directed, one per flight; no self-loops; $a \to b$ and $b \to a$ may coexist (a 2-cycle is a valid trip: ``one or more other cities'') \\
Weights & none \\
Question & find any directed cycle \\
Size & $|V| \le 10^5$, $|E| \le 2\cdot10^5$: \VE{}, iterative DFS \\
\end{tabularx}
\end{modelbox}

\subsubsection*{3. Key observations}

\begin{itemize}[leftmargin=*]
    \item \textbf{Visited is not enough.} City 4 is reached twice ($3 \to 4$ and $2 \to 4$) without any cycle. In a directed graph a revisit means ``two paths lead here'', not ``I came back''.
    \item \textbf{Gray is exactly right.} An edge to a vertex that is still on the current DFS path closes a cycle; by the cycle theorem (Lesson~\ref{ssec:ch2-cycles}) DFS meets such an edge iff any cycle exists.
    \item \textbf{No parent skip.} Unlike Problem~\ref{ssec:ch2-round-trip}, the edge $v \to u$ back to the parent is a different flight from $u \to v$ and genuinely forms the trip $u, v, u$.
    \item \textbf{Direction of printing.} Parent pointers point \emph{against} the flights; the walk from $u$ up to $v$ must be reversed.
\end{itemize}

\subsubsection*{4. Approaches}

\begin{bruteforcebox}
\textbf{Idea: test every city as a possible start.} City $s$ lies on a cycle iff $s$ can be reached again from one of its successors. Run a BFS from the successors of $s$; if it arrives at $s$, the BFS parents spell the trip.

\begin{lstlisting}[style=cpp]
#include <bits/stdc++.h>
using namespace std;

// Brute force: city s lies on a cycle iff s can be reached again from one of its successors.
// For every s run a BFS that starts at the successors of s; if it reaches s, rebuild the path.
int main() {
    int n, m;
    cin >> n >> m;
    vector<vector<int>> adj(n + 1);
    for (int i = 0; i < m; i++) {
        int a, b;
        cin >> a >> b;
        adj[a].push_back(b);
    }
    for (int s = 1; s <= n; s++) {
        vector<int> par(n + 1, 0);
        queue<int> q;
        for (int v : adj[s])
            if (!par[v]) { par[v] = s; q.push(v); }
        while (!q.empty() && !par[s]) {
            int u = q.front(); q.pop();
            for (int v : adj[u])
                if (!par[v]) { par[v] = u; q.push(v); }
        }
        if (!par[s]) continue;                        // s is on no cycle
        vector<int> cyc;                              // walk back from s along par
        int x = s;
        do { cyc.push_back(x); x = par[x]; } while (x != s);
        cyc.push_back(s);
        reverse(cyc.begin(), cyc.end());
        cout << cyc.size() << "\n";
        for (int y : cyc) cout << y << " ";
        cout << "\n";
        return 0;
    }
    cout << "IMPOSSIBLE\n";
}
\end{lstlisting}

As a bonus the BFS finds a \emph{shortest} cycle through $s$. On the sample, $s = 1$ succeeds at once and the program prints \texttt{4 / 1 3 2 1}.
\begin{complexitybox}
\textbf{Time:} \bigO{n(n+m)} --- on a DAG every one of the $10^5$ searches fails after up to $3\cdot10^5$ steps. \textbf{Space:} \bigO{n+m}. \textbf{Verdict:} \tle; a reference for testing.
\end{complexitybox}
\end{bruteforcebox}

\begin{algorithmbox}[Approach 2: DFS with three colours]
\textbf{Idea.} Colour 0 = never entered, 1 = on the current path, 2 = finished. An edge $u \to v$ with colour 1 at $v$ is a back edge: the trip is $v \to \dots \to u \to v$ along the parent pointers. Iterative, so a chain of $10^5$ flights is no problem.

\begin{lstlisting}[style=cpp]
#include <bits/stdc++.h>
using namespace std;

// Round Trip II (CSES 1678): print a directed cycle or IMPOSSIBLE.
// Iterative DFS with colours 0 white / 1 gray / 2 black; a gray target closes a cycle.
int main() {
    ios::sync_with_stdio(false);
    cin.tie(nullptr);
    int n, m;
    cin >> n >> m;
    vector<vector<int>> adj(n + 1);
    for (int i = 0; i < m; i++) {
        int a, b;
        cin >> a >> b;
        adj[a].push_back(b);
    }
    vector<int> color(n + 1, 0), par(n + 1, 0);
    vector<size_t> it(n + 1, 0);
    for (int s = 1; s <= n; s++) {
        if (color[s]) continue;
        color[s] = 1;
        vector<int> st(1, s);
        while (!st.empty()) {
            int u = st.back();
            if (it[u] == adj[u].size()) {              // all edges of u examined
                color[u] = 2;
                st.pop_back();
                continue;
            }
            int v = adj[u][it[u]++];
            if (color[v] == 0) {                       // tree edge
                color[v] = 1; par[v] = u;
                st.push_back(v);
            } else if (color[v] == 1) {                // back edge u -> v: cycle v -> ... -> u -> v
                vector<int> cyc;
                for (int x = u; x != v; x = par[x]) cyc.push_back(x);
                cyc.push_back(v);
                reverse(cyc.begin(), cyc.end());
                cyc.push_back(v);
                cout << cyc.size() << "\n";
                for (int x : cyc) cout << x << " ";
                cout << "\n";
                return 0;
            }                                          // black: forward / cross edge, ignore
        }
    }
    cout << "IMPOSSIBLE\n";
}
\end{lstlisting}

\begin{complexitybox}
\textbf{Time:} \VE. \textbf{Space:} \bigO{n + m}. \textbf{Verdict:} \ac{}. On a DAG with $10^5$ cities and $2\cdot10^5$ flights (a long chain plus random forward flights) it prints \texttt{IMPOSSIBLE} in a fraction of a second; on the chain $1 \to 2 \to \dots \to 10^5 \to 1$ it prints the full trip of $100\,001$ cities.
\end{complexitybox}
\end{algorithmbox}

\begin{invariantbox}[Why gray means cycle and black means nothing]
\textbf{Invariant.} The gray vertices are exactly the stack $x_0, \dots, x_k = u$, with an edge $x_i \to x_{i+1}$ for every $i$; a black vertex has all its out-edges examined and \emph{no cycle was found through them}.

\emph{Gray target.} $v = x_j$: the flights $x_j \to x_{j+1} \to \dots \to x_k \to x_j$ form a trip with distinct intermediate cities.

\emph{Black target.} Every vertex reachable from a black vertex $w$ is black too: when $w$ finished, all its successors were black or had just been finished below it; if some vertex reachable from $w$ were gray, the path to it would have produced a back edge before $w$ finished. So from $v$ there is no way back onto the current stack, and $u \to v$ cannot be part of a cycle through gray vertices. Ignoring it is safe.

\emph{Completeness.} If a cycle exists, the cycle theorem guarantees a back edge, so the search cannot end with \texttt{IMPOSSIBLE}. \hfill$\blacksquare$
\end{invariantbox}

\begin{alternativebox}[Approach 3: peel the sources (Kahn), then walk backwards]
Repeatedly delete cities with no incoming flight (Kahn's algorithm, Chapter~5). If every city gets deleted, the graph is acyclic. Otherwise, every survivor still has an incoming flight from another survivor --- it was never deleted because its in-degree never dropped to 0. Start anywhere among the survivors and keep stepping to \emph{some surviving predecessor}. With finitely many cities a city must repeat, and the repeated stretch, read in flight direction, is a cycle.

\begin{lstlisting}[style=cpp]
#include <bits/stdc++.h>
using namespace std;

// Alternative: peel vertices of in-degree 0 (Kahn). If everything is peeled, there is no cycle.
// Otherwise every survivor has an incoming edge from another survivor: walk BACKWARDS along such
// edges until a vertex repeats; the repeated part, read forwards, is a cycle.
int main() {
    ios::sync_with_stdio(false);
    cin.tie(nullptr);
    int n, m;
    cin >> n >> m;
    vector<vector<int>> adj(n + 1), radj(n + 1);
    vector<int> indeg(n + 1, 0);
    for (int i = 0; i < m; i++) {
        int a, b;
        cin >> a >> b;
        adj[a].push_back(b);
        radj[b].push_back(a);
        indeg[b]++;
    }
    vector<char> alive(n + 1, 1);
    queue<int> q;
    for (int v = 1; v <= n; v++)
        if (indeg[v] == 0) q.push(v);
    int peeled = 0;
    while (!q.empty()) {
        int v = q.front(); q.pop();
        alive[v] = 0; peeled++;
        for (int w : adj[v])
            if (--indeg[w] == 0) q.push(w);
    }
    if (peeled == n) { cout << "IMPOSSIBLE\n"; return 0; }
    int start = 1;
    while (!alive[start]) start++;
    vector<int> pos(n + 1, -1), walk;                 // pos[v] = index of v in walk
    int v = start;
    while (pos[v] == -1) {
        pos[v] = walk.size();
        walk.push_back(v);
        for (int u : radj[v])                          // any alive predecessor will do
            if (alive[u]) { v = u; break; }
    }
    vector<int> cyc(walk.begin() + pos[v], walk.end());   // v <- ... <- (back to v)
    reverse(cyc.begin(), cyc.end());                       // now in edge direction
    cyc.insert(cyc.begin(), v);                            // v -> ... -> v
    cout << cyc.size() << "\n";
    for (int x : cyc) cout << x << " ";
    cout << "\n";
}
\end{lstlisting}
On the sample no city has in-degree 0, so nothing is peeled; the backward walk is $1 \leftarrow 2 \leftarrow 3 \leftarrow 1$ and the output is again \texttt{4 / 1 3 2 1}.
\begin{complexitybox}
\textbf{Time:} \VE{} (the backward walk scans each city's predecessor list at most once). \textbf{Space:} \bigO{n+m}. \textbf{Verdict:} \ac. No DFS at all, no recursion; the peeled order is a topological order of the deleted part, which Chapter~5 develops.
\end{complexitybox}
\end{alternativebox}

\begin{tricksbox}[Approach 4: Implementation Refinements]
\begin{itemize}[leftmargin=*]
    \item \textbf{One colour array for all starts.} Black vertices from an earlier start are never explored again; a fresh array per start would make the search quadratic.
    \item \textbf{Stop at the first cycle.} Returning immediately keeps the code short; the colours of unfinished vertices are never needed again.
    \item \textbf{Recursive version.} Possible, but a chain of $10^5$ flights makes it $10^5$ deep (Lesson~\ref{ssec:ch2-call-stack}).
    \item \textbf{Which cycle?} DFS finds the cycle closed by the first back edge in DFS order; BFS from every start (Approach~1) would find a shortest one. The problem accepts any.
\end{itemize}
\end{tricksbox}

\subsubsection*{5. Trace}

\begin{tracebox}[{Sample: adjacency lists 1: 3,\ 2: 1 4,\ 3: 2 4,\ 4: (none)}]
\begin{center}
\begin{tikzpicture}[x=0.95cm, y=0.95cm, scale=0.92, transform shape]
    \tikzset{gv/.append style={minimum size=6.5mm, font=\scriptsize\bfseries}}
    \foreach \k/\t/\sa/\sb/\sc/\sd in {
        0/{1. enter 1 (gray)}/gvcur/gv/gv/gv,
        1/{2. 1$\to$3 tree; enter 3}/gvseen/gv/gvcur/gv,
        2/{3. 3$\to$2 tree; enter 2}/gvseen/gvcur/gvseen/gv,
        3/{4. 2$\to$1: 1 is gray!}/gvsel/gvcur/gvsel/gv}{
        \begin{scope}[xshift=\k*4.1cm]
            \node[\sa] (t1) at (0,0) {1}; \node[\sc] (t3) at (1.6,0.8) {3};
            \node[\sb] (t2) at (1.6,-0.8) {2}; \node[\sd] (t4) at (3.0,0) {4};
            \draw[gde] (t2) -- (t4); \draw[gde] (t3) -- (t4);
            \ifnum\k>0 \draw[gtree, -{Stealth[length=2.2mm]}] (t1) -- (t3); \else \draw[gde] (t1) -- (t3); \fi
            \ifnum\k>1 \draw[gtree, -{Stealth[length=2.2mm]}] (t3) -- (t2); \else \draw[gde] (t3) -- (t2); \fi
            \ifnum\k>2 \draw[gback] (t2) -- (t1); \else \draw[gde] (t2) -- (t1); \fi
            \node[font=\scriptsize\bfseries] at (1.5,1.55) {\t};
        \end{scope}
    }
\end{tikzpicture}
\end{center}
\begin{center}\small
\begin{tabular}{@{}clll@{}}
\toprule
\textbf{step} & \textbf{at $u$} & \textbf{examine} & \textbf{action} \\
\midrule
1 & 1 & $1 \to 3$, colour 0 & tree, \texttt{par[3] = 1}, stack \texttt{1 3} \\
2 & 3 & $3 \to 2$, colour 0 & tree, \texttt{par[2] = 3}, stack \texttt{1 3 2} \\
3 & 2 & $2 \to 1$, colour 1 & \textbf{back edge}: walk $2, \texttt{par}[2] = 3$, stop at $1$ $\Rightarrow$ \texttt{2 3 1} \\
\bottomrule
\end{tabular}
\end{center}
Reversed: \texttt{1 3 2}; closed: \texttt{1 3 2 1}, $k = 4$. City 4 is never entered. If the adjacency list of 3 had been \texttt{4 2}, DFS would first enter 4, finish it (black), and then the edge $2 \to 4$ would find a \emph{black} vertex --- ignored, correctly.
\end{tracebox}

\subsubsection*{6. Test suite}

\begin{testbox}
\centering\small
\begin{tabular}{@{}p{5.6cm}lp{5.6cm}@{}}
\toprule
\textbf{Input} & \textbf{Our output} & \textbf{Why} \\
\midrule
\texttt{3 3 / 1 2 / 1 3 / 2 3} & \texttt{IMPOSSIBLE} & DAG with a vertex of in-degree 2: the black-target trap \\
\texttt{2 2 / 1 2 / 2 1} & \texttt{3 / 1 2 1} & a 2-cycle is valid here (no parent skip!) \\
\texttt{4 3 / 1 2 / 3 4 / 4 3} & \texttt{3 / 3 4 3} & the cycle is not reachable from 1 \\
\texttt{4 4 / 1 2 / 2 3 / 3 4 / 4 2} & \texttt{4 / 2 3 4 2} & the cycle does not contain the start \\
\texttt{5 5 / 1 2 / 2 3 / 1 3 / 3 4 / 4 5} & \texttt{IMPOSSIBLE} & a forward edge $1 \to 3$ \\
chain $1 \to \dots \to 10^5 \to 1$ & $k = 100001$ & deepest search \\
\bottomrule
\end{tabular}
\end{testbox}
All three approaches were checked on $3000$ random directed graphs with $n \le 9$: a checker validated every printed trip, and ``cycle or not'' was compared with Kahn's algorithm.

\subsubsection*{7. Common mistakes}

\begin{warningbox}[Mistake 1 --- copying the undirected solution]
We ran Problem~\ref{ssec:ch2-round-trip}'s code (``visited and not the parent $\Rightarrow$ cycle'') with one-way adjacency lists. On the first test above it treats the black city 3, reached again by $1 \to 3$, as a cycle; its parent walk from 1 never meets 3, and the program died with \texttt{bad\_alloc}. On the second test it skips the flight $2 \to 1$ as ``the parent edge'' and prints \texttt{IMPOSSIBLE} although $1, 2, 1$ is a valid trip.
\end{warningbox}

\begin{warningbox}[Mistake 2 --- not reversing]
\texttt{2 3 1 2} uses the flights $2 \to 3$ and $3 \to 1$, which do not exist. Reverse the parent walk.
\end{warningbox}

\begin{warningbox}[Mistake 3 --- resetting colours per start]
Re-initialising \texttt{color} inside the loop over $s$ makes every start re-explore the whole graph: \bigO{n(n+m)}, \tle{} on a DAG.
\end{warningbox}

\subsubsection*{8. Takeaways}

\begin{itemize}[leftmargin=*]
    \item Directed cycle $=$ edge to a \emph{gray} vertex. Black vertices are dead ends for cycles.
    \item Undirected and directed cycle detection differ in two details: the parent skip and the meaning of ``visited''. Know both by heart.
    \item An acyclic answer comes with a bonus: the order in which vertices turn black, reversed, is a topological order (Chapter~5).
    \item Kahn's peeling is a DFS-free alternative that both detects and prints a cycle.
\end{itemize}

\begin{variantsbox}
\variant{Only yes/no: can all courses be finished?}{\varnone}
  {report only whether a cycle exists.}
  {the same graph.}
  {either approach without printing; Kahn's version is the usual one (Chapter~5).}
\variant{Contradictory orders}{\varin{https://codeforces.com/problemset/problem/1931/F}}
  {several people report the order of a chat list from their own point of view; is there one order consistent with all reports?}
  {one arrow per pair of consecutive names in a report.}
  {consistent iff the arrow graph has no directed cycle (Exercise E1.5).}
\variant{Which edge to remove}{\varin{https://leetcode.com/problems/redundant-connection-ii/}}
  {a rooted tree plus one extra directed edge; remove an edge so that a rooted tree remains.}
  {either a vertex with two parents, or a directed cycle, or both.}
  {find the in-degree-2 vertex and the cycle; combine the cases (Exercise E1.6).}
\variant{Every city on a cycle}{\varnone}
  {list all cities that lie on some round trip.}
  {a city lies on a cycle iff its strongly connected component has $\ge 2$ cities.}
  {strongly connected components (Chapter~19); for functional graphs the simpler three-state walk of Chapter~1, Pattern~5.}
\end{variantsbox}

%============================================================
% CH2 · PATTERN 1 — EXERCISES, HINTS, ANSWERS
%============================================================
\subsection{Pattern 1: Exercises and Going Further}
\label{ssec:ch2-p1-ex}

\begin{patternbox}{Pattern 1 Summary Card}
\renewcommand{\arraystretch}{1.25}
\begin{tabularx}{\linewidth}{@{}>{\bfseries}l X@{}}
Recognise & ``round trip'', ``loop'', ``deadlock'', ``contradiction in the order'', ``is it a tree / exactly one cycle?'' \\
Directed & colours $0/1/2$; an edge to a \textbf{gray} vertex closes a cycle; black targets are harmless \\
Undirected & skip the edge you came in by (the parent vertex is enough in simple graphs); any other visited neighbour is an ancestor \\
Printing & parent walk from $u$ up to $v$; reverse it for directed graphs; close the loop if the format asks for it \\
Counting & undirected yes/no: cycle iff $m > n - c$; one cycle and connected iff $m = n$ \\
Top traps & ``visited'' instead of gray; no parent skip; start only from vertex 1; recursion on $10^5$ vertices \\
\end{tabularx}
\end{patternbox}

\exlegend
\begin{enumerate}[leftmargin=*, label=\textbf{E1.\arabic*}, itemsep=6pt]
    \item \exeasy\ \pconcept \exown{The stack is the cycle.} Rewrite the directed cycle finder without a parent array: keep the current path in a vector \texttt{path} and an index \texttt{pos[v]} (position of $v$ in \texttt{path}, or $-1$). When an edge reaches a vertex with $\texttt{pos}[v] \ne -1$, cut the cycle out of \texttt{path}. Test it on $D$ (Part~A) and on a graph with a 2-cycle.
    \item \exeasy\ \pclassic \cf{1500}\ \exprob{Cthulhu}{Codeforces 103B}{https://codeforces.com/problemset/problem/103/B}. An undirected simple graph ($n \le 100$) is ``Cthulhu'' if it consists of three or more rooted trees whose roots are joined by a simple cycle. Decide whether the graph is Cthulhu.
    \item \exmedium\ \pclassic \cf{1500}\ \exprob{Fox And Two Dots}{Codeforces 510B}{https://codeforces.com/problemset/problem/510/B}. An $n \times m$ board ($n, m \le 50$) of coloured dots. Is there a cycle of at least 4 distinct, side-adjacent dots of the same colour?
    \item \exmedium\ \ptricky \cf{1600}\ \exprob{New Reform}{Codeforces 659E}{https://codeforces.com/problemset/problem/659/E}. Orient every road of an undirected simple graph ($n, m \le 10^5$) so that as few cities as possible have no incoming road. Print that minimum.
    \item \exhard\ \precent \cf{1700}\ \exprob{Chat Screenshots}{Codeforces 1931F}{https://codeforces.com/problemset/problem/1931/F}. $n$ chat members are ordered by activity, but everyone sees himself on top. $k$ members post screenshots of the order they see. Is there one order consistent with all screenshots? ($n \cdot k \le 2\cdot10^5$.)
    \item \exhard\ \ptricky \exprob{Redundant Connection II}{LeetCode 685}{https://leetcode.com/problems/redundant-connection-ii/}\ \lc{Hard}. A rooted tree on $n \le 1000$ vertices received one extra directed edge. Return an edge whose removal leaves a rooted tree; if several work, the one that appears last in the input.
\end{enumerate}

\subsubsection*{Hints}
\begin{enumerate}[leftmargin=*, label=\textbf{E1.\arabic*}]
    \item \texttt{pos[v] != -1} is exactly ``$v$ is gray''. Reset \texttt{pos[u] = -1} when $u$ is left.
    \item ``Trees hanging from a cycle'' is a connected graph with exactly one cycle. Count instead of searching.
    \item Each colour is a separate grid graph. A grid graph is bipartite, so its cycles are automatically even --- how short can one be?
    \item In a component with a cycle, can every city get an incoming road? In a tree, how many roads are there per city?
    \item A screenshot $a_1, a_2, \dots, a_n$ says nothing about its author $a_1$, but fixes the order of $a_2, \dots, a_n$.
    \item Either some vertex has two parents, or there is a directed cycle, or both. Try deleting the later of the two parent edges and test what remains.
\end{enumerate}

\subsubsection*{Answers}

\begin{solutionbox}[E1.1 The stack is the cycle]
The path vector replaces both the gray colour and the parent array: the cycle is simply the tail of \texttt{path} starting at \texttt{pos[v]}, already in edge direction --- no reversal needed.
\begin{lstlisting}[style=cpp]
#include <bits/stdc++.h>
using namespace std;

// E1.1: the stack IS the cycle. Keep the current path in a vector; pos[v] = index of v in it,
// or -1 if v is not on the path. "Gray" is now simply pos[v] != -1.
vector<vector<int>> adj;
vector<int> path, pos;
vector<bool> done;                                // black
vector<int> cycleFound;

bool dfs(int u) {
    pos[u] = path.size();
    path.push_back(u);
    for (int v : adj[u]) {
        if (pos[v] != -1) {                       // v is on the path: cut the cycle out of it
            cycleFound.assign(path.begin() + pos[v], path.end());
            return true;
        }
        if (!done[v] && dfs(v)) return true;
    }
    path.pop_back();
    pos[u] = -1;
    done[u] = true;
    return false;
}

int main() {
    int n, m;
    cin >> n >> m;
    adj.assign(n + 1, vector<int>());
    for (int i = 0; i < m; i++) {
        int a, b;
        cin >> a >> b;
        adj[a].push_back(b);
    }
    pos.assign(n + 1, -1);
    done.assign(n + 1, false);
    for (int s = 1; s <= n; s++)
        if (!done[s] && dfs(s)) break;
    if (cycleFound.empty()) cout << "acyclic\n";
    else {
        for (int x : cycleFound) cout << x << " ";
        cout << "\n";
    }
}
\end{lstlisting}
On $D$ it prints \textbf{1 2 3}; on \texttt{2 2 / 1 2 / 2 1} it prints \textbf{1 2}; on the DAG \texttt{3 3 / 1 2 / 1 3 / 2 3} it prints \textbf{acyclic}; on \texttt{5 6 / 1 2 / 2 3 / 3 4 / 4 5 / 5 3 / 1 4} it prints \textbf{3 4 5} --- the vertices 1 and 2 before the cycle stay on \texttt{path} but are not part of the answer.
\end{solutionbox}

\begin{solutionbox}[E1.2 Cthulhu]
A connected graph with $n$ vertices has at least $n - 1$ edges; $n - 1$ edges make a tree (no cycle) and every further edge adds one independent cycle. So ``connected with exactly one cycle'' $\iff$ connected and $m = n$. The ``three or more trees'' condition is automatic: in a simple graph a cycle has at least 3 vertices, each the root of its own tree. One search for connectivity, \bigO{n+m}.
\begin{lstlisting}[style=cpp]
#include <bits/stdc++.h>
using namespace std;

// Cthulhu: a connected graph with exactly one cycle  <=>  connected and m == n.
int main() {
    int n, m;
    cin >> n >> m;
    vector<vector<int>> adj(n + 1);
    for (int i = 0; i < m; i++) {
        int x, y;
        cin >> x >> y;
        adj[x].push_back(y);
        adj[y].push_back(x);
    }
    vector<bool> seen(n + 1, false);
    vector<int> st(1, 1);
    seen[1] = true;
    int reached = 1;
    while (!st.empty()) {                       // any search will do: we only count
        int u = st.back(); st.pop_back();
        for (int v : adj[u])
            if (!seen[v]) { seen[v] = true; reached++; st.push_back(v); }
    }
    cout << (reached == n && m == n ? "FHTAGN!" : "NO") << "\n";
}
\end{lstlisting}
Samples: \textbf{FHTAGN!}, \textbf{NO}. Checked against counting all simple cycles explicitly on 400 random graphs with $n \le 6$. The trap: $m = n$ alone is not enough --- a triangle plus a separate edge plus an isolated vertex has $n = m = 5$ but is disconnected.
\end{solutionbox}

\begin{solutionbox}[E1.3 Fox And Two Dots]
Cells of one colour with side adjacencies form an undirected simple graph; run the undirected cycle test of the pattern on it, skipping the cell we came from. A grid graph is bipartite (colour cells like a chessboard), so every cycle has even length; the smallest possible is 4 --- exactly the required minimum, so no length check is needed.
\begin{lstlisting}[style=cpp]
#include <bits/stdc++.h>
using namespace std;

// Fox And Two Dots: undirected cycle detection on the grid graph of each colour.
// A same-coloured visited neighbour that is not the cell we came from closes a cycle.
// Grid graphs are bipartite, so every cycle has even length >= 4 automatically.
int n, m;
vector<string> g;
vector<vector<bool>> seen;
const int dr[4] = {1, -1, 0, 0}, dc[4] = {0, 0, 1, -1};

bool dfs(int r, int c, int pr, int pc) {        // (pr, pc) = the cell we came from
    seen[r][c] = true;
    for (int d = 0; d < 4; d++) {
        int nr = r + dr[d], nc = c + dc[d];
        if (nr < 0 || nr >= n || nc < 0 || nc >= m || g[nr][nc] != g[r][c]) continue;
        if (nr == pr && nc == pc) continue;      // the tree edge back to the parent
        if (seen[nr][nc]) return true;           // back edge: a cycle
        if (dfs(nr, nc, r, c)) return true;
    }
    return false;
}

int main() {
    cin >> n >> m;
    g.resize(n);
    for (auto& row : g) cin >> row;
    seen.assign(n, vector<bool>(m, false));
    for (int r = 0; r < n; r++)
        for (int c = 0; c < m; c++)
            if (!seen[r][c] && dfs(r, c, -1, -1)) { cout << "Yes\n"; return 0; }
    cout << "No\n";
}
\end{lstlisting}
All five samples reproduced (\textbf{Yes, No, Yes, Yes, No}); checked against explicit cycle enumeration on 400 random $2$--$4 \times 2$--$4$ boards with two colours. Recursion is safe here: at most $50 \cdot 50 = 2500$ levels.
\end{solutionbox}

\begin{solutionbox}[E1.4 New Reform]
The model twist: this looks like an orientation problem but is a cycle test per component. A tree component with $k$ cities has $k - 1$ roads, so at least one city gets no incoming road --- and one suffices (orient every road away from any root). A component with a cycle can reach zero: orient the cycle around, then all other roads away from the cycle. So the answer is the \textbf{number of components without a cycle}, i.e.\ with $\#\text{edges} = \#\text{vertices} - 1$.
\begin{lstlisting}[style=cpp]
#include <bits/stdc++.h>
using namespace std;

// New Reform: a component with a cycle can be oriented with no separate city;
// a tree component always keeps exactly one. Answer = number of tree components.
int main() {
    ios::sync_with_stdio(false);
    cin.tie(nullptr);
    int n, m;
    cin >> n >> m;
    vector<vector<int>> adj(n + 1);
    for (int i = 0; i < m; i++) {
        int x, y;
        cin >> x >> y;
        adj[x].push_back(y);
        adj[y].push_back(x);
    }
    vector<bool> seen(n + 1, false);
    int answer = 0;
    for (int s = 1; s <= n; s++) {
        if (seen[s]) continue;
        long long verts = 0, degSum = 0;            // degSum = 2 * (edges of the component)
        vector<int> st(1, s);
        seen[s] = true;
        while (!st.empty()) {
            int u = st.back(); st.pop_back();
            verts++;
            degSum += adj[u].size();
            for (int v : adj[u])
                if (!seen[v]) { seen[v] = true; st.push_back(v); }
        }
        if (degSum / 2 == verts - 1) answer++;      // a tree: no cycle in this component
    }
    cout << answer << "\n";
}
\end{lstlisting}
Samples: \textbf{1, 0, 1}. Checked against trying all $2^m$ orientations on 400 random graphs with $m \le 8$. Isolated cities are tree components too (1 vertex, 0 roads) and count --- they are always separate.
\end{solutionbox}

\begin{solutionbox}[E1.5 Chat Screenshots]
Build an arrow $a_j \to a_{j+1}$ for every pair of consecutive non-authors $j \ge 2$ of every screenshot. A valid order exists iff these arrows contain no directed cycle: if they are acyclic, any topological order works (and puts each author where the others agree); a cycle means two screenshots contradict each other. Iterative colour DFS per test, \bigO{nk}.
\begin{lstlisting}[style=cpp]
#include <bits/stdc++.h>
using namespace std;

// Chat Screenshots: screenshot a1, a2, ..., an says a2 before a3 before ... before an
// (a1 is the author, always shown on top). An order exists iff these arrows have no cycle.
int main() {
    ios::sync_with_stdio(false);
    cin.tie(nullptr);
    int t;
    cin >> t;
    while (t--) {
        int n, k;
        cin >> n >> k;
        vector<vector<int>> adj(n + 1);
        for (int i = 0; i < k; i++) {
            vector<int> a(n);
            for (auto& x : a) cin >> x;
            for (int j = 1; j + 1 < n; j++) adj[a[j]].push_back(a[j + 1]);   // skip the author a[0]
        }
        vector<int> color(n + 1, 0);
        vector<size_t> it(n + 1, 0);
        bool cyclic = false;
        for (int s = 1; s <= n && !cyclic; s++) {
            if (color[s]) continue;
            color[s] = 1;
            vector<int> st(1, s);
            while (!st.empty() && !cyclic) {
                int u = st.back();
                if (it[u] == adj[u].size()) { color[u] = 2; st.pop_back(); continue; }
                int v = adj[u][it[u]++];
                if (color[v] == 1) cyclic = true;            // gray: contradiction
                else if (color[v] == 0) { color[v] = 1; st.push_back(v); }
            }
        }
        cout << (cyclic ? "NO" : "YES") << "\n";
    }
}
\end{lstlisting}
The official sample (ten tests) is reproduced: \textbf{YES YES YES YES NO YES YES YES YES NO}. Checked against trying every order of the members on 1500 random tests with $n \le 5$. Why ignoring the author is right: in a screenshot the author's own position carries no information about the real order.
\end{solutionbox}

\begin{solutionbox}[E1.6 Redundant Connection II]
Every vertex of a rooted tree has exactly one parent, so the extra edge $u \to v$ either gives $v$ a second parent, or points into the root and closes a directed cycle --- or both (it gives $v$ a second parent \emph{and} the old parent edge of $v$ lies on a cycle).
\begin{itemize}[leftmargin=*]
    \item \textbf{No vertex with two parents:} the edges of the unique cycle are the candidates; return the one latest in the input.
    \item \textbf{$v$ has parent edges $A$ (earlier) and $B$ (later):} delete $B$. If no directed cycle remains, $B$ is the answer (it is the later candidate); otherwise $A$ must go.
\end{itemize}
\begin{lstlisting}[style=cpp]
#include <bits/stdc++.h>
using namespace std;

// Redundant Connection II: a rooted tree plus one extra directed edge.
// Case A: some vertex has two parents -> one of its two incoming edges must go.
// Case B: no such vertex -> the extra edge closed a directed cycle; remove its last edge in input order.
class Solution {
    // directed cycle among the given edges (skipping edge 'skip'); returns the edge ids on it
    vector<int> cycleEdges(int n, const vector<vector<int>>& edges, int skip) {
        vector<vector<pair<int, int>>> adj(n + 1);   // (target, edge id)
        for (int i = 0; i < (int)edges.size(); i++)
            if (i != skip) adj[edges[i][0]].push_back({edges[i][1], i});
        vector<int> color(n + 1, 0), parEdge(n + 1, -1);
        vector<size_t> it(n + 1, 0);
        for (int s = 1; s <= n; s++) {
            if (color[s]) continue;
            color[s] = 1;
            vector<int> st(1, s);
            while (!st.empty()) {
                int u = st.back();
                if (it[u] == adj[u].size()) { color[u] = 2; st.pop_back(); continue; }
                int v = adj[u][it[u]].first, id = adj[u][it[u]].second;
                it[u]++;
                if (color[v] == 0) { color[v] = 1; parEdge[v] = id; st.push_back(v); }
                else if (color[v] == 1) {            // back edge id closes the cycle v ... u -> v
                    vector<int> ids(1, id);
                    for (int x = u; x != v; x = edges[parEdge[x]][0]) ids.push_back(parEdge[x]);
                    return ids;
                }
            }
        }
        return {};
    }
public:
    vector<int> findRedundantDirectedConnection(vector<vector<int>>& edges) {
        int n = edges.size();
        vector<int> parentEdge(n + 1, -1);
        int first = -1, second = -1;                 // the two edges into a vertex with 2 parents
        for (int i = 0; i < n; i++) {
            int v = edges[i][1];
            if (parentEdge[v] == -1) parentEdge[v] = i;
            else { first = parentEdge[v]; second = i; }
        }
        if (second == -1) {                          // case B: only a cycle
            vector<int> ids = cycleEdges(n, edges, -1);
            return edges[*max_element(ids.begin(), ids.end())];
        }
        // case A: drop the later edge; if no cycle remains it was the right one
        if (cycleEdges(n, edges, second).empty()) return edges[second];
        return edges[first];
    }
};

int main() {
    Solution s;
    vector<vector<int>> a = {{1, 2}, {1, 3}, {2, 3}}, b = {{1, 2}, {2, 3}, {3, 4}, {4, 1}, {1, 5}};
    vector<int> x = s.findRedundantDirectedConnection(a), y = s.findRedundantDirectedConnection(b);
    cout << x[0] << "," << x[1] << " " << y[0] << "," << y[1] << "\n";   // 2,3 4,1
}
\end{lstlisting}
Examples: \textbf{[2,3]} and \textbf{[4,1]}. Checked against ``try deleting each edge from the last to the first and test for a rooted tree'' on 1500 random trees with one extra edge ($n \le 7$). The trap is the mixed case: when both a second parent and a cycle exist, the answer is the parent edge \emph{on the cycle}, which may be the earlier one even though the problem prefers later edges.
\end{solutionbox}


%============================================================
\clearpage
\section{Pattern 2: Subtree Aggregation and Push-Down}
\label{sec:ch2-p2}
%============================================================
%============================================================
% CH2 · PATTERN 2 — SUBTREE AGGREGATION AND PUSH-DOWN: THEORY
%============================================================
\subsection{Pattern Theory}
\label{ssec:ch2-p2-theory}

% the tree T of the toolbox lesson as a pic with per-vertex styles (prefix g2p2-)
\tikzset{
    g2p2-v1/.style=gv, g2p2-v2/.style=gv, g2p2-v3/.style=gv, g2p2-v4/.style=gv,
    g2p2-v5/.style=gv, g2p2-v6/.style=gv, g2p2-v7/.style=gv, g2p2-v8/.style=gv,
    g2p2-e/.style=gtree,
    g2p2-T/.pic={
        \node[g2p2-v1] (t1) at (0,0) {1};
        \node[g2p2-v2] (t2) at (-1.8,-1) {2}; \node[g2p2-v3] (t3) at (0,-1) {3}; \node[g2p2-v4] (t4) at (1.8,-1) {4};
        \node[g2p2-v5] (t5) at (-2.5,-2) {5}; \node[g2p2-v6] (t6) at (-1.1,-2) {6}; \node[g2p2-v7] (t7) at (1.8,-2) {7};
        \node[g2p2-v8] (t8) at (1.8,-3) {8};
        \foreach \a/\b in {1/2, 1/3, 1/4, 2/5, 2/6, 4/7, 7/8} \draw[g2p2-e] (t\a) -- (t\b);
    }
}

\begin{patternbox}{Pattern 2 at a Glance}
Root the tree. Information flows in two directions along the DFS: \textbf{up}, at the moment a vertex is \emph{left} (post-order) --- every child is finished, so the vertex combines its children into a value for its whole subtree; and \textbf{down}, at the moment a vertex is \emph{entered} (pre-order) --- its parent is finished, so the vertex extends the parent's value by one step. One pass, \VE{} $= \bigO{n}$.

\begin{center}
\begin{tikzpicture}[x=1.05cm, y=0.95cm]
    \tikzset{g2p2-v1/.style=gvsrc, g2p2-v4/.style=gvcur, g2p2-v7/.style=gvdone, g2p2-v8/.style=gvdone}
    \pic {g2p2-T};
    \draw[->, thick, accentorange] (0.45,0.25) to[bend left=30] node[pos=0.5, above right, font=\scriptsize, text=accentorange] {down} (2.15,-0.6);
    \draw[->, thick, primaryblue] (2.45,-3.0) -- node[pos=0.5, right, font=\scriptsize, text=primaryblue] {up} (2.45,-1.3);
    \node[font=\scriptsize, anchor=west, align=left] at (4.0,-0.5) {\textcolor{accentorange}{\textbf{enter 4:}} $\mathrm{dep}[4] = \mathrm{dep}[1] + 1 = 1$\\[1pt]\quad the parent 1 (green) is already final};
    \node[font=\scriptsize, anchor=west, align=left] at (4.0,-2.1) {\textcolor{primaryblue}{\textbf{leave 4:}} $\mathrm{sz}[4] = 1 + \mathrm{sz}[7] = 3$\\[1pt]\quad the subtree $\{7, 8\}$ (blue) is already final};
\end{tikzpicture}
\end{center}
\begin{recurrencebox}
\centering\small
\textbf{up} (leave $v$): $\ \mathrm{agg}[v] = \mathrm{base}(v) \oplus \textstyle\bigoplus_{c} \mathrm{agg}[c]$ over the children $c$
\qquad e.g.\ $\mathrm{sz}[v] = 1 + \sum_c \mathrm{sz}[c]$\\[3pt]
the part above $v$ has $n - \mathrm{sz}[v]$ vertices\\[3pt]
\textbf{down} (enter $v$): $\ \mathrm{acc}[v] = \mathrm{acc}[\mathrm{par}(v)] \otimes \mathrm{step}(v)$
\qquad e.g.\ $\mathrm{dep}[v] = \mathrm{dep}[\mathrm{par}(v)] + 1$\\[3pt]
no recursion: one BFS gives \texttt{order}; forward over \texttt{order} $=$ down, backward $=$ up
\end{recurrencebox}
\end{patternbox}

%------------------------------------------------------------
\subsubsection*{From the previous pattern to this one}

Pattern~1 looked at the edges DFS meets that are \emph{not} tree edges. A tree has none: every edge is a tree edge, the only visited neighbour of a vertex is its parent, and the DFS from the root \emph{is} the rooted tree. So this pattern ignores edge types and uses the other thing Lesson~\ref{ssec:ch2-call-stack} showed: the two hooks. Code in the post-order hook runs when a vertex's whole subtree is done; code in the pre-order hook runs when the whole path from the root is done. The toolbox lesson (\S\ref{ssec:ch2-toolbox}) computed depth and size this way on the tree $T$; here we make that program into a pattern and remove its recursion.

Three things change compared with a DFS on a general graph:
\begin{itemize}[leftmargin=*]
    \item \textbf{No \texttt{visited} array.} Pass the parent instead: in a tree, ``$v \ne \mathrm{par}(u)$'' is the same as ``$v$ not visited''. This is the only way the input's $n - 1$ undirected edges become a rooted tree.
    \item \textbf{Values, not marks.} Each vertex stores a number (or a small record) computed from its children or from its parent.
    \item \textbf{The search order is irrelevant.} Only the tree matters, so any order in which parents come before children works --- including BFS order, or simply the labels when the input guarantees $p_i < i$.
\end{itemize}

%------------------------------------------------------------
\subsubsection*{When to recognise it (and when not)}

\begin{itemize}[leftmargin=*]
    \item ``\textbf{for every vertex}, how many / what sum / which maximum \textbf{in its subtree}'' --- post-order aggregation (Problem~\ref{ssec:ch2-white-black});
    \item ``remove a vertex / cut an edge: how big are the pieces?'' --- the pieces are the child subtrees and the part above, of size $n - \mathrm{sz}[v]$ (Problem~\ref{ssec:ch2-highest-score});
    \item ``\textbf{add to the whole subtree} of $p$'', ``a value accumulated along the path from the root'', depths, ancestors' properties --- pre-order push-down (Problem~\ref{ssec:ch2-ki});
    \item ``everything that passes through the edge above $v$'' --- traffic, cars, fuel, people: it is a function of the subtree of $v$ (E2.3);
    \item a greedy that cuts a subtree as soon as it is ``complete'' (even size, sum divisible by $k$): decide at the moment $v$ is left (E2.2, E2.5).
\end{itemize}
\textbf{Not this pattern:}
\begin{itemize}[leftmargin=*]
    \item the answer is wanted \textbf{for every possible root} (``sum of distances from every vertex'') --- that is \emph{rerooting}, a second pass that moves the root across each edge: \dpbook{19}{Rerooting};
    \item the value of a vertex is a \textbf{choice} that depends on choices in the children (take / skip a vertex, colourings, matchings, knapsack over subtrees) --- general tree DP: \dpbook{11}{DP on Trees};
    \item many \textbf{queries} or \textbf{updates on subtrees mixed in time} (``add to the subtree of $p$, then ask the sum of the subtree of $q$, then \dots'') --- flatten the tree with entry times (Pattern~3) and use a Fenwick or segment tree (Chapter~15);
    \item questions about \textbf{two vertices} (distance, lowest common ancestor) --- Chapter~14.
\end{itemize}

%------------------------------------------------------------
\subsubsection*{The algorithm and its invariant}

\begin{center}
\small
\renewcommand{\arraystretch}{1.3}
\begin{tabularx}{\linewidth}{@{}>{\bfseries}l X X@{}}
\toprule
 & \textbf{up (aggregation)} & \textbf{down (push-down)} \\
\midrule
hook & post-order: after the loop over the children & pre-order: before the loop \\
needs & every child $c$ is finished & the parent is finished \\
typical values & size, number of white vertices, sum, max, product of child sizes, a top-3 list & depth, sum along the root path, ``best path ending here'', total added by ancestors \\
iterative form & loop over \texttt{order} \textbf{backwards} & loop over \texttt{order} \textbf{forwards} \\
\bottomrule
\end{tabularx}
\end{center}

\begin{invariantbox}[When $v$ is left, every value of its subtree is final]
\textbf{Invariant (up).} At the moment the DFS leaves $v$, it has entered and left every descendant of $v$ and no other vertex since entering $v$; hence $\mathrm{agg}[c]$ is final for every child $c$, and $\mathrm{agg}[v]$ computed now is final too.

\emph{Proof by induction on the height of $v$.} A leaf has no children: $\mathrm{agg}[v] = \mathrm{base}(v)$. Otherwise the loop of $v$ calls the DFS on each child $c$ and returns only after $c$ is left; by induction $\mathrm{agg}[c]$ is final then, and the subtrees of different children are disjoint (a tree has no cross edges), so nothing is counted twice. The subtree of $v$ is $\{v\} \cup \bigcup_c \mathrm{subtree}(c)$, which is exactly what the formula combines.

\textbf{Invariant (down).} At the moment the DFS enters $v$, the value of every ancestor of $v$ is final --- they are exactly the vertices on the stack (Lesson~\ref{ssec:ch2-call-stack}), each entered before $v$.

\emph{Iterative form.} The same two statements hold for any \texttt{order} in which $\mathrm{par}(v)$ precedes $v$: going backwards, all children of $v$ appear after $v$ in \texttt{order}, so they are processed before $v$; going forwards, the parent is processed before $v$. \hfill$\blacksquare$
\end{invariantbox}

\begin{diagrambox}
\begin{tikzpicture}[x=0.9cm, y=0.85cm]
    \begin{scope}
        \pic {g2p2-T};
        \foreach \v/\s/\b/\w in {1/8/36/right, 2/3/13/left, 3/1/3/below, 4/3/19/right, 5/1/5/below, 6/1/6/below, 7/2/15/right, 8/1/8/right}
            \node[gdist, text=primaryblue, fill=white, \w=1pt of t\v] {\s\,|\,\b};
        \node[font=\small\bfseries] at (0,1.0) {up: $\mathrm{sz}\,|\,\mathrm{sub}$};
        \node[note, text width=5.2cm] at (0,-4.1) {values $\mathrm{val}(v) = v$; $\mathrm{sub}[4] = 4 + \mathrm{sub}[7] = 4 + 15 = 19$};
    \end{scope}
    \begin{scope}[xshift=7.6cm]
        \pic {g2p2-T};
        \foreach \v/\d/\p/\w in {1/0/1/right, 2/1/3/left, 3/1/4/below, 4/1/5/right, 5/2/8/below, 6/2/9/below, 7/2/12/right, 8/3/20/right}
            \node[gdist, text=accentorange, fill=white, \w=1pt of t\v] {\d\,|\,\p};
        \node[font=\small\bfseries] at (0,1.0) {down: $\mathrm{dep}\,|\,\mathrm{path}$};
        \node[note, text width=5.2cm] at (0,-4.1) {$\mathrm{path}[8] = \mathrm{path}[7] + 8 = 12 + 8 = 20 = 1 + 4 + 7 + 8$};
    \end{scope}
\end{tikzpicture}
\figcap{Both directions on $T$ with $\mathrm{val}(v) = v$ (numbers printed by the templates below). Left: subtree size and subtree sum, computed bottom-up. Right: depth and the sum along the root path, computed top-down.}
\end{diagrambox}

%------------------------------------------------------------
\subsubsection*{The part above $v$: $n - \mathrm{sz}[v]$}

Deleting a vertex $v$ splits the tree into the subtrees of its children and \emph{one} more piece: everything that is not in the subtree of $v$. That piece has $n - \mathrm{sz}[v]$ vertices (zero for the root). So after one post-order pass we know the size of every piece of every deletion, although the DFS never ``looked upwards''. Products of piece sizes (Problem~\ref{ssec:ch2-highest-score}), the largest piece (centroids, Chapter~14), and the number of paths through $v$ are all built from these numbers.

\begin{diagrambox}
\begin{tikzpicture}[x=0.9cm, y=0.85cm]
    \begin{scope}
        \tikzset{g2p2-v2/.style=gvdead, g2p2-v5/.style=gvA, g2p2-v6/.style=gvB,
                 g2p2-v1/.style=gvsel, g2p2-v3/.style=gvsel, g2p2-v4/.style=gvsel, g2p2-v7/.style=gvsel, g2p2-v8/.style=gvsel}
        \pic {g2p2-T};
        \node[font=\small\bfseries] at (0,1.0) {delete 2};
        \node[note, text width=5.6cm] at (0,-4.5) {pieces $\{5\}$, $\{6\}$ (children, size 1 each) and the purple rest of size $n - \mathrm{sz}[2] = 8 - 3 = 5$; largest piece 5};
    \end{scope}
    \begin{scope}[xshift=7.6cm]
        \tikzset{g2p2-v1/.style=gvdead, g2p2-v2/.style=gvA, g2p2-v5/.style=gvA, g2p2-v6/.style=gvA,
                 g2p2-v3/.style=gvB, g2p2-v4/.style=gvsel, g2p2-v7/.style=gvsel, g2p2-v8/.style=gvsel}
        \pic {g2p2-T};
        \node[font=\small\bfseries] at (0,1.0) {delete the root 1};
        \node[note, text width=5.6cm] at (0,-4.5) {only child pieces: $3, 1, 3$; the part above is empty ($8 - \mathrm{sz}[1] = 0$); largest piece 3 --- the smallest ``worst'' of all vertices};
    \end{scope}
\end{tikzpicture}
\figcap{Deleting a vertex: the child subtrees plus the part above. Template (4) below prints the largest piece for every $v$: $3, 5, 7, 5, 7, 7, 6, 7$.}
\end{diagrambox}

%------------------------------------------------------------
\subsubsection*{No recursion: the order array}

With $n = 2\cdot10^5$ a path-shaped tree makes a recursive DFS go $2\cdot10^5$ levels deep, and on a 2\,MB stack that crashes near $65\,000$ (Lesson~\ref{ssec:ch2-call-stack}). Neither direction needs the DFS rhythm, only ``parents before children'', so:
\begin{enumerate}[leftmargin=*]
    \item run a BFS from the root, recording \texttt{order} and \texttt{parent} (skip the parent instead of checking \texttt{visited});
    \item \textbf{forward} over \texttt{order}: push-down values;
    \item \textbf{backward} over \texttt{order}: aggregate into the parent.
\end{enumerate}
If the input is a parent array with $p_i < i$ (Codeforces loves this format), even the BFS is unnecessary: the labels $1, 2, \dots, n$ already are such an order.

\begin{diagrambox}
\begin{tikzpicture}[x=1.15cm]
    \node[font=\scriptsize\bfseries, anchor=east] at (-0.6,0) {order};
    \foreach \v [count=\i from 0] in {1,2,3,4,5,6,7,8}
        \node[cell, minimum width=9mm] (o\i) at (\i,0) {\v};
    \node[font=\scriptsize\bfseries, anchor=east] at (-0.6,-0.75) {parent};
    \foreach \p [count=\i from 0] in {--,1,1,1,2,2,4,7}
        \node[idx] at (\i,-0.75) {\p};
    \draw[->, thick, accentorange] (-0.3,1.25) -- node[above, font=\scriptsize]{forward: $\mathrm{dep}[v] = \mathrm{dep}[\mathrm{par}]+1$ (parent already done)} (7.3,1.25);
    \draw[->, thick, primaryblue] (7.3,-1.35) -- node[below, font=\scriptsize]{backward: $\mathrm{sz}[\mathrm{par}] \mathrel{+}= \mathrm{sz}[v]$ (children already done)} (-0.3,-1.35);
    \draw[dep] (o7.north) to[bend right=40] (o6.north);
    \draw[dep] (o6.north) to[bend right=40] (o3.north);
    \node[note, text width=3.2cm, anchor=west] at (8.0,-0.4) {BFS order of $T$ from 1; orange arcs: 8 hands its size to 7, then 7 to 4};
\end{tikzpicture}
\figcap{The order-array trick. BFS gives every parent before its children, so one forward and one backward loop replace both hooks of the recursive DFS.}
\end{diagrambox}

%------------------------------------------------------------
\subsubsection*{Variations map}

\begin{diagrambox}
\begin{tikzpicture}[
    var/.style={rectangle, rounded corners=2pt, draw=patternborder, fill=patternbg,
                font=\scriptsize, align=center, text width=3.25cm, minimum height=1.3cm, inner sep=3pt}]
    \node[var, fill=answerfill, draw=accentorange, text width=5.4cm, font=\scriptsize\bfseries] (root) at (0,0)
        {rooted tree; combine children at leave, extend the parent at enter};
    \node[var] (v1) at (-5.4,-2) {\textbf{count / sum per subtree}\\colour balance\\(White-Black, 2.2.1)};
    \node[var] (v2) at (-1.8,-2) {\textbf{pieces after deletion}\\child sizes and $n - \mathrm{sz}$\\(Highest Score, 2.2.2)};
    \node[var] (v3) at (1.8,-2)  {\textbf{add to a subtree}\\mark the top, push down\\(Ki, 2.2.3)};
    \node[var] (v4) at (5.4,-2)  {\textbf{per-label counts}\\counter before / after\\(E2.1)};
    \node[var] (v5) at (-5.4,-3.6) {\textbf{greedy cuts}\\cut when complete\\(E2.2, E2.5)};
    \node[var] (v6) at (-1.8,-3.6) {\textbf{traffic on an edge}\\a function of $\mathrm{sz}$\\(E2.3)};
    \node[var] (v7) at (1.8,-3.6)  {\textbf{best path ending here}\\push-down with reset\\(E2.4)};
    \node[var] (v8) at (5.4,-3.6)  {\textbf{small summary per subtree}\\top-3 / bottom-2\\(E2.6)};
    \node[draw=patternborder, dashed, rounded corners=4pt, fit=(v1)(v4)(v5)(v8), inner sep=5pt] (grp) {};
    \draw[flow] (root.south) -- (grp.north);
    \node[note, below=2mm of grp] {all roots: rerooting (DP book Ch.~19) $\cdot$ choices: tree DP (DP book Ch.~11) $\cdot$ mixed queries: Chapter~15};
\end{tikzpicture}
\end{diagrambox}

\textbf{Twists and how to handle them}
\begin{itemize}[leftmargin=*]
    \item \textbf{Edges given as pairs $a_i < b_i$:} that does \emph{not} make $a_i$ the parent. Always root the tree yourself by a search from the root (Problem~\ref{ssec:ch2-ki}).
    \item \textbf{Parent array with $p_i < i$:} loop the labels backwards --- no adjacency lists at all (Problem~\ref{ssec:ch2-white-black}).
    \item \textbf{Parent array without that guarantee} (LeetCode's \texttt{parents[]}): build children lists and a BFS order, or peel leaves (Problem~\ref{ssec:ch2-highest-score}).
    \item \textbf{Products and sums overflow:} three pieces of $3\cdot10^4$ vertices multiply to $2.7\cdot10^{13}$; sums of $2\cdot10^5$ values up to $10^9$ need 64 bits.
    \item \textbf{Merging big summaries is expensive:} keep only what the parent can use (5 numbers in E2.6, one counter difference in E2.1) instead of whole multisets.
    \item \textbf{Updates to subtrees, answers at the end:} mark only the top vertex and push down once (Problem~\ref{ssec:ch2-ki}); the ``difference on pre-order intervals'' form of the same idea becomes a Fenwick tree when queries are interleaved (Chapter~15).
\end{itemize}

%------------------------------------------------------------
\subsubsection*{Code templates}

\begin{lstlisting}[style=cpp]
#include <bits/stdc++.h>
using namespace std;
typedef long long ll;

// (1) RECURSIVE, tree given as undirected edges: the parent parameter replaces visited[].
//     Down (enter): depth, root-path sum.  Up (leave): size, subtree sum.  Depth of recursion = height.
struct TreeRec {
    vector<vector<int>> adj;
    vector<ll> val, sub, path;
    vector<int> sz, dep;
    void dfs(int u, int p) {
        dep[u] = (p == 0) ? 0 : dep[p] + 1;          // PRE-ORDER: p is final, push down
        path[u] = (p == 0 ? 0 : path[p]) + val[u];
        sz[u] = 1; sub[u] = val[u];
        for (int v : adj[u])
            if (v != p) {
                dfs(v, u);
                sz[u] += sz[v];                       // POST-ORDER: v is final, pull up
                sub[u] += sub[v];
            }
    }
    void run(int n, const vector<pair<int, int>>& edges, const vector<ll>& w, int root) {
        adj.assign(n + 1, vector<int>());
        for (auto& e : edges) { adj[e.first].push_back(e.second); adj[e.second].push_back(e.first); }
        val = w; sub.assign(n + 1, 0); path.assign(n + 1, 0); sz.assign(n + 1, 0); dep.assign(n + 1, 0);
        dfs(root, 0);
    }
};

// (2) ITERATIVE with an order array: one BFS gives parent[] and an order with parents first.
//     Forward pass = push-down, backward pass = aggregation. Safe for n = 10^6.
struct TreeIter {
    vector<int> parent, order, sz, dep;
    vector<ll> sub, path;
    void run(int n, const vector<pair<int, int>>& edges, const vector<ll>& val, int root) {
        vector<vector<int>> adj(n + 1);
        for (auto& e : edges) { adj[e.first].push_back(e.second); adj[e.second].push_back(e.first); }
        parent.assign(n + 1, 0); order.assign(1, root);
        for (size_t i = 0; i < order.size(); i++) {
            int u = order[i];
            for (int v : adj[u]) if (v != parent[u]) { parent[v] = u; order.push_back(v); }
        }
        dep.assign(n + 1, 0); path.assign(n + 1, 0);
        for (int u : order) {                         // parents before children: push down
            int p = parent[u];
            dep[u] = (p == 0) ? 0 : dep[p] + 1;
            path[u] = (p == 0 ? 0 : path[p]) + val[u];
        }
        sz.assign(n + 1, 1); sub.assign(val.begin(), val.end());
        for (int i = n - 1; i >= 1; i--) {            // children before parents: pull up
            int u = order[i];
            sz[parent[u]] += sz[u];
            sub[parent[u]] += sub[u];
        }
    }
};

// (3) PARENT ARRAY with p[i] < i (vertex 1 is the root): the labels are already an order.
vector<int> sizesFromParents(int n, const vector<int>& p) {
    vector<int> sz(n + 1, 1);
    for (int i = n; i >= 2; i--) sz[p[i]] += sz[i];  // i's children all have labels > i
    return sz;
}

// (4) THE PART ABOVE v: deleting v leaves its child subtrees and one piece of n - sz[v] vertices.
vector<int> largestPieceAfterDeleting(int n, const TreeIter& t) {
    vector<int> worst(n + 1, 0);
    for (int v = 1; v <= n; v++) worst[v] = n - t.sz[v];          // the piece above v (0 for the root)
    for (int v = 1; v <= n; v++)
        if (t.parent[v]) worst[t.parent[v]] = max(worst[t.parent[v]], t.sz[v]);   // a piece below
    return worst;
}

int main() {
    // the tree T of the toolbox lesson; the value of vertex v is v
    int n = 8;
    vector<pair<int, int>> edges = {{1, 2}, {1, 3}, {1, 4}, {2, 5}, {2, 6}, {4, 7}, {7, 8}};
    vector<ll> val(n + 1);
    for (int v = 0; v <= n; v++) val[v] = v;
    TreeRec r; r.run(n, edges, val, 1);
    TreeIter t; t.run(n, edges, val, 1);
    vector<int> worst = largestPieceAfterDeleting(n, t);
    vector<int> p = {0, 0, 1, 1, 1, 2, 2, 4, 7};                     // the same tree as a parent array
    vector<int> sz3 = sizesFromParents(n, p);
    cout << "v  sz sub dep path worst\n";
    for (int v = 1; v <= n; v++) {
        assert(r.sz[v] == t.sz[v] && r.sub[v] == t.sub[v] && r.dep[v] == t.dep[v] && r.path[v] == t.path[v] && sz3[v] == t.sz[v]);
        cout << v << "  " << t.sz[v] << "  " << t.sub[v] << "  " << t.dep[v] << "  " << t.path[v] << "  " << worst[v] << "\n";
    }
    cout << "BFS order:";
    for (int u : t.order) cout << " " << u;
    cout << "\n";
}
\end{lstlisting}

\begin{lstlisting}[style=trace]
output: v  sz sub dep path worst          5  1  5  2  8  7
        1  8  36  0  1  3                 6  1  6  2  9  7
        2  3  13  1  3  5                 7  2  15  2  12  6
        3  1  3  1  4  7                  8  1  8  3  20  7
        4  3  19  1  5  5                 BFS order: 1 2 3 4 5 6 7 8
\end{lstlisting}

\begin{insightbox}[Which template when?]
\renewcommand{\arraystretch}{1.25}
\begin{tabularx}{\linewidth}{@{}>{\bfseries}l X@{}}
(1) \texttt{TreeRec} & short and natural; use when the height is small ($n \le 5\cdot10^4$, or a judge with a big stack) \\
(2) \texttt{TreeIter} & the default for $n$ up to $10^6$: edges in any orientation, any root, no recursion \\
(3) \texttt{sizesFromParents} & input is a parent array with $p_i < i$: no adjacency lists, no search \\
(4) \texttt{largestPieceAfterDeleting} & anything about the pieces left by deleting a vertex: child sizes and $n - \mathrm{sz}[v]$ \\
\end{tabularx}
\end{insightbox}

%------------------------------------------------------------
\subsubsection*{Common mistakes}

\begin{warningbox}[Mistake 1 --- treating the first endpoint as the parent]
\begin{lstlisting}[style=snippet]
for (each edge a b) par[b] = a;          // WRONG: "a < b" or "listed first" says nothing about rooting
// RIGHT: root the tree by a BFS/DFS from the root, skipping the parent
\end{lstlisting}
On the edges $1{-}3$, $2{-}3$ with root 1, vertex 2 is a child of 3, not the other way round. Problem~\ref{ssec:ch2-ki} shows the wrong output this produces.
\end{warningbox}

\begin{warningbox}[Mistake 2 --- using a value before it is final]
Aggregation reads the children, so it must run \emph{after} the recursive calls (or in the \emph{backward} loop). Writing \texttt{sz[u] += sz[v]} before \texttt{dfs(v, u)} adds the child's initial 0. Push-down reads the parent, so it runs \emph{before} the calls (or in the \emph{forward} loop).
\end{warningbox}

\begin{warningbox}[Mistake 3 --- forgetting the part above]
The pieces left by deleting $v$ are its child subtrees \textbf{and} $n - \mathrm{sz}[v]$. Forgetting the last one (or multiplying by it for the root, where it is $0$) is the classic bug of Problem~\ref{ssec:ch2-highest-score}.
\end{warningbox}

\begin{warningbox}[Mistake 4 --- recursion on a path of $2\cdot10^5$]
Trees in tests are often paths or ``caterpillars''. Our recursive push-down for Problem~\ref{ssec:ch2-ki} died without printing anything on a path of $2\cdot10^5$ vertices; template (2) handles it instantly.
\end{warningbox}

\begin{warningbox}[Mistake 5 --- \texttt{int} for sizes multiplied or values summed]
Sizes fit in \texttt{int}, their products do not: the first tree whose best deletion score exceeds $2^{31} - 1$ has only $n = 3872$ vertices ($1290 \cdot 1290 \cdot 1291$).
\end{warningbox}

\begin{checklistbox}[``I'm stuck'' checklist for Pattern 2]
\begin{itemize}[leftmargin=*]
    \item[$\square$] Is the answer for vertex $v$ about its \emph{subtree} (up) or about its \emph{root path} (down)? Or both?
    \item[$\square$] What does a vertex need from its children: is one number enough, or a small record (counts per colour, top 3)?
    \item[$\square$] Do I need the part above $v$? It is $n - \mathrm{sz}[v]$, or ``total minus subtree'' for sums.
    \item[$\square$] Who is the root, and is the input really oriented? Root it myself.
    \item[$\square$] How deep can the tree be? Above $5\cdot10^4$: order array.
    \item[$\square$] Do products or sums need \texttt{long long}?
    \item[$\square$] Do I need the answer for \emph{every} root, or with interleaved queries? Then this is the wrong pattern (rerooting, Chapter~15).
\end{itemize}
\end{checklistbox}

%============================================================
% CH2 · P2 · PROBLEM 1 — WHITE-BLACK BALANCED SUBTREES (CF 1676G)
%============================================================
\subsection{Problem 2.2.1: White-Black Balanced Subtrees}
\problemlabel{2.2.1}{ssec:ch2-white-black}

% sample 1 tree; vertex colours W / B, per-vertex style override g2p2-s<k>
\tikzset{
    g2p2-W/.style={gv},
    g2p2-B/.style={gv, fill=darkgray, text=white, draw=black},
    g2p2-s1/.style=g2p2-W, g2p2-s2/.style=g2p2-B, g2p2-s3/.style=g2p2-B, g2p2-s4/.style=g2p2-W,
    g2p2-s5/.style=g2p2-W, g2p2-s6/.style=g2p2-B, g2p2-s7/.style=g2p2-W,
    g2p2-S/.pic={
        \node[g2p2-s1] (s1) at (0,0) {1};
        \node[g2p2-s2] (s2) at (-1.2,-1) {2}; \node[g2p2-s3] (s3) at (1.2,-1) {3};
        \node[g2p2-s4] (s4) at (-1.2,-2) {4}; \node[g2p2-s5] (s5) at (0.6,-2) {5}; \node[g2p2-s6] (s6) at (1.8,-2) {6};
        \node[g2p2-s7] (s7) at (0.6,-3) {7};
        \foreach \a/\b in {1/2, 1/3, 2/4, 3/5, 3/6, 5/7} \draw[gtree] (s\a) -- (s\b);
    }
}

\problemcard{2.2.1}{White-Black Balanced Subtrees}{Codeforces 1676G}{https://codeforces.com/problemset/problem/1676/G}{\exeasy\ \cf{1300}\ (the purest subtree aggregation)}{Post-order aggregation of one number per vertex, $+1$ for white and $-1$ for black; a subtree is balanced iff its sum is 0; parents $a_i < i$ give the order for free}

\subsubsection*{1. Problem statement}

A tree with $n$ vertices is rooted at vertex 1; vertex $i \ge 2$ has parent $a_i$ with $a_i < i$. Each vertex is white (\texttt{W}) or black (\texttt{B}), given as a string $s$. The subtree of $u$ consists of $u$ and all its descendants. A subtree is \textbf{balanced} if it contains as many white as black vertices. Count the balanced subtrees (one per vertex, so the answer is the number of vertices whose subtree is balanced).

\textbf{Constraints:} $t \le 10^4$ test cases, $2 \le n \le 4000$, $\sum n \le 2\cdot10^5$, $1 \le a_i < i$.

\begin{table}[H]
\centering\small
\begin{tabular}{@{}llcl@{}}
\toprule
\textbf{$n$ / parents $a_2..a_n$} & \textbf{$s$} & \textbf{Output} & \textbf{Balanced subtrees} \\
\midrule
\texttt{7 / 1 1 2 3 3 5} & \texttt{WBBWWBW} & \texttt{2} & of 2 and 3 \\
\texttt{2 / 1} & \texttt{BW} & \texttt{1} & of 1 \\
\texttt{8 / 1 2 3 4 5 6 7} & \texttt{BWBWBWBW} & \texttt{4} & of 1, 3, 5, 7 \\
\bottomrule
\end{tabular}
\end{table}

\begin{diagrambox}
\begin{tikzpicture}[x=1cm, y=0.9cm]
    \begin{scope}
        \pic {g2p2-S};
        \node[font=\scriptsize\bfseries] at (0,0.8) {sample 1};
    \end{scope}
    \begin{scope}[xshift=4.4cm]
        \pic {g2p2-S};
        \node[draw=successgreen, line width=1.2pt, rounded corners=6pt, dashed, fit=(s2)(s4), inner sep=3pt] {};
        \node[draw=successgreen, line width=1.2pt, rounded corners=6pt, dashed, fit=(s3)(s5)(s6)(s7), inner sep=3pt] {};
        \node[font=\scriptsize\bfseries] at (0,0.8) {the two balanced subtrees};
    \end{scope}
    \node[note, anchor=west, text width=4.6cm, align=left] at (7.2,-1.4) {subtree of 2: $\{2, 4\}$: 1 W, 1 B\\subtree of 3: $\{3, 5, 6, 7\}$: 2 W, 2 B\\subtree of 1: 4 W, 3 B --- not balanced};
    \begin{scope}[yshift=-4.3cm, xshift=-1.2cm]
        \node[font=\scriptsize\bfseries, anchor=east] at (-0.5,0) {sample 3};
        \foreach \i in {1,...,8}{
            \ifodd\i \node[g2p2-B] (q\i) at (\i*1.25-1.25,0) {\i};
            \else \node[g2p2-W] (q\i) at (\i*1.25-1.25,0) {\i}; \fi
        }
        \foreach \i [evaluate=\i as \j using int(\i+1)] in {1,...,7} \draw[gtree] (q\i) -- (q\j);
        \foreach \i in {1,3,5,7} \node[gdist, below=1pt of q\i] {balanced};
    \end{scope}
\end{tikzpicture}
\figcap{Sample 1 (black disks are black vertices) and sample 3, a path rooted at 1 that alternates colours: the subtree of $i$ is the suffix $i..8$, balanced exactly when its length is even.}
\end{diagrambox}

\subsubsection*{2. The graph model}

\begin{modelbox}[Graph Model --- White-Black Balanced Subtrees]
\renewcommand{\arraystretch}{1.35}
\begin{tabularx}{\linewidth}{@{}>{\bfseries}l X@{}}
Vertices & the $n$ tree vertices, each with a weight $w(v) = +1$ (white) or $-1$ (black) \\
Edges & tree edges $a_i \to i$, given as a parent array; $a_i < i$, so every parent has a smaller label than its children \\
Weights & the vertex weights $\pm 1$ \\
Question & for how many $v$ is $\mathrm{bal}[v] = \sum_{x \in \mathrm{subtree}(v)} w(x)$ equal to 0? \\
Size & $\sum n \le 2\cdot10^5$: \bigO{n} per test; $n \le 4000$ even allows \bigO{n^2} per test --- but $t$ tests of $n = 4000$ cost $50 \cdot 4000^2 / 2 = 4\cdot10^8$ \\
\end{tabularx}
\end{modelbox}

\begin{diagrambox}
\begin{tikzpicture}[x=1cm, y=0.9cm]
    \begin{scope}
        \node[font=\scriptsize\bfseries, anchor=east] at (-0.6,0) {$i$};
        \foreach \i in {2,...,7} \node[cell, minimum width=8mm] (c\i) at (\i*0.9-1.8,0) {\i};
        \node[font=\scriptsize\bfseries, anchor=east] at (-0.6,-0.8) {$a_i$};
        \foreach \p [count=\i from 2] in {1,1,2,3,3,5} \node[cell, minimum width=8mm, fill=theorybg] at (\i*0.9-1.8,-0.8) {\p};
        \node[note, text width=5.2cm] at (2.2,-2.0) {every $a_i < i$: reading $i = n, n-1, \dots$ meets every child before its parent};
    \end{scope}
    \begin{scope}[xshift=7.6cm]
        \pic {g2p2-S};
        \foreach \v/\w/\d in {1/+1/right, 2/-1/left, 3/-1/right, 4/+1/left, 5/+1/left, 6/-1/right, 7/+1/right}
            \node[gdist, text=deepblue, \d=1pt of s\v] {$\w$};
        \node[note, anchor=west, text width=3.1cm, align=left] at (2.6,-1.5) {white $= +1$, black $= -1$:\\balanced $\iff$ sum $= 0$};
    \end{scope}
\end{tikzpicture}
\figcap{The model: the parent array is already a valid order, and one signed number per vertex replaces the two colour counts.}
\end{diagrambox}

\subsubsection*{3. Key observations}

\begin{itemize}[leftmargin=*]
    \item \textbf{One number instead of two.} ``$\#W = \#B$'' is ``$\#W - \#B = 0$''. With $w = \pm 1$ the subtree sum $\mathrm{bal}[v]$ is all we need, and it aggregates: $\mathrm{bal}[v] = w(v) + \sum_c \mathrm{bal}[c]$.
    \item \textbf{The labels are an order.} Since $a_i < i$, all children of $v$ have labels $> v$. Looping $v = n, n-1, \dots, 1$ processes every child before its parent: at the moment we reach $v$, $\mathrm{bal}[v]$ is final (the pattern's invariant) and can be tested and handed up.
    \item \textbf{No graph is built.} No adjacency lists, no search, no recursion: two arrays.
    \item \textbf{Odd subtrees are never balanced.} A balanced subtree has even size --- a quick sanity check on traces (in sample 1 the balanced ones have sizes 2 and 4).
\end{itemize}

\subsubsection*{4. Approaches}

\begin{bruteforcebox}
\textbf{Idea: list every subtree and count.} For each $u$, collect its subtree by walking only from parents to children, and compare the two counts.

\begin{lstlisting}[style=cpp]
#include <bits/stdc++.h>
using namespace std;

// Brute force: for every vertex u, list its subtree explicitly (a search that only walks
// from a vertex to its children) and count the two colours. O(n^2) per test.
int main() {
    int t;
    cin >> t;
    while (t--) {
        int n;
        cin >> n;
        vector<vector<int>> children(n + 1);
        for (int i = 2; i <= n; i++) {
            int p;
            cin >> p;
            children[p].push_back(i);
        }
        string s;
        cin >> s;
        int balanced = 0;
        for (int u = 1; u <= n; u++) {
            int white = 0, black = 0;
            vector<int> todo(1, u);                // the subtree of u = everything below u
            while (!todo.empty()) {
                int x = todo.back(); todo.pop_back();
                if (s[x - 1] == 'W') white++; else black++;
                for (int c : children[x]) todo.push_back(c);
            }
            if (white == black) balanced++;
        }
        cout << balanced << "\n";
    }
}
\end{lstlisting}

\begin{complexitybox}
\textbf{Time:} $\sum_u \mathrm{sz}[u]$ = $\sum_v (\mathrm{dep}[v] + 1)$, which is $\Theta(n^2)$ on a path. Fifty paths of $4000$ vertices ($\sum n = 2\cdot10^5$) took about 3\,s here, against $0.26$\,s for Approach~2 including process start-up. \textbf{Space:} \bigO{n}. \textbf{Verdict:} borderline at these limits, hopeless if $n$ were $2\cdot10^5$ --- and the reference for testing.
\end{complexitybox}
\end{bruteforcebox}

\begin{algorithmbox}[Approach 2: one backward pass over the labels]
\textbf{Idea.} $\mathrm{bal}[v] = \pm1$ initially. For $v = n$ down to $1$: $\mathrm{bal}[v]$ is now final; count it if it is 0; add it to $\mathrm{bal}[a_v]$. This is template (3) of the pattern with $\pm1$ in place of $1$.

\begin{lstlisting}[style=cpp]
#include <bits/stdc++.h>
using namespace std;

// White-Black Balanced Subtrees (CF 1676G). Parents satisfy a_i < i, so the labels n, n-1, ..., 2
// already list every child before its parent: bal[v] = (#white - #black) in the subtree of v.
int main() {
    ios::sync_with_stdio(false);
    cin.tie(nullptr);
    int t;
    cin >> t;
    while (t--) {
        int n;
        cin >> n;
        vector<int> par(n + 1, 0);
        for (int i = 2; i <= n; i++) cin >> par[i];
        string s;
        cin >> s;
        vector<int> bal(n + 1, 0);
        for (int v = 1; v <= n; v++) bal[v] = (s[v - 1] == 'W') ? 1 : -1;
        int balanced = 0;
        for (int v = n; v >= 1; v--) {             // children (larger labels) are finished first
            if (bal[v] == 0) balanced++;           // the subtree of v is complete: test it
            if (v > 1) bal[par[v]] += bal[v];      // hand the finished value to the parent
        }
        cout << balanced << "\n";
    }
}
\end{lstlisting}

\begin{complexitybox}
\textbf{Time:} \bigO{n} per test, \bigO{\sum n} in total. \textbf{Space:} \bigO{n}. \textbf{Verdict:} reproduces all samples; no recursion, so depth is irrelevant.
\end{complexitybox}
\end{algorithmbox}

\begin{invariantbox}[Why the test happens at the right moment]
\textbf{Invariant.} When the loop reaches $v$, $\mathrm{bal}[v] = \sum_{x \in \mathrm{subtree}(v)} w(x)$.

Every descendant $x$ of $v$ has a label $> v$ (labels grow along every root path because $a_i < i$), so it was processed earlier, and when it was processed it added its own final value to its parent. By induction from the largest label down, each child $c$ of $v$ added exactly $\mathrm{bal}[c] = $ the sum over its subtree, and nothing else was ever added to $\mathrm{bal}[v]$ (only children write into it). The children's subtrees are disjoint and together with $v$ form the subtree of $v$. Hence the value is final and complete when $v$ is reached, and testing it there counts each balanced subtree exactly once. \hfill$\blacksquare$
\end{invariantbox}

\begin{alternativebox}[Approach 3: a real post-order DFS]
Build children lists and let a recursive DFS return the subtree balance. The depth is at most $n \le 4000$ frames --- far below the $\approx 65\,000$ that a 2\,MB stack survives (Lesson~\ref{ssec:ch2-call-stack}), so recursion is safe \emph{here}. This version does not depend on $a_i < i$ and works for any parent array.

\begin{lstlisting}[style=cpp]
#include <bits/stdc++.h>
using namespace std;

// Alternative: a real post-order DFS on the children lists. It returns (#white - #black)
// of the subtree and counts balanced subtrees on the way out. Depth <= n <= 4000: safe.
vector<vector<int>> children;
string s;
int balanced;

int dfs(int u) {
    int bal = (s[u - 1] == 'W') ? 1 : -1;
    for (int c : children[u]) bal += dfs(c);  // every child subtree is finished here
    if (bal == 0) balanced++;
    return bal;
}

int main() {
    ios::sync_with_stdio(false);
    cin.tie(nullptr);
    int t;
    cin >> t;
    while (t--) {
        int n;
        cin >> n;
        children.assign(n + 1, vector<int>());
        for (int i = 2; i <= n; i++) {
            int p;
            cin >> p;
            children[p].push_back(i);
        }
        cin >> s;
        balanced = 0;
        dfs(1);
        cout << balanced << "\n";
    }
}
\end{lstlisting}
\begin{complexitybox}
\textbf{Time:} \bigO{n} per test. \textbf{Space:} \bigO{n} plus a stack of height $\le n$. \textbf{Verdict:} reproduces the samples; agrees with Approaches 1 and 2 on all random tests.
\end{complexitybox}
\end{alternativebox}

\begin{tricksbox}[Approach 4: Implementation Refinements]
\begin{itemize}[leftmargin=*]
    \item \textbf{Reset per test.} With $t = 10^4$ tests, allocate arrays of size $n + 1$ per test (as above), never of size $2\cdot10^5$ per test --- that alone would be $2\cdot10^9$ writes.
    \item \textbf{Two counters work too.} Keeping \texttt{white[v]} and \texttt{black[v]} is equally correct; the signed sum halves the memory and is the form that generalises (``sum of $\pm1$ weights $= 0$'').
    \item \textbf{The root is a subtree too.} The loop includes $v = 1$; sample 2 (\texttt{BW}) has its only balanced subtree at the root.
\end{itemize}
\end{tricksbox}

\subsubsection*{5. Trace}

\begin{tracebox}[{Sample 1: $a = (1,1,2,3,3,5)$, $s = \texttt{WBBWWBW}$, initial $\mathrm{bal} = (+1,-1,-1,+1,+1,-1,+1)$}]
\begin{center}
\begin{tikzpicture}[x=0.85cm, y=0.8cm]
    \tikzset{gv/.append style={minimum size=6.5mm, font=\scriptsize\bfseries}}
    \begin{scope}
        \tikzset{g2p2-s7/.style={g2p2-W, draw=tabulationblue, line width=1.3pt}, g2p2-s6/.style={g2p2-B, draw=tabulationblue, line width=1.3pt},
                 g2p2-s5/.style={g2p2-W, draw=accentorange, line width=1.6pt}}
        \pic {g2p2-S};
        \foreach \v/\b/\d in {1/+1/right, 2/-1/left, 3/-2/right, 4/+1/left, 5/+2/left, 6/-1/right, 7/+1/right} \node[gdist, text=deepblue, \d=0pt of s\v] {$\b$};
        \node[font=\scriptsize\bfseries] at (0,0.9) {1. $v = 7, 6$ done; at $v = 5$};
    \end{scope}
    \begin{scope}[xshift=5.4cm]
        \tikzset{g2p2-s3/.style={g2p2-B, draw=successgreen, line width=1.8pt}, g2p2-s2/.style={g2p2-B, draw=accentorange, line width=1.6pt}}
        \pic {g2p2-S};
        \foreach \v/\b/\d in {1/+1/right, 2/0/left, 3/0/right, 4/+1/left, 5/+2/left, 6/-1/right, 7/+1/right} \node[gdist, text=deepblue, \d=0pt of s\v] {$\b$};
        \node[font=\scriptsize\bfseries] at (0,0.9) {2. after $v = 5, 4$: 3 and 2 at 0};
    \end{scope}
    \begin{scope}[xshift=10.8cm]
        \tikzset{g2p2-s3/.style={g2p2-B, draw=successgreen, line width=1.8pt}, g2p2-s2/.style={g2p2-B, draw=successgreen, line width=1.8pt},
                 g2p2-s1/.style={g2p2-W, draw=accentorange, line width=1.6pt}}
        \pic {g2p2-S};
        \foreach \v/\b/\d in {1/+1/right, 2/0/left, 3/0/right, 4/+1/left, 5/+2/left, 6/-1/right, 7/+1/right} \node[gdist, text=deepblue, \d=0pt of s\v] {$\b$};
        \node[font=\scriptsize\bfseries] at (0,0.9) {3. end: count $= 2$};
    \end{scope}
\end{tikzpicture}
\end{center}
\begin{center}\small
\begin{tabular}{@{}cccl@{}}
\toprule
$v$ & $\mathrm{bal}[v]$ (final) & balanced? & hand up \\
\midrule
7 & $+1$ & no & $\mathrm{bal}[5]$: $+1 \to +2$ \\
6 & $-1$ & no & $\mathrm{bal}[3]$: $-1 \to -2$ \\
5 & $+2$ & no & $\mathrm{bal}[3]$: $-2 \to 0$ \\
4 & $+1$ & no & $\mathrm{bal}[2]$: $-1 \to 0$ \\
3 & $0$ & \textbf{yes} (count 1) & $\mathrm{bal}[1]$: $+1 \to +1$ \\
2 & $0$ & \textbf{yes} (count 2) & $\mathrm{bal}[1]$: $+1 \to +1$ \\
1 & $+1$ & no & --- \\
\bottomrule
\end{tabular}
\end{center}
Printed by a copy of Approach~2 with one trace line per step. Mistake 1 below was also run: the forward loop prints \texttt{1}. Orange ring: the vertex being processed; green ring: a balanced subtree; blue rings: already handed up.
\end{tracebox}

\subsubsection*{6. Test suite}

\begin{testbox}
\centering\small
\begin{tabular}{@{}p{5.2cm}cp{6.4cm}@{}}
\toprule
\textbf{Input ($n$ / parents / $s$)} & \textbf{Output} & \textbf{Why} \\
\midrule
\texttt{2 / 1 / WW} & \texttt{0} & no subtree can balance \\
\texttt{2 / 1 / BW} & \texttt{1} & only the root (sample 2) \\
\texttt{5 / 1 1 1 1 / WBBBB} & \texttt{0} & star: leaves have size 1, root has 1 W vs 4 B \\
\texttt{4 / 1 1 1 / BWBW} & \texttt{1} & star: only the root \\
\texttt{6 / 1 2 3 4 5 / WBWBWB} & \texttt{3} & path: the suffixes of even length \\
50 paths, $n = 4000$, \texttt{WBWB\dots} & \texttt{2000} each & largest input; deepest tree \\
\bottomrule
\end{tabular}
\end{testbox}
All outputs from Approach~2; Approach~1 printed the same. In addition, Approaches 1--3 agreed on $3000$ random inputs (1--3 test cases each, $n \le 12$, random / path / star / caterpillar shapes).

\subsubsection*{7. Common mistakes}

\begin{warningbox}[Mistake 1 --- the forward loop]
Looping $v = 2, \dots, n$ and adding $\mathrm{bal}[v]$ to the parent hands up values \emph{before} their own children have arrived. On sample 1, vertex 5 would hand up $+1$ instead of $+2$, $\mathrm{bal}[3]$ would end at $-1$, and the answer would be 1 instead of 2.
\end{warningbox}

\begin{warningbox}[Mistake 2 --- testing before the subtree is complete]
Counting zeros \emph{while} adding (``if \texttt{bal[par[v]]} became 0, count it'') counts a parent that passes through 0 only temporarily. On \texttt{3 / 1 1 / BBW} the root goes $-1 \to 0$ (vertex 3 arrives) $\to -1$ (vertex 2 arrives): that version prints \texttt{1}, the correct answer is \texttt{0}. Test each $v$ once, at the moment it is reached.
\end{warningbox}

\begin{warningbox}[Mistake 3 --- off-by-one in the string]
Vertices are 1-based, the string is 0-based: the colour of $v$ is \texttt{s[v-1]}.
\end{warningbox}

\subsubsection*{8. Takeaways}

\begin{itemize}[leftmargin=*]
    \item Turn ``two counts equal'' into ``a signed sum is 0'': one aggregated number.
    \item A parent array with $a_i < i$ is an order: one backward loop is a post-order traversal.
    \item Test a subtree's value exactly once, when the vertex is left (reached in the backward loop).
\end{itemize}

\begin{variantsbox}
\variant{Count per label}{\varin{https://leetcode.com/problems/number-of-nodes-in-the-sub-tree-with-the-same-label/}}
  {for each vertex, how many vertices in its subtree have its own label (26 labels).}
  {the same tree, a counter per label instead of one signed number.}
  {a global counter, read at entry and at exit of $v$ (Exercise E2.1).}
\variant{Three colours, all equal}{\varnone}
  {colours R, G, B; a subtree is balanced if the three counts are equal.}
  {two signed sums: $\#R - \#G$ and $\#G - \#B$.}
  {aggregate the pair; balanced iff both are 0. Matches an explicit-count brute force on 5000 random trees.}
\variant{Arbitrary edge list}{\varnone}
  {the tree is given as $n - 1$ undirected edges, root 1.}
  {no label order any more.}
  {template (2): BFS order, backward loop. Matches the label loop on 5000 randomly relabelled trees.}
\variant{Even subtrees to cut}{\varin{https://codeforces.com/problemset/problem/982/C}}
  {cut as many edges as possible so that every piece has even size.}
  {aggregate sizes; an edge above $v$ can be cut iff $\mathrm{sz}[v]$ is even.}
  {the same backward loop with a parity test (Exercise E2.2).}
\end{variantsbox}

%============================================================
% CH2 · P2 · PROBLEM 2 — COUNT NODES WITH THE HIGHEST SCORE (LeetCode 2049)
%============================================================
\subsection{Problem 2.2.2: Count Nodes With the Highest Score}
\problemlabel{2.2.2}{ssec:ch2-highest-score}

% example 1 tree (0-based labels), per-vertex styles g2p2-h<k>
\tikzset{
    g2p2-h0/.style=gv, g2p2-h1/.style=gv, g2p2-h2/.style=gv, g2p2-h3/.style=gv, g2p2-h4/.style=gv,
    g2p2-H/.pic={
        \node[g2p2-h0] (h0) at (0,0) {0};
        \node[g2p2-h2] (h2) at (-1.1,-1) {2}; \node[g2p2-h4] (h4) at (1.1,-1) {4};
        \node[g2p2-h1] (h1) at (-1.7,-2) {1}; \node[g2p2-h3] (h3) at (-0.5,-2) {3};
        \foreach \a/\b in {0/2, 0/4, 2/1, 2/3} \draw[gtree] (h\a) -- (h\b);
    }
}

\problemcard{2.2.2}{Count Nodes With the Highest Score}{LeetCode 2049}{https://leetcode.com/problems/count-nodes-with-the-highest-score/}{\exmedium\ \lc{Medium}\ (the ``part above'' $n - \mathrm{sz}[v]$)}{Subtree sizes by post-order aggregation; the pieces left by deleting $v$ are its child subtrees and the part above of size $n - \mathrm{sz}[v]$; multiply in 64 bits}

\subsubsection*{1. Problem statement}

A \textbf{binary} tree has $n$ nodes $0..n-1$, rooted at 0, given by \texttt{parents} with $\texttt{parents}[0] = -1$. The \textbf{score} of a node: delete the node and its incident edges; the tree falls into one or more non-empty pieces; the score is the product of the piece sizes. Return how many nodes have the highest score.

\textbf{Constraints:} $2 \le n \le 10^5$; \texttt{parents} describes a valid binary tree (each node has at most two children). Note: $\texttt{parents}[i]$ may be larger than $i$.

\begin{table}[H]
\centering\small
\begin{tabular}{@{}lcl@{}}
\toprule
\textbf{parents} & \textbf{Output} & \textbf{Scores of nodes $0, 1, 2, \dots$} \\
\midrule
$[-1,2,0,2,0]$ & \texttt{3} & $3, 4, 2, 4, 4$ \\
$[-1,2,0]$ & \texttt{2} & $2, 2, 1$ \\
\bottomrule
\end{tabular}
\end{table}

\begin{diagrambox}
\begin{tikzpicture}[x=1cm, y=0.95cm]
    \begin{scope}
        \tikzset{g2p2-h1/.style=gvdst, g2p2-h3/.style=gvdst, g2p2-h4/.style=gvdst}
        \pic {g2p2-H};
        \foreach \v/\s/\d in {0/3/right, 1/4/left, 2/2/left, 3/4/right, 4/4/right} \node[gdist, text=deepblue, \d=0pt of h\v] {score \s};
        \node[font=\scriptsize\bfseries] at (0,0.9) {example 1: three nodes score 4};
    \end{scope}
    \begin{scope}[xshift=5.2cm]
        \tikzset{g2p2-h2/.style=gvdead, g2p2-h1/.style=gvA, g2p2-h3/.style=gvB, g2p2-h0/.style=gvsel, g2p2-h4/.style=gvsel}
        \pic {g2p2-H};
        \node[font=\scriptsize\bfseries] at (0,0.9) {delete 2: pieces $1 \cdot 1 \cdot 2 = 2$};
    \end{scope}
    \begin{scope}[xshift=10.4cm]
        \tikzset{g2p2-h0/.style=gvdead, g2p2-h2/.style=gvA, g2p2-h1/.style=gvA, g2p2-h3/.style=gvA, g2p2-h4/.style=gvB}
        \pic {g2p2-H};
        \node[font=\scriptsize\bfseries] at (0,0.9) {delete the root: $3 \cdot 1 = 3$};
    \end{scope}
\end{tikzpicture}
\figcap{Example 1. Each colour is one piece. Deleting a leaf leaves one piece of $n - 1 = 4$ nodes, which here is the maximum.}
\end{diagrambox}

\subsubsection*{2. The graph model}

\begin{modelbox}[Graph Model --- Count Nodes With the Highest Score]
\renewcommand{\arraystretch}{1.35}
\begin{tabularx}{\linewidth}{@{}>{\bfseries}l X@{}}
Vertices & the $n$ nodes; root 0 \\
Edges & tree edges $\texttt{parents}[v] \to v$; at most two children per node; labels give \emph{no} order ($\texttt{parents}[1] = 2$ in example 1) \\
Weights & none; every node counts 1 \\
Question & $\mathrm{score}(v) = \prod_{c \text{ child of } v} \mathrm{sz}[c] \cdot (n - \mathrm{sz}[v])$, the last factor omitted for the root; count the $v$ attaining the maximum \\
Size & $n \le 10^5$: \bigO{n}; a path of $10^5$ nodes is legal, so the depth can be $10^5$; scores reach $3.7\cdot10^{13}$ \\
\end{tabularx}
\end{modelbox}

\begin{diagrambox}
\begin{tikzpicture}[x=1cm, y=1cm]
    \node[gvcur] (v) at (0,0) {$v$};
    \coordinate (p) at (0,1.0);
    \draw[gfaded, dashed] (p) -- (v);
    \draw[draw=softpurple, fill=optimizedbg, rounded corners=8pt] (-2.6,1.0) rectangle (2.6,2.4);
    \node[font=\scriptsize, text width=4.8cm, align=center] at (0,2.05) {everything outside the subtree of $v$: $n - \mathrm{sz}[v]$ nodes};
    \draw[draw=primaryblue, fill=primaryblue!12] (-1.2,-0.6) -- (-2.1,-2.0) -- (-0.3,-2.0) -- cycle;
    \draw[draw=recursionred, fill=dupfill] (1.2,-0.6) -- (0.3,-2.0) -- (2.1,-2.0) -- cycle;
    \draw[gfaded, dashed] (v) -- (-1.2,-0.6); \draw[gfaded, dashed] (v) -- (1.2,-0.6);
    \node[font=\scriptsize] at (-1.2,-1.55) {$\mathrm{sz}[\ell]$};
    \node[font=\scriptsize] at (1.2,-1.55) {$\mathrm{sz}[r]$};
    \node[note, text width=5.6cm, anchor=west, align=left] at (3.2,0) {$\mathrm{score}(v) = \mathrm{sz}[\ell] \cdot \mathrm{sz}[r] \cdot (n - \mathrm{sz}[v])$\\[2pt]missing child: factor omitted (not 0)\\root: no part above (factor omitted, not 0)\\all three sizes add up to $n - 1$};
\end{tikzpicture}
\figcap{The model: at most three pieces, all of them readable from the subtree sizes.}
\end{diagrambox}

\subsubsection*{3. Key observations}

\begin{itemize}[leftmargin=*]
    \item \textbf{The pieces are known from sizes alone.} The children's pieces are their subtrees; the remaining piece is everything else, $n - \mathrm{sz}[v]$ nodes. One post-order pass computes all $\mathrm{sz}$, then every score is \bigO{1}.
    \item \textbf{Empty pieces do not count.} ``One or more non-empty subtrees'': a missing child and the empty part above the root contribute nothing --- multiplying by $0$ would make the score 0.
    \item \textbf{64 bits are needed.} The three sizes add up to $n - 1$, so the product is at most $\left(\frac{n-1}{3}\right)^3 \approx 3.7\cdot10^{13}$ for $n = 10^5$ (AM--GM), above $2^{31} - 1 \approx 2.1\cdot10^9$ but far below $2^{63}$. The bound is attained: a chain of $33\,333$ nodes from the root ending at $m$, and two chains of $33\,333$ below $m$, give $\mathrm{score}(m) = 33\,333^3 = 37\,035\,925\,937\,037$. The smallest tree that overflows \texttt{int} has $n = 3872$.
    \item \textbf{The labels are no order.} $\texttt{parents}[i] > i$ is allowed, so a backward loop over the labels is wrong; build children lists and a BFS order.
\end{itemize}

\subsubsection*{4. Approaches}

\begin{bruteforcebox}
\textbf{Idea: really delete each node.} For every $v$, block it and run a BFS from every remaining unvisited node; multiply the component sizes.

\begin{lstlisting}[style=cpp]
#include <bits/stdc++.h>
using namespace std;
typedef long long ll;

class SolutionBrute {
public:
    int countHighestScoreNodes(vector<int>& parents) {
        int n = parents.size();
        vector<vector<int>> adj(n);
        for (int v = 1; v < n; v++) { adj[v].push_back(parents[v]); adj[parents[v]].push_back(v); }
        vector<ll> score(n, 1);
        for (int del = 0; del < n; del++) {
            vector<char> seen(n, 0);
            seen[del] = 1;                              // the deleted node blocks the search
            for (int s = 0; s < n; s++) {
                if (seen[s]) continue;
                ll piece = 0;
                vector<int> q(1, s);
                seen[s] = 1;
                for (size_t i = 0; i < q.size(); i++) {
                    piece++;
                    for (int w : adj[q[i]]) if (!seen[w]) { seen[w] = 1; q.push_back(w); }
                }
                score[del] *= piece;
            }
        }
        ll best = *max_element(score.begin(), score.end());
        return count(score.begin(), score.end(), best);
    }
};
\end{lstlisting}

\begin{complexitybox}
\textbf{Time:} $n$ searches of \bigO{n}: $\Theta(n^2) = 10^{10}$ for $n = 10^5$. \textbf{Space:} \bigO{n}. \textbf{Verdict:} \tle; it is the reference: it does not know anything about sizes or the part above.
\end{complexitybox}
\end{bruteforcebox}

\begin{algorithmbox}[Approach 2: BFS order, sizes backwards, scores from sizes]
\textbf{Idea.} Children lists; a BFS order from 0 (parents first); sizes by the backward loop (template (2) of the pattern); then $\mathrm{score}(v)$ = product of the children's sizes, times $n - \mathrm{sz}[v]$ unless $v = 0$.

\begin{lstlisting}[style=cpp]
#include <bits/stdc++.h>
using namespace std;
typedef long long ll;

class Solution {
public:
    int countHighestScoreNodes(vector<int>& parents) {
        int n = parents.size();
        vector<vector<int>> children(n);
        for (int v = 1; v < n; v++) children[parents[v]].push_back(v);
        vector<int> order(1, 0);                       // BFS order: parents before children
        for (size_t i = 0; i < order.size(); i++)
            for (int c : children[order[i]]) order.push_back(c);
        vector<ll> sz(n, 1);
        for (int i = n - 1; i >= 1; i--) sz[parents[order[i]]] += sz[order[i]];
        ll best = -1;
        int count = 0;
        for (int v = 0; v < n; v++) {
            ll score = 1;
            for (int c : children[v]) score *= sz[c];  // the parts below v
            if (v != 0) score *= n - sz[v];            // the part above v
            if (score > best) { best = score; count = 1; }
            else if (score == best) count++;
        }
        return count;
    }
};
\end{lstlisting}

\begin{complexitybox}
\textbf{Time:} \bigO{n}. \textbf{Space:} \bigO{n}, no recursion --- the path of $10^5$ nodes is handled (answer 2: the two ends score $n - 1$). \textbf{Verdict:} returns 3 and 2 on the two examples.
\end{complexitybox}
\end{algorithmbox}

\begin{invariantbox}[Why the scores are right]
\textbf{Sizes.} In BFS order every node appears after its parent, so in the backward loop every node is handled after all its children: $\mathrm{sz}[v]$ is final when $v$ hands it up (the pattern's invariant).

\textbf{Pieces.} Removing $v$ removes the edges to its children and to its parent. Each child $c$ keeps its whole subtree, which no longer touches anything else: one piece of $\mathrm{sz}[c]$ nodes. All other nodes --- the nodes outside the subtree of $v$ --- are still connected to each other, because the path between any two of them in the tree does not pass through $v$ (it would have to enter and leave the subtree of $v$ through the single edge above $v$). That is one piece of $n - \mathrm{sz}[v]$ nodes if this number is positive, i.e.\ if $v$ is not the root. These are all the pieces, so the product is the score. \hfill$\blacksquare$
\end{invariantbox}

\begin{alternativebox}[Approach 3: peel the leaves]
Another order with ``children first'': keep, for every node, the number of children whose size is still unknown; start a queue with the leaves; when a node is taken from the queue its size is final, it multiplies its parent's score by its size, and the parent joins the queue when its last child is done. This is Kahn's algorithm (Chapter~5) on the tree directed towards the root, and it needs no BFS from the root.

\begin{lstlisting}[style=cpp]
#include <bits/stdc++.h>
using namespace std;
typedef long long ll;

class SolutionPeel {
public:
    int countHighestScoreNodes(vector<int>& parents) {
        int n = parents.size();
        vector<int> pending(n, 0);                     // children whose size is not yet known
        for (int v = 1; v < n; v++) pending[parents[v]]++;
        vector<ll> sz(n, 1), score(n, 1);
        vector<int> q;
        for (int v = 0; v < n; v++) if (pending[v] == 0) q.push_back(v);
        for (size_t i = 0; i < q.size(); i++) {
            int v = q[i];                              // all children of v are finished
            if (v == 0) continue;
            score[v] *= n - sz[v];
            int p = parents[v];
            sz[p] += sz[v];
            score[p] *= sz[v];
            if (--pending[p] == 0) q.push_back(p);
        }
        ll best = *max_element(score.begin(), score.end());
        return count(score.begin(), score.end(), best);
    }
};
\end{lstlisting}
\begin{complexitybox}
\textbf{Time:} \bigO{n}. \textbf{Space:} \bigO{n}. \textbf{Verdict:} returns 3 and 2 on the examples; agrees with Approaches 1 and 2 on all random trees. The queue order on example 1 is $1, 3, 4, 2, 0$ --- a ``reverse BFS'' from the leaves.
\end{complexitybox}
\end{alternativebox}

\begin{tricksbox}[Approach 4: Implementation Refinements]
\begin{itemize}[leftmargin=*]
    \item \textbf{Recursive DFS.} \texttt{sz[v] = 1 + dfs(left) + dfs(right)} is the shortest code, but a path of $10^5$ nodes recurses $10^5$ deep; LeetCode's stack is not documented, so the BFS order is the safe choice.
    \item \textbf{Count while computing.} Keep \texttt{best} and \texttt{count} in one pass; ties are compared exactly because the scores are integers (never use \texttt{double} for products of this size).
    \item \textbf{Children lists are optional.} For a binary tree, two arrays \texttt{left}, \texttt{right} with $-1$ for ``none'' are enough.
\end{itemize}
\end{tricksbox}

\subsubsection*{5. Trace}

\begin{tracebox}[{Example 1: $\texttt{parents} = [-1,2,0,2,0]$, $n = 5$}]
\begin{center}
\begin{tikzpicture}[x=1.05cm]
    \node[font=\scriptsize\bfseries, anchor=east] at (-0.6,0) {order};
    \foreach \v [count=\i from 0] in {0,2,4,1,3} \node[cell, minimum width=9mm] (o\i) at (\i,0) {\v};
    \node[font=\scriptsize\bfseries, anchor=east] at (-0.6,-0.75) {parent};
    \foreach \p [count=\i from 0] in {--,0,0,2,2} \node[idx] at (\i,-0.75) {\p};
    \draw[->, thick, primaryblue] (4.3,-1.25) -- node[below, font=\scriptsize]{backward: $\mathrm{sz}[\mathrm{par}] \mathrel{+}= \mathrm{sz}[v]$} (-0.3,-1.25);
    \node[note, anchor=west, text width=3.8cm, align=left] at (5.0,-0.4) {BFS from 0 with \texttt{children[0] = 2 4}, \texttt{children[2] = 1 3}};
\end{tikzpicture}
\end{center}
\begin{center}
\begin{tikzpicture}[x=0.95cm, y=0.85cm]
    \tikzset{gv/.append style={minimum size=6.5mm, font=\scriptsize\bfseries}}
    \begin{scope}
        \tikzset{g2p2-h3/.style=gvdone, g2p2-h1/.style=gvdone, g2p2-h4/.style=gvdone}
        \pic {g2p2-H};
        \foreach \v/\s/\d in {0/1/right, 1/1/left, 2/1/left, 3/1/right, 4/1/right} \node[gdist, text=primaryblue, \d=0pt of h\v] {\s};
        \node[font=\scriptsize\bfseries] at (0,0.85) {1. leaves 3, 1, 4: final 1};
    \end{scope}
    \begin{scope}[xshift=4.6cm]
        \tikzset{g2p2-h3/.style=gvdone, g2p2-h1/.style=gvdone, g2p2-h4/.style=gvdone, g2p2-h2/.style=gvcur}
        \pic {g2p2-H};
        \foreach \v/\s/\d in {0/2/right, 1/1/left, 2/3/left, 3/1/right, 4/1/right} \node[gdist, text=primaryblue, \d=0pt of h\v] {\s};
        \node[font=\scriptsize\bfseries] at (0,0.85) {2. node 2: $1 + 1 + 1 = 3$};
    \end{scope}
    \begin{scope}[xshift=9.2cm]
        \tikzset{g2p2-h3/.style=gvdone, g2p2-h1/.style=gvdone, g2p2-h4/.style=gvdone, g2p2-h2/.style=gvdone, g2p2-h0/.style=gvdone}
        \pic {g2p2-H};
        \foreach \v/\s/\d in {0/5/right, 1/1/left, 2/3/left, 3/1/right, 4/1/right} \node[gdist, text=primaryblue, \d=0pt of h\v] {\s};
        \node[font=\scriptsize\bfseries] at (0,0.85) {3. root: $\mathrm{sz}[0] = 5$};
    \end{scope}
\end{tikzpicture}
\end{center}
\begin{center}\small
\begin{tabular}{@{}cccccl@{}}
\toprule
$v$ & children (sizes) & $\mathrm{sz}[v]$ & $n - \mathrm{sz}[v]$ & score & best / count \\
\midrule
0 & 2 (3), 4 (1) & 5 & --- (root) & $3 \cdot 1 = 3$ & 3 / 1 \\
1 & --- & 1 & 4 & $4$ & 4 / 1 \\
2 & 1 (1), 3 (1) & 3 & 2 & $1 \cdot 1 \cdot 2 = 2$ & 4 / 1 \\
3 & --- & 1 & 4 & $4$ & 4 / 2 \\
4 & --- & 1 & 4 & $4$ & 4 / 3 \\
\bottomrule
\end{tabular}
\end{center}
Return \textbf{3}. The backward loop hands up in the order $3, 1, 4, 2$; the sizes after each step were $\mathrm{sz} = (1,1,2,1,1)$, $(1,1,3,1,1)$, $(2,1,3,1,1)$, $(5,1,3,1,1)$. In snapshot 2 the root shows the partial size 2 (node 4 has arrived, node 2 not yet).
\end{tracebox}

\subsubsection*{6. Test suite}

\begin{testbox}
\centering\small
\begin{tabular}{@{}p{5.6cm}cp{6.0cm}@{}}
\toprule
\textbf{parents} & \textbf{Output} & \textbf{Why} \\
\midrule
$[-1,0]$ & \texttt{2} & both nodes score 1 \\
$[-1,0,0]$ & \texttt{2} & root $1 \cdot 1 = 1$, leaves $2$ \\
$[-1,0,1,2,3]$ (path) & \texttt{3} & scores $4,3,4,3,4$: ends and middle \\
$[-1,0,0,1,1,2,2]$ (complete) & \texttt{1} & root $3 \cdot 3 = 9$ beats leaves (6) \\
path, $n = 10^5$ & \texttt{2} & depth $10^5$: the two ends score $99\,999$ \\
three arms of $33\,333$, $n = 10^5$ & \texttt{1} & $\mathrm{score}(m) = 37\,035\,925\,937\,037$; with \texttt{int} the program returns \texttt{3} \\
\bottomrule
\end{tabular}
\end{testbox}
All outputs from Approach~2 (Approach~3 and the brute force agree on the small ones). In addition, the three approaches agreed on $5000$ random binary trees with $n \le 13$ and shuffled labels (so that parents can have larger labels than children).

\subsubsection*{7. Common mistakes}

\begin{warningbox}[Mistake 1 --- \texttt{int} scores]
Replacing \texttt{long long} by \texttt{int} still passes both examples and every small test --- and returns \texttt{3} instead of \texttt{1} on the three-arm tree, where the products wrap around. The first tree where it can happen has only $3872$ nodes.
\end{warningbox}

\begin{warningbox}[Mistake 2 --- multiplying by an empty piece]
\begin{lstlisting}[style=snippet]
ll score = sz[l] * sz[r] * (n - sz[v]);     // WRONG: an absent child or the root gives 0
ll score = 1;                               // RIGHT: multiply only the non-empty pieces
for (int c : children[v]) score *= sz[c];
if (v != 0) score *= n - sz[v];
\end{lstlisting}
\end{warningbox}

\begin{warningbox}[Mistake 3 --- looping over the labels]
\texttt{for (v = n-1; v >= 1; v--) sz[parents[v]] += sz[v];} is template (3) of the pattern, valid only when $\texttt{parents}[v] < v$. In example 1, node 1 hands its size to node 2 \emph{after} node 2 has already handed its own size to the root, so the root ends with $\mathrm{sz}[0] = 4$ instead of 5 (the answer survives by luck). On the path $[-1,2,3,0]$ ($0 - 3 - 2 - 1$) the wrong sizes change the answer: \texttt{1} instead of \texttt{2}. Among $5000$ random small binary trees with shuffled labels, the label loop gave a wrong answer on $881$.
\end{warningbox}

\subsubsection*{8. Takeaways}

\begin{itemize}[leftmargin=*]
    \item Deleting $v$ leaves the child subtrees and one piece of $n - \mathrm{sz}[v]$: one size pass answers every deletion.
    \item Bound products before choosing the type: three parts summing to $n - 1$ give at most $((n-1)/3)^3$.
    \item When the labels are not an order, make one: BFS from the root, or peel leaves.
\end{itemize}

\begin{variantsbox}
\variant{Smallest largest piece}{\varnone}
  {find the node whose deletion leaves the smallest largest piece (a centroid).}
  {the same pieces; maximum instead of product.}
  {template (4) of the pattern: $\max(\max_c \mathrm{sz}[c],\ n - \mathrm{sz}[v])$; on the tree $T$ it picks vertex 1 (largest piece 3).}
\variant{Any tree}{\varnone}
  {nodes may have any number of children.}
  {up to $\deg(v)$ pieces; the product can exceed 64 bits.}
  {the same loop; compare $\sum \log$ of the piece sizes (with a tolerance) instead of products. The log comparison picked the same count as exact products on $3000$ random trees with $n \le 40$.}
\variant{Number of separated pairs}{\varnone}
  {for every node, how many pairs of other nodes lose their connection when it is deleted.}
  {pairs taken from two different pieces.}
  {$\binom{n-1}{2} - \sum_{\text{pieces}} \binom{s}{2}$, all pieces from sizes; matched a brute force (delete, BFS, count pairs) on the same $3000$ random trees.}
\variant{Fuel through an edge}{\varin{https://leetcode.com/problems/minimum-fuel-cost-to-report-to-the-capital/}}
  {everybody drives to the root in cars of $s$ seats.}
  {the edge above $v$ carries $\mathrm{sz}[v]$ people.}
  {sum $\lceil \mathrm{sz}[v] / s \rceil$ over all edges (Exercise E2.3).}
\end{variantsbox}

%============================================================
% CH2 · P2 · PROBLEM 3 — KI (AtCoder ABC138 D)
%============================================================
\subsection{Problem 2.2.3: Ki}
\problemlabel{2.2.3}{ssec:ch2-ki}

% sample 1 tree, per-vertex styles g2p2-k<k>
\tikzset{
    g2p2-k1/.style=gv, g2p2-k2/.style=gv, g2p2-k3/.style=gv, g2p2-k4/.style=gv,
    g2p2-K/.pic={
        \node[g2p2-k1] (k1) at (0,0) {1};
        \node[g2p2-k2] (k2) at (0,-1.1) {2};
        \node[g2p2-k3] (k3) at (-0.9,-2.2) {3}; \node[g2p2-k4] (k4) at (0.9,-2.2) {4};
        \foreach \a/\b in {1/2, 2/3, 2/4} \draw[gtree] (k\a) -- (k\b);
    }
}

\problemcard{2.2.3}{Ki}{AtCoder ABC138 D}{https://atcoder.jp/contests/abc138/tasks/abc138_d}{\exeasy\ (AtCoder ABC D, 400 points; the push-down direction)}{Record each operation only at the top vertex $p$; one pre-order pass pushes the totals down: $\mathrm{cnt}[v] \mathrel{+}= \mathrm{cnt}[\mathrm{par}(v)]$; iterative because $N = 2\cdot10^5$}

\subsubsection*{1. Problem statement}

A tree with $N$ vertices is rooted at vertex 1; the $i$-th edge joins $a_i$ and $b_i$. Every vertex has a counter, initially 0. Operation $j$ adds $x_j$ to the counter of \textbf{every vertex in the subtree of $p_j$}. Print all counters after the $Q$ operations.

\textbf{Constraints:} $2 \le N \le 2\cdot10^5$, $1 \le Q \le 2\cdot10^5$, $1 \le a_i < b_i \le N$, $1 \le p_j \le N$, $1 \le x_j \le 10^4$.

\begin{table}[H]
\centering\small
\begin{tabular}{@{}lll@{}}
\toprule
\textbf{Input} & \textbf{Official output} & \textbf{Our output} \\
\midrule
\texttt{4 3 / 1 2 / 2 3 / 2 4 / 2 10 / 1 100 / 3 1} & \texttt{100 110 111 110} & \texttt{100 110 111 110} \\
\texttt{6 2 / 1 2 / 1 3 / 2 4 / 3 6 / 2 5 / 1 10 / 1 10} & \texttt{20 20 20 20 20 20} & \texttt{20 20 20 20 20 20} \\
\bottomrule
\end{tabular}
\end{table}

\begin{diagrambox}
\begin{tikzpicture}[x=1cm, y=0.95cm]
    \begin{scope}
        \pic {g2p2-K};
        \node[font=\scriptsize\bfseries] at (0,0.85) {sample 1};
        \node[note, anchor=west, text width=3.9cm, align=left] at (1.4,-1.1) {op 1: $+10$ on $\{2,3,4\}$\\op 2: $+100$ on all\\op 3: $+1$ on $\{3\}$};
    \end{scope}
    \begin{scope}[xshift=6.4cm]
        \tikzset{g2p2-k1/.style=gvdone, g2p2-k2/.style=gvdone, g2p2-k3/.style=gvdst, g2p2-k4/.style=gvdone}
        \pic {g2p2-K};
        \foreach \v/\c in {1/100, 2/110, 3/111, 4/110} \node[gdist, text=deepblue, right=1pt of k\v] {\c};
        \node[font=\scriptsize\bfseries] at (0,0.85) {final counters};
    \end{scope}
    \begin{scope}[xshift=11.2cm, y=0.95cm]
        \node[gv] (m1) at (0,0) {1};
        \node[gv] (m2) at (-0.8,-1.1) {2}; \node[gv] (m3) at (0.8,-1.1) {3};
        \node[gv] (m4) at (-1.3,-2.2) {4}; \node[gv] (m5) at (-0.3,-2.2) {5}; \node[gv] (m6) at (0.8,-2.2) {6};
        \foreach \a/\b in {1/2, 1/3, 2/4, 2/5, 3/6} \draw[gtree] (m\a) -- (m\b);
        \node[font=\scriptsize\bfseries] at (0,0.85) {sample 2: $2 \times (+10)$ at 1};
    \end{scope}
\end{tikzpicture}
\figcap{The samples. In sample 2 both operations hit the root, so every counter ends at 20.}
\end{diagrambox}

\subsubsection*{2. The graph model}

\begin{modelbox}[Graph Model --- Ki]
\renewcommand{\arraystretch}{1.35}
\begin{tabularx}{\linewidth}{@{}>{\bfseries}l X@{}}
Vertices & the $N$ vertices, root 1; $\mathrm{add}[p] = \sum_{j : p_j = p} x_j$, the total put on $p$ itself \\
Edges & $N - 1$ undirected edges; $a_i < b_i$ is only a normalisation of the pair, \emph{not} ``$a_i$ is the parent'' \\
Weights & the vertex weights $\mathrm{add}[v]$ \\
Question & $\mathrm{cnt}[v] = \sum_{u \text{ ancestor of } v \text{ (or } v)} \mathrm{add}[u]$ --- a sum along the \textbf{root path}, i.e.\ push-down: $\mathrm{cnt}[v] = \mathrm{add}[v] + \mathrm{cnt}[\mathrm{par}(v)]$ \\
Size & $N, Q \le 2\cdot10^5$: \bigO{N + Q}; a path is legal, so the depth is up to $2\cdot10^5$; counters up to $Q \cdot \max x = 2\cdot10^9$ \\
\end{tabularx}
\end{modelbox}

\begin{diagrambox}
\begin{tikzpicture}[x=1cm, y=0.95cm]
    \begin{scope}
        \tikzset{g2p2-k1/.style=gvsel, g2p2-k2/.style=gvsel, g2p2-k3/.style=gvsel}
        \pic {g2p2-K};
        \node[gdist, text=softpurple, left=1pt of k1] {$+100$};
        \node[gdist, text=softpurple, left=1pt of k2] {$+10$};
        \node[gdist, text=softpurple, left=1pt of k3] {$+1$};
        \node[font=\scriptsize\bfseries] at (0,0.85) {step 1: mark only the tops};
    \end{scope}
    \draw[flow] (2.0,-1.1) -- (3.6,-1.1);
    \begin{scope}[xshift=5.6cm]
        \pic {g2p2-K};
        \draw[->, thick, accentorange] (k1) to[bend left=35] node[lbl, right, text=accentorange]{$+100$} (k2);
        \draw[->, thick, accentorange] (k2) to[bend right=30] node[lbl, left, text=accentorange]{$+110$} (k3);
        \draw[->, thick, accentorange] (k2) to[bend left=30] node[lbl, right, text=accentorange]{$+110$} (k4);
        \node[font=\scriptsize\bfseries] at (0,0.85) {step 2: push down once};
    \end{scope}
    \node[note, anchor=west, text width=4.2cm, align=left] at (8.3,-1.1) {a vertex $v$ is in the subtree of $p$ $\iff$ $p$ is on the root path of $v$; so $v$ receives exactly the marks on its root path};
\end{tikzpicture}
\figcap{The model: ``add to a subtree'' turned around becomes ``sum over the root path'', which is what a push-down computes.}
\end{diagrambox}

\subsubsection*{3. Key observations}

\begin{itemize}[leftmargin=*]
    \item \textbf{Turn the operation around.} $v \in \mathrm{subtree}(p) \iff p \text{ is an ancestor of } v$ (or $v$ itself). So the final counter of $v$ is the sum of the marks $\mathrm{add}[u]$ over its root path --- the pre-order direction of the pattern.
    \item \textbf{Operations commute.} All operations are additions and nothing is asked in between, so we may postpone all the work: record \bigO{1} per operation, then one pass.
    \item \textbf{Root it yourself.} $a_i < b_i$ does not say who is the parent: with edges $1{-}3$, $2{-}3$ vertex 3 is the parent of 2. A program that sets $\mathrm{par}[b_i] = a_i$ prints \texttt{0 0 1} instead of \texttt{0 1 1} for the operation ``$+1$ on 3''.
    \item \textbf{Depth.} A path of $2\cdot10^5$ vertices is a legal tree; our recursive push-down died on it without output. BFS order makes the pass iterative.
\end{itemize}

\subsubsection*{4. Approaches}

\begin{bruteforcebox}
\textbf{Idea: do what the statement says.} Root the tree once, then for each operation walk the subtree of $p_j$ (downwards only) and add $x_j$.

\begin{lstlisting}[style=cpp]
#include <bits/stdc++.h>
using namespace std;
typedef long long ll;

// Brute force: perform every operation literally. Root the tree once (parent array), then for each
// operation walk the whole subtree of p (only downwards: never back to the parent) and add x.
int main() {
    int n, q;
    cin >> n >> q;
    vector<vector<int>> adj(n + 1);
    for (int i = 0; i < n - 1; i++) {
        int a, b;
        cin >> a >> b;
        adj[a].push_back(b);
        adj[b].push_back(a);
    }
    vector<int> parent(n + 1, 0), order(1, 1);
    for (size_t i = 0; i < order.size(); i++)
        for (int v : adj[order[i]])
            if (v != parent[order[i]]) { parent[v] = order[i]; order.push_back(v); }
    vector<ll> cnt(n + 1, 0);
    for (int j = 0; j < q; j++) {
        int p, x;
        cin >> p >> x;
        vector<int> st(1, p);
        while (!st.empty()) {
            int u = st.back(); st.pop_back();
            cnt[u] += x;
            for (int v : adj[u]) if (v != parent[u]) st.push_back(v);
        }
    }
    for (int v = 1; v <= n; v++) cout << cnt[v] << (v == n ? '\n' : ' ');
}
\end{lstlisting}

\begin{complexitybox}
\textbf{Time:} \bigO{\sum_j \mathrm{sz}[p_j]} $= \bigO{NQ}$ $= 4\cdot10^{10}$ when every operation hits the root. On a path of $2\cdot10^5$ vertices with $2\cdot10^5$ operations at the root it had not finished after three minutes. \textbf{Space:} \bigO{N}. \textbf{Verdict:} \tle; reference for testing.
\end{complexitybox}
\end{bruteforcebox}

\begin{algorithmbox}[Approach 2: mark the top, push down in BFS order]
\textbf{Idea.} $\mathrm{cnt}[p_j] \mathrel{+}= x_j$ for each operation. Then a BFS from 1 that skips the parent: when it discovers $v$ from $u$, $\mathrm{cnt}[u]$ is final, and $\mathrm{cnt}[v] \mathrel{+}= \mathrm{cnt}[u]$. This is the forward loop of template (2), fused with the BFS itself.

\begin{lstlisting}[style=cpp]
#include <bits/stdc++.h>
using namespace std;
typedef long long ll;

// Ki (ABC138 D): put x on the counter of p only, then push the counters DOWN the tree:
// final[v] = add[v] + final[parent(v)], in an order where every parent comes first (BFS order).
int main() {
    ios::sync_with_stdio(false);
    cin.tie(nullptr);
    int n, q;
    cin >> n >> q;
    vector<vector<int>> adj(n + 1);
    for (int i = 0; i < n - 1; i++) {
        int a, b;
        cin >> a >> b;
        adj[a].push_back(b);
        adj[b].push_back(a);
    }
    vector<ll> cnt(n + 1, 0);
    for (int j = 0; j < q; j++) {
        int p, x;
        cin >> p >> x;
        cnt[p] += x;                               // only the top of the subtree, O(1)
    }
    vector<int> parent(n + 1, 0), order(1, 1);
    parent[1] = -1;
    for (size_t i = 0; i < order.size(); i++) {    // BFS from the root, no recursion
        int u = order[i];
        for (int v : adj[u])
            if (v != parent[u]) {
                parent[v] = u;
                cnt[v] += cnt[u];                  // pre-order push-down: u is already final
                order.push_back(v);
            }
    }
    for (int v = 1; v <= n; v++) cout << cnt[v] << (v == n ? '\n' : ' ');
}
\end{lstlisting}

\begin{complexitybox}
\textbf{Time:} \bigO{N + Q}. \textbf{Space:} \bigO{N}. \textbf{Verdict:} reproduces both samples; on the path of $2\cdot10^5$ vertices with all operations at the root it prints $2\cdot10^9$ everywhere in $0.9$\,s wall time here, start-up included.
\end{complexitybox}
\end{algorithmbox}

\begin{invariantbox}[Why one push-down is enough]
\textbf{Claim.} When $v$ is discovered from $u = \mathrm{par}(v)$, $\mathrm{cnt}[u] = \sum_{w \in \mathrm{path}(1, u)} \mathrm{add}[w]$, and after the update $\mathrm{cnt}[v] = \sum_{w \in \mathrm{path}(1, v)} \mathrm{add}[w]$.

\emph{Induction along the BFS order.} The root is never updated: $\mathrm{cnt}[1] = \mathrm{add}[1]$, its root path is $\{1\}$. A vertex $u$ is taken from the queue only after it was discovered (so its own update is done), and $\mathrm{cnt}[u]$ never changes afterwards: only its parent writes into it, once. So $\mathrm{cnt}[v] = \mathrm{add}[v] + \mathrm{cnt}[u]$ is the sum over $\mathrm{path}(1, u) \cup \{v\} = \mathrm{path}(1, v)$.

\emph{The answer.} $\mathrm{path}(1, v)$ is the set of $p$ whose subtree contains $v$, so $\mathrm{cnt}[v] = \sum_{j : v \in \mathrm{subtree}(p_j)} x_j$, the required counter. \hfill$\blacksquare$
\end{invariantbox}

\begin{alternativebox}[Approach 3: a difference array on pre-order intervals]
Number the vertices in DFS pre-order: $\mathrm{pos}[v]$. The subtree of $p$ is then the block of positions $\mathrm{pos}[p], \dots, \mathrm{pos}[p] + \mathrm{sz}[p] - 1$ (this is the parenthesis theorem of Lesson~\ref{ssec:ch2-timestamps}, the core of Pattern~3). ``Add $x$ to a block'' on an array is a difference update: $+x$ at the start, $-x$ just after the end; one prefix sum at the end gives every value. This version is the one that generalises when queries are mixed with updates: replace the difference array by a Fenwick tree (Chapter~15).

\begin{lstlisting}[style=cpp]
#include <bits/stdc++.h>
using namespace std;
typedef long long ll;

// Alternative: number the vertices in DFS pre-order. The subtree of p is the block of positions
// pos[p] .. pos[p] + sz[p] - 1, so "add x to the subtree" becomes a range update on an array:
// diff[pos[p]] += x, diff[pos[p] + sz[p]] -= x, and one prefix sum at the end.
int main() {
    ios::sync_with_stdio(false);
    cin.tie(nullptr);
    int n, q;
    cin >> n >> q;
    vector<vector<int>> adj(n + 1);
    for (int i = 0; i < n - 1; i++) {
        int a, b;
        cin >> a >> b;
        adj[a].push_back(b);
        adj[b].push_back(a);
    }
    // iterative DFS: pre-order positions, then sizes from the reversed pre-order
    vector<int> pos(n + 1), parent(n + 1, 0), pre, st(1, 1);
    while (!st.empty()) {
        int u = st.back(); st.pop_back();
        pos[u] = pre.size();
        pre.push_back(u);
        for (int v : adj[u])
            if (v != parent[u]) { parent[v] = u; st.push_back(v); }
    }
    vector<int> sz(n + 1, 1);
    for (int i = n - 1; i >= 1; i--) sz[parent[pre[i]]] += sz[pre[i]];
    vector<ll> diff(n + 1, 0);
    for (int j = 0; j < q; j++) {
        int p, x;
        cin >> p >> x;
        diff[pos[p]] += x;                         // the block of p starts here ...
        diff[pos[p] + sz[p]] -= x;                 // ... and ends just before here
    }
    for (int i = 1; i < n; i++) diff[i] += diff[i - 1];
    for (int v = 1; v <= n; v++) cout << diff[pos[v]] << (v == n ? '\n' : ' ');
}
\end{lstlisting}
The stack here pushes all children of the popped vertex at once. In a general graph that is the ``shortcut that is not DFS'' of Lesson~\ref{ssec:ch2-call-stack}; in a tree every vertex is pushed exactly once (by its parent), and the last-pushed child's whole subtree is popped before its siblings, so the subtrees still occupy contiguous blocks --- a genuine pre-order, with the children in reverse order.
\begin{complexitybox}
\textbf{Time:} \bigO{N + Q}. \textbf{Space:} \bigO{N}. \textbf{Verdict:} reproduces both samples and prints the same output as Approach~2 on the large tests (identical checksums).
\end{complexitybox}
\end{alternativebox}

\begin{tricksbox}[Approach 4: Implementation Refinements]
\begin{itemize}[leftmargin=*]
    \item \textbf{\texttt{int} or \texttt{long long}?} The largest counter is $Q \cdot \max x = 2\cdot10^5 \cdot 10^4 = 2\cdot10^9 < 2^{31} - 1 = 2\,147\,483\,647$: \texttt{int} fits, with 7\,\% to spare. One more zero in any limit breaks it; \texttt{long long} costs nothing here.
    \item \textbf{Fast I/O.} $8\cdot10^5$ numbers in, $2\cdot10^5$ out: \texttt{sync\_with\_stdio(false)} and one \texttt{'\textbackslash n'} at the end.
    \item \textbf{Fusing passes.} The push-down can be done inside the BFS (as above) because a vertex is final at the moment it is dequeued. Aggregation (the up direction) cannot be fused this way: it needs the backward loop after the BFS.
\end{itemize}
\end{tricksbox}

\subsubsection*{5. Trace}

\begin{tracebox}[{Sample 1: edges $1{-}2$, $2{-}3$, $2{-}4$; operations $(2, 10)$, $(1, 100)$, $(3, 1)$}]
\begin{center}
\begin{tikzpicture}[x=0.95cm, y=0.85cm]
    \tikzset{gv/.append style={minimum size=6.5mm, font=\scriptsize\bfseries}}
    \begin{scope}
        \pic {g2p2-K};
        \foreach \v/\c in {1/100, 2/10, 3/1, 4/0} \node[gdist, text=softpurple, right=1pt of k\v] {\c};
        \node[font=\scriptsize\bfseries] at (0,0.8) {1. marks after reading};
    \end{scope}
    \begin{scope}[xshift=4.2cm]
        \tikzset{g2p2-k1/.style=gvdone, g2p2-k2/.style=gvcur}
        \pic {g2p2-K};
        \foreach \v/\c in {1/100, 2/110, 3/1, 4/0} \node[gdist, text=deepblue, right=1pt of k\v] {\c};
        \node[font=\scriptsize\bfseries] at (0,0.8) {2. from 1: $\mathrm{cnt}[2] = 10 + 100$};
    \end{scope}
    \begin{scope}[xshift=8.4cm]
        \tikzset{g2p2-k1/.style=gvdone, g2p2-k2/.style=gvdone, g2p2-k3/.style=gvcur, g2p2-k4/.style=gvcur}
        \pic {g2p2-K};
        \foreach \v/\c in {1/100, 2/110, 3/111, 4/110} \node[gdist, text=deepblue, right=1pt of k\v] {\c};
        \node[font=\scriptsize\bfseries] at (0,0.8) {3. from 2: 3 and 4};
    \end{scope}
\end{tikzpicture}
\end{center}
\begin{center}\small
\begin{tabular}{@{}clll@{}}
\toprule
\textbf{dequeue $u$} & \textbf{discover} & \textbf{update} & \textbf{$\mathrm{cnt}[1..4]$ after} \\
\midrule
--- & --- & marks: $\mathrm{cnt}[2] = 10$, $\mathrm{cnt}[1] = 100$, $\mathrm{cnt}[3] = 1$ & $100, 10, 1, 0$ \\
1 & 2 & $\mathrm{cnt}[2] = 10 + 100$ & $100, 110, 1, 0$ \\
2 & 3, 4 & $\mathrm{cnt}[3] = 1 + 110$, $\mathrm{cnt}[4] = 0 + 110$ & $100, 110, 111, 110$ \\
3, 4 & --- & leaves & $100, 110, 111, 110$ \\
\bottomrule
\end{tabular}
\end{center}
\textbf{Approach 3 on the same sample.} The stack DFS gives the pre-order $1, 2, 4, 3$ (children of 2 pushed as 3, 4; 4 is popped first), so $\mathrm{pos} = (0, 1, 3, 2)$ for vertices $1..4$ and $\mathrm{sz} = (4, 3, 1, 1)$.
\begin{center}
\begin{tikzpicture}[x=1.15cm]
    \node[font=\scriptsize\bfseries, anchor=east] at (-0.6,0.75) {vertex};
    \foreach \v [count=\i from 0] in {1,2,4,3} \node[idx] at (\i,0.75) {\v};
    \node[font=\scriptsize\bfseries, anchor=east] at (-0.6,0) {diff};
    \foreach \d [count=\i from 0] in {100,10,0,1,$-111$} \node[cell, minimum width=10mm] (d\i) at (\i,0) {\d};
    \node[font=\scriptsize\bfseries, anchor=east] at (-0.6,-0.75) {prefix};
    \foreach \d [count=\i from 0] in {100,110,110,111} \node[cell, minimum width=10mm, fill=computedfill] at (\i,-0.75) {\d};
    \draw[decorate, decoration={brace, amplitude=4pt, mirror}, primaryblue] (0.6,-1.2) -- (3.4,-1.2) node[midway, below=4pt, font=\scriptsize]{subtree(2) = positions 1..3: $+10$ at 1, $-10$ at 4};
    \node[note, anchor=west, text width=3.6cm, align=left] at (5.4,0) {op (1,100): $+100$ at 0, $-100$ at 4\\op (3,1): $+1$ at 3, $-1$ at 4};
\end{tikzpicture}
\end{center}
Read at $\mathrm{pos}[v]$: vertex 1 $\to 100$, 2 $\to 110$, 3 $\to 111$ (position 3), 4 $\to 110$ (position 2). The arrays $\mathrm{pos}$, $\mathrm{sz}$, diff and prefix were printed by an instrumented copy of Approach~3.
\end{tracebox}

\subsubsection*{6. Test suite}

\begin{testbox}
\centering\small
\begin{tabular}{@{}p{5.4cm}p{3.2cm}p{4.8cm}@{}}
\toprule
\textbf{Input} & \textbf{Output} & \textbf{Why} \\
\midrule
\texttt{2 1 / 1 2 / 2 5} & \texttt{0 5} & operation at a leaf \\
\texttt{4 2 / 1 4 / 2 4 / 3 4 / 4 7 / 1 1} & \texttt{1 8 8 8} & star centred at 4, but the root 1 is a leaf \\
\texttt{5 2 / 1 2 / 2 3 / 3 4 / 4 5 / 3 10 / 5 1} & \texttt{0 0 10 10 11} & path, operations in the middle \\
\texttt{3 2 / 1 3 / 2 3 / 3 1 / 2 2} & \texttt{0 3 1} & $a_i < b_i$ but 3 is the parent of 2 \\
path $1..2\cdot10^5$, $2\cdot10^5 \times (1, 10^4)$ & $2\cdot10^9$ everywhere & deepest tree, largest counter \\
random tree, $N = Q = 2\cdot10^5$ & (same checksum) & Approaches 2 and 3 agree \\
\bottomrule
\end{tabular}
\end{testbox}
All outputs from Approach~2; Approaches 1 and 3 printed the same on the small tests. In addition, the three approaches agreed on $3000$ random inputs ($N \le 10$, $Q \le 8$, random / path / star / caterpillar trees with shuffled labels and randomly ordered edges).

\subsubsection*{7. Common mistakes}

\begin{warningbox}[Mistake 1 --- trusting $a_i < b_i$]
\begin{lstlisting}[style=snippet]
for (each edge a b) par[b] = a;                       // WRONG: smaller label != parent
for (int v = 2; v <= n; v++) cnt[v] += cnt[par[v]];
\end{lstlisting}
Both samples happen to have $a_i = \mathrm{par}(b_i)$, so this passes them. On \texttt{3 1 / 1 3 / 2 3 / 3 1} it prints \texttt{0 0 1} instead of \texttt{0 1 1}; our random test found a mismatch on its very first input.
\end{warningbox}

\begin{warningbox}[Mistake 2 --- recursion on a path]
A recursive \texttt{dfs(v, p)} doing the same push-down crashed (no output at all) on the path of $2\cdot10^5$ vertices with the default 2\,MB stack; it works on random trees, whose depth is small. Use the BFS order.
\end{warningbox}

\begin{warningbox}[Mistake 3 --- simulating the operations]
Walking the subtree per operation looks harmless on the samples but is \bigO{NQ} (Approach~1). The point of the problem is to defer the work.
\end{warningbox}

\subsubsection*{8. Takeaways}

\begin{itemize}[leftmargin=*]
    \item ``Add to the subtree of $p$'' $=$ ``every descendant of $p$ receives it'' $=$ ``$v$ receives the marks on its root path'': mark the top, push down once.
    \item Push-down is pre-order: forward over a BFS order, and it can be fused into the BFS.
    \item Never derive the parent from the input order of an edge; root the tree yourself.
    \item Subtrees are contiguous in pre-order: subtree updates become range updates (Pattern~3, Chapter~15).
\end{itemize}

\begin{variantsbox}
\variant{Queries between updates}{\varnone}
  {updates ``add $x$ to subtree($p$)'' and queries ``value of $v$'' are interleaved.}
  {the same pre-order blocks.}
  {a Fenwick tree over positions: range add, point query, \bigO{\log N} each (Chapter~15). Our Fenwick version matched a per-operation subtree walk on $3000$ random trees with 15 interleaved operations each.}
\variant{Add along a root path}{\varnone}
  {operation: add $x$ to every vertex on the path from $p$ up to the root.}
  {the transpose: $v$ receives $x$ iff $p$ is in the subtree of $v$.}
  {mark $p$, then \emph{aggregate up} (backward loop): $\mathrm{cnt}[v] = \sum_{\mathrm{subtree}(v)} \mathrm{add}$. Matched a per-operation walk up the parents on $3000$ random trees.}
\variant{Best path ending here}{\varin{https://codeforces.com/problemset/problem/682/C}}
  {for every vertex, the largest sum of an upward path ending at it, with negative edges.}
  {push-down of a maximum, restarted when the parent's value is negative.}
  {$\mathrm{best}[v] = \max(\mathrm{best}[\mathrm{par}], 0) + c$ (Exercise E2.4).}
\variant{Counts per label}{\varin{https://leetcode.com/problems/number-of-nodes-in-the-sub-tree-with-the-same-label/}}
  {how many vertices of the subtree share $v$'s label.}
  {a pre-order position interval per subtree.}
  {a global counter read at entry and exit: the before/after trick (Exercise E2.1).}
\end{variantsbox}

%============================================================
% CH2 · PATTERN 2 — EXERCISES, HINTS, ANSWERS
%============================================================
\subsection{Pattern 2: Exercises and Going Further}
\label{ssec:ch2-p2-ex}

\begin{patternbox}{Pattern 2 Summary Card}
\renewcommand{\arraystretch}{1.25}
\begin{tabularx}{\linewidth}{@{}>{\bfseries}l X@{}}
Recognise & ``for every vertex, \dots{} in its subtree'' (up); ``\dots{} along the path from the root'', ``add to a whole subtree'' (down); ``delete a vertex / cut an edge'' (pieces) \\
Root it & pass the parent instead of \texttt{visited}; never trust the orientation of the input edges \\
Up & at leave: $\mathrm{agg}[v] = \mathrm{base}(v) \oplus \mathrm{agg}[\text{children}]$; the part above has $n - \mathrm{sz}[v]$ vertices \\
Down & at enter: $\mathrm{acc}[v] = \mathrm{acc}[\mathrm{par}] \otimes \mathrm{step}(v)$; for subtree updates mark the top and push down once \\
No recursion & BFS order: forward $=$ down, backward $=$ up; parent array with $p_i < i$: the labels are the order \\
Top traps & wrong parent from $a_i < b_i$; a value used before it is final; forgetting $n - \mathrm{sz}$; \texttt{int} products; deep recursion \\
\end{tabularx}
\end{patternbox}

\exlegend
\begin{enumerate}[leftmargin=*, label=\textbf{E2.\arabic*}, itemsep=6pt]
    \item \exeasy\ \pconcept \lc{Medium}\ \exprob{Number of Nodes in the Sub-Tree With the Same Label}{LeetCode 1519}{https://leetcode.com/problems/number-of-nodes-in-the-sub-tree-with-the-same-label/}. A tree on $0..n-1$ rooted at 0, each node labelled with a lowercase letter. For every node $i$, return how many nodes of its subtree carry the same label as $i$. ($n \le 10^5$.)
    \item \exeasy\ \pclassic \cf{1500}\ \exprob{Cut 'em all!}{Codeforces 982C}{https://codeforces.com/problemset/problem/982/C}. Remove as many edges of a tree as possible so that every remaining component has an even number of vertices; print that number, or $-1$ if impossible. ($n \le 10^5$.)
    \item \exmedium\ \pclassic \lc{Medium}\ \exprob{Minimum Fuel Cost to Report to the Capital}{LeetCode 2477}{https://leetcode.com/problems/minimum-fuel-cost-to-report-to-the-capital/}. A tree of cities with capital 0; one representative and one car with \texttt{seats} seats in every city; driving a road costs one litre per car; representatives may change cars. Minimum total fuel for everybody to reach the capital. ($n, \texttt{seats} \le 10^5$.)
    \item \exmedium\ \ptricky \cf{1600}\ \exprob{Alyona and the Tree}{Codeforces 682C}{https://codeforces.com/problemset/problem/682/C}. A tree rooted at 1 with a number $a_u$ on every vertex and a (possibly negative) number on every edge. Vertex $v$ is \emph{sad} if some $u$ in its subtree has $\mathrm{dist}(v, u) > a_u$. Remove leaves one at a time until no vertex is sad; minimise the number removed. ($n \le 10^5$, $|c| \le 10^9$.) The trap: the edge weights can be negative, so the worst ancestor is not the root.
    \item \exmedium\ \pconcept \lc{Hard}\ \exprob{Maximum Number of K-Divisible Components}{LeetCode 2872}{https://leetcode.com/problems/maximum-number-of-k-divisible-components/}. Remove any set of edges of a tree with node values so that every component's value sum is divisible by $k$ (the total is). Maximise the number of components. ($n \le 3\cdot10^4$.)
    \item \exhard\ \precent \lc{Hard}\ \exprob{Find Number of Coins to Place in Tree Nodes}{LeetCode 2973}{https://leetcode.com/problems/find-number-of-coins-to-place-in-tree-nodes/}. A tree rooted at 0 with a non-zero cost on every node. For each node: 1 coin if its subtree has fewer than 3 nodes, otherwise the maximum product of the costs of 3 distinct nodes of its subtree, or 0 if that product is negative. ($n \le 2\cdot10^4$, $|\mathrm{cost}| \le 10^4$.)
\end{enumerate}

\subsubsection*{Hints}
\begin{enumerate}[leftmargin=*, label=\textbf{E2.\arabic*}]
    \item Merging 26 counters per node works, but there is a cheaper way: a \emph{single} global counter per letter. What does it hold when you enter $i$, and when you leave $i$?
    \item With $n$ odd the total is odd: $-1$. Otherwise, which edges can be cut on their own? Does cutting one change whether another can be cut?
    \item Look at one road: how many representatives must cross it, at least? Can all roads reach their lower bound at the same time?
    \item A vertex $u$ that makes some ancestor sad must go, and with it its whole subtree. Which ancestor maximises $\mathrm{dist}(v, u)$? Think of the best upward path ending at $u$.
    \item Process children before parents. When the remaining sum of a subtree is divisible by $k$, is it ever wrong to cut it off?
    \item A maximum product of three numbers uses either the three largest, or the two smallest (both negative) and the largest. How many numbers must a subtree pass to its parent?
\end{enumerate}

\subsubsection*{Answers}

\begin{solutionbox}[E2.1 Number of Nodes in the Sub-Tree With the Same Label]
\textbf{The before/after trick.} Keep one global counter per letter, incremented when a node is \emph{entered}. Between entering and leaving $i$, exactly the nodes of $i$'s subtree are entered (Lesson~\ref{ssec:ch2-call-stack}), so $\mathrm{ans}[i] = \mathrm{cnt}_{\text{leave}}[\ell_i] - \mathrm{cnt}_{\text{enter}}[\ell_i]$, where the entry value is read \emph{before} $i$ itself is counted. This needs a real DFS (the BFS order does not keep subtrees together), so it is the iterative DFS with \texttt{it[u]}; $n = 10^5$ rules out recursion on LeetCode.
\begin{center}
\begin{tikzpicture}[x=1.15cm, y=0.9cm]
    \node[gvsrc] (a0) at (0,0) {0}; \node[gv] (a1) at (-1,-1) {1}; \node[gv] (a3) at (1,-1) {3}; \node[gv] (a2) at (-1,-2) {2};
    \foreach \a/\b in {0/1, 1/2, 0/3} \draw[gtree] (a\a) -- (a\b);
    \node[note, anchor=west, text width=8.2cm, align=left] at (2.2,-1) {labels \texttt{bbbb}. Counter of \texttt{b}: enter 0 (before 0), enter 1 (before 1), enter 2 (before 2), leave 2 at 3 $\Rightarrow 1$, leave 1 at 3 $\Rightarrow 2$, enter 3 (before 3), leave 3 at 4 $\Rightarrow 1$, leave 0 at 4 $\Rightarrow 4$};
\end{tikzpicture}
\end{center}
\begin{lstlisting}[style=cpp]
#include <bits/stdc++.h>
using namespace std;

class Solution {
public:
    vector<int> countSubTrees(int n, vector<vector<int>>& edges, string labels) {
        vector<vector<int>> adj(n);
        for (auto& e : edges) { adj[e[0]].push_back(e[1]); adj[e[1]].push_back(e[0]); }
        vector<int> ans(n, 0), before(n, 0), parent(n, -1), cnt(26, 0);
        vector<size_t> it(n, 0);
        vector<int> st(1, 0);
        before[0] = cnt[labels[0] - 'a']++;            // enter 0: remember the old count
        while (!st.empty()) {
            int u = st.back();
            if (it[u] < adj[u].size()) {
                int v = adj[u][it[u]++];
                if (v == parent[u]) continue;
                parent[v] = u;
                before[v] = cnt[labels[v] - 'a']++;    // enter v
                st.push_back(v);
            } else {                                   // leave u: its whole subtree was counted
                ans[u] = cnt[labels[u] - 'a'] - before[u];
                st.pop_back();
            }
        }
        return ans;
    }
};
\end{lstlisting}
Examples: $[2,1,1,1,1,1,1]$, $[4,2,1,1]$, $[3,2,1,1,1]$. On a path of $10^5$ nodes all labelled \texttt{z} it returns $10^5$ for the root and 1 for the last node. Matched a brute force (walk every subtree and count) on $5000$ random trees with $n \le 12$ and 3 letters. The alternative --- an array of 26 counts per node, added into the parent in the backward BFS loop --- is also correct and \bigO{26n}, but needs $26n$ integers of memory; the trick above needs 26 in total.
\end{solutionbox}

\begin{solutionbox}[E2.2 Cut 'em all!]
If $n$ is odd, some component is odd: $-1$. Otherwise cut the edge above $v$ exactly when $\mathrm{sz}[v]$ is even. \emph{Why this is allowed:} cutting such an edge leaves an even piece below and an even remainder above, so parities elsewhere do not change --- every such edge can be cut independently. \emph{Why no other edge:} after any valid set of cuts, the vertices below the edge above $v$ form a union of whole components iff that edge is cut; a union of even components has even size, so a cut edge always has $\mathrm{sz}[v]$ even. Hence the answer is the number of non-root $v$ with even $\mathrm{sz}[v]$.
\begin{center}
\begin{tikzpicture}[x=1.2cm, y=0.9cm]
    \node[gvsrc] (c1) at (0,0) {1}; \node[gv] (c4) at (-0.8,-1) {4}; \node[gv] (c3) at (0.8,-1) {3}; \node[gv] (c2) at (-0.8,-2) {2};
    \draw[gtree] (c1) -- (c3); \draw[gtree] (c4) -- (c2); \draw[gcut] (c1) -- (c4);
    \node[gdist, text=primaryblue, left=1pt of c4] {sz 2}; \node[gdist, text=primaryblue, right=1pt of c3] {sz 1}; \node[gdist, text=primaryblue, left=1pt of c2] {sz 1};
    \node[note, anchor=west, text width=6cm, align=left] at (2.0,-1) {sample 1, rooted at 1: only $\mathrm{sz}[4] = 2$ is even $\Rightarrow$ cut $1{-}4$, answer 1};
\end{tikzpicture}
\end{center}
\begin{lstlisting}[style=cpp]
#include <bits/stdc++.h>
using namespace std;

// Cut 'em all! (CF 982C): an edge (parent, v) can be cut iff sz[v] is even; with n odd nothing works.
int main() {
    int n;
    cin >> n;
    vector<vector<int>> adj(n + 1);
    for (int i = 0; i < n - 1; i++) {
        int a, b;
        cin >> a >> b;
        adj[a].push_back(b);
        adj[b].push_back(a);
    }
    if (n % 2 == 1) { cout << -1 << "\n"; return 0; }
    vector<int> parent(n + 1, 0), order(1, 1), sz(n + 1, 1);
    for (size_t i = 0; i < order.size(); i++)
        for (int v : adj[order[i]])
            if (v != parent[order[i]]) { parent[v] = order[i]; order.push_back(v); }
    int cuts = 0;
    for (int i = n - 1; i >= 1; i--) {             // reverse BFS order: children first
        int v = order[i];
        if (sz[v] % 2 == 0) cuts++;                // sz[v] is final: cut the edge above v
        sz[parent[v]] += sz[v];
    }
    cout << cuts << "\n";
}
\end{lstlisting}
Samples: \texttt{1}, \texttt{-1}, \texttt{4}, \texttt{0}. Matched a brute force that tries all $2^{n-1}$ edge subsets on $3000$ random trees with $n \le 12$ (including $n = 1$, answer $-1$).
\end{solutionbox}

\begin{solutionbox}[E2.3 Minimum Fuel Cost to Report to the Capital]
\textbf{Lower bound.} The road above $v$ must be crossed by all $\mathrm{sz}[v]$ representatives of the subtree of $v$, so it is driven by at least $\lceil \mathrm{sz}[v] / \texttt{seats} \rceil$ cars. \textbf{Achievable.} Move everybody up level by level: at $v$, the $\mathrm{sz}[v]$ people present pack into $\lceil \mathrm{sz}[v]/\texttt{seats} \rceil$ cars --- there are enough cars, because every person arrived with a car or in one ($\ge$ that many cars are at $v$). So the answer is $\sum_{v \ne 0} \lceil \mathrm{sz}[v] / \texttt{seats} \rceil$: aggregation plus one term per edge.
\begin{lstlisting}[style=cpp]
#include <bits/stdc++.h>
using namespace std;
typedef long long ll;

class Solution {
public:
    long long minimumFuelCost(vector<vector<int>>& roads, int seats) {
        int n = roads.size() + 1;
        vector<vector<int>> adj(n);
        for (auto& r : roads) { adj[r[0]].push_back(r[1]); adj[r[1]].push_back(r[0]); }
        vector<int> parent(n, -1), order(1, 0);
        for (size_t i = 0; i < order.size(); i++)
            for (int v : adj[order[i]])
                if (v != parent[order[i]] && v != 0) { parent[v] = order[i]; order.push_back(v); }
        vector<ll> sz(n, 1);
        ll fuel = 0;
        for (int i = n - 1; i >= 1; i--) {
            int v = order[i];
            fuel += (sz[v] + seats - 1) / seats;       // sz[v] people cross the road above v
            sz[parent[v]] += sz[v];
        }
        return fuel;
    }
};
\end{lstlisting}
Examples: \textbf{3}, \textbf{7}, \textbf{0}. The answer needs 64 bits: a path of $10^5$ cities with one seat costs $\sum_{k=1}^{99\,999} k = 4\,999\,950\,000$ (printed by the program). Checked on $5000$ random trees ($n \le 14$, $\texttt{seats} \le 4$) against an independent computation of the crossings (every representative walks its own path to the capital) and against a simulation that moves the cars level by level and asserts that enough cars are present. The simulation verifies feasibility; optimality is the lower bound above.
\end{solutionbox}

\begin{solutionbox}[E2.4 Alyona and the Tree]
\textbf{Which vertices must go.} If $u$ makes an ancestor $v$ sad, $u$ cannot stay (removing leaves never shortens the path $v \leadsto u$), and removing $u$ means removing its whole subtree first. Conversely, after removing exactly the subtrees of such ``bad'' $u$, nobody is sad. So the answer is the number of vertices with a bad vertex on their root path (itself included).

\textbf{The trap.} $\max_v \mathrm{dist}(v, u)$ over ancestors $v$ is \emph{not} $\mathrm{dist}(1, u)$: with negative edges, a later start can be better. It is the best sum of an upward path ending at $u$, a push-down with reset (Kadane's algorithm along every root path): $\mathrm{best}[u] = \max(\mathrm{best}[\mathrm{par}], 0) + c(\mathrm{par}, u)$. The ``$\max(\cdot, 0)$'' starts a fresh path at the parent. Do not descend below a bad vertex; the answer is $n - $ (vertices reached). Values reach $10^5 \cdot 10^9$: \texttt{long long}.
\begin{lstlisting}[style=cpp]
#include <bits/stdc++.h>
using namespace std;
typedef long long ll;

// Alyona and the Tree (CF 682C). best[u] = max over ancestors v of dist(v, u) = the largest sum of a
// path that ENDS at u and goes up: best[u] = max(best[p], 0) + c(p, u). A vertex with best[u] > a[u]
// must go, and with it its whole subtree. Answer = n - (vertices kept).
int main() {
    ios::sync_with_stdio(false);
    cin.tie(nullptr);
    int n;
    cin >> n;
    vector<ll> a(n + 1);
    for (int i = 1; i <= n; i++) cin >> a[i];
    vector<vector<pair<int, ll>>> adj(n + 1);
    for (int i = 2; i <= n; i++) {
        int p; ll c;
        cin >> p >> c;
        adj[p].push_back({i, c});
        adj[i].push_back({p, c});
    }
    vector<int> parent(n + 1, 0), order(1, 1);
    vector<ll> best(n + 1, 0);
    int kept = 0;
    for (size_t i = 0; i < order.size(); i++) {    // BFS = pre-order push-down, no recursion
        int u = order[i];
        kept++;
        for (auto& e : adj[u]) {
            int v = e.first;
            if (v == parent[u]) continue;
            parent[v] = u;
            best[v] = max(best[u], 0LL) + e.second; // restart the path at u if that is better
            if (best[v] > a[v]) continue;           // v makes an ancestor sad: drop v and its subtree
            order.push_back(v);
        }
    }
    cout << n - kept << "\n";
}
\end{lstlisting}
Sample: \textbf{5}. Matched a brute force (for every $u$, walk all ancestors and test $\mathrm{dist}(v,u) > a_u$; then count vertices below a bad one) on $3000$ random trees with $n \le 10$, edge weights in $[-20, 20]$ and parents that may carry larger labels than their children --- the input lists ``edge between $i+1$ and $p_i$'', and $p_i$ is the parent but not necessarily smaller.
\end{solutionbox}

\begin{solutionbox}[E2.5 Maximum Number of K-Divisible Components]
Children first (backward BFS order), keep $\mathrm{rem}[v]$ = the sum of what is still attached to $v$, modulo $k$. If $\mathrm{rem}[v] = 0$ when $v$ is left, cut the edge above $v$ (for the root: the last component); otherwise hand $\mathrm{rem}[v]$ to the parent. \emph{Why greedy is safe:} if the attached part of the subtree of $v$ has sum $\equiv 0$, cutting it off changes neither the residue of anything above (it removes $0 \bmod k$) nor any choice below; and an edge above a $v$ with $\mathrm{rem}[v] \ne 0$ can never be cut, because the part below it would have a sum $\not\equiv 0$. So the greedy cuts every edge that any valid split could cut.
\begin{center}
\begin{tikzpicture}[x=1.2cm, y=0.9cm]
    \node[gvsrc] (k0) at (0,0) {0}; \node[gv] (k2) at (0,-1) {2}; \node[gv] (k1) at (-0.9,-2) {1}; \node[gv] (k4) at (0.9,-2) {4}; \node[gv] (k3) at (-0.9,-3) {3};
    \draw[gtree] (k0) -- (k2); \draw[gtree] (k2) -- (k4); \draw[gtree] (k1) -- (k3); \draw[gcut] (k2) -- (k1);
    \foreach \v/\x in {0/1, 2/1, 1/8, 4/4, 3/4} \node[gdist, text=deepblue, right=1pt of k\v] {\x};
    \node[note, anchor=west, text width=6.4cm, align=left] at (2.0,-1.5) {example 1, $k = 6$ (values in blue). Subtree of 1: $8 + 4 = 12 \equiv 0$: cut. Node 2: $1 + 4 = 5$, handed up; root: $1 + 5 = 6 \equiv 0$: the last piece. Answer 2};
\end{tikzpicture}
\end{center}
\begin{lstlisting}[style=cpp]
#include <bits/stdc++.h>
using namespace std;
typedef long long ll;

class Solution {
public:
    int maxKDivisibleComponents(int n, vector<vector<int>>& edges, vector<int>& values, int k) {
        vector<vector<int>> adj(n);
        for (auto& e : edges) { adj[e[0]].push_back(e[1]); adj[e[1]].push_back(e[0]); }
        vector<int> parent(n, -1), order(1, 0);
        for (size_t i = 0; i < order.size(); i++)
            for (int v : adj[order[i]])
                if (v != parent[order[i]] && v != 0) { parent[v] = order[i]; order.push_back(v); }
        vector<ll> rem(n);                             // subtree sum mod k (of what is still attached)
        for (int v = 0; v < n; v++) rem[v] = values[v] % k;
        int comps = 0;
        for (int i = n - 1; i >= 0; i--) {
            int v = order[i];
            if (rem[v] == 0) comps++;                  // cut above v (the root: the last piece)
            else rem[parent[v]] = (rem[parent[v]] + rem[v]) % k;
        }
        return comps;
    }
};
\end{lstlisting}
Examples: \textbf{2}, \textbf{3}. Matched a brute force over all $2^{n-1}$ edge subsets on $5000$ random trees ($n \le 11$, $k \le 5$, values $0..6$ adjusted so that the total is divisible by $k$). Reducing modulo $k$ at every step keeps the numbers small; without it the sums still fit in \texttt{long long} ($3\cdot10^4 \cdot 10^9$), but not in \texttt{int}.
\end{solutionbox}

\begin{solutionbox}[E2.6 Find Number of Coins to Place in Tree Nodes]
The best product of three numbers from a multiset is $\max(\text{top}_1\text{top}_2\text{top}_3,\ \text{low}_1\text{low}_2\,\text{top}_1)$: either the three largest, or the two most negative times the largest. So a subtree needs to pass up only its \textbf{3 largest and 2 smallest} costs --- at most 5 numbers. Each node merges its own cost with its children's lists, sorts, drops everything between the 2nd smallest and the 3rd largest, and evaluates the formula if its subtree has $\ge 3$ nodes. Since every child passes at most 5 numbers, node $v$ sorts at most $5\deg(v) + 1$ numbers: \bigO{n \log n} in the worst case (a star), \bigO{n} on bounded degrees.
\begin{lstlisting}[style=cpp]
#include <bits/stdc++.h>
using namespace std;
typedef long long ll;

class Solution {
public:
    vector<long long> placedCoins(vector<vector<int>>& edges, vector<int>& cost) {
        int n = cost.size();
        vector<vector<int>> adj(n);
        for (auto& e : edges) { adj[e[0]].push_back(e[1]); adj[e[1]].push_back(e[0]); }
        vector<int> parent(n, -1), order(1, 0);
        for (size_t i = 0; i < order.size(); i++)
            for (int v : adj[order[i]])
                if (v != parent[order[i]] && v != 0) { parent[v] = order[i]; order.push_back(v); }
        // keep[v] = at most 5 costs of v's subtree: its 3 largest and its 2 smallest
        vector<vector<int>> keep(n);
        vector<int> sz(n, 1);
        vector<ll> coin(n, 1);
        for (int v = 0; v < n; v++) keep[v].push_back(cost[v]);
        for (int i = n - 1; i >= 0; i--) {
            int v = order[i];
            vector<int>& c = keep[v];                  // all children of v were merged into c
            sort(c.begin(), c.end());
            if (c.size() > 5) c.erase(c.begin() + 2, c.end() - 3);   // drop the middle
            if (sz[v] >= 3) {
                int m = c.size();
                ll top3 = (ll)c[m - 1] * c[m - 2] * c[m - 3];
                ll two_neg = (ll)c[0] * c[1] * c[m - 1];
                coin[v] = max(0LL, max(top3, two_neg));
            }
            if (i > 0) {
                int p = parent[v];
                sz[p] += sz[v];
                keep[p].insert(keep[p].end(), c.begin(), c.end());
            }
        }
        return coin;
    }
};
\end{lstlisting}
Examples: $[120,1,1,1,1,1]$, $[280,140,32,1,1,1,1,1,1]$, $[0,1,1]$. Matched a brute force that collects every subtree and tries all triples on $5000$ random trees ($n \le 13$, costs in $\pm[1, 9]$). Two details: with only 3 or 4 kept values the ``top'' and ``low'' entries overlap, which is fine --- both formulas still use three different positions of the sorted list, i.e.\ three distinct nodes; and the products reach $10^{12}$, so they are computed in \texttt{long long}.
\end{solutionbox}


%============================================================
\clearpage
\section{Pattern 3: Entry and Exit Times}
\label{sec:ch2-p3}
%============================================================
%============================================================
% CH2 · PATTERN 3 — ENTRY AND EXIT TIMES: THEORY
%============================================================
\subsection{Pattern Theory}
\label{ssec:ch2-p3-theory}

% the running tree T of this pattern: per-vertex styles can be overridden in a scope
\tikzset{
    g2p3-v1/.style=gv, g2p3-v2/.style=gv, g2p3-v3/.style=gv, g2p3-v4/.style=gv, g2p3-v5/.style=gv,
    g2p3-v6/.style=gv, g2p3-v7/.style=gv, g2p3-v8/.style=gv, g2p3-v9/.style=gv,
    g2p3-e/.style=gtree,
    g2p3-T/.pic={
        \node[g2p3-v1] (t1) at (3,0) {1};
        \node[g2p3-v2] (t2) at (1.2,-1.2) {2};
        \node[g2p3-v3] (t3) at (4.8,-1.2) {3};
        \node[g2p3-v4] (t4) at (0.4,-2.4) {4};
        \node[g2p3-v5] (t5) at (2.0,-2.4) {5};
        \node[g2p3-v6] (t6) at (4.0,-2.4) {6};
        \node[g2p3-v7] (t7) at (5.6,-2.4) {7};
        \node[g2p3-v8] (t8) at (3.4,-3.6) {8};
        \node[g2p3-v9] (t9) at (4.6,-3.6) {9};
        \foreach \a/\b in {1/2, 1/3, 2/4, 2/5, 3/6, 3/7, 6/8, 6/9} \draw[g2p3-e] (t\a) -- (t\b);
    }
}

\begin{patternbox}{Pattern 3 at a Glance}
Run \textbf{one} DFS and write down, for every vertex, \emph{when} it was entered and \emph{when} it was left. Afterwards the tree is no longer needed: ``$v$ is below $u$'' is two comparisons, ``the subtree of $u$'' is one contiguous block of an array, ``the descendants of $u$ at depth $d$'' is a binary search, and the order in which vertices \emph{finish} is an order in which they can be peeled off without disconnecting the rest.

\begin{center}
\begin{tikzpicture}[x=0.9cm, y=0.9cm]
    \tikzset{gv/.append style={minimum size=6.5mm, font=\scriptsize\bfseries}}
    \begin{scope}[g2p3-v3/.style=gvcur, g2p3-v6/.style=gvseen, g2p3-v7/.style=gvseen,
                  g2p3-v8/.style=gvseen, g2p3-v9/.style=gvseen]
        \pic {g2p3-T};
    \end{scope}
    \foreach \v/\a/\b/\p in {1/0/17/above, 2/1/6/left, 3/7/16/right, 4/2/3/below, 5/4/5/below,
                             6/8/13/left, 7/14/15/below, 8/9/10/below, 9/11/12/below}
        \node[gdist, text=darkgray, \p=0.4mm of t\v] {\tiny[\a,\b]};
    % the pre-order array
    \begin{scope}[xshift=7.3cm, yshift=-1.0cm]
        \foreach \v/\i in {1/0, 2/1, 4/2, 5/3, 3/4, 6/5, 8/6, 9/7, 7/8}{
            \ifnum\i<4 \node[gcell, minimum size=6.5mm, font=\scriptsize\bfseries] (c\i) at (\i*0.75,0) {\v};
            \else \node[gcell, fill=hitfill, minimum size=6.5mm, font=\scriptsize\bfseries] (c\i) at (\i*0.75,0) {\v}; \fi
            \node[font=\tiny, text=darkgray] at (\i*0.75,-0.55) {\i};
        }
        \node[gcell, fill=currentfill, draw=accentorange, line width=1pt, minimum size=6.5mm, font=\scriptsize\bfseries] at (4*0.75,0) {3};
        \draw[accentorange, line width=1.2pt, decorate, decoration={brace, amplitude=4pt}] (2.75,0.45) -- (6.25,0.45)
            node[midway, above=4pt, font=\scriptsize, text=accentorange] {pos$[3] = 4$, sz$[3] = 5$};
        \node[font=\scriptsize, anchor=west] at (-0.4,-1.2) {pre-order \texttt{order[]} (index $=$ pos)};
        \node[font=\scriptsize, anchor=west, text width=6.4cm] at (-0.4,-2.0) {$3$ is above $9$: $7 \le 11$ and $12 \le 16$.\quad $2$ is not above $6$: $\tout{6} = 13 > 6$.};
    \end{scope}
\end{tikzpicture}
\end{center}
\begin{recurrencebox}
\centering\small
enter $u$: $\tin{u} = \texttt{timer++}$, $\mathrm{pos}[u] = \texttt{cnt++}$ \qquad leave $u$: $\tout{u} = \texttt{timer++}$, $\mathrm{sz}[u] = \texttt{cnt} - \mathrm{pos}[u]$, append $u$ to \texttt{post}\\[3pt]
$u$ above $v \iff \tin{u} \le \tin{v} \wedge \tout{v} \le \tout{u}$ \qquad subtree$(u) = \texttt{order}[\mathrm{pos}[u] \,..\, \mathrm{pos}[u]+\mathrm{sz}[u]-1]$\\[3pt]
one DFS \VE, then \bigO{1} per ancestor query and \bigO{\log n} per ``at depth $d$'' query
\end{recurrencebox}
\end{patternbox}

%------------------------------------------------------------
\subsubsection*{From the lessons to the pattern}

Lesson~\ref{ssec:ch2-timestamps} introduced the two numberings and proved the parenthesis theorem; Pattern~2 used DFS to \emph{compute} something for every subtree while the search was running. This pattern turns the point of view around: the DFS runs \textbf{once, first}, and records only times. All questions are answered \emph{afterwards}, from the arrays, often many thousands of them. The single new ingredient is the observation that time is an \emph{order}, and orders can be searched:
\begin{itemize}[leftmargin=*]
    \item the interval $[\tin{u}, \tout{u}]$ contains exactly the entry times of $u$'s descendants (parenthesis theorem), so ``descendant'' becomes ``number in an interval'';
    \item in pre-order numbering the descendants are not only inside an interval, they \emph{fill} it: the subtree is a block of consecutive positions;
    \item the finishing order (post-order) lists every vertex after all its descendants --- the order for peeling, and, on DAGs, the reverse of a topological order.
\end{itemize}

%------------------------------------------------------------
\subsubsection*{When to recognise it (and when not)}

\begin{itemize}[leftmargin=*]
    \item many queries ``is $u$ an ancestor of $v$?'', ``does the path from the root to $v$ pass through $u$?'', ``is $v$ in the subtree of $u$?'' --- the \bigO{1} test;
    \item ``the $k$-th vertex visited when a DFS starts at $u$'', ``list the subtree of $u$ in DFS order'' --- the pre-order block (Problem~\ref{ssec:ch2-military});
    \item ``how many descendants of $u$ have depth $d$ / colour $c$ / label $\ell$'' --- per group a sorted list of entry times and a binary search, or an offline counter difference (Problem~\ref{ssec:ch2-count-descendants});
    \item ``remove $k$ vertices (cells) so that the rest stays connected'', ``in which order can the pieces be dismantled'' --- the finishing order (Problem~\ref{ssec:ch2-maze});
    \item ``the vertices are numbered so that every subtree is consecutive'' --- the input \emph{is} a pre-order: use the block directly (E3.5).
\end{itemize}
\textbf{Not this pattern:} if a subtree value can be computed bottom-up during the DFS itself (sizes, sums, counts), Pattern~2 is simpler. If the queries \emph{change} values (``add 5 to vertex $v$, then ask the sum of a subtree''), the block is still the key, but the range structure on top of it (Fenwick tree, segment tree) belongs to Chapter~15. Lowest common ancestors and distances between two arbitrary vertices need binary lifting or an Euler tour with range minimum (Chapter~14). Ancestor tests in a general directed graph speak about the DFS tree, \emph{not} about reachability (Lesson~\ref{ssec:ch2-timestamps}).

%------------------------------------------------------------
\subsubsection*{The framework}

\begin{center}
\small
\renewcommand{\arraystretch}{1.3}
\begin{tabularx}{\linewidth}{@{}>{\bfseries}l X@{}}
\toprule
Record & on entry: $\tin{u}$, $\mathrm{pos}[u]$, $\mathrm{dep}[u]$, append $u$ to \texttt{order}; on exit: $\tout{u}$, $\mathrm{sz}[u]$, append $u$ to \texttt{post} \\
Ancestor & $\tin{u} \le \tin{v}$ and $\tout{v} \le \tout{u}$ (or with the block: $\mathrm{pos}[u] \le \mathrm{pos}[v] < \mathrm{pos}[u] + \mathrm{sz}[u]$) \\
Subtree list & $\texttt{order}[\mathrm{pos}[u]], \dots, \texttt{order}[\mathrm{pos}[u] + \mathrm{sz}[u] - 1]$; the first element is $u$ itself \\
Group count & \texttt{byDepth[d]} $=$ entry times of depth-$d$ vertices (sorted for free); count those in $[\tin{u}, \tout{u}]$ \\
Peel order & \texttt{post}: every vertex after all of its descendants \\
Cost & one DFS \VE{} (iterative), $\bigO{n}$ memory per array; queries \bigO{1} or \bigO{\log n} \\
\bottomrule
\end{tabularx}
\end{center}

\begin{invariantbox}[Why a subtree is a block, and why peeling in finishing order is safe]
\textbf{Invariant (pre-order).} While \texttt{dfs(u)} is open, every vertex that receives a position is a descendant of $u$, and every descendant of $u$ receives its position while \texttt{dfs(u)} is open (white-path theorem: at $\tin{u}$ all of $u$'s subtree is white). Positions are handed out by one counter, so the positions given during \texttt{dfs(u)} are the consecutive numbers $\mathrm{pos}[u], \mathrm{pos}[u]+1, \dots, \texttt{cnt}-1$ at the moment $u$ is left --- exactly $\mathrm{sz}[u]$ of them. That is the block.

\textbf{Invariant (post-order).} At any moment, the set $F$ of finished vertices is closed under taking descendants: a vertex finishes only after all its children did. Hence for every \emph{unfinished} vertex $v \ne$ root, its parent is unfinished too (a finished parent would have finished $v$ before). So the unfinished vertices, together with their tree edges, form a connected subtree containing the root. Removing the first $k$ vertices of \texttt{post} therefore never disconnects the rest --- in any connected graph, since the DFS tree spans it.

\textbf{Depth lists are sorted.} Vertices are appended to \texttt{byDepth[d]} at their entry time, which only grows; so each list is increasing and binary search applies with no sort. \hfill$\blacksquare$
\end{invariantbox}

\begin{diagrambox}
\begin{tikzpicture}[x=0.62cm, y=0.5cm]
    \draw[->, darkgray] (-0.3,0) -- (18.4,0) node[right, font=\scriptsize] {time};
    \foreach \t in {0,...,17} {\draw[darkgray] (\t,0.12) -- (\t,-0.12); \node[font=\tiny, below] at (\t,-0.12) {\t};}
    \foreach \v/\a/\b/\l/\c in {1/0/17/5/primaryblue, 2/1/6/4/successgreen, 3/7/16/4/successgreen,
                                4/2/3/3/accentorange, 5/4/5/3/accentorange, 6/8/13/3/accentorange, 7/14/15/3/accentorange,
                                8/9/10/2/softpurple, 9/11/12/2/softpurple}{
        \draw[\c, line width=2.2pt] (\a,\l*0.95) -- (\b,\l*0.95);
        \draw[\c, line width=1pt] (\a,\l*0.95-0.25) -- (\a,\l*0.95+0.25);
        \draw[\c, line width=1pt] (\b,\l*0.95-0.25) -- (\b,\l*0.95+0.25);
        \node[font=\scriptsize\bfseries, fill=white, inner sep=1pt, text=\c] at ({(\a+\b)/2},\l*0.95+0.45) {\v};
    }
    \fill[accentorange, opacity=0.12] (7,-0.9) rectangle (16,5.5);
    \node[font=\scriptsize, text=accentorange, anchor=north] at (11.5,-0.95) {window of 3: every interval inside it is a descendant of 3};
\end{tikzpicture}
\figcap{The intervals $[\tin{u}, \tout{u}]$ of $T$, one row per depth. The ancestor test asks ``is $v$'s interval inside $u$'s window?''; because intervals never overlap partially, it is enough to look at the entry time $\tin{v}$ alone: $\tin{u} \le \tin{v} \le \tout{u}$.}
\end{diagrambox}

%------------------------------------------------------------
\subsubsection*{Combining the block with groups: per-depth lists}

``Descendants of $u$ with depth $d$'' mixes a \emph{subtree} condition with a \emph{group} condition. Keep one sorted list per group (here: per depth) of entry times. The subtree condition becomes an interval of times, and an interval of a sorted list is counted by two binary searches:
\[
    \#\{v \text{ below } u,\ \mathrm{dep}[v] = d\} \;=\; \texttt{upper\_bound}(L_d, \tout{u}) - \texttt{lower\_bound}(L_d, \tin{u}).
\]

\begin{diagrambox}
\begin{tikzpicture}[x=1cm, y=0.8cm]
    \foreach \d/\lst in {0/{0/1}, 1/{1/2, 7/3}, 2/{2/4, 4/5, 8/6, 14/7}, 3/{9/8, 11/9}}{
        \node[font=\scriptsize\bfseries, anchor=east] at (-0.2,-\d) {$L_{\d}$};
        \foreach \t/\v [count=\i from 0] in \lst {
            \node[gcell, minimum width=9mm, minimum height=6mm, font=\scriptsize] (L\d-\i) at (\i*1.0+0.3,-\d) {\t};
            \node[font=\tiny, text=darkgray, anchor=west] at (\i*1.0+0.3+0.32,-\d-0.28) {v\v};
        }
    }
    \node[gcell, minimum width=9mm, minimum height=6mm, font=\scriptsize\bfseries, fill=answerfill, draw=accentorange] at (0.3,-3) {9};
    \node[gcell, minimum width=9mm, minimum height=6mm, font=\scriptsize\bfseries, fill=answerfill, draw=accentorange] at (1.3,-3) {11};
    \node[gcell, minimum width=9mm, minimum height=6mm, font=\scriptsize\bfseries, fill=answerfill, draw=accentorange] at (2.3,-2) {8};
    \node[gcell, minimum width=9mm, minimum height=6mm, font=\scriptsize\bfseries, fill=answerfill, draw=accentorange] at (3.3,-2) {14};
    \node[note, text width=6.8cm, align=left, anchor=west] at (4.6,-1.0) {query $(u = 3, d = 3)$: window $[7, 16]$; in $L_3 = (9, 11)$ the first entry $\ge 7$ has index 0, the first $> 16$ has index 2: answer $2$ (vertices 8, 9).\\[3pt] query $(u = 3, d = 2)$: in $L_2 = (2, 4, 8, 14)$ the window covers indices $2..3$: answer $2$ (vertices 6, 7).\\[3pt] query $(u = 4, d = 1)$: $d < \mathrm{dep}[4]$, answer $0$.};
\end{tikzpicture}
\figcap{Per-depth lists of entry times of $T$ (small labels: the vertex). Each list is filled in time order, hence sorted. Values from the template program below.}
\end{diagrambox}

The same trick works for any partition of the vertices into groups (colour, label, letter): one list per group, $n$ entries in total. When the queries can be collected first, an offline variant needs no binary search at all: keep a counter per group, read it when entering $u$ and again when leaving $u$; the difference counts the group inside the subtree (Approach~3 of Problem~\ref{ssec:ch2-count-descendants}).

%------------------------------------------------------------
\subsubsection*{The finishing order: peeling and topological order}

The post-order is the third array a DFS gives for free. Two facts make it useful:
\begin{enumerate}[leftmargin=*]
    \item \textbf{Peeling.} The first vertices to finish are leaves of the DFS tree; in general every prefix of \texttt{post} is a union of complete subtrees, so removing a prefix leaves a connected graph (invariant above). This is the whole solution of Problem~\ref{ssec:ch2-maze}. The mirror statement for the \emph{pre}-order (and for a BFS order) is: every prefix is connected, because each vertex is discovered from an earlier one.
    \item \textbf{Reverse topological order.} In a DAG every edge $u \to v$ satisfies $\tout{v} < \tout{u}$: $v$ is white (tree edge, finishes inside $u$), or black (already finished), and it can never be gray (that would be a back edge, a cycle). So listing the vertices by \emph{decreasing} exit time puts every edge forwards: a topological order. Chapter~5 builds on this; here it is one \texttt{reverse}.
\end{enumerate}

\begin{diagrambox}
\begin{tikzpicture}[x=0.9cm, y=0.9cm]
    \tikzset{gv/.append style={minimum size=6.5mm, font=\scriptsize\bfseries}}
    \begin{scope}[g2p3-v4/.style=gvdead, g2p3-v5/.style=gvdead, g2p3-v2/.style=gvdead, g2p3-v8/.style=gvdead]
        \pic {g2p3-T};
        \foreach \a/\b in {1/2, 2/4, 2/5, 6/8} \draw[draw=white, line width=3pt] (t\a) -- (t\b);
        \foreach \a/\b in {1/2, 2/4, 2/5, 6/8} \draw[gfaded, dashed] (t\a) -- (t\b);
        \node[font=\scriptsize\bfseries] at (3,0.9) {remove the first 4 of \texttt{post}};
        \node[font=\scriptsize\ttfamily] at (3,-4.5) {post = \textcolor{midgray}{4 5 2 8} 9 6 7 3 1};
    \end{scope}
    \begin{scope}[xshift=7.6cm, yshift=-0.6cm]
        \node[font=\scriptsize\bfseries] at (2.0,1.5) {a DAG and its finishing order};
        \node[gvdone] (a1) at (0,0) {1}; \node[gvdone] (a2) at (1.4,0.7) {2}; \node[gvdone] (a3) at (1.4,-0.7) {3};
        \node[gvdone] (a4) at (2.8,0) {4}; \node[gvdone] (a5) at (4.2,0) {5}; \node[gvdone] (a6) at (0,-1.6) {6};
        \foreach \a/\b in {1/2, 1/3, 2/4, 3/4, 4/5, 6/3} \draw[gde] (a\a) -- (a\b);
        \foreach \v/\f/\p in {1/4/left, 2/2/above, 3/3/below, 4/1/above, 5/0/above, 6/5/left}
            \node[gdist, text=deepblue, \p=0.5mm of a\v] {\tiny\#\f};
        \node[font=\scriptsize\ttfamily, anchor=west] at (-0.3,-2.6) {finish: 5 4 2 3 1 6};
        \node[font=\scriptsize\ttfamily, anchor=west] at (-0.3,-3.15) {reverse: 6 1 3 2 4 5};
        \node[note, anchor=west, text width=4.6cm, align=left] at (-0.3,-3.95) {every edge goes left to right in the reversed list: a topological order (Chapter~5)};
    \end{scope}
\end{tikzpicture}
\figcap{Left: the four earliest finishers of $T$ are removed; what remains ($1, 3, 6, 7, 9$) is still a tree. Right: \#$i$ $=$ position in the finishing order; reversing it orders the DAG. Both orders printed by the template program.}
\end{diagrambox}

%------------------------------------------------------------
\subsubsection*{Computing the times iteratively}

With $n = 2\cdot10^5$ and a tree that may be a path, a recursive DFS needs $2\cdot10^5$ frames --- far beyond the about $65{,}000$ that the default 2\,MB stack survives (Lesson~\ref{ssec:ch2-call-stack}). The iterative form keeps, for every vertex on the stack, the index \texttt{it[u]} of the next neighbour to try. ``Enter'' happens when a vertex is pushed, ``leave'' when its iterator runs out and it is popped --- exactly the two moments where the recursive version writes \texttt{tin} and \texttt{tout}:

\begin{diagrambox}
\begin{tikzpicture}[node distance=4mm]
    \node[fcproc=3.0cm] (top) {$u =$ top of stack};
    \node[fcdec=3.0cm, below=of top] (d) {\texttt{it[u]} $<$ \texttt{deg(u)}?};
    \node[fcproc=3.6cm, right=14mm of d] (enter) {$v = \texttt{adj[u][it[u]++]}$; if $v \ne$ parent: \textbf{enter} $v$ (tin, pos, dep), push};
    \node[fcproc=3.6cm, below=of d] (leave) {\textbf{leave} $u$: tout, $\mathrm{sz}[u] = \texttt{cnt} - \mathrm{pos}[u]$, \texttt{post}, pop};
    \draw[fcarrow] (top) -- (d);
    \draw[fcarrow] (d) -- node[lbl, above]{yes} (enter);
    \draw[fcarrow] (d) -- node[lbl, right]{no} (leave);
    \draw[fcarrow] (enter.north) |- (top.east);
    \draw[fcarrow] (leave.west) -- ++(-0.9,0) |- (top.west);
\end{tikzpicture}
\figcap{One step of the iterative timestamp DFS. Each vertex is entered once and left once; each adjacency entry is read once: \VE{} and no recursion.}
\end{diagrambox}

When the input gives parents with $p_i < i$ (Problems~\ref{ssec:ch2-military} and~\ref{ssec:ch2-count-descendants}), sizes and depths need no DFS at all: depths by a loop $i = 2..n$, sizes by a loop $i = n..2$. The pre-order still needs the children order, but even it can be computed without a stack by handing out blocks top-down (Approach~3 of Problem~\ref{ssec:ch2-military}).

%------------------------------------------------------------
\subsubsection*{Variations map}

\begin{diagrambox}
\begin{tikzpicture}[
    var/.style={rectangle, rounded corners=2pt, draw=patternborder, fill=patternbg,
                font=\scriptsize, align=center, text width=3.25cm, minimum height=1.3cm, inner sep=3pt}]
    \node[var, fill=answerfill, draw=accentorange, text width=5.4cm, font=\scriptsize\bfseries] (root) at (0,0)
        {one DFS records tin / tout / pos / sz / post; answer from the arrays};
    \node[var] (v1) at (-5.4,-2) {\textbf{$k$-th in DFS order}\\\texttt{order[pos[u]+k-1]}\\(Military, 2.3.1)};
    \node[var] (v2) at (-1.8,-2) {\textbf{count at depth $d$}\\sorted tin per depth\\(Count Descendants, 2.3.2)};
    \node[var] (v3) at (1.8,-2)  {\textbf{peel keeping connectivity}\\first $k$ of post-order\\(Maze, 2.3.3)};
    \node[var] (v4) at (5.4,-2)  {\textbf{ancestor queries}\\two comparisons\\(E3.1)};
    \node[var] (v5) at (-5.4,-3.6) {\textbf{Euler walk}\\print on entry and on return\\(E3.2)};
    \node[var] (v6) at (-1.8,-3.6) {\textbf{edge types after the fact}\\tree / back / forward / cross\\(E3.3)};
    \node[var] (v7) at (1.8,-3.6)  {\textbf{which branch holds $x$?}\\child $c$ with $c$ above $x$\\(E3.4)};
    \node[var] (v8) at (5.4,-3.6)  {\textbf{input already in pre-order}\\edges cross the block twice\\(E3.5)};
    \node[draw=patternborder, dashed, rounded corners=4pt, fit=(v1)(v4)(v5)(v8), inner sep=5pt] (grp) {};
    \draw[flow] (root.south) -- (grp.north);
    \node[note, below=2mm of grp] {range sums / updates on the block: Chapter 15 \quad flattening as a technique: Chapter 13 \quad LCA: Chapter 14};
\end{tikzpicture}
\end{diagrambox}

\textbf{Twists and how to handle them}
\begin{itemize}[leftmargin=*]
    \item \textbf{``Children in increasing order''.} The pre-order depends on the order of the adjacency lists. Sort them, or build them so that they come out sorted (with $p_i < i$, appending $i$ to \texttt{ch[p\_i]} for $i = 2..n$ already sorts them).
    \item \textbf{The root is not vertex 1.} Run the DFS from the given root; for a question about ``the side of $x$ as seen from $y$'', simply root at $y$ (E3.4).
    \item \textbf{Strict or non-strict ancestor?} The test with $\le$ counts $u = v$ as an ancestor. For ``proper'' add $u \ne v$.
    \item \textbf{Both conventions at once.} One clock ($0..2n-1$) is best for comparisons; the pre-order counter ($0..n-1$) is best for indexing arrays. Recording both costs nothing.
    \item \textbf{A forest or a disconnected graph.} Keep the timer running across the outer loop; intervals of different trees are disjoint, so the tests stay correct.
    \item \textbf{The depth may exceed the list count.} Guard \texttt{d < byDepth.size()} (and $d \ge \mathrm{dep}[u]$) before indexing.
\end{itemize}

%------------------------------------------------------------
\subsubsection*{Code templates}

\begin{lstlisting}[style=cpp]
#include <bits/stdc++.h>
using namespace std;

// (1) Rooted tree (undirected adjacency lists, 1-indexed), ITERATIVE: safe for n = 2e5 on a path.
//     One clock:   tin[u] < tout[u] in 0..2n-1           -> ancestor test
//     Pre-order:   pos[u] in 0..n-1, order[pos[u]] = u   -> subtree = order[pos[u] .. pos[u]+sz[u]-1]
//     Post-order:  post = vertices in the order they finish (leaves of the DFS tree come early)
struct Timestamps {
    vector<int> tin, tout, pos, order, sz, dep, par, post;
    vector<vector<int>> byDepth;                  // byDepth[d] = tin of all depth-d vertices, sorted

    Timestamps(const vector<vector<int>>& adj, int root) {
        int n = adj.size() - 1, timer = 0, cnt = 0;
        tin.assign(n + 1, -1); tout = pos = sz = dep = par = tin;
        order.assign(n, 0);
        vector<size_t> it(n + 1, 0);
        vector<int> st;
        auto enter = [&](int v, int p) {
            par[v] = p; dep[v] = (p ? dep[p] + 1 : 0);
            tin[v] = timer++; pos[v] = cnt; order[cnt++] = v;
            if ((int)byDepth.size() <= dep[v]) byDepth.push_back(vector<int>());
            byDepth[dep[v]].push_back(tin[v]);    // pushed in increasing tin: already sorted
            st.push_back(v);
        };
        enter(root, 0);
        while (!st.empty()) {
            int u = st.back();
            if (it[u] < adj[u].size()) {
                int v = adj[u][it[u]++];
                if (v != par[u]) enter(v, u);
            } else {                                // u finishes
                tout[u] = timer++;
                sz[u] = cnt - pos[u];               // vertices entered since u
                post.push_back(u);
                st.pop_back();
            }
        }
    }
    bool isAnc(int u, int v) const { return tin[u] <= tin[v] && tout[v] <= tout[u]; }
    int kth(int u, int k) const { return k <= sz[u] ? order[pos[u] + k - 1] : -1; }  // k-th in u's subtree
    int countAtDepth(int u, int d) const {          // descendants of u (u included) at depth d
        if (d >= (int)byDepth.size() || d < dep[u]) return 0;
        const vector<int>& L = byDepth[d];
        return upper_bound(L.begin(), L.end(), tout[u]) - lower_bound(L.begin(), L.end(), tin[u]);
    }
};

// (2) Finishing order of a DIRECTED graph (all start vertices), iterative.
//     On a DAG the reverse of this order is a topological order (Chapter 5).
vector<int> finishOrder(const vector<vector<int>>& adj) {
    int n = adj.size() - 1;
    vector<int> color(n + 1, 0), res, st;
    vector<size_t> it(n + 1, 0);
    for (int s = 1; s <= n; s++) {
        if (color[s]) continue;
        color[s] = 1; st.push_back(s);
        while (!st.empty()) {
            int u = st.back();
            if (it[u] < adj[u].size()) {
                int v = adj[u][it[u]++];
                if (!color[v]) { color[v] = 1; st.push_back(v); }
            } else { color[u] = 2; res.push_back(u); st.pop_back(); }
        }
    }
    return res;
}

int main() {
    // the tree T: 1-2, 1-3, 2-4, 2-5, 3-6, 3-7, 6-8, 6-9 (neighbours in increasing order)
    int n = 9;
    vector<vector<int>> adj(n + 1);
    int E[8][2] = {{1, 2}, {1, 3}, {2, 4}, {2, 5}, {3, 6}, {3, 7}, {6, 8}, {6, 9}};
    for (auto& e : E) { adj[e[0]].push_back(e[1]); adj[e[1]].push_back(e[0]); }
    Timestamps T(adj, 1);
    for (int u = 1; u <= n; u++)
        cout << u << ": tin " << T.tin[u] << " tout " << T.tout[u] << " pos " << T.pos[u]
             << " sz " << T.sz[u] << " dep " << T.dep[u] << "\n";
    cout << "pre-order :"; for (int v : T.order) cout << " " << v; cout << "\n";
    cout << "post-order:"; for (int v : T.post) cout << " " << v; cout << "\n";
    cout << T.isAnc(3, 9) << T.isAnc(2, 6) << T.isAnc(6, 6) << "\n";              // 101
    cout << T.kth(3, 3) << " " << T.kth(3, 5) << " " << T.kth(2, 4) << "\n";        // 8 7 -1
    cout << T.countAtDepth(1, 2) << " " << T.countAtDepth(3, 2) << " "
         << T.countAtDepth(3, 3) << " " << T.countAtDepth(4, 1) << "\n";            // 4 2 2 0
    // a DAG: 1->2, 1->3, 2->4, 3->4, 4->5, 6->3
    vector<vector<int>> dag(7);
    dag[1] = {2, 3}; dag[2] = {4}; dag[3] = {4}; dag[4] = {5}; dag[6] = {3};
    vector<int> f = finishOrder(dag);
    cout << "finish:"; for (int v : f) cout << " " << v; cout << "\n";
    reverse(f.begin(), f.end());
    cout << "topo  :"; for (int v : f) cout << " " << v; cout << "\n";
}
\end{lstlisting}

\begin{lstlisting}[style=trace]
output: 1: tin 0 tout 17 pos 0 sz 9 dep 0      6: tin 8 tout 13 pos 5 sz 3 dep 2
        2: tin 1 tout 6 pos 1 sz 3 dep 1       7: tin 14 tout 15 pos 8 sz 1 dep 2
        3: tin 7 tout 16 pos 4 sz 5 dep 1      8: tin 9 tout 10 pos 6 sz 1 dep 3
        4: tin 2 tout 3 pos 2 sz 1 dep 2       9: tin 11 tout 12 pos 7 sz 1 dep 3
        5: tin 4 tout 5 pos 3 sz 1 dep 2       (one per line in the real output)
        pre-order : 1 2 4 5 3 6 8 9 7
        post-order: 4 5 2 8 9 6 7 3 1
        101
        8 7 -1
        4 2 2 0
        finish: 5 4 2 3 1 6
        topo  : 6 1 3 2 4 5
\end{lstlisting}
On a path of $2\cdot10^5$ vertices the iterative constructor runs without any stack problem; the recursive version of the same code crashes at that depth.

\begin{insightbox}[Which template when?]
\renewcommand{\arraystretch}{1.25}
\begin{tabularx}{\linewidth}{@{}>{\bfseries}l X@{}}
\texttt{isAnc} & ``is $u$ above $v$'', ``does the root path of $v$ pass $u$'', ``which child of $y$ leads to $x$'' \\
\texttt{kth} / \texttt{order} & ``$k$-th vertex of a DFS from $u$'', ``list / iterate the subtree of $u$'', any per-subtree array question (range queries on it: Chapter~15) \\
\texttt{countAtDepth} & ``how many descendants of $u$ in group $g$'' online; offline, the counter difference is shorter \\
\texttt{post} / \texttt{finishOrder} & peel vertices keeping the rest connected; reverse it for a topological order of a DAG \\
\end{tabularx}
\end{insightbox}

%------------------------------------------------------------
\subsubsection*{Common mistakes}

\begin{warningbox}[Mistake 1 --- mixing the two conventions]
\begin{lstlisting}[style=snippet]
int kth = order[tin[u] + k - 1];       // WRONG: tin uses one clock for entry AND exit (0..2n-1)
int kth = order[pos[u] + k - 1];       // RIGHT: pos counts entries only (0..n-1)
\end{lstlisting}
With one clock the subtree is \emph{inside} the interval but does not \emph{fill} it: on $T$, $\tin{3} = 7$ while $\mathrm{pos}[3] = 4$.
\end{warningbox}

\begin{warningbox}[Mistake 2 --- unsorted children when the order matters]
If the problem prescribes ``children by increasing index'' and the adjacency lists were filled from unordered edges, the pre-order --- and every \texttt{kth} answer --- is wrong, while ancestor tests and counts stay right. Sort the lists (or build them sorted).
\end{warningbox}

\begin{warningbox}[Mistake 3 --- binary search on the wrong side of the window]
\begin{lstlisting}[style=snippet]
lower_bound(L, tout[u]) - lower_bound(L, tin[u])      // WRONG only if tout values are in L
upper_bound(L, tout[u]) - lower_bound(L, tin[u])      // RIGHT and robust
\end{lstlisting}
With one clock no entry time equals $\tout{u}$, so both forms happen to agree --- but with the pre-order convention ($[\mathrm{pos}[u], \mathrm{pos}[u]+\mathrm{sz}[u]-1]$) the last position \emph{is} an entry and \texttt{lower\_bound} on the right end drops it. Use the half-open form consistently.
\end{warningbox}

\begin{warningbox}[Mistake 4 --- recursion on a path]
$n = 2\cdot10^5$ with $p_i = i - 1$ is a legal input of all three solved problems; a grid of $500 \times 500$ can force a DFS $125{,}000$ cells deep. Use the iterative template.
\end{warningbox}

\begin{warningbox}[Mistake 5 --- removing in entry order]
The pre-order is the order for \emph{keeping} (every prefix is connected), the post-order the order for \emph{removing} (every prefix is a union of subtrees). Removing the first $k$ vertices of the pre-order removes the root first and can cut the graph into pieces.
\end{warningbox}

\begin{checklistbox}[``I'm stuck'' checklist for Pattern 3]
\begin{itemize}[leftmargin=*]
    \item[$\square$] Can the question be phrased as ``$v$ in the subtree of $u$'' plus something simple? Then the interval $[\tin{u}, \tout{u}]$ or the block $\mathrm{pos}[u]..\mathrm{pos}[u]+\mathrm{sz}[u]-1$ is the answer's skeleton.
    \item[$\square$] Is there a second condition (depth, colour, label)? One sorted list of times per group, or an offline counter difference.
    \item[$\square$] Does the answer depend on the visiting order (``smallest index first'')? Sort the children.
    \item[$\square$] Is it a peeling or dismantling question? Finishing order removes, entry or BFS order keeps.
    \item[$\square$] Where is the root? Re-root at the vertex the question is ``seen from''.
    \item[$\square$] Depth up to $n$? Iterative DFS. Products of sizes? \texttt{long long}.
\end{itemize}
\end{checklistbox}

%============================================================
% CH2 · P3 · PROBLEM 1 — MILITARY PROBLEM (CF 1006E)
%============================================================
\subsection{Problem 2.3.1: Military Problem}
\problemlabel{2.3.1}{ssec:ch2-military}

\problemcard{2.3.1}{Military Problem}{Codeforces 1006E}{https://codeforces.com/problemset/problem/1006/E}{\exmedium\ \cf{1600}}{One pre-order of the whole tree; the command list of $u$ is the block \texttt{order[pos[u] .. pos[u]+sz[u]-1]}}

% local layout of the sample tree
\tikzset{
    g2p3m-v1/.style=gv, g2p3m-v2/.style=gv, g2p3m-v3/.style=gv, g2p3m-v4/.style=gv, g2p3m-v5/.style=gv,
    g2p3m-v6/.style=gv, g2p3m-v7/.style=gv, g2p3m-v8/.style=gv, g2p3m-v9/.style=gv,
    g2p3m-e/.style={gde},
    g2p3m-T/.pic={
        \node[g2p3m-v1] (m1) at (2.4,0) {1};
        \node[g2p3m-v2] (m2) at (0.6,-1.2) {2};
        \node[g2p3m-v3] (m3) at (2.4,-1.2) {3};
        \node[g2p3m-v4] (m4) at (4.2,-1.2) {4};
        \node[g2p3m-v5] (m5) at (1.6,-2.4) {5};
        \node[g2p3m-v7] (m7) at (3.2,-2.4) {7};
        \node[g2p3m-v6] (m6) at (1.0,-3.6) {6};
        \node[g2p3m-v8] (m8) at (2.2,-3.6) {8};
        \node[g2p3m-v9] (m9) at (3.2,-3.6) {9};
        \foreach \a/\b in {1/2, 1/3, 1/4, 3/5, 3/7, 5/6, 5/8, 7/9} \draw[g2p3m-e] (m\a) -- (m\b);
    }
}

\subsubsection*{1. Problem statement}

An army of $n$ officers forms a tree rooted at officer 1; officer $i \ge 2$ has the direct superior $p_i < i$. When officer $u$ issues a command it spreads by DFS: the current officer hands it to his direct subordinate with the \textbf{smallest index} that has not received it yet, waits until that subordinate has finished spreading it through his own subtree, then picks the next one; when no subordinate is left, he is done. Answer $q$ independent queries $(u_i, k_i)$: the $k_i$-th officer to receive the command issued by $u_i$ ($u_i$ himself is the first), or $-1$ if fewer than $k_i$ officers receive it.

\textbf{Constraints:} $2 \le n \le 2\cdot10^5$, $1 \le q \le 2\cdot10^5$, $1 \le p_i < i$, $1 \le u_i, k_i \le n$.

\begin{table}[H]
\centering\small
\begin{tabular}{@{}lll@{}}
\toprule
\textbf{Input} & \textbf{Official output} & \textbf{Our output} \\
\midrule
\texttt{9 6 / 1 1 1 3 5 3 5 7 / 3 1 / 1 5 / 3 4 / 7 3 / 1 8 / 1 9} & \texttt{3 6 8 -1 9 4} & \texttt{3 6 8 -1 9 4} \\
\bottomrule
\end{tabular}
\end{table}

\begin{diagrambox}
\begin{tikzpicture}[x=1cm, y=0.9cm]
    \begin{scope}
        \pic {g2p3m-T};
        \node[font=\scriptsize\bfseries] at (2.4,0.8) {the sample: $p = (1,1,1,3,5,3,5,7)$};
    \end{scope}
    \begin{scope}[xshift=5.5cm, g2p3m-v3/.style=gvsrc, g2p3m-v5/.style=gvseen, g2p3m-v6/.style=gvseen,
                  g2p3m-v8/.style=gvdst, g2p3m-v7/.style=gvseen, g2p3m-v9/.style=gvseen]
        \pic {g2p3m-T};
        \node[font=\scriptsize\bfseries] at (2.4,0.8) {command from 3: $3, 5, 6, \mathbf{8}, 7, 9$};
        \foreach \v/\k/\p in {3/1/right, 5/2/left, 6/3/below, 8/4/below, 7/5/right, 9/6/below}
            \node[gdist, text=accentorange, \p=0.6mm of m\v] {\tiny\k};
        \node[note, text width=3.3cm, align=left, anchor=west] at (4.9,-2.6) {query $(3, 4)$: the 4th officer is \textbf{8}; query $(7, 3)$: only $7, 9$ receive it, answer $-1$};
    \end{scope}
\end{tikzpicture}
\figcap{The sample tree and the spreading from officer 3 (small numbers: receiving order). Children are always taken by increasing index.}
\end{diagrambox}

\subsubsection*{2. The graph model}

\begin{modelbox}[Graph Model --- Military Problem]
\renewcommand{\arraystretch}{1.35}
\begin{tabularx}{\linewidth}{@{}>{\bfseries}l X@{}}
Vertices & the $n$ officers \\
Edges & tree edges $p_i \to i$, rooted at 1; children ordered by index \\
Weights & none \\
Question & the $k$-th vertex of the pre-order of the subtree of $u$ (with this child order) \\
Size & $n, q \le 2\cdot10^5$: \bigO{n} once plus \bigO{1} per query; the tree can be a path ($2\cdot10^5$ deep) \\
\end{tabularx}
\end{modelbox}

\begin{diagrambox}
\begin{tikzpicture}[x=1cm, y=1cm]
    \foreach \v/\i in {1/0, 2/1, 3/2, 5/3, 6/4, 8/5, 7/6, 9/7, 4/8}{
        \node[gcell, minimum size=8mm, font=\small\bfseries] (c\i) at (\i*0.95,0) {\v};
        \node[font=\tiny, text=darkgray] at (\i*0.95,-0.6) {pos \i};
    }
    \draw[primaryblue, line width=1.4pt, decorate, decoration={brace, amplitude=5pt}] (-0.45,0.5) -- (8.05,0.5)
        node[midway, above=5pt, font=\scriptsize, text=primaryblue] {command from 1: all 9 (sz 9)};
    \draw[accentorange, line width=1.3pt, decorate, decoration={brace, amplitude=4pt, mirror}] (1.45,-0.85) -- (7.1,-0.85)
        node[midway, below=4pt, font=\scriptsize, text=accentorange] {command from 3: pos $2..7$ (sz 6)};
    \draw[softpurple, line width=1.3pt, decorate, decoration={brace, amplitude=4pt}] (5.25,0.45) -- (7.1,0.45)
        node[midway, above=4pt, font=\scriptsize, text=softpurple] {from 7};
\end{tikzpicture}
\figcap{The pre-order of the whole tree, \texttt{order[]}. Every officer's command list is a contiguous block of it, starting at the officer himself.}
\end{diagrambox}

\subsubsection*{3. Key observations}

\begin{itemize}[leftmargin=*]
    \item \textbf{The spreading is a DFS pre-order} with children in increasing order. A command from $u$ reaches exactly the subtree of $u$, in the same relative order in which a DFS from the root would visit it --- the DFS from the root, once it enters $u$, does exactly what a DFS from $u$ does until it leaves $u$.
    \item \textbf{So all $n$ lists are blocks of one array.} By the pre-order invariant of the pattern, $u$'s list is $\texttt{order}[\mathrm{pos}[u]], \dots, \texttt{order}[\mathrm{pos}[u]+\mathrm{sz}[u]-1]$. The answer is $\texttt{order}[\mathrm{pos}[u]+k-1]$ if $k \le \mathrm{sz}[u]$, else $-1$.
    \item \textbf{Children come sorted for free.} Reading $p_2, p_3, \dots$ and appending $i$ to \texttt{ch[$p_i$]} puts the children in increasing order.
    \item \textbf{$p_i < i$ makes sizes trivial:} for $i = n$ down to $2$, every child is finished before its parent is read.
    \item \textbf{Depth.} $p_i = i - 1$ gives a path of $2\cdot10^5$ officers: the DFS must be iterative.
\end{itemize}

\subsubsection*{4. Approaches}

\begin{bruteforcebox}
\textbf{Idea: simulate.} For each query run the spreading from $u_i$ literally and stop after $k_i$ officers.

\begin{lstlisting}[style=cpp]
#include <bits/stdc++.h>
using namespace std;

// Brute force: simulate the spreading for every query, stop after k officers.
vector<vector<int>> ch;
vector<int> got;

void spread(int a, int k) {
    if ((int)got.size() == k) return;              // the k-th officer is known: stop early
    got.push_back(a);
    for (int b : ch[a]) spread(b, k);              // children are stored in increasing order
}

int main() {
    int n, q;
    cin >> n >> q;
    ch.assign(n + 1, vector<int>());
    for (int i = 2; i <= n; i++) {
        int p;
        cin >> p;
        ch[p].push_back(i);
    }
    while (q--) {
        int u, k;
        cin >> u >> k;
        got.clear();
        spread(u, k);
        cout << ((int)got.size() == k ? got[k - 1] : -1) << "\n";
    }
}
\end{lstlisting}

\begin{complexitybox}
\textbf{Time:} \bigO{n} per query, \bigO{nq} $= 4\cdot10^{10}$ in the worst case (a star with $u = 1$, $k = n$ on every query). \textbf{Space:} \bigO{n}, but the recursion is as deep as the tree --- on the path input it also crashes. \textbf{Verdict:} \tle{}; it is the reference for the random tests.
\end{complexitybox}
\end{bruteforcebox}

\begin{algorithmbox}[Approach 2: one pre-order, one block per officer]
\textbf{Idea.} Compute the pre-order of the whole tree once, iteratively: pop a vertex, give it the next position, push its children in \emph{reverse} order so that the smallest is popped next. Then sizes bottom-up, and each query is one array access.

\begin{lstlisting}[style=cpp]
#include <bits/stdc++.h>
using namespace std;

// Military Problem (CF 1006E): the command from u spreads in DFS pre-order of u's subtree
// (children by increasing index). One pre-order of the whole tree answers every query:
// u's list is the block order[pos[u] .. pos[u] + sz[u] - 1].
int main() {
    ios::sync_with_stdio(false);
    cin.tie(nullptr);
    int n, q;
    cin >> n >> q;
    vector<vector<int>> ch(n + 1);
    for (int i = 2; i <= n; i++) {
        int p;
        cin >> p;
        ch[p].push_back(i);                    // i increases: children lists come out sorted
    }
    vector<int> order, pos(n + 1), sz(n + 1, 1);
    vector<int> st(1, 1);
    while (!st.empty()) {                      // iterative pre-order (a path of 2e5 is legal)
        int u = st.back(); st.pop_back();
        pos[u] = order.size();
        order.push_back(u);
        for (int j = (int)ch[u].size() - 1; j >= 0; j--)
            st.push_back(ch[u][j]);            // reverse: the smallest child is popped first
    }
    for (int u = n; u >= 1; u--)               // p_i < i: every child is finished before u
        for (int c : ch[u]) sz[u] += sz[c];
    while (q--) {
        int u, k;
        cin >> u >> k;
        cout << (k <= sz[u] ? order[pos[u] + k - 1] : -1) << "\n";
    }
}
\end{lstlisting}

\begin{complexitybox}
\textbf{Time:} \bigO{n} for the pre-order and the sizes, \bigO{1} per query: \bigO{n + q}. \textbf{Space:} \bigO{n}, no recursion. \textbf{Verdict:} \ac{}-ready: on the path $p_i = i-1$ with $2\cdot10^5$ random queries the whole run, input and output included, takes under a second on this machine.
\end{complexitybox}
\end{algorithmbox}

\begin{invariantbox}[Why the block is the command list]
\textbf{Claim.} Popping from the stack produces a pre-order of the tree with children in increasing order, and the vertices popped from the moment $u$ is popped until the stack is back to the height it had below $u$ are exactly $u$'s subtree.

When $u$ is popped, its children are pushed above everything that was on the stack, smallest on top. Everything pushed later lies above them, so nothing below $u$'s children is popped before all of $u$'s children --- and, by induction, all their descendants --- have been popped. Hence $u$'s subtree is popped as one consecutive run that starts with $u$, and inside it the children's subtrees come one after another in increasing order of the child: the spreading order. The run has length $\mathrm{sz}[u]$, which the bottom-up loop computes ($p_i < i$, so when $u$ is processed all its children already hold final sizes). \hfill$\blacksquare$
\end{invariantbox}

\begin{alternativebox}[Approach 3: hand out the blocks top-down --- no stack at all]
Because $p_i < i$, every vertex is processed after its parent in the loop $u = 1..n$. If $u$ already knows where its block starts, its children's blocks follow right after $u$, one after another in increasing index: child $c_1$ starts at $\mathrm{pos}[u]+1$, child $c_2$ at $\mathrm{pos}[u]+1+\mathrm{sz}[c_1]$, and so on. Sizes come first, bottom-up.

\begin{lstlisting}[style=cpp]
#include <bits/stdc++.h>
using namespace std;

// Alternative without any stack: because p_i < i, compute the sizes bottom-up (i = n..2),
// then hand out the blocks top-down (u = 1..n): u's children get consecutive blocks
// starting right after u, in increasing index order.
int main() {
    ios::sync_with_stdio(false);
    cin.tie(nullptr);
    int n, q;
    cin >> n >> q;
    vector<int> p(n + 1, 0);
    vector<vector<int>> ch(n + 1);
    for (int i = 2; i <= n; i++) {
        cin >> p[i];
        ch[p[i]].push_back(i);
    }
    vector<int> sz(n + 1, 1), pos(n + 1, 0), order(n);
    for (int i = n; i >= 2; i--) sz[p[i]] += sz[i];     // children have larger indices
    pos[1] = 0;
    for (int u = 1; u <= n; u++) {                      // pos[u] is known: its parent came earlier
        order[pos[u]] = u;
        int next = pos[u] + 1;                          // first free slot inside u's block
        for (int c : ch[u]) { pos[c] = next; next += sz[c]; }
    }
    while (q--) {
        int u, k;
        cin >> u >> k;
        cout << (k <= sz[u] ? order[pos[u] + k - 1] : -1) << "\n";
    }
}
\end{lstlisting}

\begin{diagrambox}
\begin{tikzpicture}[x=0.95cm, y=1cm]
    \foreach \i in {0,...,8} \node[gcell, minimum size=7mm, font=\scriptsize] (b\i) at (\i,0) {};
    \foreach \i in {0,...,8} \node[font=\tiny, text=darkgray] at (\i,-0.55) {\i};
    \node[gcell, minimum size=7mm, font=\scriptsize\bfseries, fill=currentfill] at (0,0) {1};
    \draw[successgreen, line width=1.2pt] (0.55,0.45) -- (1.45,0.45) node[midway, above, font=\scriptsize] {2 (sz 1)};
    \draw[accentorange, line width=1.2pt] (1.55,0.45) -- (7.45,0.45) node[midway, above, font=\scriptsize] {3 (sz 6)};
    \draw[softpurple, line width=1.2pt] (7.55,0.45) -- (8.45,0.45) node[midway, above, font=\scriptsize] {4 (sz 1)};
    \node[note, anchor=west, text width=5.2cm, align=left] at (9.0,0.1) {$u = 1$: pos$[2] = 1$, pos$[3] = 1+1 = 2$, pos$[4] = 2+6 = 8$. Later $u = 3$ splits $2..7$ into $3 \mid 5,6,8 \mid 7,9$.};
\end{tikzpicture}
\figcap{Approach 3 on the sample: a vertex cuts its own block into the blocks of its children, left to right.}
\end{diagrambox}

\begin{complexitybox}
\textbf{Time:} \bigO{n + q}. \textbf{Space:} \bigO{n}. \textbf{Verdict:} \ac{}-ready; it gives the same \texttt{pos} and \texttt{order} as Approach~2 (identical outputs on the $2\cdot10^5$ path test). It works only because $p_i < i$ gives a parents-first order for free; for a tree given by an edge list, use Approach~2 or the template.
\end{complexitybox}
\end{alternativebox}

\begin{tricksbox}[Approach 4: Implementation Refinements]
\begin{itemize}[leftmargin=*]
    \item \textbf{One clock would not do.} $\tin{}$ from the one-clock convention skips numbers (exits use them too); indexing \texttt{order} needs the entries-only counter \texttt{pos}.
    \item \textbf{Fast I/O.} $2\cdot10^5$ answers: \texttt{"\textbackslash n"}, never \texttt{endl}; \texttt{sync\_with\_stdio(false)}.
    \item \textbf{Reverse push.} Forgetting to push the children in reverse produces the pre-order with children by \emph{decreasing} index --- wrong answers, but only on vertices with $\ge 2$ children.
    \item \textbf{Recursive alternative.} A recursive pre-order is three lines, but at depth $2\cdot10^5$ it crashes on a 2\,MB stack; our test of exactly that case terminated abnormally.
\end{itemize}
\end{tricksbox}

\subsubsection*{5. Trace}

\begin{tracebox}[{Sample: children 1: 2 3 4,\ 3: 5 7,\ 5: 6 8,\ 7: 9}]
\begin{center}
\begin{tikzpicture}[x=0.72cm, y=0.72cm]
    \tikzset{gv/.append style={minimum size=6mm, font=\scriptsize\bfseries}}
    \begin{scope}[g2p3m-v1/.style=gvdone, g2p3m-v2/.style=gvcur, g2p3m-v3/.style=gvseen, g2p3m-v4/.style=gvseen]
        \pic {g2p3m-T};
        \node[font=\scriptsize\bfseries] at (2.4,0.9) {pop 2 (pos 1)};
        \node[font=\scriptsize\ttfamily] at (2.4,-4.5) {stack: 4 3};
    \end{scope}
    \begin{scope}[xshift=4.6cm, g2p3m-v1/.style=gvdone, g2p3m-v2/.style=gvdone, g2p3m-v3/.style=gvdone,
                  g2p3m-v5/.style=gvcur, g2p3m-v4/.style=gvseen, g2p3m-v7/.style=gvseen, g2p3m-v6/.style=gvseen, g2p3m-v8/.style=gvseen]
        \pic {g2p3m-T};
        \node[font=\scriptsize\bfseries] at (2.4,0.9) {pop 5 (pos 3)};
        \node[font=\scriptsize\ttfamily] at (2.4,-4.5) {stack: 4 7 8 6};
    \end{scope}
    \begin{scope}[xshift=9.2cm, g2p3m-v1/.style=gvdone, g2p3m-v2/.style=gvdone, g2p3m-v3/.style=gvdone,
                  g2p3m-v5/.style=gvdone, g2p3m-v6/.style=gvdone, g2p3m-v8/.style=gvdone, g2p3m-v7/.style=gvcur,
                  g2p3m-v4/.style=gvseen, g2p3m-v9/.style=gvseen]
        \pic {g2p3m-T};
        \node[font=\scriptsize\bfseries] at (2.4,0.9) {pop 7 (pos 6)};
        \node[font=\scriptsize\ttfamily] at (2.4,-4.5) {stack: 4 9};
    \end{scope}
\end{tikzpicture}
\end{center}
\begin{center}\small
\begin{tabular}{@{}clll@{}}
\toprule
\textbf{pop} & \textbf{pos} & \textbf{push (reversed children)} & \textbf{stack after} \\
\midrule
1 & 0 & 4, 3, 2 & \texttt{4 3 2} \\
2 & 1 & --- & \texttt{4 3} \\
3 & 2 & 7, 5 & \texttt{4 7 5} \\
5 & 3 & 8, 6 & \texttt{4 7 8 6} \\
6 & 4 & --- & \texttt{4 7 8} \\
8 & 5 & --- & \texttt{4 7} \\
7 & 6 & 9 & \texttt{4 9} \\
9 & 7 & --- & \texttt{4} \\
4 & 8 & --- & (empty) \\
\bottomrule
\end{tabular}
\qquad
\begin{tabular}{@{}cccc@{}}
\toprule
\textbf{query} & \textbf{pos, sz} & \textbf{index} & \textbf{answer} \\
\midrule
$(3,1)$ & 2, 6 & 2 & 3 \\
$(1,5)$ & 0, 9 & 4 & 6 \\
$(3,4)$ & 2, 6 & 5 & 8 \\
$(7,3)$ & 6, 2 & $3 > 2$ & $-1$ \\
$(1,8)$ & 0, 9 & 7 & 9 \\
$(1,9)$ & 0, 9 & 8 & 4 \\
\bottomrule
\end{tabular}
\end{center}
Sizes from the bottom-up loop: $\mathrm{sz} = (9, 1, 6, 1, 3, 1, 2, 1, 1)$ for officers $1..9$. Orange $=$ just popped, light orange $=$ on the stack, blue $=$ already numbered.
\end{tracebox}

\subsubsection*{6. Test suite}

\begin{testbox}
\centering\small
\begin{tabular}{@{}p{6.6cm}lp{4.8cm}@{}}
\toprule
\textbf{Input} & \textbf{Our output} & \textbf{Why} \\
\midrule
\texttt{2 3 / 1 / 1 1 / 1 2 / 2 2} & \texttt{1 2 -1} & smallest tree; $k$ beyond a leaf \\
\texttt{5 5 / 1 2 3 4 / 1 5 / 2 4 / 3 3 / 5 1 / 5 2} & \texttt{5 5 5 5 -1} & a path: every block ends at 5 \\
\texttt{5 4 / 1 1 1 1 / 1 2 / 1 5 / 2 1 / 3 2} & \texttt{2 5 2 -1} & a star: children by index \\
\texttt{6 3 / 1 1 2 2 3 / 1 6 / 2 2 / 3 2} & \texttt{6 4 6} & the subtree of 2 comes before 3 \\
path $p_i = i-1$, $n = q = 2\cdot10^5$ & (same for 2 and 3) & depth $2\cdot10^5$, iterative only \\
\bottomrule
\end{tabular}
\end{testbox}
All outputs come from running Approach~2. Approaches 1--3 reproduce the official sample and agree on $1000$ random trees with $n \le 10$ (paths, stars and random parents, $q \le 10$ random queries each).

\subsubsection*{7. Common mistakes}

\begin{warningbox}[Mistake 1 --- one DFS per query]
Re-running the spreading (even iteratively) for every query is the brute force: \bigO{nq}. The point of the problem is that all $n$ lists are pieces of one list.
\end{warningbox}

\begin{warningbox}[Mistake 2 --- \texttt{tin} instead of \texttt{pos}]
With the one-clock times of the lesson, $\tin{3} = 3$ on the sample, but officer 3 is at position 2 of the pre-order; \texttt{order[tin[u]+k-1]} even runs past the end of the array for large $k$.
\end{warningbox}

\begin{warningbox}[Mistake 3 --- off-by-one in the bound]
$u$ himself is the first officer: the block has $\mathrm{sz}[u]$ elements and the answer exists iff $k \le \mathrm{sz}[u]$, at index $\mathrm{pos}[u] + k - 1$. Query $(7, 3)$ of the sample ($\mathrm{sz}[7] = 2$) catches \texttt{k <= sz[u] + 1}.
\end{warningbox}

\begin{warningbox}[Mistake 4 --- children in the wrong order]
Pushing the children in increasing order onto a stack pops the \emph{largest} first; the sample then prints \texttt{3 9 5 -1 6 2} instead of \texttt{3 6 8 -1 9 4} (our run of that bug). Only the first answer, which never leaves officer 3 himself, survives. Reverse the push, or use the iterator form of the template.
\end{warningbox}

\subsubsection*{8. Takeaways}

\begin{itemize}[leftmargin=*]
    \item A DFS started at $u$ visits exactly the part of the global pre-order that belongs to $u$'s subtree: compute the global order once.
    \item Subtree $=$ block \texttt{order[pos[u] .. pos[u]+sz[u]-1]}: positions count entries only.
    \item With $p_i < i$, sizes need only a reverse loop, and even the positions can be handed out top-down without a stack.
    \item $n = 2\cdot10^5$ and a parent array allow a path: iterative.
\end{itemize}

\begin{variantsbox}
\variant{The $k$-th officer at a given depth}{\varnone}
  {query $(u, d, k)$: the $k$-th officer of depth $d$ that receives $u$'s command.}
  {the depth-$d$ vertices of $u$'s subtree are a contiguous run of the depth-$d$ list sorted by $\mathrm{pos}$.}
  {binary search the first entry $\ge \mathrm{pos}[u]$ in that list and step $k-1$ further; check that it is still $< \mathrm{pos}[u]+\mathrm{sz}[u]$ (the idea of Problem~\ref{ssec:ch2-count-descendants}).}
\variant{Where in the list is officer $v$?}{\varnone}
  {query $(u, v)$: at which step does $v$ receive $u$'s command, or never?}
  {$v$ receives it iff $v$ is in the block of $u$.}
  {answer $\mathrm{pos}[v] - \mathrm{pos}[u] + 1$ if $\mathrm{pos}[u] \le \mathrm{pos}[v] < \mathrm{pos}[u]+\mathrm{sz}[u]$, else ``never''; checked against the brute force on random trees.}
\variant{Command returns to the issuer}{\varin{https://atcoder.jp/contests/abc213/tasks/abc213_d}}
  {every officer reports back to his superior when done; list every visit, including returns.}
  {the Euler walk: a vertex appears on entry and after each child.}
  {print on push and again after each pop (Exercise E3.2).}
\variant{Children by decreasing seniority}{\varnone}
  {each officer passes the command to his subordinates by a given priority instead of index.}
  {same tree, different child order.}
  {sort each \texttt{ch[u]} by the priority before the pre-order; queries unchanged.}
\end{variantsbox}

%============================================================
% CH2 · P3 · PROBLEM 2 — COUNT DESCENDANTS (AtCoder ABC202 E)
%============================================================
\subsection{Problem 2.3.2: Count Descendants}
\problemlabel{2.3.2}{ssec:ch2-count-descendants}

\problemcard{2.3.2}{Count Descendants}{AtCoder ABC202 E}{https://atcoder.jp/contests/abc202/tasks/abc202_e}{\exmedium\ (AtCoder ABC 500 points)}{Entry times per depth, sorted for free; count those inside $[\tin{U}, \tout{U}]$ with two binary searches --- or offline, a counter difference between entering and leaving $U$}

\tikzset{
    g2p3c-v1/.style=gv, g2p3c-v2/.style=gv, g2p3c-v3/.style=gv, g2p3c-v4/.style=gv,
    g2p3c-v5/.style=gv, g2p3c-v6/.style=gv, g2p3c-v7/.style=gv,
    g2p3c-T/.pic={
        \node[g2p3c-v1] (c1) at (2.4,0) {1};
        \node[g2p3c-v2] (c2) at (1.2,-1.2) {2};
        \node[g2p3c-v3] (c3) at (3.6,-1.2) {3};
        \node[g2p3c-v4] (c4) at (0.2,-2.4) {4};
        \node[g2p3c-v5] (c5) at (1.2,-2.4) {5};
        \node[g2p3c-v7] (c7) at (2.2,-2.4) {7};
        \node[g2p3c-v6] (c6) at (0.2,-3.6) {6};
        \foreach \a/\b in {1/2, 1/3, 2/4, 2/5, 2/7, 4/6} \draw[gtree, -{Stealth[length=2.2mm]}] (c\a) -- (c\b);
    }
}

\subsubsection*{1. Problem statement}

A tree with $N$ vertices is rooted at vertex 1; vertex $i \ge 2$ has parent $P_i < i$. Answer $Q$ queries $(U_i, D_i)$: how many vertices $u$ have $U_i$ on their path to the root (endpoints included) \emph{and} exactly $D_i$ edges between $u$ and the root?

\textbf{Constraints:} $2 \le N \le 2\cdot10^5$, $1 \le P_i < i$, $1 \le Q \le 2\cdot10^5$, $1 \le U_i \le N$, $0 \le D_i \le N-1$.

\begin{table}[H]
\centering\small
\begin{tabular}{@{}lll@{}}
\toprule
\textbf{Input} & \textbf{Official output} & \textbf{Our output} \\
\midrule
\texttt{7 / 1 1 2 2 4 2 / 4 / 1 2 / 7 2 / 4 1 / 5 5} & \texttt{3 1 0 0} & \texttt{3 1 0 0} \\
\bottomrule
\end{tabular}
\end{table}

\begin{diagrambox}
\begin{tikzpicture}[x=1cm, y=0.9cm]
    \begin{scope}
        \pic {g2p3c-T};
        \foreach \d in {0,1,2,3} \node[font=\scriptsize, text=darkgray, anchor=east] at (-0.4,-1.2*\d) {depth \d};
        \node[font=\scriptsize\bfseries] at (2.0,0.8) {the sample tree};
    \end{scope}
    \begin{scope}[xshift=6.2cm, g2p3c-v1/.style=gvsrc, g2p3c-v4/.style=gvdst, g2p3c-v5/.style=gvdst, g2p3c-v7/.style=gvdst]
        \pic {g2p3c-T};
        \draw[accentorange, dashed, rounded corners=3pt] (-0.4,-2.0) rectangle (2.7,-2.8);
        \node[font=\scriptsize\bfseries] at (2.0,0.8) {query $(1, 2)$: answer 3};
        \node[note, anchor=west, text width=3.6cm, align=left] at (3.2,-2.6) {every vertex of depth 2 lies below the root: $4, 5, 7$};
    \end{scope}
\end{tikzpicture}
\figcap{``$U$ is on the path from $u$ to the root'' means ``$u$ is a descendant of $U$ (or $U$ itself)''. Query $(7, 2)$ counts only 7; query $(4, 1)$ asks for depth 1 below a vertex of depth 2: 0; depth 5 does not exist: 0.}
\end{diagrambox}

\subsubsection*{2. The graph model}

\begin{modelbox}[Graph Model --- Count Descendants]
\renewcommand{\arraystretch}{1.35}
\begin{tabularx}{\linewidth}{@{}>{\bfseries}l X@{}}
Vertices & the $N$ tree vertices, each with a depth $\mathrm{dep}[v]$ \\
Edges & $P_i \to i$, rooted at 1 \\
Weights & none \\
Question & $\#\{v : v \text{ in the subtree of } U,\ \mathrm{dep}[v] = D\}$ --- a subtree condition and a group condition \\
Size & $N, Q \le 2\cdot10^5$: \bigO{N} or \bigO{N \log N} preprocessing, \bigO{\log N} per query; depth up to $N-1$ \\
\end{tabularx}
\end{modelbox}

\begin{diagrambox}
\begin{tikzpicture}[x=0.75cm, y=0.5cm]
    \draw[->, darkgray] (-0.3,0) -- (14.4,0) node[right, font=\scriptsize] {time};
    \foreach \t in {0,...,13} {\draw[darkgray] (\t,0.12) -- (\t,-0.12); \node[font=\tiny, below] at (\t,-0.12) {\t};}
    \foreach \v/\a/\b/\l/\c in {1/0/13/4/primaryblue, 2/1/10/3/successgreen, 3/11/12/3/successgreen,
                                4/2/5/2/accentorange, 5/6/7/2/accentorange, 7/8/9/2/accentorange, 6/3/4/1/softpurple}{
        \draw[\c, line width=2.2pt] (\a,\l*1.0) -- (\b,\l*1.0);
        \draw[\c, line width=1pt] (\a,\l*1.0-0.25) -- (\a,\l*1.0+0.25);
        \draw[\c, line width=1pt] (\b,\l*1.0-0.25) -- (\b,\l*1.0+0.25);
        \node[font=\scriptsize\bfseries, fill=white, inner sep=1pt, text=\c] at ({(\a+\b)/2},\l*1.0+0.45) {\v};
    }
    \foreach \d/\l in {0/4, 1/3, 2/2, 3/1} \node[font=\scriptsize, text=darkgray, anchor=east] at (-0.4,\l*1.0) {dep \d};
\end{tikzpicture}
\figcap{The model in time: intervals $[\tin{v}, \tout{v}]$ of the sample, one row per depth. A query is ``how many intervals of row $D$ start inside the interval of $U$''.}
\end{diagrambox}

\subsubsection*{3. Key observations}

\begin{itemize}[leftmargin=*]
    \item \textbf{Path to the root through $U$ $=$ descendant of $U$.} Then the parenthesis theorem turns the subtree condition into $\tin{U} \le \tin{v} \le \tout{U}$.
    \item \textbf{Group by depth.} For each depth $d$ keep the list $L_d$ of entry times of the depth-$d$ vertices. The DFS appends to $L_d$ in time order, so every list is sorted without sorting.
    \item \textbf{Count $=$ two binary searches.} $\texttt{upper\_bound}(L_D, \tout{U}) - \texttt{lower\_bound}(L_D, \tin{U})$. A depth below $U$'s own depth gives 0 automatically (no shallower vertex is inside the window), and so does a depth that does not occur (empty list).
    \item \textbf{Why not per-vertex tables?} Storing ``count per depth'' for every vertex is \bigO{N^2} memory on a path.
    \item \textbf{Depth up to $2\cdot10^5$}: iterative DFS.
\end{itemize}

\subsubsection*{4. Approaches}

\begin{bruteforcebox}
\textbf{Idea: test every vertex.} For a query, take each vertex of depth $D$ and climb to the root; count it if $U$ is met.

\begin{lstlisting}[style=cpp]
#include <bits/stdc++.h>
using namespace std;

// Brute force: for every query, look at every vertex v of depth D and walk from v up to the root;
// count v if U is met on the way.
int main() {
    int n;
    cin >> n;
    vector<int> p(n + 1, 0), dep(n + 1, 0);
    for (int i = 2; i <= n; i++) {
        cin >> p[i];
        dep[i] = dep[p[i]] + 1;                  // p_i < i: the parent's depth is known
    }
    int q;
    cin >> q;
    while (q--) {
        int u, d, ans = 0;
        cin >> u >> d;
        for (int v = 1; v <= n; v++) {
            if (dep[v] != d) continue;
            for (int x = v; x != 0; x = p[x])
                if (x == u) { ans++; break; }
        }
        cout << ans << "\n";
    }
}
\end{lstlisting}

\begin{complexitybox}
\textbf{Time:} \bigO{N \cdot \text{depth}} per query --- on a path with $D = N-1$ that is $N$ steps per query, $4\cdot10^{10}$ in total. \textbf{Space:} \bigO{N}. \textbf{Verdict:} \tle{}; reference only.
\end{complexitybox}
\end{bruteforcebox}

\begin{algorithmbox}[Approach 2: sorted entry times per depth + binary search]
\textbf{Idea.} One iterative DFS records $\tin{}$, $\tout{}$, $\mathrm{dep}$ and appends $\tin{v}$ to \texttt{at[dep[v]]} at the moment $v$ is entered. Each query is two binary searches in one list.

\begin{lstlisting}[style=cpp]
#include <bits/stdc++.h>
using namespace std;

// Count Descendants (ABC202 E): count v at depth D with U on the path root..v,
// i.e. descendants of U (U included) at depth D. Iterative DFS with one clock;
// at[d] = entry times of the depth-d vertices, sorted because they are appended in time order.
int main() {
    ios::sync_with_stdio(false);
    cin.tie(nullptr);
    int n;
    cin >> n;
    vector<vector<int>> ch(n + 1);
    for (int i = 2; i <= n; i++) {
        int p;
        cin >> p;
        ch[p].push_back(i);
    }
    vector<int> tin(n + 1), tout(n + 1), dep(n + 1, 0);
    vector<size_t> it(n + 1, 0);
    vector<vector<int>> at(n);
    int timer = 0;
    vector<int> st(1, 1);
    tin[1] = timer++; at[0].push_back(tin[1]);
    while (!st.empty()) {
        int u = st.back();
        if (it[u] < ch[u].size()) {
            int v = ch[u][it[u]++];
            dep[v] = dep[u] + 1;
            tin[v] = timer++;
            at[dep[v]].push_back(tin[v]);
            st.push_back(v);
        } else {
            tout[u] = timer++;
            st.pop_back();
        }
    }
    int q;
    cin >> q;
    while (q--) {
        int u, d;
        cin >> u >> d;
        const vector<int>& L = at[d];             // d <= n - 1, so the index is valid
        int ans = upper_bound(L.begin(), L.end(), tout[u]) - lower_bound(L.begin(), L.end(), tin[u]);
        cout << ans << "\n";
    }
}
\end{lstlisting}

\begin{complexitybox}
\textbf{Time:} \bigO{N} for the DFS (the lists hold $N$ numbers in total), \bigO{\log N} per query: \bigO{N + Q \log N}. \textbf{Space:} \bigO{N}. \textbf{Verdict:} \ac{}-ready; on the path $P_i = i-1$ with $2\cdot10^5$ random queries the run takes under a second here.
\end{complexitybox}
\end{algorithmbox}

\begin{invariantbox}[Why the count is right]
\textbf{Lists are sorted.} \texttt{timer} only grows, and $\tin{v}$ is appended to \texttt{at[dep[v]]} at the moment it is assigned; so each list is strictly increasing.

\textbf{The window selects the subtree.} By the parenthesis theorem, $v$ is a descendant of $U$ (or $U$ itself) iff $\tin{U} \le \tin{v} \le \tout{U}$. Restricted to the list of depth $D$, the vertices that qualify are exactly the entries in $[\tin{U}, \tout{U}]$; in a sorted list they form the index range $[\texttt{lower\_bound}(\tin{U}), \texttt{upper\_bound}(\tout{U}))$, whose length is the answer. Depths smaller than $\mathrm{dep}[U]$ have no entry in the window (no ancestor of a vertex can be its descendant), and missing depths are empty lists. \hfill$\blacksquare$
\end{invariantbox}

\begin{alternativebox}[Approach 3: offline --- counter difference between entering and leaving $U$]
Collect the queries at their vertex $U$. During the DFS keep \texttt{cnt[d]} $=$ number of depth-$d$ vertices entered so far. Everything entered after ``enter $U$'' and before ``leave $U$'' is in $U$'s subtree, so
\[
    \text{answer}(U, D) = \texttt{cnt}[D]\big|_{\text{leaving } U} - \texttt{cnt}[D]\big|_{\text{entering } U,\ \text{before counting } U}.
\]
No sorting, no binary search: \bigO{N + Q}.

\begin{lstlisting}[style=cpp]
#include <bits/stdc++.h>
using namespace std;

// Offline alternative: cnt[d] = number of depth-d vertices entered so far.
// Everything entered between "enter U" and "leave U" is in U's subtree, so
// answer(U, D) = cnt[D] when leaving U  -  cnt[D] when entering U.
int main() {
    ios::sync_with_stdio(false);
    cin.tie(nullptr);
    int n;
    cin >> n;
    vector<vector<int>> ch(n + 1);
    for (int i = 2; i <= n; i++) {
        int p;
        cin >> p;
        ch[p].push_back(i);
    }
    int q;
    cin >> q;
    vector<vector<pair<int, int>>> qs(n + 1);     // qs[u] = (D, query index)
    vector<long long> ans(q, 0);
    for (int i = 0; i < q; i++) {
        int u, d;
        cin >> u >> d;
        qs[u].push_back({d, i});
    }
    vector<int> cnt(n, 0), dep(n + 1, 0);
    vector<size_t> it(n + 1, 0);
    vector<int> st;
    auto enter = [&](int v) {
        for (auto& pr : qs[v]) ans[pr.second] -= cnt[pr.first];   // before v itself is counted
        cnt[dep[v]]++;
        st.push_back(v);
    };
    enter(1);
    while (!st.empty()) {
        int u = st.back();
        if (it[u] < ch[u].size()) {
            int v = ch[u][it[u]++];
            dep[v] = dep[u] + 1;
            enter(v);
        } else {
            for (auto& pr : qs[u]) ans[pr.second] += cnt[pr.first];
            st.pop_back();
        }
    }
    for (int i = 0; i < q; i++) cout << ans[i] << "\n";
}
\end{lstlisting}

\begin{complexitybox}
\textbf{Time:} \bigO{N + Q}. \textbf{Space:} \bigO{N + Q}. \textbf{Verdict:} \ac{}-ready, identical output to Approach~2 on the $2\cdot10^5$ path test. It needs all queries in advance; Approach~2 is \emph{online} (it could answer a query before reading the next). The difference trick generalises to any group key (colour, letter) and is the seed of the ``DSU on tree'' / small-to-large techniques of Chapter~15.
\end{complexitybox}
\end{alternativebox}

\begin{diagrambox}
\begin{tikzpicture}[x=0.78cm, y=0.8cm]
    \foreach \t/\ev in {0/{in 1}, 1/{in 2}, 2/{in 4}, 3/{in 6}, 4/{out 6}, 5/{out 4}, 6/{in 5}, 7/{out 5}, 8/{in 7}, 9/{out 7}, 10/{out 2}, 11/{in 3}, 12/{out 3}, 13/{out 1}}{
        \node[font=\tiny, rotate=60, anchor=west] at (\t,0.25) {\ev};
        \draw[darkgray] (\t,0.12) -- (\t,-0.12);
    }
    \draw[darkgray] (-0.3,0) -- (13.3,0);
    \node[font=\scriptsize, anchor=east] at (-0.4,-0.7) {cnt[2] after};
    \foreach \t/\c in {0/0, 1/0, 2/1, 3/1, 4/1, 5/1, 6/2, 7/2, 8/3, 9/3, 10/3, 11/3, 12/3, 13/3}
        \node[font=\scriptsize\ttfamily] at (\t,-0.7) {\c};
    \draw[accentorange, line width=1.2pt, -{Stealth[length=2mm]}] (7,-1.2) -- (7,-0.95);
    \draw[accentorange, line width=1.2pt, -{Stealth[length=2mm]}] (9,-1.2) -- (9,-0.95);
    \node[font=\scriptsize, text=accentorange, anchor=north, text width=8cm, align=center] at (8,-1.25) {query $(7, 2)$: on entering 7, before 7 itself is counted, \texttt{cnt[2]} is still 2 (the value after event 7); on leaving 7 it is 3: $3 - 2 = 1$};
\end{tikzpicture}
\figcap{Approach 3 on the sample: the depth-2 counter along the 14 events of the DFS. Query $(1,2)$ reads $0$ and $3$; query $(4,1)$ reads \texttt{cnt[1]} $= 1$ on both sides: $0$.}
\end{diagrambox}

\begin{tricksbox}[Approach 4: Implementation Refinements]
\begin{itemize}[leftmargin=*]
    \item \textbf{Size the depth table by $N$.} $D \le N - 1$, so \texttt{at} with $N$ lists can be indexed directly; with \texttt{at} sized by the maximum depth, guard the index.
    \item \textbf{Pre-order positions work too.} With $\mathrm{pos}$ and $\mathrm{sz}$ the window is $[\mathrm{pos}[U], \mathrm{pos}[U]+\mathrm{sz}[U]-1]$ --- then use \texttt{upper\_bound} on the right end (Mistake 3 of the pattern).
    \item \textbf{No DFS needed for depths} ($P_i < i$), but the entry times need the DFS; the recursion depth can be $N$: iterative.
    \item \textbf{Memory.} $N$ entries in all lists together plus the children lists: a few MB.
\end{itemize}
\end{tricksbox}

\subsubsection*{5. Trace}

\begin{tracebox}[{Sample: children 1: 2 3,\ 2: 4 5 7,\ 4: 6}]
\begin{center}
\begin{tikzpicture}[x=0.72cm, y=0.72cm]
    \tikzset{gv/.append style={minimum size=6mm, font=\scriptsize\bfseries}}
    \begin{scope}[g2p3c-v1/.style=gvseen, g2p3c-v2/.style=gvseen, g2p3c-v4/.style=gvseen, g2p3c-v6/.style=gvcur]
        \pic {g2p3c-T};
        \node[font=\scriptsize\bfseries] at (2.0,0.9) {t = 3: enter 6};
        \node[font=\scriptsize\ttfamily] at (2.0,-4.5) {stack 1 2 4 6};
    \end{scope}
    \begin{scope}[xshift=4.4cm, g2p3c-v1/.style=gvseen, g2p3c-v2/.style=gvseen, g2p3c-v4/.style=gvdone, g2p3c-v6/.style=gvdone, g2p3c-v5/.style=gvdone, g2p3c-v7/.style=gvcur]
        \pic {g2p3c-T};
        \node[font=\scriptsize\bfseries] at (2.0,0.9) {t = 8: enter 7};
        \node[font=\scriptsize\ttfamily] at (2.0,-4.5) {stack 1 2 7};
    \end{scope}
    \begin{scope}[xshift=8.8cm, g2p3c-v1/.style=gvcur, g2p3c-v2/.style=gvdone, g2p3c-v3/.style=gvdone, g2p3c-v4/.style=gvdone, g2p3c-v6/.style=gvdone, g2p3c-v5/.style=gvdone, g2p3c-v7/.style=gvdone]
        \pic {g2p3c-T};
        \node[font=\scriptsize\bfseries] at (2.0,0.9) {t = 13: leave 1};
        \node[font=\scriptsize\ttfamily] at (2.0,-4.5) {stack (empty)};
    \end{scope}
\end{tikzpicture}
\end{center}
\begin{center}\small
\begin{tabular}{@{}lccccccc@{}}
\toprule
vertex & 1 & 2 & 3 & 4 & 5 & 6 & 7 \\
\midrule
$\tin{}$ & 0 & 1 & 11 & 2 & 6 & 3 & 8 \\
$\tout{}$ & 13 & 10 & 12 & 5 & 7 & 4 & 9 \\
depth & 0 & 1 & 1 & 2 & 2 & 3 & 2 \\
\bottomrule
\end{tabular}
\qquad
\begin{tabular}{@{}ll@{}}
\toprule
\textbf{depth} & \textbf{\texttt{at[d]}} \\
\midrule
0 & 0 \\
1 & 1, 11 \\
2 & 2, 6, 8 \\
3 & 3 \\
\bottomrule
\end{tabular}
\end{center}
\begin{center}\small
\begin{tabular}{@{}ccll@{}}
\toprule
\textbf{query} & \textbf{window} & \textbf{list} & \textbf{answer} \\
\midrule
$(1,2)$ & $[0, 13]$ & $2, 6, 8$: all inside & 3 \\
$(7,2)$ & $[8, 9]$ & only 8 inside & 1 \\
$(4,1)$ & $[2, 5]$ & $1, 11$: none inside & 0 \\
$(5,5)$ & $[6, 7]$ & \texttt{at[5]} empty & 0 \\
\bottomrule
\end{tabular}
\end{center}
\end{tracebox}

\subsubsection*{6. Test suite}

\begin{testbox}
\centering\small
\begin{tabular}{@{}p{6.4cm}lp{5cm}@{}}
\toprule
\textbf{Input} & \textbf{Our output} & \textbf{Why} \\
\midrule
\texttt{2 / 1 / 3 / 1 0 / 1 1 / 2 0} & \texttt{1 1 0} & $D = 0$ counts the root itself; depth above $U$ \\
\texttt{5 / 1 2 3 4 / 3 / 1 4 / 3 2 / 5 4} & \texttt{1 1 1} & a path: $U$ counts itself when $D = \mathrm{dep}[U]$ \\
\texttt{5 / 1 1 1 1 / 3 / 1 1 / 2 1 / 1 4} & \texttt{4 1 0} & a star; a missing depth \\
\texttt{7 / 1 1 2 2 4 2 / 3 / 2 2 / 3 2 / 6 3} & \texttt{3 0 1} & the sample tree, other queries \\
path $P_i = i-1$, $N = Q = 2\cdot10^5$ & (same for 2 and 3) & depth $2\cdot10^5$; iterative only \\
\bottomrule
\end{tabular}
\end{testbox}
All outputs come from running Approach~2. Approaches 1--3 reproduce the official sample and agree on $1000$ random trees with $N \le 10$ and $Q \le 10$ random queries (paths and stars included).

\subsubsection*{7. Common mistakes}

\begin{warningbox}[Mistake 1 --- forgetting $U$ itself]
``Endpoints included'': with $D = \mathrm{dep}[U]$ the answer is 1 ($U$ alone). The window $[\tin{U}, \tout{U}]$ includes $\tin{U}$, so the formula handles it; a ``strict descendants'' implementation answers \texttt{1 0 0} instead of \texttt{1 1 1} on the second test.
\end{warningbox}

\begin{warningbox}[Mistake 2 --- sorting what is already sorted, or sorting by vertex number]
The lists must be sorted by \emph{time}. Building them in a loop over $v = 1..N$ (``\texttt{at[dep[v]].push\_back(tin[v])}'' after the DFS) gives lists ordered by vertex number, and binary search on them is wrong unless you sort. Appending during the DFS needs no sort.
\end{warningbox}

\begin{warningbox}[Mistake 3 --- the offline version counts $U$ before subtracting]
In Approach 3, \texttt{ans -= cnt[D]} must happen \emph{before} \texttt{cnt[dep[U]]++}; otherwise $U$ itself is never counted, and $(7, 2)$ of the sample becomes 0.
\end{warningbox}

\begin{warningbox}[Mistake 4 --- recursion]
$P_i = i - 1$ is allowed: a recursive DFS goes $2\cdot10^5$ deep and crashes on a 2\,MB stack.
\end{warningbox}

\subsubsection*{8. Takeaways}

\begin{itemize}[leftmargin=*]
    \item ``$U$ on the path from $u$ to the root'' is the subtree of $U$: an interval of times.
    \item Subtree $\times$ group $=$ one sorted list of times per group, count inside the window by two binary searches.
    \item Offline: a counter read on entering and leaving is the same count without searching.
    \item Lists filled during the DFS are sorted by time for free.
\end{itemize}

\begin{variantsbox}
\variant{Depth relative to $U$}{\varnone}
  {count the descendants exactly $K$ levels below $U$.}
  {the same query with $D = \mathrm{dep}[U] + K$.}
  {one addition before the lookup; $D > N-1$ gives 0.}
\variant{Count by colour instead of depth}{\varnone}
  {every vertex has a colour; count descendants of $U$ with colour $c$.}
  {group the vertices by colour instead of depth.}
  {lists of entry times per colour, the same two binary searches; or the offline counter per colour (checked against the brute force on random coloured trees).}
\variant{The $k$-th descendant at depth $D$}{\varnone}
  {in DFS order, the $k$-th vertex of depth $D$ below $U$.}
  {the qualifying entries are a contiguous run of $L_D$.}
  {index \texttt{lower\_bound}$(\tin{U}) + k - 1$, valid if it is still $\le \tout{U}$; store the vertex next to its time.}
\variant{Ancestor test only}{\varnone}
  {just ``is $U$ on the root path of $v$''.}
  {no group at all.}
  {two comparisons of times (Exercise E3.1).}
\end{variantsbox}

%============================================================
% CH2 · P3 · PROBLEM 3 — MAZE (CF 377A)
%============================================================
\subsection{Problem 2.3.3: Maze}
\problemlabel{2.3.3}{ssec:ch2-maze}

\problemcard{2.3.3}{Maze}{Codeforces 377A}{https://codeforces.com/problemset/problem/377/A}{\exmedium\ \cf{1600}}{Wall up the first $k$ cells of the DFS finishing order (they form whole subtrees), or keep the first $s-k$ cells of any traversal order (a connected prefix)}

\tikzset{
    g2p3z-e/.style={gcell, minimum size=6.5mm, font=\scriptsize\bfseries},
    g2p3z-w/.style={gwall, minimum size=6.5mm},
    g2p3z-x/.style={gcell, minimum size=6.5mm, fill=dupfill, draw=recursionred, font=\scriptsize\bfseries, text=recursionred},
    g2p3z-c/.style={gcell, minimum size=6.5mm, fill=currentfill, draw=accentorange, line width=1pt, font=\scriptsize\bfseries},
    g2p3z-s/.style={gcell, minimum size=6.5mm, fill=hitfill, font=\scriptsize\bfseries},
    g2p3z-d/.style={gcell, minimum size=6.5mm, fill=computedfill, font=\scriptsize\bfseries},
}
% \gtpz{row}{col}{style}{text}: one cell of a maze drawing (row 0 on top, 0.65 cm per cell)
\newcommand{\gtpz}[4]{\node[#3] at (#2*0.65, -#1*0.65) {#4};}

\subsubsection*{1. Problem statement}

A maze is an $n \times m$ grid of empty cells (\texttt{.}) and walls (\texttt{\#}); moves go between side-adjacent empty cells, and all empty cells form one connected area. Turn \textbf{exactly} $k$ empty cells into walls, marked \texttt{X}, so that the remaining empty cells are still connected. Print the new maze; a solution is guaranteed and any valid one is accepted.

\textbf{Constraints:} $1 \le n, m \le 500$, $0 \le k < s$ where $s$ is the number of empty cells.

\begin{table}[H]
\centering\small
\begin{tabular}{@{}p{5.2cm}p{4.3cm}p{4.3cm}@{}}
\toprule
\textbf{Input} & \textbf{Official output} & \textbf{Our output} \\
\midrule
\texttt{3 4 2 / \#..\# / ..\#. / \#...} & \texttt{\#.X\# / X.\#. / \#...} & \texttt{\#..\# / ..\#X / \#..X} \\
\texttt{5 4 5 / \#... / \#.\#. / .\#.. / ...\# / .\#.\#} & \texttt{\#XXX / \#X\#. / X\#.. / ...\# / .\#.\#} & \texttt{\#... / \#X\#. / X\#.. / X..\# / X\#X\#} \\
\bottomrule
\end{tabular}
\end{table}

\begin{diagrambox}
\begin{tikzpicture}
    \begin{scope}
        \foreach \r/\c/\s in {0/0/w, 0/1/e, 0/2/e, 0/3/w, 1/0/e, 1/1/e, 1/2/w, 1/3/e, 2/0/w, 2/1/e, 2/2/e, 2/3/e}
            \gtpz{\r}{\c}{g2p3z-\s}{}
        \node[font=\scriptsize\bfseries] at (0.98,0.6) {sample 1, $k = 2$};
    \end{scope}
    \begin{scope}[xshift=3.7cm]
        \foreach \r/\c/\s/\t in {0/0/w/, 0/1/e/, 0/2/x/X, 0/3/w/, 1/0/x/X, 1/1/e/, 1/2/w/, 1/3/e/, 2/0/w/, 2/1/e/, 2/2/e/, 2/3/e/}
            \gtpz{\r}{\c}{g2p3z-\s}{\t}
        \node[font=\scriptsize\bfseries] at (0.98,0.6) {official answer};
    \end{scope}
    \begin{scope}[xshift=7.4cm]
        \foreach \r/\c/\s/\t in {0/0/w/, 0/1/e/, 0/2/e/, 0/3/w/, 1/0/e/, 1/1/e/, 1/2/w/, 1/3/x/X, 2/0/w/, 2/1/e/, 2/2/e/, 2/3/x/X}
            \gtpz{\r}{\c}{g2p3z-\s}{\t}
        \node[font=\scriptsize\bfseries] at (0.98,0.6) {our answer};
    \end{scope}
    \node[note, text width=3.6cm, anchor=west, align=left] at (10.6,-0.65) {two different answers, both valid: the 6 remaining empty cells are connected in each};
\end{tikzpicture}
\figcap{Sample 1. Dark cells are walls, red cells the $k$ new walls \texttt{X}. The official answer removes two dead ends at the top; ours removes the tail $(1,3), (2,3)$ on the right.}
\end{diagrambox}

\subsubsection*{2. The graph model}

\begin{modelbox}[Graph Model --- Maze]
\renewcommand{\arraystretch}{1.35}
\begin{tabularx}{\linewidth}{@{}>{\bfseries}l X@{}}
Vertices & the $s$ empty cells \\
Edges & between side-adjacent empty cells, implicit (computed from the grid) \\
Weights & none \\
Question & delete exactly $k$ vertices of a connected graph so that the rest stays connected --- equivalently, choose $s - k$ vertices inducing a connected subgraph \\
Size & $s \le 2.5\cdot10^5$, $|E| < 5\cdot10^5$: one traversal; a DFS can be $1.25\cdot10^5$ cells deep \\
\end{tabularx}
\end{modelbox}

\begin{diagrambox}
\begin{tikzpicture}[x=1.1cm, y=1.1cm]
    \tikzset{gv/.append style={minimum size=8mm, font=\tiny\bfseries}}
    \foreach \r/\c in {0/1, 0/2, 1/0, 1/1, 1/3, 2/1, 2/2, 2/3}
        \node[gv] (g\r\c) at (\c,-\r) {\r,\c};
    \foreach \a/\b in {01/02, 01/11, 10/11, 11/21, 21/22, 22/23, 13/23} \draw[ge] (g\a) -- (g\b);
    \node[note, anchor=west, text width=5.6cm, align=left] at (4.0,-1.0) {the 8 empty cells of sample 1 as a graph (labels: row, column). Deleting a vertex keeps it connected iff the vertex is not a cut vertex; leaves of any spanning tree are never cut vertices.};
\end{tikzpicture}
\figcap{The model: cells are vertices, shared sides are edges. Here the graph happens to be a tree, so the safe cells to remove first are exactly its leaves $(0,2)$, $(1,0)$, $(1,3)$.}
\end{diagrambox}

\subsubsection*{3. Key observations}

\begin{itemize}[leftmargin=*]
    \item \textbf{Greedy removal one by one works} --- a connected graph always has a vertex whose removal keeps it connected (a leaf of a spanning tree) --- but finding such a vertex anew $k$ times is slow.
    \item \textbf{All $k$ at once with the finishing order.} Every prefix of the DFS post-order is a union of complete DFS subtrees; the unfinished cells always contain the root and the tree path from each of them to the root (invariant of the pattern). So the first $k$ cells to finish can be walled up together.
    \item \textbf{The mirror view: keep a connected prefix.} In a BFS (or DFS pre-) order every cell except the start is discovered from an earlier cell, so the first $s - k$ cells are connected. Walling up the \emph{last} $k$ cells of a BFS order is equally correct.
    \item \textbf{Exactly $k$, not at most.} Both rules produce exactly $k$ walls, since $k < s$.
    \item \textbf{Depth.} A snake-shaped $500 \times 500$ maze has $125{,}250$ empty cells in one path: the DFS must be iterative (or use BFS).
\end{itemize}

\subsubsection*{4. Approaches}

\begin{bruteforcebox}
\textbf{Idea: remove one safe cell at a time.} $k$ times, scan the empty cells in row-major order and wall up the first one whose removal leaves the other empty cells connected (checked by a BFS).

\begin{lstlisting}[style=cpp]
#include <bits/stdc++.h>
using namespace std;

// Brute force: k times, try the empty cells in row-major order and wall up the first one
// whose removal leaves the other empty cells connected (checked by a BFS).
// A connected graph always has such a cell (a leaf of any spanning tree), so this never gets stuck.
int n, m;
vector<string> g;

bool connectedWithout(int skip) {
    int start = -1, total = 0;
    for (int i = 0; i < n * m; i++)
        if (i != skip && g[i / m][i % m] == '.') { total++; if (start < 0) start = i; }
    if (total == 0) return true;
    vector<char> seen(n * m, 0);
    queue<int> q;
    q.push(start); seen[start] = 1;
    int reached = 0;
    const int dr[4] = {-1, 1, 0, 0}, dc[4] = {0, 0, -1, 1};
    while (!q.empty()) {
        int c = q.front(); q.pop(); reached++;
        for (int d = 0; d < 4; d++) {
            int nr = c / m + dr[d], nc = c % m + dc[d];
            if (nr < 0 || nr >= n || nc < 0 || nc >= m) continue;
            int x = nr * m + nc;
            if (x == skip || g[nr][nc] != '.' || seen[x]) continue;
            seen[x] = 1; q.push(x);
        }
    }
    return reached == total;
}

int main() {
    int k;
    cin >> n >> m >> k;
    g.resize(n);
    for (auto& row : g) cin >> row;
    for (int t = 0; t < k; t++)
        for (int i = 0; i < n * m; i++)
            if (g[i / m][i % m] == '.' && connectedWithout(i)) { g[i / m][i % m] = 'X'; break; }
    for (auto& row : g) cout << row << "\n";
}
\end{lstlisting}
On sample 1 it walls up $(0,2)$ (removing $(0,1)$ first would isolate $(0,2)$) and then $(0,1)$, printing \texttt{\#XX\# / ..\#. / \#...} --- a third valid answer.

\begin{complexitybox}
\textbf{Time:} up to $k \cdot s$ connectivity checks of \bigO{nm} each; for an open $500 \times 500$ grid that is far beyond $10^{15}$. \textbf{Space:} \bigO{nm}. \textbf{Verdict:} \tle{}; the reference for random tests (together with a checker, since answers are not unique).
\end{complexitybox}
\end{bruteforcebox}

\begin{algorithmbox}[Approach 2: wall up the first $k$ cells to finish (iterative DFS)]
\textbf{Idea.} DFS from any empty cell. Each stack entry remembers the next of its four directions to try (\texttt{dir[c]}); when all four are tried, the cell \emph{finishes}: if fewer than $k$ cells have been walled so far, mark it \texttt{X}.

\begin{lstlisting}[style=cpp]
#include <bits/stdc++.h>
using namespace std;

// Maze (CF 377A): wall up the FIRST k cells in DFS finishing order.
// A finished set is always a union of whole DFS subtrees, so the unfinished rest stays connected.
// Iterative: a snake-shaped maze of 500 x 500 makes the DFS 125,000+ cells deep.
int main() {
    ios::sync_with_stdio(false);
    cin.tie(nullptr);
    int n, m, k;
    cin >> n >> m >> k;
    vector<string> g(n);
    for (auto& row : g) cin >> row;
    int start = -1;
    for (int i = 0; i < n * m && start < 0; i++)
        if (g[i / m][i % m] == '.') start = i;
    const int dr[4] = {-1, 1, 0, 0}, dc[4] = {0, 0, -1, 1};
    vector<char> seen(n * m, 0), dir(n * m, 0);      // dir[c] = next direction to try from c
    vector<int> st(1, start);
    seen[start] = 1;
    int walled = 0;
    while (!st.empty()) {
        int c = st.back(), r = c / m, col = c % m;
        if (dir[c] == 4) {                           // c finishes
            st.pop_back();
            if (walled < k) { g[r][col] = 'X'; walled++; }
            continue;
        }
        int d = dir[c]++;
        int nr = r + dr[d], nc = col + dc[d];
        if (nr < 0 || nr >= n || nc < 0 || nc >= m || g[nr][nc] != '.' || seen[nr * m + nc]) continue;
        seen[nr * m + nc] = 1;
        st.push_back(nr * m + nc);
    }
    for (auto& row : g) cout << row << "\n";
}
\end{lstlisting}

\begin{complexitybox}
\textbf{Time:} \bigO{nm}: each cell is pushed once and tries four directions. \textbf{Space:} \bigO{nm}. \textbf{Verdict:} \ac{}-ready: about $0.2$\,s here on the $500 \times 500$ snake ($125{,}250$ cells, DFS depth $125{,}250$) and on the open $500 \times 500$ grid with $k = s - 1$; the checker accepts both outputs.
\end{complexitybox}
\end{algorithmbox}

\begin{invariantbox}[Why the remaining cells stay connected]
Let $F$ be the set of finished cells at any moment. A cell finishes only after every cell it discovered has finished, so $F$ is closed under taking DFS-tree descendants. Take any unfinished cell $v$ other than the start: its DFS parent finishes after $v$, so it is unfinished too. Repeating, the whole tree path from $v$ to the start avoids $F$. Hence every unfinished cell is joined to the start inside the unfinished set.

The program walls up exactly the cells that finish while \texttt{walled < k}, i.e.\ the first $k$ cells of the finishing order --- the set $F$ at the moment the $k$-th cell finished. The remaining $s-k \ge 1$ cells are unfinished at that moment and therefore connected. Writing \texttt{X} during the search does not disturb it: a cell is marked only after it has finished, and a finished cell is already \texttt{seen}, so it is never pushed again. \hfill$\blacksquare$
\end{invariantbox}

\begin{alternativebox}[Approach 3: keep a connected prefix --- wall up the last $k$ cells of a BFS order]
The other side of the same coin: the first $j$ cells of a BFS order are connected for every $j$ (each was discovered from an earlier one). Keep the first $s - k$, wall up the rest. BFS needs no direction counters and no deep stack; the queue is the order.

\begin{lstlisting}[style=cpp]
#include <bits/stdc++.h>
using namespace std;

// Alternative: keep a connected PREFIX. In BFS order every cell (except the start) is
// discovered from an earlier cell, so the first s - k cells are connected: wall up the last k.
int main() {
    ios::sync_with_stdio(false);
    cin.tie(nullptr);
    int n, m, k;
    cin >> n >> m >> k;
    vector<string> g(n);
    for (auto& row : g) cin >> row;
    int start = -1;
    for (int i = 0; i < n * m && start < 0; i++)
        if (g[i / m][i % m] == '.') start = i;
    const int dr[4] = {-1, 1, 0, 0}, dc[4] = {0, 0, -1, 1};
    vector<char> seen(n * m, 0);
    vector<int> order;                               // BFS order = the queue itself
    order.push_back(start);
    seen[start] = 1;
    for (size_t h = 0; h < order.size(); h++) {
        int r = order[h] / m, c = order[h] % m;
        for (int d = 0; d < 4; d++) {
            int nr = r + dr[d], nc = c + dc[d];
            if (nr < 0 || nr >= n || nc < 0 || nc >= m || g[nr][nc] != '.' || seen[nr * m + nc]) continue;
            seen[nr * m + nc] = 1;
            order.push_back(nr * m + nc);
        }
    }
    int s = order.size();
    for (int i = s - k; i < s; i++) g[order[i] / m][order[i] % m] = 'X';
    for (auto& row : g) cout << row << "\n";
}
\end{lstlisting}

\begin{diagrambox}
\begin{tikzpicture}
    \begin{scope}
        \foreach \r/\c/\s/\t in {0/0/w/, 0/1/d/0, 0/2/d/2, 0/3/w/, 1/0/d/4, 1/1/d/1, 1/2/w/, 1/3/x/7, 2/0/w/, 2/1/d/3, 2/2/d/5, 2/3/x/6}
            \gtpz{\r}{\c}{g2p3z-\s}{\t}
        \node[font=\scriptsize\bfseries] at (0.98,0.6) {BFS order, keep $0..5$};
    \end{scope}
    \begin{scope}[xshift=4.2cm]
        \foreach \r/\c/\s/\t in {0/0/w/, 0/1/d/8, 0/2/d/7, 0/3/w/, 1/0/d/5, 1/1/d/6, 1/2/w/, 1/3/x/1, 2/0/w/, 2/1/d/4, 2/2/d/3, 2/3/x/2}
            \gtpz{\r}{\c}{g2p3z-\s}{\t}
        \node[font=\scriptsize\bfseries] at (0.98,0.6) {DFS finish \#, remove 1, 2};
    \end{scope}
    \node[note, text width=5.6cm, anchor=west, align=left] at (7.6,-0.65) {On sample 1 both rules wall up the same two cells; on sample 2 they differ (\texttt{\#... / \#X\#. / X\#.. / X..\# / X\#X\#} by DFS, \texttt{\#... / \#.\#. / X\#.. / XX.\# / X\#X\#} by BFS), and both pass the checker.};
\end{tikzpicture}
\figcap{The two orders on sample 1. Left: BFS positions; the last $k = 2$ become walls. Right: DFS finishing numbers; the first $k = 2$ become walls.}
\end{diagrambox}

\begin{complexitybox}
\textbf{Time:} \bigO{nm}. \textbf{Space:} \bigO{nm}. \textbf{Verdict:} \ac{}-ready, also about $0.2$\,s on the two $500 \times 500$ tests. The DFS \emph{pre}-order would work just as well as BFS for ``keep a prefix''; the post-order is the natural order for ``remove a prefix''.
\end{complexitybox}
\end{alternativebox}

\begin{tricksbox}[Approach 4: Implementation Refinements]
\begin{itemize}[leftmargin=*]
    \item \textbf{Recursive DFS is a trap here.} The snake maze sends a recursive DFS $125{,}250$ frames deep --- about twice what a 2\,MB stack survives (Lesson~\ref{ssec:ch2-call-stack}). Either the iterative form above, or BFS.
    \item \textbf{One index per cell.} \texttt{c = r * m + col} keeps the stack and the arrays one-dimensional; \texttt{dir} fits in a \texttt{char}.
    \item \textbf{Mark \texttt{X} in place.} After a cell finishes it is never examined again, so writing \texttt{X} immediately is safe; the BFS version writes after the search.
    \item \textbf{Start anywhere.} Any empty cell can be the root; the first in row-major order is simplest.
\end{itemize}
\end{tricksbox}

\subsubsection*{5. Trace}

\begin{tracebox}[{Sample 1, directions tried up, down, left, right; start $(0,1)$, $k = 2$}]
\begin{center}
\begin{tikzpicture}
    \begin{scope}
        \foreach \r/\c/\s/\t in {0/0/w/, 0/1/s/0, 0/2/e/, 0/3/w/, 1/0/e/, 1/1/s/1, 1/2/w/, 1/3/c/5, 2/0/w/, 2/1/s/2, 2/2/s/3, 2/3/s/4}
            \gtpz{\r}{\c}{g2p3z-\s}{\t}
        \node[font=\scriptsize\bfseries] at (0.98,0.6) {enter $(1,3)$, depth 5};
        \draw[primaryblue, line width=1.4pt, -{Stealth[length=2mm]}] (0.65,-0.2) -- (0.65,-0.45);
        \draw[primaryblue, line width=1.4pt, -{Stealth[length=2mm]}] (0.65,-0.85) -- (0.65,-1.1);
        \draw[primaryblue, line width=1.4pt, -{Stealth[length=2mm]}] (0.85,-1.3) -- (1.1,-1.3);
        \draw[primaryblue, line width=1.4pt, -{Stealth[length=2mm]}] (1.5,-1.3) -- (1.75,-1.3);
        \draw[primaryblue, line width=1.4pt, -{Stealth[length=2mm]}] (1.95,-1.1) -- (1.95,-0.85);
    \end{scope}
    \begin{scope}[xshift=3.9cm]
        \foreach \r/\c/\s/\t in {0/0/w/, 0/1/s/, 0/2/e/, 0/3/w/, 1/0/e/, 1/1/s/, 1/2/w/, 1/3/x/X, 2/0/w/, 2/1/s/, 2/2/c/, 2/3/x/X}
            \gtpz{\r}{\c}{g2p3z-\s}{\t}
        \node[font=\scriptsize\bfseries] at (0.98,0.6) {2 finished: $k$ reached};
    \end{scope}
    \begin{scope}[xshift=7.8cm]
        \foreach \r/\c/\s/\t in {0/0/w/, 0/1/d/, 0/2/d/, 0/3/w/, 1/0/d/, 1/1/d/, 1/2/w/, 1/3/x/X, 2/0/w/, 2/1/d/, 2/2/d/, 2/3/x/X}
            \gtpz{\r}{\c}{g2p3z-\s}{\t}
        \node[font=\scriptsize\bfseries] at (0.98,0.6) {end: 6 cells kept};
    \end{scope}
    \node[note, text width=3.4cm, anchor=west, align=left] at (11.0,-0.65) {orange $=$ top of stack, light orange $=$ on the stack, blue $=$ finished and kept, red $=$ walled};
\end{tikzpicture}
\end{center}
\begin{center}\small
\begin{tabular}{@{}clp{5.6cm}l@{}}
\toprule
\textbf{\#} & \textbf{event} & \textbf{cell (entered from)} & \textbf{action} \\
\midrule
1 & enter 0 & $(0,1)$ & start \\
2 & enter 1--5 & $(1,1), (2,1), (2,2), (2,3), (1,3)$, each from the previous & down, down, right, right, up \\
3 & finish 1 & $(1,3)$ & \texttt{X} (walled 1) \\
4 & finish 2 & $(2,3)$ & \texttt{X} (walled 2 $= k$) \\
5 & finish 3, 4 & $(2,2), (2,1)$ & kept \\
6 & enter 6, finish 5 & $(1,0)$ (from $(1,1)$) & kept \\
7 & finish 6 & $(1,1)$ & kept \\
8 & enter 7, finish 7, 8 & $(0,2)$ (from $(0,1)$), then $(0,1)$ & kept \\
\bottomrule
\end{tabular}
\end{center}
Finishing order: $(1,3), (2,3), (2,2), (2,1), (1,0), (1,1), (0,2), (0,1)$. The first two are a whole subtree (the tail hanging below $(2,2)$).
\end{tracebox}

\subsubsection*{6. Test suite}

\begin{testbox}
\centering\small
\begin{tabular}{@{}p{4.8cm}p{3.6cm}p{5.2cm}@{}}
\toprule
\textbf{Input} & \textbf{Our output} & \textbf{Why} \\
\midrule
\texttt{1 1 0 / .} & \texttt{.} & one cell, nothing to do \\
\texttt{1 5 3 / .....} & \texttt{..XXX} & a path: the far end finishes first \\
\texttt{2 2 3 / .. / ..} & \texttt{.X / XX} & a 4-cycle, $k = s - 1$ \\
\texttt{3 3 4 / ... / .\#. / ...} & \texttt{.XX / .\#X / ..X} & a ring: removing a contiguous arc \\
$500 \times 500$ snake, $k = 62625$ & valid (checker) & DFS depth $125{,}250$ \\
$500 \times 500$ open, $k = s - 1$ & valid (checker) & only the start survives \\
\bottomrule
\end{tabular}
\end{testbox}
All outputs come from running Approach~2 (the BFS version prints \texttt{..X / .\#X / .XX} on the ring). Because answers are not unique, a checker verifies every output: exactly $k$ cells changed, each from \texttt{.} to \texttt{X}, nothing else changed, and the remaining \texttt{.} cells connected. Approaches 1--3 pass it on both official samples and on $1000$ random mazes up to $6 \times 6$ (random walls, then only the largest component kept, random $k < s$).

\subsubsection*{7. Common mistakes}

\begin{warningbox}[Mistake 1 --- walling the first $k$ cells \emph{entered}]
Pre-order is the order for keeping, not for removing. On sample 1 the first two cells entered are $(0,1)$ and $(1,1)$; walling them up cuts $(0,2)$ off from everything else.
\end{warningbox}

\begin{warningbox}[Mistake 2 --- any $k$ cells with degree 1]
``Remove dead ends'' is right one cell at a time, but $k$ may exceed the number of dead ends (on an open grid there are none), and removing one dead end creates new ones. The finishing order handles all of that in one pass.
\end{warningbox}

\begin{warningbox}[Mistake 3 --- recursion]
The snake maze is a legal test; recursive DFS crashes at depth $125{,}250$.
\end{warningbox}

\begin{warningbox}[Mistake 4 --- changing walls or printing \texttt{\#}]
Only empty cells may change, and they must be printed as \texttt{X}, not \texttt{\#}. A checker rejects both.
\end{warningbox}

\subsubsection*{8. Takeaways}

\begin{itemize}[leftmargin=*]
    \item Every prefix of a DFS finishing order is a union of complete subtrees: removing it keeps the rest connected.
    \item Every prefix of a BFS (or DFS entry) order is connected: keeping it is the mirror statement.
    \item ``Exactly $k$, any valid answer'' asks for a construction, and a traversal order is often the construction.
    \item Grids up to $500 \times 500$ need an iterative DFS or a BFS.
\end{itemize}

\begin{variantsbox}
\variant{Keep exactly $t$ cells}{\varnone}
  {keep a connected region of exactly $t$ cells that contains a given cell $c$.}
  {start the traversal at $c$.}
  {BFS from $c$, keep the first $t$ cells (or remove the first $s - t$ finishers of a DFS rooted at $c$, which never removes $c$).}
\variant{Remove vertices of a general graph}{\varnone}
  {a connected graph (not a grid); delete $k$ vertices keeping the rest connected.}
  {same model, explicit adjacency lists.}
  {the first $k$ vertices of the finishing order of a DFS (template \texttt{Timestamps::post} or \texttt{finishOrder}).}
\variant{Dismantle a tree leaf by leaf}{\varnone}
  {print an order in which the vertices of a tree can be deleted so that the remaining tree stays connected after every step.}
  {the tree itself.}
  {the post-order of any root is such an order (each prefix is a union of subtrees); the reverse of a BFS order works too.}
\variant{Topological order of a DAG}{\varnone}
  {order the vertices of a DAG so that all edges point forwards.}
  {the finishing order again, now in a directed graph.}
  {reverse the finishing order (pattern template \texttt{finishOrder}; Chapter~5).}
\end{variantsbox}

%============================================================
% CH2 · PATTERN 3 — EXERCISES, HINTS, ANSWERS
%============================================================
\subsection{Pattern 3: Exercises and Going Further}
\label{ssec:ch2-p3-ex}

\begin{patternbox}{Pattern 3 Summary Card}
\renewcommand{\arraystretch}{1.25}
\begin{tabularx}{\linewidth}{@{}>{\bfseries}l X@{}}
Recognise & many queries about ``below $u$'', ``the subtree of $u$'', ``the $k$-th visited from $u$''; a peeling order; vertices numbered so that subtrees are consecutive \\
Record & one iterative DFS: $\tin{}$/$\tout{}$ (one clock), $\mathrm{pos}$/$\mathrm{sz}$/\texttt{order} (pre-order), \texttt{post} (finishing order), depth \\
Ancestor & $\tin{u} \le \tin{v} \wedge \tout{v} \le \tout{u}$ --- also ``which child of $y$ leads to $x$'' \\
Subtree & block \texttt{order[pos[u] .. pos[u]+sz[u]-1]}; with groups: sorted entry times per group + binary search, or the offline counter difference \\
Orders & pre-order / BFS prefixes are connected; post-order prefixes are unions of subtrees; reversed post-order of a DAG is topological \\
Top traps & \texttt{tin} used as an array index; unsorted children; $U$ itself forgotten; recursion at depth $2\cdot10^5$; products of sizes in \texttt{int} \\
\end{tabularx}
\end{patternbox}

\exlegend
\begin{enumerate}[leftmargin=*, label=\textbf{E3.\arabic*}, itemsep=6pt]
    \item \exeasy\ \pconcept \exown{Ancestor queries.} A tree with $n$ vertices is given by its $n-1$ edges and rooted at 1. Answer $q$ queries ``$u$ $v$'': does $u$ lie on the path from the root to $v$ ($u = v$ counts)? ($n, q \le 2\cdot10^5$.)
    \item \exeasy\ \pclassic \exprob{Takahashi Tour}{AtCoder ABC213 D}{https://atcoder.jp/contests/abc213/tasks/abc213_d}. A tree of $N$ cities. Starting in city 1, Takahashi always moves to the unvisited neighbour with the smallest number; if there is none, he goes back to the city he came from (and stops when that is impossible in city 1). Print every city he is in, in order, including city 1 at both ends. ($N \le 2\cdot10^5$.)
    \item \exmedium\ \pconcept \exown{Edge types after the fact.} Run an iterative DFS on a directed graph and store only $\tin{}$, $\tout{}$ and, for every vertex, the id of the edge through which it was entered. Afterwards classify every edge as tree, back, forward or cross --- without colours --- and verify the result against the colour method of Lesson~\ref{ssec:ch2-dfs-tree} on random graphs with self-loops and parallel edges.
    \item \exmedium\ \ptricky \cf{1600}\ \exprob{Kuro and Walking Route}{Codeforces 979C}{https://codeforces.com/problemset/problem/979/C}. A tree with $n$ towns and two special towns $x \ne y$. Count the ordered pairs $(u, v)$, $u \ne v$, whose path does \emph{not} visit $x$ and later $y$. ($n \le 3\cdot10^5$.)
    \item \exmedium\ \precent \cf{1800}\ \exprob{Iris and the Tree}{Codeforces 2006B}{https://codeforces.com/problemset/problem/2006/B}. A rooted tree whose vertices are numbered in DFS order (every subtree is a range of consecutive numbers); edge $(i, p_i)$ has an unknown non-negative weight $t_i$, and $\sum t_i = w$. After each of $n-1$ events ``$t_x = y$'' print $\sum_{i=1}^{n} \max \operatorname{dist}(i, i \bmod n + 1)$, where each maximum is taken independently over the still unknown weights. ($\sum n \le 2\cdot10^5$, $w \le 10^{12}$.)
\end{enumerate}

\subsubsection*{Hints}
\begin{enumerate}[leftmargin=*, label=\textbf{E3.\arabic*}]
    \item ``On the path from the root to $v$'' is ``ancestor of $v$''. One DFS, two comparisons per query.
    \item The walk prints a vertex when it is entered and once more each time a child's subtree is finished. In the iterative DFS these are exactly the push and the moment the top is popped and the parent becomes the top again.
    \item A tree edge is the edge recorded for its head. For every other edge $u \to v$ the intervals of $u$ and $v$ are nested one way, the other way, or disjoint.
    \item Root the tree at $y$. A bad path starts in the subtree of $x$ and ends outside the branch of $y$ that contains $x$. Which child of $y$ is that branch? The answer can exceed $2^{31}$.
    \item The subtree of $x$ is the range $x..x+\mathrm{sz}[x]-1$. A path $i \to i+1$ uses the edge above $x$ iff it enters or leaves that range. How many paths do that? Then think about where to put the unknown weight.
\end{enumerate}

\subsubsection*{Answers}

\begin{solutionbox}[E3.1 Ancestor queries]
$u$ is on the root path of $v$ iff $u$ is an ancestor of $v$ (or $v$ itself): the parenthesis test. The DFS is iterative because the tree may be a path; the parent is skipped since the input is an undirected edge list.
\begin{lstlisting}[style=cpp]
#include <bits/stdc++.h>
using namespace std;

// E3.1: tree rooted at 1, q queries "u v": does u lie on the path from the root to v?
// One iterative DFS (one clock), then two comparisons per query.
int main() {
    ios::sync_with_stdio(false);
    cin.tie(nullptr);
    int n;
    cin >> n;
    vector<vector<int>> adj(n + 1);
    for (int i = 0; i < n - 1; i++) {
        int a, b;
        cin >> a >> b;
        adj[a].push_back(b);
        adj[b].push_back(a);
    }
    vector<int> tin(n + 1), tout(n + 1), par(n + 1, 0), st(1, 1);
    vector<size_t> it(n + 1, 0);
    int timer = 0;
    tin[1] = timer++;
    while (!st.empty()) {
        int u = st.back();
        if (it[u] < adj[u].size()) {
            int v = adj[u][it[u]++];
            if (v == par[u]) continue;
            par[v] = u; tin[v] = timer++;
            st.push_back(v);
        } else {
            tout[u] = timer++;
            st.pop_back();
        }
    }
    int q;
    cin >> q;
    while (q--) {
        int u, v;
        cin >> u >> v;
        cout << (tin[u] <= tin[v] && tout[v] <= tout[u] ? "YES" : "NO") << "\n";
    }
}
\end{lstlisting}
On the tree $T$ of the pattern theory with the queries $(3,9)$, $(2,6)$, $(6,6)$, $(1,7)$, $(8,6)$ it prints \textbf{YES NO YES YES NO}. It agrees with a brute force (BFS parents from the root, climb from $v$ looking for $u$) on $500$ random trees with $n \le 10$ and random labels.
\end{solutionbox}

\begin{solutionbox}[E3.2 Takahashi Tour]
Sort the neighbour lists. The tour is the \textbf{Euler walk} of the DFS tree: a vertex is printed on entry and again after each of its children returns, $2N - 1$ numbers in total. With the iterator-style iterative DFS, ``entering $v$'' is the push, and ``coming back'' is the pop after which the new top is the parent.
\begin{lstlisting}[style=cpp]
#include <bits/stdc++.h>
using namespace std;

// E3.2 Takahashi Tour (ABC213 D): print the vertex when it is entered and again
// every time the walk comes back to it from a child. Neighbours sorted; iterative DFS.
int main() {
    ios::sync_with_stdio(false);
    cin.tie(nullptr);
    int n;
    cin >> n;
    vector<vector<int>> adj(n + 1);
    for (int i = 0; i < n - 1; i++) {
        int a, b;
        cin >> a >> b;
        adj[a].push_back(b);
        adj[b].push_back(a);
    }
    for (auto& l : adj) sort(l.begin(), l.end());
    vector<int> par(n + 1, 0), st(1, 1), walk(1, 1);
    vector<size_t> it(n + 1, 0);
    while (!st.empty()) {
        int u = st.back();
        if (it[u] < adj[u].size()) {
            int v = adj[u][it[u]++];
            if (v == par[u]) continue;
            par[v] = u;
            st.push_back(v);
            walk.push_back(v);                     // entering v
        } else {
            st.pop_back();
            if (!st.empty()) walk.push_back(st.back());   // back in the parent
        }
    }
    for (size_t i = 0; i < walk.size(); i++) cout << walk[i] << (i + 1 < walk.size() ? ' ' : '\n');
}
\end{lstlisting}
Both official samples are reproduced: \texttt{1 2 4 2 1 3 1} and \texttt{1 2 1 3 1 4 1 5 1}. A literal simulation of the statement (scan all neighbours at every step) gives the same walk on $500$ random trees with $N \le 10$; the path of $2\cdot10^5$ cities (walk of $399{,}999$ numbers) runs without stack trouble.
\end{solutionbox}

\begin{solutionbox}[E3.3 Edge types after the fact]
After the DFS, for an edge $i = (u \to v)$:
\begin{center}\small
\begin{tabular}{@{}ll@{}}
\toprule
\textbf{test (in this order)} & \textbf{type} \\
\midrule
$v$ was entered through edge $i$ & tree \\
$\tin{v} \le \tin{u}$ and $\tout{u} \le \tout{v}$ ($v$ ancestor of $u$, or $u = v$) & back \\
$\tin{u} < \tin{v}$ and $\tout{v} < \tout{u}$ ($v$ a proper descendant) & forward \\
otherwise (disjoint intervals; then $\tout{v} < \tin{u}$) & cross \\
\bottomrule
\end{tabular}
\end{center}
\emph{Why it matches the colour method.} When the colour method examines $u \to v$, $u$ is gray. A gray $v$ is an ancestor (back); a black $v$ with later entry is a descendant (forward); a black $v$ with earlier entry is finished and unrelated (cross); a white $v$ becomes a child --- but a \emph{second} parallel edge to the same child, examined later, finds it black and descendant: forward. That is why ``tree'' must be decided by the edge id, not by \texttt{par[v] == u}.
\begin{lstlisting}[style=cpp]
#include <bits/stdc++.h>
using namespace std;

// E3.3: classify the edges of a directed graph AFTER the DFS, from tin/tout and the tree-edge ids
// only, and compare with the colour method of Lesson "The DFS Tree" on random graphs.
// Edge i = (a, b). Types: 0 tree, 1 back, 2 forward, 3 cross.
typedef vector<pair<int, int>> Edges;

// one iterative DFS (outer loop over 1..n, edges in input order); returns tin, tout, treeEdge[v]
void stamps(int n, const Edges& e, vector<int>& tin, vector<int>& tout, vector<int>& treeEdge) {
    vector<vector<int>> out(n + 1);                     // edge ids leaving each vertex
    for (int i = 0; i < (int)e.size(); i++) out[e[i].first].push_back(i);
    tin.assign(n + 1, -1); tout.assign(n + 1, -1); treeEdge.assign(n + 1, -1);
    vector<size_t> it(n + 1, 0);
    int timer = 0;
    for (int s = 1; s <= n; s++) {
        if (tin[s] >= 0) continue;
        tin[s] = timer++;
        vector<int> st(1, s);
        while (!st.empty()) {
            int u = st.back();
            if (it[u] < out[u].size()) {
                int id = out[u][it[u]++], v = e[id].second;
                if (tin[v] < 0) { tin[v] = timer++; treeEdge[v] = id; st.push_back(v); }
            } else { tout[u] = timer++; st.pop_back(); }
        }
    }
}

vector<int> classifyByTimes(int n, const Edges& e) {
    vector<int> tin, tout, treeEdge, type(e.size());
    stamps(n, e, tin, tout, treeEdge);
    for (int i = 0; i < (int)e.size(); i++) {
        int u = e[i].first, v = e[i].second;
        if (treeEdge[v] == i) type[i] = 0;                                  // v was entered through i
        else if (tin[v] <= tin[u] && tout[u] <= tout[v]) type[i] = 1;      // v ancestor of u (or u = v)
        else if (tin[u] < tin[v] && tout[v] < tout[u]) type[i] = 2;        // v descendant of u
        else type[i] = 3;                                                   // disjoint intervals
    }
    return type;
}

// reference: the colour method, classifying each edge at the moment it is examined (recursive, small n)
vector<vector<int>> outR;
vector<int> colR, tinR, typeR;
Edges eR;
int timerR;
void dfsR(int u) {
    colR[u] = 1; tinR[u] = timerR++;
    for (int id : outR[u]) {
        int v = eR[id].second;
        if (colR[v] == 0) { typeR[id] = 0; dfsR(v); }
        else if (colR[v] == 1) typeR[id] = 1;
        else typeR[id] = (tinR[u] < tinR[v]) ? 2 : 3;
    }
    colR[u] = 2; timerR++;
}
vector<int> classifyByColours(int n, const Edges& e) {
    eR = e; outR.assign(n + 1, vector<int>()); colR.assign(n + 1, 0); tinR.assign(n + 1, -1);
    typeR.assign(e.size(), -1); timerR = 0;
    for (int i = 0; i < (int)e.size(); i++) outR[e[i].first].push_back(i);
    for (int s = 1; s <= n; s++) if (!colR[s]) dfsR(s);
    return typeR;
}

int main() {
    const char* name[4] = {"tree", "back", "forward", "cross"};
    // the graph D of Part A, edges in adjacency order
    Edges D = {{1, 2}, {1, 5}, {1, 6}, {2, 3}, {2, 4}, {3, 1}, {4, 3}, {5, 4}, {5, 6}, {6, 7}, {7, 6}, {8, 7}};
    vector<int> t = classifyByTimes(8, D);
    for (int i = 0; i < (int)D.size(); i++) cout << D[i].first << "->" << D[i].second << " " << name[t[i]] << "\n";
    // random check, self-loops and parallel edges included
    mt19937 rng(2024);
    int bad = 0;
    for (int test = 0; test < 20000; test++) {
        int n = rng() % 7 + 1, m = rng() % 13;
        Edges e;
        for (int i = 0; i < m; i++) e.push_back({(int)(rng() % n) + 1, (int)(rng() % n) + 1});
        if (classifyByTimes(n, e) != classifyByColours(n, e)) bad++;
    }
    cout << "random graphs: 20000, mismatches: " << bad << "\n";
}
\end{lstlisting}
\begin{lstlisting}[style=trace]
output: 1->2 tree      1->5 tree      1->6 forward   2->3 tree
        2->4 tree      3->1 back      4->3 cross     5->4 cross
        5->6 tree      6->7 tree      7->6 back      8->7 cross
        random graphs: 20000, mismatches: 0          (one edge per line in the real output)
\end{lstlisting}
The classification of $D$ is exactly the one of Lesson~\ref{ssec:ch2-dfs-tree}, and the $20{,}000$ random graphs ($n \le 7$, up to 12 edges, loops and parallel edges allowed) produce no mismatch.
\end{solutionbox}

\begin{solutionbox}[E3.4 Kuro and Walking Route]
\textbf{Model.} Root the tree at $y$. The path $u \to v$ visits $x$ and \emph{then} $y$ iff $u$ lies in the subtree of $x$ (then the path from $u$ climbs through $x$ towards $y$) and $v$ lies outside the branch of $y$ that contains $x$ (so the path must climb all the way up to $y$). That branch is the child $c$ of $y$ that is an ancestor of $x$ --- one ancestor test per child. Bad pairs: $\mathrm{sz}[x] \cdot (n - \mathrm{sz}[c])$; answer $n(n-1) - \mathrm{sz}[x](n - \mathrm{sz}[c])$.

\textbf{The trap.} $n(n-1)$ reaches $9\cdot10^{10}$: every factor must be \texttt{long long} (the star with centre $x$ gives $(n-1)^2 = 89{,}999{,}400{,}001$ for $n = 3\cdot10^5$).
\begin{lstlisting}[style=cpp]
#include <bits/stdc++.h>
using namespace std;

// E3.4 Kuro and Walking Route (CF 979C). Root the tree at y. A pair (u, v) is bad iff the path
// u -> v meets x and then y: u lies in x's subtree, v lies outside the branch of y that contains x.
// That branch is the child c of y with c an ancestor of x (tin/tout test).
// Answer = n(n-1) - sz[x] * (n - sz[c]), in 64-bit.
int main() {
    ios::sync_with_stdio(false);
    cin.tie(nullptr);
    int n, x, y;
    cin >> n >> x >> y;
    vector<vector<int>> adj(n + 1);
    for (int i = 0; i < n - 1; i++) {
        int a, b;
        cin >> a >> b;
        adj[a].push_back(b);
        adj[b].push_back(a);
    }
    vector<int> tin(n + 1), tout(n + 1), sz(n + 1, 1), par(n + 1, 0), st(1, y);
    vector<size_t> it(n + 1, 0);
    int timer = 0;
    tin[y] = timer++;
    while (!st.empty()) {                       // iterative: n = 3e5 and the tree may be a path
        int u = st.back();
        if (it[u] < adj[u].size()) {
            int v = adj[u][it[u]++];
            if (v == par[u]) continue;
            par[v] = u; tin[v] = timer++;
            st.push_back(v);
        } else {
            tout[u] = timer++;
            if (par[u]) sz[par[u]] += sz[u];
            st.pop_back();
        }
    }
    int c = -1;
    for (int v : adj[y])
        if (tin[v] <= tin[x] && tout[x] <= tout[v]) c = v;   // the branch of y holding x
    long long N = n;
    cout << N * (N - 1) - (long long)sz[x] * (N - sz[c]) << "\n";
}
\end{lstlisting}
Official samples: \textbf{5} and \textbf{4}. A brute force over all ordered pairs (BFS path, check the positions of $x$ and $y$) agrees on $500$ random trees with $n \le 9$. Large checks: the path $1 - 2 - \dots - 3\cdot10^5$ with $x = 1$, $y = n$ gives $89{,}999{,}699{,}999 = n(n-1) - 1$, the star gives the value above.
\end{solutionbox}

\begin{solutionbox}[E3.5 Iris and the Tree]
\textbf{The block.} The numbering is a pre-order, so the subtree of $x$ is the range $B_x = [x, x + \mathrm{sz}[x] - 1]$, with sizes from the loop $i = n..2$ ($p_i < i$). The edge $e_x = (x, p_x)$ lies on the path between $a$ and $b$ iff exactly one of them is in $B_x$. The $n$ paths $i \to i \bmod n + 1$ walk once around the numbers $1, 2, \dots, n, 1$; they cross the boundary of the range $B_x$ exactly twice: path $x - 1$ (entering) and path $x + \mathrm{sz}[x] - 1$ (leaving; if that is $n$, the path $n \to 1$, and $1 \notin B_x$ since $x \ge 2$). So \textbf{every edge lies on exactly two paths}.

\textbf{The maximum.} Let $K$ be the sum of the known weights. A path that still contains an unknown edge can receive the whole rest $w - K$ on it; a path without unknown edges has a fixed length. Summing, the known weights contribute $2K$, and every ``open'' path contributes $w - K$:
\[
    \text{answer} = 2K + (\#\text{open paths}) \cdot (w - K).
\]
Keep, per path, the number of unknown edges on it; an event decrements the two counters of $e_x$, and a counter reaching 0 closes a path. \bigO{n} per test case.
\begin{lstlisting}[style=cpp]
#include <bits/stdc++.h>
using namespace std;

// E3.5 Iris and the Tree (CF 2006B). Vertices are numbered in DFS pre-order, so the subtree of x
// is the block x .. x+sz[x]-1. The edge (x, p_x) lies on exactly two of the n paths i -> i mod n + 1:
// path x-1 (enters the block) and path x+sz[x]-1 (leaves it; index n means n -> 1).
// Answer = 2 * known + (#paths that still have an unknown edge) * (w - known).
int main() {
    ios::sync_with_stdio(false);
    cin.tie(nullptr);
    int T;
    cin >> T;
    while (T--) {
        int n;
        long long w;
        cin >> n >> w;
        vector<int> p(n + 1, 0), sz(n + 1, 1), unknown(n + 1, 0);
        for (int i = 2; i <= n; i++) cin >> p[i];
        for (int i = n; i >= 2; i--) sz[p[i]] += sz[i];          // p_i < i
        for (int x = 2; x <= n; x++) { unknown[x - 1]++; unknown[x + sz[x] - 1]++; }
        int open = 0;                                              // paths with an unknown edge
        for (int i = 1; i <= n; i++) open += (unknown[i] > 0);
        long long known = 0;
        for (int e = 0; e < n - 1; e++) {
            int x;
            long long y;
            cin >> x >> y;
            known += y;
            if (--unknown[x - 1] == 0) open--;
            if (--unknown[x + sz[x] - 1] == 0) open--;
            cout << 2 * known + (long long)open * (w - known) << (e + 2 < n ? ' ' : '\n');
        }
    }
}
\end{lstlisting}
The official sample is reproduced: \texttt{2000000000000}, \texttt{25 18 18}, \texttt{449 302 247 200 200}, \texttt{4585 4473 2681 1567 1454 1322 1094 1022 1022}. A brute force that, after every event, walks each of the $n$ paths by climbing parents and adds ``known on the path $+ (w - K)$ if an unknown edge is on it'' agrees on $500$ random inputs (up to 3 test cases, $n \le 9$, random DFS-numbered trees). On a path of $2\cdot10^5$ vertices with $w = 10^{12}$ the first answer is $1.99999\cdot10^{17}$ --- \texttt{long long} is required.
\end{solutionbox}


%============================================================
\clearpage
\section{Pattern 4: Backtracking over Paths}
\label{sec:ch2-p4}
%============================================================
%============================================================
% CH2 · PATTERN 4 — BACKTRACKING OVER PATHS: THEORY
%============================================================
\subsection{Pattern Theory}
\label{ssec:ch2-p4-theory}

\begin{patternbox}{Pattern 4 at a Glance}
Run DFS, but \textbf{remove the mark of a vertex when you leave it}. The marks now mean ``on the current path'' (the gray vertices), and the search walks through \textbf{every simple path} from the start --- one recursive call per path. Keep whatever the question needs (the path itself, its sum, its set of colours) \textbf{on the path}: add it when you enter a vertex, undo it when you leave. The cost is the number of simple paths, which is exponential: affordable for about $n \le 10$--$15$ vertices or $\le 25$ cells, or when pruning, a cap or a DAG keeps the count small.

\begin{center}
\begin{tikzpicture}[x=1cm, y=1cm]
    \begin{scope}
        \tikzset{g2-u1/.style=gvsrc, g2-u6/.style=gvdst}
        \pic {g2-G3};
        \node[font=\scriptsize\bfseries] at (2.15,-1.75) {$G_3$, start 1, target 6};
    \end{scope}
    \begin{scope}[xshift=6.3cm, yshift=1.3cm,
        tn/.style={rectangle, rounded corners=2pt, draw=darkgray!70, fill=white, font=\scriptsize\ttfamily,
                   inner sep=2pt, minimum height=4.2mm},
        tg/.style={tn, draw=accentorange, fill=answerfill, line width=0.9pt},
        tr/.style={tn, draw=successgreen, fill=basefill}]
        \node[tr] (n1) at (3.6,0) {1};
        \node[tn] (n12) at (1.5,-0.75) {1 2};
        \node[tn] (n14) at (5.7,-0.75) {1 4};
        \node[tn] (n123) at (0.2,-1.5) {1 2 3};
        \node[tn] (n125) at (2.4,-1.5) {1 2 5};
        \node[tn] (n145) at (5.7,-1.5) {1 4 5};
        \node[tn] (n1254) at (1.5,-2.25) {1 2 5 4};
        \node[tg] (n1256) at (3.3,-2.25) {1 2 5 6};
        \node[tn] (n1452) at (4.8,-2.25) {1 4 5 2};
        \node[tg] (n1456) at (6.6,-2.25) {1 4 5 6};
        \node[tn] (n14523) at (4.8,-3.0) {1 4 5 2 3};
        \foreach \a/\b in {n1/n12, n1/n14, n12/n123, n12/n125, n14/n145, n125/n1254, n125/n1256,
                           n145/n1452, n145/n1456, n1452/n14523}
            \draw[darkgray!80, -{Stealth[length=1.6mm]}] (\a) -- (\b);
        \node[note, text width=5.6cm, anchor=north] at (3.6,-3.35) {the search tree: one node per simple path from 1 (11 calls). Vertex 5 appears in two nodes --- marks are removed on exit};
    \end{scope}
\end{tikzpicture}
\end{center}
\begin{recurrencebox}
\centering\small
\texttt{go(u):}\ \ \textcolor{successgreen}{\textbf{enter}}: \texttt{onPath[u] = 1}, add $u$ to the state \quad record the path if it is an answer\\[2pt]
\texttt{for v in adj[u]: if !onPath[v] and promising(v): go(v)}\\[2pt]
\textcolor{accentorange}{\textbf{exit}}: \texttt{onPath[u] = 0}, undo exactly what \emph{enter} did\\[4pt]
time $=$ (number of search-tree nodes) $\times$ (degree) \qquad memory $=$ \bigO{n} (one path)
\end{recurrencebox}
\end{patternbox}

%------------------------------------------------------------
\subsubsection*{From the lessons to the pattern}

Tool~4 of Lesson~\ref{ssec:ch2-toolbox} made a single change to DFS: \texttt{onPath[u] = false} when $u$ is left. Plain DFS visits every \emph{vertex} once and is linear; backtracking visits every \emph{simple path} once and is exponential. The lesson's program printed the two paths $1, 2, 5, 6$ and $1, 4, 5, 6$ of $G_3$; the picture above shows everything it looked at on the way --- the \textbf{search tree}, whose nodes are the simple paths starting at 1 and whose edges are ``extend the path by one vertex''. The call stack at any moment is one root-to-node path of this tree, and the gray vertices are exactly the vertices of the current graph path.

Three things are new compared with the other patterns of this chapter:
\begin{itemize}[leftmargin=*]
    \item \textbf{counting the cost before coding} --- the search tree, not the graph, is what the program walks;
    \item \textbf{state on the path} --- a sum, a set of colours, a list of vertices --- maintained by ``add on enter, undo on exit'', the same discipline as the \texttt{onPath} marks themselves;
    \item \textbf{cutting the tree} --- pruning, a cap on the count, or memoising when two branches have the same future.
\end{itemize}

%------------------------------------------------------------
\subsubsection*{How many simple paths are there?}

\begin{center}
\small
\renewcommand{\arraystretch}{1.2}
\begin{tabular}{@{}lrrrrrr@{}}
\toprule
$k \times k$ grid graph, $k =$ & 1 & 2 & 3 & 4 & 5 & 6 \\
\midrule
paths corner $\to$ opposite corner & 1 & 2 & 12 & 184 & 8\,512 & 1\,262\,816 \\
all simple paths from a corner & 1 & 7 & 79 & 2\,111 & 153\,745 & 31\,811\,177 \\
\bottomrule
\end{tabular}

\smallskip
\begin{tabular}{@{}lrrrrrr@{}}
\toprule
complete graph $K_n$, $n =$ & 4 & 6 & 8 & 10 & 11 & 12 \\
\midrule
simple paths from one vertex & 16 & 326 & 13\,700 & 986\,410 & 9\,864\,101 & 108\,505\,112 \\
\bottomrule
\end{tabular}
\end{center}
Every number was counted by the backtracking program itself (the $6 \times 6$ row took 0.7\,s, $K_{12}$ took 2.3\,s). The complete graph is the worst case for $n$ vertices: a path from a fixed start is an ordered choice of distinct vertices, $\sum_{k} \frac{(n-1)!}{(n-1-k)!} \approx e\,(n-1)!$ paths. So:

\begin{diagrambox}
\begin{tikzpicture}[x=0.95cm, y=0.55cm]
    \draw[->, darkgray] (0,0) -- (0,9.2) node[above, font=\scriptsize]{$\log_{10}$(paths)};
    \foreach \y in {0,2,4,6,8} { \draw[darkgray] (-0.08,\y) -- (0.08,\y); \node[font=\tiny, anchor=east] at (-0.1,\y) {$10^{\y}$}; }
    \foreach \x/\v/\l in {1/1.08/3x3, 2/2.26/4x4, 3/3.93/5x5, 4/6.10/6x6}
        { \fill[primaryblue!60] (\x-0.3,0) rectangle (\x+0.3,\v); \node[font=\tiny, below] at (\x,0) {\l}; }
    \foreach \x/\v/\l in {6/1.20/$K_4$, 7/2.51/$K_6$, 8/4.14/$K_8$, 9/5.99/$K_{10}$, 10/8.04/$K_{12}$}
        { \fill[accentorange!70] (\x-0.3,0) rectangle (\x+0.3,\v); \node[font=\tiny, below] at (\x,0) {\l}; }
    \draw[warningred, dashed, thick] (0,8) -- (10.6,8) node[right, font=\scriptsize, align=left]{$\approx 10^8$ calls:\\2.3\,s here};
    \node[font=\scriptsize, text=primaryblue] at (2.5,7.0) {grid, corner to corner};
    \node[font=\scriptsize, text=accentorange] at (7.2,8.9) {$K_n$, from one vertex};
\end{tikzpicture}
\figcap{The two tables on a log scale. Each step of the grid size multiplies the count by 6 to 150; each extra vertex of $K_n$ multiplies it by $n$.}
\end{diagrambox}

%------------------------------------------------------------
\subsubsection*{When to recognise it (and when not)}

\textbf{Use backtracking when}
\begin{itemize}[leftmargin=*]
    \item the statement asks for \textbf{all} paths, or the best path under a rule that depends on the \emph{whole} path (``no vertex twice'', ``collect each cell's gold once'', ``a word along adjacent cells''), and
    \item the instance is \textbf{tiny}: $n \le 10$--$15$ vertices, at most $25$ cells of a grid, a word of $\le 15$ letters, a 7$\times$7 board with strong pruning; or
    \item the output itself has to be listed (then its size is the lower bound anyway: Problem~\ref{ssec:ch2-all-paths}); or
    \item the answer is \textbf{capped} (``print $\min(K, 10^6)$'') --- every call is one path, so the cap bounds the work (Exercise E4.4).
\end{itemize}
A special, linear case: on a \textbf{tree} there is exactly one path from the root to each vertex, so ``state along the path'' costs nothing extra --- the search tree \emph{is} the tree. The path-stack technique (template 3) is then a plain DFS.

\textbf{Do NOT use it when}
\begin{itemize}[leftmargin=*]
    \item only the \emph{number} of paths, or the best path, in a \textbf{DAG} is wanted: the future of a path depends only on its last vertex, so memoise --- DP on DAGs, \dpbook{12}{DP on DAGs \& Graphs};
    \item the shortest path is wanted: BFS (Chapter~3) or Dijkstra (Chapter~7);
    \item the question is ``visit every vertex'' (Hamiltonian, TSP) with $n \approx 15$--$20$: backtracking is $n!$, bitmask DP over (set, last vertex) is $2^n n^2$ --- \dpbook{13}{Bitmask DP} and Chapter~23;
    \item $n$ is $10^5$: no exponential search survives; look for a tree, a DAG or a DP.
\end{itemize}

%------------------------------------------------------------
\subsubsection*{The framework and its invariant}

\begin{center}
\small
\renewcommand{\arraystretch}{1.25}
\begin{tabularx}{\linewidth}{@{}>{\bfseries}l X@{}}
\toprule
Search tree & node $=$ simple path $s = x_0, x_1, \dots, x_k$; child $=$ the path extended by one neighbour of $x_k$ not on it \\
State & \texttt{onPath[]} plus whatever the answer needs: the vertex list, a running sum, colour counts, a bitmask \\
Enter $u$ & \texttt{onPath[u] = 1}; push $u$; add its contribution (sum $+{=}$ value, \texttt{cnt[colour]++}) \\
Exit $u$ & the exact inverse, in reverse order: \texttt{cnt[colour]--}, sum $-{=}$ value, pop, \texttt{onPath[u] = 0} \\
Answer & checked at every node (``best path anywhere'') or only at goal nodes (``paths to $t$'') \\
Cost & $\#\text{nodes} \times \deg$ time, \bigO{n} memory; the recursion depth is the path length $\le n$ \\
\bottomrule
\end{tabularx}
\end{center}

\begin{invariantbox}[Why backtracking sees every simple path exactly once]
\textbf{Invariant.} When \texttt{go(u)} is called, the call stack holds a simple path $x_0 = s, \dots, x_k = u$; \texttt{onPath} is true exactly for $x_0, \dots, x_k$, and every piece of path state equals its value for this path. When \texttt{go(u)} returns, \texttt{onPath} and the state are back to their values for $x_0, \dots, x_{k-1}$.

\emph{Proof} by induction on the call depth. Entering $u$ sets exactly the mark and the contributions of $u$; each recursive call returns with them unchanged (induction hypothesis); the exit undoes exactly the entry. So between two sibling calls the state is that of $x_0, \dots, x_k$.

\emph{Every path once.} A simple path $s = y_0, \dots, y_m$ is reached by the chain of calls $\texttt{go}(y_0) \to \dots \to \texttt{go}(y_m)$: each $y_{i+1}$ is a neighbour of $y_i$ that is not on the path $y_0..y_i$, so the loop of \texttt{go}$(y_i)$ calls it. It is reached only once, because the sequence of calls \emph{is} the path. Hence the search tree has one node per simple path from $s$. \hfill$\blacksquare$
\end{invariantbox}

%------------------------------------------------------------
\subsubsection*{State along the path: the path stack}

The marks are only the simplest piece of path state. Any quantity that is a ``sum'' of the vertices of the path can be kept the same way, \textbf{in \bigO{1} per step}, instead of being recomputed from the path at every node:

\begin{diagrambox}
\begin{tikzpicture}[x=0.95cm, y=0.95cm]
    \tikzset{gv/.append style={minimum size=6.8mm, font=\scriptsize\bfseries}}
    \node[gvseen] (r1) at (0,0) {1};
    \node[gvseen] (r4) at (1.6,-1) {4}; \node[gvdone] (r2) at (-1.6,-1) {2}; \node[gvdone] (r3) at (0,-1) {3};
    \node[gvdone] (r5) at (-2.3,-2) {5}; \node[gvdone] (r6) at (-0.9,-2) {6};
    \node[gvseen] (r7) at (1.6,-2) {7}; \node[gvcur] (r8) at (1.6,-3) {8};
    \foreach \a/\b in {1/2, 1/3, 2/5, 2/6} \draw[ge] (r\a) -- (r\b);
    \foreach \a/\b in {1/4, 4/7, 7/8} \draw[gtree] (r\a) -- (r\b);
    \foreach \v/\c in {1/a, 2/b, 3/a, 4/c, 5/b, 6/c, 7/a, 8/b}
        \node[font=\tiny\bfseries, text=softpurple, right=0pt of r\v] {\c};
    \node[note, text width=3.6cm] at (-1.6,-3.1) {black: entered and left again --- their colours were removed};
    % the stack
    \begin{scope}[xshift=3.7cm]
        \node[font=\scriptsize\bfseries] at (0.6,0.55) {call stack};
        \foreach \v/\c/\y in {1/a/0, 4/c/-0.75, 7/a/-1.5, 8/b/-2.25}
            \node[draw=memoorange, fill=hitfill, minimum width=1.3cm, minimum height=6.5mm, font=\scriptsize] at (0.6,\y) {\v\ (\c)};
        \draw[flow] (1.5,-2.25) -- (2.3,-2.25) node[right, font=\scriptsize, align=left]{enter 8: \texttt{cnt[b]++}};
        \draw[flow, accentorange] (2.3,-2.85) -- (1.5,-2.85) node[pos=0, right, font=\scriptsize, align=left]{exit 8: \texttt{cnt[b]--}};
    \end{scope}
    \begin{scope}[xshift=8.5cm]
        \node[font=\scriptsize\bfseries] at (0.8,0.55) {counts on the path 1, 4, 7, 8};
        \foreach \c/\k/\x in {a/2/0, b/1/0.8, c/1/1.6}{
            \node[cell] at (\x,-0.3) {\c};
            \node[curcell] at (\x,-1.0) {\k};
        }
        \node[font=\scriptsize, align=left, anchor=north west] at (-0.4,-1.6) {distinct colours $= 3$\\(the program prints 3 for 8)};
    \end{scope}
\end{tikzpicture}
\figcap{The tree $T$ of Lesson~\ref{ssec:ch2-toolbox} with colours a/b/c (template 3 uses $1,2,3$). While the DFS sits at 8, the stack is the root path $1, 4, 7, 8$ and \texttt{cnt[]} describes exactly that path.}
\end{diagrambox}

On a tree this is a \emph{linear} algorithm: every vertex is entered once and left once, so a question like ``for every vertex, does its colour already occur on the path from the root?'' (Exercise E4.5) or ``for every vertex, binary-search the prefix sums of its root path'' (Exercise E4.6) costs \VE{} plus the work per step. In a general graph the same discipline runs inside the exponential search; the state then has to be undone just as carefully as the marks.

%------------------------------------------------------------
\subsubsection*{Cutting the search tree}

Pruning means: at a node, prove that no extension can matter, and return. It is only allowed when the proof is airtight --- a wrong prune loses answers silently.

\begin{diagrambox}
\begin{tikzpicture}[x=0.62cm, y=0.62cm]
    \tikzset{c/.style={gcell, minimum size=6mm, font=\scriptsize}, v/.style={c, fill=hitfill}, cu/.style={c, fill=currentfill, draw=accentorange, line width=0.9pt},
             f/.style={c, fill=basefill}}
    % (a) bound
    \begin{scope}
        \node[font=\scriptsize\bfseries, align=center] at (1.5,1.05) {(a) bound};
        \foreach \x/\y/\s/\t in {0/0/v/2, 1/0/v/3, 2/0/cu/1, 3/0/c/1, 0/-1/c/1, 1/-1/c/2, 2/-1/c/1, 3/-1/c/1}
            \node[\s] at (\x,\y) {\t};
        \node[font=\scriptsize, align=center, text width=3.4cm] at (1.5,-2.5) {sum 6, unvisited 6;\\best so far 12 $\ge 6 + 6$:\\cut};
    \end{scope}
    % (b) split
    \begin{scope}[xshift=4.9cm]
        \node[font=\scriptsize\bfseries] at (1.5,1.05) {(b) the free cells split};
        \foreach \x in {0,1,2,3} \node[v] at (\x,0) {};
        \node[v] at (3,-1) {}; \node[cu] at (3,-2) {};
        \node[f] at (2,-2) {L}; \node[f] at (4,-2) {R};
        \node[gwall, minimum size=6mm] at (3,-3) {};
        \foreach \x/\y in {0/-1, 1/-1, 2/-1, 4/-1, 4/0, 0/-2, 1/-2} \node[c] at (\x,\y) {};
        \draw[flow, accentorange] (3,-1.35) -- (3,-1.65);
        \node[font=\scriptsize, align=center, text width=3.6cm] at (2,-4.75) {moving down, the wall is ahead, L and R are both free: the path can reach only one side};
    \end{scope}
    % (c) parity
    \begin{scope}[xshift=10.2cm]
        \node[font=\scriptsize\bfseries] at (1.5,1.05) {(c) parity};
        \foreach \x in {0,...,3} \foreach \y in {0,-1,-2} {
            \pgfmathtruncatemacro{\p}{mod(\x-\y+10,2)}
            \ifnum\p=0 \node[c, fill=primaryblue!12] at (\x,\y) {}; \else \node[c] at (\x,\y) {}; \fi }
        \node[cu] at (0,0) {s}; \node[c, fill=answerfill, draw=accentorange, line width=0.9pt] at (3,-2) {t};
        \node[font=\scriptsize, align=center, text width=3.6cm] at (1.5,-3.75) {every step changes the colour: a walk $s \to t$ has odd length here, so ``exactly 6 steps'' is impossible};
    \end{scope}
\end{tikzpicture}
\figcap{Three proofs that cut a branch. (a) An optimistic \textbf{bound} --- current value plus everything still available --- cannot beat the best: template 2. (b) A \textbf{dead-end test} --- some free cell can no longer be reached by the rest of the path: Exercise E4.7. (c) A \textbf{parity} argument on the bipartite grid.}
\end{diagrambox}

\begin{itemize}[leftmargin=*]
    \item \textbf{Bounds.} ``value so far $+$ optimistic value of the rest $\le$ best'' $\Rightarrow$ return. The optimistic rest can be crude (the sum of all unvisited values) or sharper (only the values still \emph{reachable} from the current vertex --- Problem~\ref{ssec:ch2-max-gold} measures the difference). Once the best is equal to the global maximum possible, stop the whole search.
    \item \textbf{Dead ends.} The target is no longer reachable from $u$ without crossing the path; or a cell that must still be visited can no longer be reached; or a cell has only one free neighbour left and therefore must be entered \emph{now}. On a DAG, ``can $v$ reach $t$ at all?'' is one reverse search (Pattern~5) and removes whole subgraphs before the search starts.
    \item \textbf{Parity and counting.} In a grid every step flips the colour of a chessboard; a path of exactly $L$ steps from $s$ to $t$ needs $L \equiv \text{dist}(s,t) \pmod 2$. Similar counting arguments: ``the remaining letters of the word occur often enough on the board'' (Exercise E4.2).
    \item \textbf{Order.} Try the most promising neighbour first (biggest value, rarest letter), so that a good ``best'' is known early and the bound bites sooner.
    \item \textbf{Caps.} If the answer is $\min(K, \text{cap})$, stop when the counter reaches the cap. Since every call is a new path, the number of calls is at most the cap.
\end{itemize}

%------------------------------------------------------------
\subsubsection*{When the future does not depend on the path: memoise}

Backtracking pays for the \emph{past}: two different paths that end in the same vertex are two different nodes of the search tree, because their sets of used vertices differ. In a \textbf{DAG} the past does not matter --- a path can never come back to a vertex it has used, so the possible continuations of a path depend only on its last vertex. The search tree then collapses into the DAG itself:

\begin{diagrambox}
\begin{tikzpicture}[x=1.05cm, y=0.9cm]
    \begin{scope}
        \node[gvsrc] (a0) at (0,0) {0}; \node[gv] (a1) at (1.4,0.8) {1}; \node[gv] (a2) at (1.4,-0.8) {2};
        \node[gvcur] (a3) at (2.8,0) {3}; \node[gv] (a4) at (4.2,0.8) {4}; \node[gvdst] (a5) at (4.2,-0.8) {5};
        \foreach \a/\b in {0/1, 0/2, 1/3, 2/3, 3/4, 3/5, 4/5} \draw[gde] (a\a) -- (a\b);
        \node[font=\scriptsize\bfseries] at (2.1,-1.6) {a DAG: 4 paths $0 \leadsto 5$};
    \end{scope}
    \begin{scope}[xshift=6.0cm, yshift=0.9cm, tn/.style={rectangle, rounded corners=2pt, draw=darkgray!70, fill=white,
                  font=\tiny\ttfamily, inner sep=1.5pt}, tz/.style={tn, fill=currentfill, draw=accentorange}]
        \node[tn] (s0) at (1.6,0) {0};
        \node[tn] (s1) at (0.7,-0.6) {01}; \node[tn] (s2) at (2.5,-0.6) {02};
        \node[tz] (s13) at (0.7,-1.2) {013}; \node[tz] (s23) at (2.5,-1.2) {023};
        \node[tn] (s134) at (0.3,-1.8) {0134}; \node[tn] (s135) at (1.1,-1.8) {0135};
        \node[tn] (s234) at (2.1,-1.8) {0234}; \node[tn] (s235) at (2.9,-1.8) {0235};
        \node[tn] (s1345) at (0.3,-2.4) {01345}; \node[tn] (s2345) at (2.1,-2.4) {02345};
        \foreach \a/\b in {s0/s1, s0/s2, s1/s13, s2/s23, s13/s134, s13/s135, s23/s234, s23/s235, s134/s1345, s234/s2345}
            \draw[darkgray!70] (\a) -- (\b);
        \draw[decorate, decoration={brace, amplitude=4pt, mirror}, accentorange] (-0.2,-2.7) -- (3.4,-2.7)
            node[midway, below=4pt, font=\scriptsize, text=accentorange, align=center]{the two subtrees under 3 are identical};
    \end{scope}
    \begin{scope}[xshift=11.0cm]
        \node[font=\scriptsize\bfseries] at (1.1,1.35) {memo: paths from $v$ to 5};
        \foreach \v/\c/\x/\y in {0/4/0/0, 1/2/0.9/0.6, 2/2/0.9/-0.6, 3/2/1.8/0, 4/1/2.7/0.6, 5/1/2.7/-0.6}{
            \node[gv, minimum size=6mm, font=\scriptsize\bfseries] (m\v) at (\x,\y) {\v};
            \node[gdist, above=0pt of m\v] {\c};
        }
        \foreach \a/\b in {0/1, 0/2, 1/3, 2/3, 3/4, 3/5, 4/5} \draw[gde] (m\a) -- (m\b);
        \node[font=\scriptsize, align=center] at (1.35,-1.55) {$P(v) = \sum_{v \to w} P(w)$, $P(5) = 1$};
    \end{scope}
\end{tikzpicture}
\figcap{Left: a DAG. Middle: its search tree from 0 --- vertex 3 is reached by two paths and everything below it is computed twice. Right: memoising on the last vertex merges the copies; the count is a DP over the DAG in \VE.}
\end{diagrambox}

\begin{itemize}[leftmargin=*]
    \item \textbf{Count or optimise on a DAG:} memoise on the last vertex --- \VE{}, even when the number of paths is astronomically large (\dpbook{12}{DP on DAGs \& Graphs}). Listing the paths still costs their total length; Problem~\ref{ssec:ch2-all-paths} compares both.
    \item \textbf{General graph, ``which vertices are used'' matters:} the honest state is (set of used vertices, last vertex) --- $2^n \cdot n$ states. For $n \le 20$ this bitmask DP beats $n!$ backtracking by far (\dpbook{13}{Bitmask DP}); for a grid with 25 gold cells it does not help, because almost every (set, last) pair is reached by only one path (Problem~\ref{ssec:ch2-max-gold}, Approach~3).
    \item \textbf{Trees:} one path per vertex --- nothing to memoise, the path stack is already linear.
\end{itemize}

%------------------------------------------------------------
\subsubsection*{Variations map}

\begin{diagrambox}
\begin{tikzpicture}[
    var/.style={rectangle, rounded corners=2pt, draw=patternborder, fill=patternbg,
                font=\scriptsize, align=center, text width=3.25cm, minimum height=1.3cm, inner sep=3pt}]
    \node[var, fill=answerfill, draw=accentorange, text width=5.6cm, font=\scriptsize\bfseries] (root) at (0,0)
        {un-mark on exit: every simple path; state added on enter, undone on exit};
    \node[var] (v1) at (-5.4,-2) {\textbf{list all paths $s \to t$}\\in a DAG, no marks needed\\(All Paths, 2.4.1)};
    \node[var] (v2) at (-1.8,-2) {\textbf{best path, any start}\\grid, temporary zeroing\\(Max Gold, 2.4.2)};
    \node[var] (v3) at (1.8,-2)  {\textbf{count paths through all}\\Hamiltonian, $n \le 8$\\(E4.1)};
    \node[var] (v4) at (5.4,-2)  {\textbf{a word along a path}\\letter-by-letter pruning\\(E4.2)};
    \node[var] (v5) at (-5.4,-3.6) {\textbf{max-weight simple path}\\$n \le 10$, every start\\(E4.3)};
    \node[var] (v6) at (-1.8,-3.6) {\textbf{count with a cap}\\$\min(K, 10^6)$, $n = 2\cdot10^5$\\(E4.4)};
    \node[var] (v7) at (1.8,-3.6)  {\textbf{path stack on a tree}\\colour counts, prefix sums\\(E4.5, E4.6)};
    \node[var] (v8) at (5.4,-3.6)  {\textbf{heavy pruning}\\$7 \times 7$ grid, 48 steps\\(E4.7)};
    \node[draw=patternborder, dashed, rounded corners=4pt, fit=(v1)(v4)(v5)(v8), inner sep=5pt] (grp) {};
    \draw[flow] (root.south) -- (grp.north);
    \node[note, below=2mm of grp] {count or optimise on a DAG: DP (DP book Ch.~12) $\cdot$ ``visit all'' with $n \approx 20$: bitmask DP (DP book Ch.~13, Chapter~23)};
\end{tikzpicture}
\end{diagrambox}

\textbf{Twists and how to handle them}
\begin{itemize}[leftmargin=*]
    \item \textbf{Any start, any end:} loop over every start vertex; record the answer at \emph{every} node, not only at a target.
    \item \textbf{Grid instead of graph:} mark a used cell by overwriting it (\texttt{0}, \texttt{'\#'}) and restore the old value on exit --- no separate \texttt{onPath} array (Problem~\ref{ssec:ch2-max-gold}).
    \item \textbf{Stop at the target:} ``paths from $s$ to $t$'' must not continue \emph{through} $t$; a path that reaches $t$ early may also be a failure (``exactly 48 steps'').
    \item \textbf{Undirected graph:} a path and its reverse are the same path; if they must be counted once, count ordered paths and divide by 2, or fix an orientation (e.g.\ start $<$ end).
    \item \textbf{Deep recursion:} the depth is only the path length, $\le n$ --- harmless for $n \le 25$, but a capped count on $2\cdot10^5$ vertices can build a path of $2\cdot10^5$ vertices: write it iteratively and reset the iterator of a vertex when it is left (Exercise E4.4).
\end{itemize}

%------------------------------------------------------------
\subsubsection*{Code templates}

\begin{lstlisting}[style=cpp]
#include <bits/stdc++.h>
using namespace std;
typedef long long ll;

// (1) ALL SIMPLE PATHS s -> t (any graph). onPath[] = the gray vertices = the current path.
//     limit: stop after that many paths (the answer may be exponential).
struct AllPaths {
    const vector<vector<int>>& adj;
    vector<char> onPath;
    vector<int> path;
    vector<vector<int>> out;
    int t;
    size_t limit;
    AllPaths(const vector<vector<int>>& g) : adj(g), onPath(g.size(), 0) {}
    void go(int u) {
        if (out.size() >= limit) return;              // cap reached: stop the whole search
        onPath[u] = 1; path.push_back(u);             // enter: extend the path
        if (u == t) out.push_back(path);              // a simple path ends here (do not continue past t)
        else
            for (int v : adj[u])
                if (!onPath[v]) go(v);
        onPath[u] = 0; path.pop_back();               // exit: undo exactly what enter did
    }
    vector<vector<int>> run(int s, int target, size_t lim) {
        t = target; limit = lim; out.clear();
        go(s);
        return out;
    }
};

// (2) BEST SIMPLE PATH FROM ANY START, vertex values val[v] >= 0, with an optimistic bound:
//     if (sum so far) + (sum of all values not on the path) <= best, no extension can win.
struct BestPath {
    const vector<vector<int>>& adj;
    const vector<int>& val;
    vector<char> onPath;
    ll best = 0, rest = 0, calls = 0;                 // rest = total value of vertices NOT on the path
    BestPath(const vector<vector<int>>& g, const vector<int>& w) : adj(g), val(w), onPath(g.size(), 0) {}
    void go(int u, ll sum) {
        calls++;
        onPath[u] = 1; sum += val[u]; rest -= val[u]; // enter
        best = max(best, sum);
        if (sum + rest > best)                        // prune: even taking everything left cannot beat best
            for (int v : adj[u])
                if (!onPath[v]) go(v, sum);
        onPath[u] = 0; rest += val[u];                // exit
    }
    ll run() {
        for (int x : val) rest += x;
        for (int s = 0; s < (int)adj.size(); s++) go(s, 0);
        return best;
    }
};

// (3) PATH-STACK DFS ON A TREE: the root path is the state. cnt[c] = how often colour c is on the
//     path root..u; add at enter, remove at exit. Iterative (trees can be 1e5+ deep).
//     Output: distinct[u] = number of different colours on the path root..u.
vector<int> distinctOnRootPath(const vector<vector<int>>& adj, const vector<int>& col, int root) {
    int n = adj.size(), cur = 0;
    vector<int> cnt(*max_element(col.begin(), col.end()) + 1, 0), distinct(n, 0), par(n, -1);
    vector<size_t> it(n, 0);
    vector<int> st;
    st.push_back(root); par[root] = root;
    if (cnt[col[root]]++ == 0) cur++;                 // enter root
    distinct[root] = cur;
    while (!st.empty()) {
        int u = st.back();
        if (it[u] < adj[u].size()) {
            int v = adj[u][it[u]++];
            if (v == par[u]) continue;
            par[v] = u;
            if (cnt[col[v]]++ == 0) cur++;            // enter v: its colour joins the path
            distinct[v] = cur;
            st.push_back(v);
        } else {
            if (--cnt[col[u]] == 0) cur--;            // exit u: its colour leaves the path
            st.pop_back();
        }
    }
    return distinct;
}

int main() {
    // G3 of the lessons, 0-indexed: edges 1-2 1-4 2-3 2-5 4-5 5-6 (vertex 7 isolated)
    vector<vector<int>> G3(7);
    int E[6][2] = {{1, 2}, {1, 4}, {2, 3}, {2, 5}, {4, 5}, {5, 6}};
    for (auto& e : E) { G3[e[0] - 1].push_back(e[1] - 1); G3[e[1] - 1].push_back(e[0] - 1); }
    AllPaths ap(G3);
    for (auto& p : ap.run(0, 5, 1000)) {               // 1 2 5 6 / 1 4 5 6
        for (int x : p) cout << x + 1 << " ";
        cout << "\n";
    }
    cout << ap.run(0, 5, 1).size() << "\n";            // 1 (capped)

    vector<int> val = {3, 1, 4, 1, 5, 9, 0};           // G3 with vertex values (7 is isolated)
    BestPath bp(G3, val);
    ll best = bp.run();
    cout << best << " " << bp.calls << "\n";           // 23 32   (63 calls without the bound)

    // the tree T of Lesson 5, 0-indexed, colours of 1..8
    vector<vector<int>> T(8);
    int F[7][2] = {{1, 2}, {1, 3}, {1, 4}, {2, 5}, {2, 6}, {4, 7}, {7, 8}};
    for (auto& e : F) { T[e[0] - 1].push_back(e[1] - 1); T[e[1] - 1].push_back(e[0] - 1); }
    vector<int> col = {1, 2, 1, 3, 2, 3, 1, 2};
    for (int d : distinctOnRootPath(T, col, 0)) cout << d << " ";   // 1 2 1 2 2 3 2 3
    cout << "\n";
}
\end{lstlisting}

On the values $3, 1, 4, 1, 5, 9$ of $G_3$ the best path is $3, 2, 1, 4, 5, 6$ with value 23, which uses the whole component; once it is found, $\text{sum} + \text{rest} = 23$ never beats it again, and the bound halves the search from 63 to 32 calls.

\begin{insightbox}[Which template when?]
\renewcommand{\arraystretch}{1.25}
\begin{tabularx}{\linewidth}{@{}>{\bfseries}l X@{}}
(1) \texttt{AllPaths} & list or count the simple paths between two vertices; with a \texttt{limit} when the count is capped; in a DAG the \texttt{onPath} test can be dropped (Problem~\ref{ssec:ch2-all-paths}) \\
(2) \texttt{BestPath} & the best simple path anywhere (max value, max length); every node is a candidate; add the strongest bound you can prove \\
(3) \texttt{distinctOnRootPath} & any ``for every vertex of a tree, something about its root path'' --- linear, iterative; replace the colour counts by a sum, a stack of prefix sums, a set \\
\end{tabularx}
\end{insightbox}

%------------------------------------------------------------
\subsubsection*{Common mistakes}

\begin{warningbox}[Mistake 1 --- forgetting to un-mark]
\begin{lstlisting}[style=snippet]
void go(int u) { onPath[u] = true; for (int v : adj[u]) if (!onPath[v]) go(v); }        // plain DFS
void go(int u) { onPath[u] = true; for (int v : adj[u]) if (!onPath[v]) go(v); onPath[u] = false; }
\end{lstlisting}
The first version is a correct DFS and a wrong backtracking: on $G_3$ it prints only $1, 2, 5, 6$ --- vertex 5 stays marked and the path $1, 4, 5, 6$ is never tried.
\end{warningbox}

\begin{warningbox}[Mistake 2 --- undoing on only one exit]
An early \texttt{return} (``target reached'', ``pruned'') that skips the un-mark leaves a vertex gray forever. Put the undo at the single end of the function, or undo before every \texttt{return}. In a grid that is marked by overwriting, restoring the wrong value (\texttt{1} instead of the saved gold) is the same bug.
\end{warningbox}

\begin{warningbox}[Mistake 3 --- continuing through the target]
For ``paths from $s$ to $t$'' a path that reaches $t$ is finished. If the search goes on, it also counts walks like $s \dots t \dots t'$ that pass through $t$ --- in a DAG this does no harm only if $t$ is a sink.
\end{warningbox}

\begin{warningbox}[Mistake 4 --- backtracking where a DP is needed]
``Count the paths from 1 to $n$ in a DAG with $n = 10^5$'' has up to $2^{n-2}$ paths --- no search can list them. The future depends only on the last vertex: memoise (and take the count modulo the given prime).
\end{warningbox}

\begin{warningbox}[Mistake 5 --- a prune that is only ``usually'' right]
``Start only at cells with few neighbours'', ``never turn twice in a row'' --- such rules can silently drop the optimum. Prove them, or test them against the unpruned search on thousands of random inputs before trusting them; Problem~\ref{ssec:ch2-max-gold} shows a rule that passed every test we could think of but has no proof, and is therefore kept optional.
\end{warningbox}

\begin{checklistbox}[``I'm stuck'' checklist for Pattern 4]
\begin{itemize}[leftmargin=*]
    \item[$\square$] How many simple paths can there be in the worst case? Count them (by the formula, or by running the search on the worst input) before coding.
    \item[$\square$] Is the graph a DAG or a tree? Then the past does not matter: no marks needed (DAG) or one path per vertex (tree) --- maybe memoise instead.
    \item[$\square$] What does the answer depend on --- only the last vertex, or the set of used vertices? That decides between DP on the last vertex, bitmask DP and backtracking.
    \item[$\square$] Is every piece of state undone on every way out of the function?
    \item[$\square$] Which bound, dead-end test or parity argument can cut a branch, and can I prove it?
    \item[$\square$] Is the answer capped, or does the search stop at the first success? Then stop the whole recursion, not just the current call.
    \item[$\square$] Can the path be $10^5$ long? Then write the search iteratively.
\end{itemize}
\end{checklistbox}

%============================================================
% CH2 · P4 · PROBLEM 1 — ALL PATHS FROM SOURCE TO TARGET (LeetCode 797)
%============================================================
\subsection{Problem 2.4.1: All Paths From Source to Target}
\problemlabel{2.4.1}{ssec:ch2-all-paths}

\problemcard{2.4.1}{All Paths From Source to Target}{LeetCode 797}{https://leetcode.com/problems/all-paths-from-source-to-target/}{\exeasy\ \lc{Medium}\ (the first backtracking program; the output itself is exponential)}{Backtracking with a path stack; in a DAG no \texttt{visited} array is needed at all; alternative: memoised path lists per vertex}

\subsubsection*{1. Problem statement}

A directed acyclic graph (DAG) has $n$ vertices $0..n-1$; \texttt{graph[i]} lists the vertices $j$ with an edge $i \to j$. Return \textbf{all} paths from vertex $0$ to vertex $n-1$, in any order.

\textbf{Constraints:} $2 \le n \le 15$; $0 \le \texttt{graph[i][j]} < n$, $\texttt{graph[i][j]} \ne i$; the entries of one list are distinct; the graph is guaranteed to be a DAG.

\begin{table}[H]
\centering\small
\begin{tabular}{@{}lll@{}}
\toprule
\textbf{graph} & \textbf{Official output} & \textbf{Our output (backtracking)} \\
\midrule
\texttt{[[1,2],[3],[3],[]]} & \texttt{[[0,1,3],[0,2,3]]} & \texttt{[[0,1,3],[0,2,3]]} \\
\texttt{[[4,3,1],[3,2,4],[3],[4],[]]} & \texttt{[[0,4],[0,3,4],[0,1,3,4],} & identical, same order \\
 & \texttt{\ [0,1,2,3,4],[0,1,4]]} & \\
\bottomrule
\end{tabular}
\end{table}

\begin{diagrambox}
\begin{tikzpicture}[x=1.15cm, y=1cm]
    \begin{scope}
        \node[gvsrc] (a0) at (0,0) {0}; \node[gv] (a1) at (1.5,0.8) {1};
        \node[gv] (a2) at (1.5,-0.8) {2}; \node[gvdst] (a3) at (3,0) {3};
        \foreach \a/\b in {0/1, 0/2, 1/3, 2/3} \draw[gde] (a\a) -- (a\b);
        \node[font=\scriptsize\bfseries] at (1.5,-1.75) {Example 1: 2 paths};
    \end{scope}
    \begin{scope}[xshift=5.6cm]
        \node[gvsrc] (b0) at (0,0) {0}; \node[gv] (b1) at (1.4,1.2) {1}; \node[gv] (b2) at (2.9,1.2) {2};
        \node[gv] (b3) at (2.9,-0.6) {3}; \node[gvdst] (b4) at (4.3,0) {4};
        \draw[gde] (b0) -- (b1); \draw[gde] (b0) -- (b3); \draw[gde] (b0) to[bend right=50] (b4);
        \draw[gde] (b1) -- (b2); \draw[gde] (b1) -- (b3); \draw[gde] (b1) -- (b4);
        \draw[gde] (b2) -- (b3); \draw[gde] (b3) -- (b4);
        \node[font=\scriptsize\bfseries] at (2.2,-1.75) {Example 2: 5 paths};
    \end{scope}
\end{tikzpicture}
\figcap{The two examples. In Example 2 vertex 3 is reached along three different paths ($0 3$, $0 1 3$, $0 1 2 3$), vertex 4 along five.}
\end{diagrambox}

\subsubsection*{2. The graph model}

\begin{modelbox}[Graph Model --- All Paths From Source to Target]
\renewcommand{\arraystretch}{1.35}
\begin{tabularx}{\linewidth}{@{}>{\bfseries}l X@{}}
Vertices & the $n \le 15$ given vertices \\
Edges & directed, given as adjacency lists; acyclic; no self-loops, no parallel edges \\
Weights & none \\
Question & list every path $0 \leadsto n-1$; in a DAG every path is automatically simple \\
Size & at most $\binom{15}{2} = 105$ edges --- the graph is tiny; the \textbf{output} is the problem: up to $2^{n-2} = 8192$ paths with $69\,632$ numbers in total \\
\end{tabularx}
\end{modelbox}

\subsubsection*{3. Key observations}

\begin{itemize}[leftmargin=*]
    \item \textbf{No marks needed.} A path in a DAG can never return to a vertex it has used: that would close a cycle. So the \texttt{onPath} test of template 1 is always false and can be dropped; only the path stack is needed.
    \item \textbf{A path is its set of middle vertices.} Fix a topological order. The vertices of any path appear in that order, so a path is determined by \emph{which} of the $n-2$ middle vertices it uses. Hence there are at most $2^{n-2}$ paths --- and the \textbf{complete DAG} (an edge $i \to j$ for all $i < j$) has exactly that many, because every subset in increasing order is a path.
    \item \textbf{Output-sensitive.} The total output size is at most $2^{13} \cdot 15$ numbers; any correct algorithm must spend at least that. The goal is to spend not much more.
    \item \textbf{Dead ends.} A branch into vertices that cannot reach $n-1$ produces no output but costs time; a reverse search from $n-1$ (Pattern~5) finds them in advance.
\end{itemize}

\begin{diagrambox}
\begin{tikzpicture}[x=1.25cm, y=1cm]
    \foreach \i in {0,...,4} \node[gv] (c\i) at (\i,0) {\i};
    \node[gvsrc] at (c0) {0}; \node[gvdst] at (c4) {4};
    \foreach \i/\j in {0/1, 1/2, 2/3, 3/4} \draw[gde] (c\i) -- (c\j);
    \foreach \i/\j/\b in {0/2/30, 1/3/30, 2/4/30} \draw[gde] (c\i) to[bend left=\b] (c\j);
    \foreach \i/\j/\b in {0/3/38, 1/4/38} \draw[gde] (c\i) to[bend right=\b] (c\j);
    \draw[gde] (c0) to[bend right=55] (c4);
    \node[font=\scriptsize\bfseries] at (2,-2.3) {complete DAG, $n = 5$};
    \begin{scope}[xshift=7.2cm]
        \node[font=\scriptsize\bfseries] at (1.5,1.05) {middle set $\to$ path};
        \foreach \s/\p/\y in {{\{\}}/0 4/0.4, {\{1\}}/0 1 4/-0.05, {\{2\}}/0 2 4/-0.5, {\{3\}}/0 3 4/-0.95,
                              {\{1,2\}}/0 1 2 4/-1.4, {\{1,3\}}/0 1 3 4/-1.85, {\{2,3\}}/0 2 3 4/-2.3, {\{1,2,3\}}/0 1 2 3 4/-2.75}{
            \node[font=\scriptsize, anchor=east] at (0.9,\y) {$\s$};
            \node[font=\scriptsize\ttfamily, anchor=west] at (1.2,\y) {\p};
        }
        \draw[->, darkgray] (0.95,-1.2) -- (1.15,-1.2);
    \end{scope}
\end{tikzpicture}
\figcap{Why $2^{n-2}$ is the maximum and is attained: in the complete DAG each of the $2^3 = 8$ subsets of $\{1,2,3\}$, read in increasing order, is a path from 0 to 4.}
\end{diagrambox}

\subsubsection*{4. Approaches}

\begin{bruteforcebox}
\textbf{Idea: try every candidate set.} By the second observation it is enough to try every subset of the middle vertices: sort it in a topological order (computed by Kahn's algorithm, Chapter~5), put 0 in front and $n-1$ at the end, and keep it if every consecutive pair is an edge.

\begin{lstlisting}[style=cpp]
#include <bits/stdc++.h>
using namespace std;

// Brute force: every subset of the middle vertices, in topological order, checked edge by edge.
class Solution {
public:
    vector<vector<int>> allPathsSourceTarget(vector<vector<int>>& graph) {
        int n = graph.size();
        vector<int> indeg(n, 0), order;                // Kahn's algorithm for a topological order
        for (auto& l : graph) for (int v : l) indeg[v]++;
        queue<int> q;
        for (int v = 0; v < n; v++) if (indeg[v] == 0) q.push(v);
        while (!q.empty()) {
            int u = q.front(); q.pop();
            order.push_back(u);
            for (int v : graph[u]) if (--indeg[v] == 0) q.push(v);
        }
        vector<vector<char>> edge(n, vector<char>(n, 0));
        for (int u = 0; u < n; u++) for (int v : graph[u]) edge[u][v] = 1;
        vector<int> mid;
        for (int v : order) if (v != 0 && v != n - 1) mid.push_back(v);
        int k = mid.size();
        vector<vector<int>> res;
        for (int mask = 0; mask < (1 << k); mask++) {
            vector<int> p(1, 0);
            for (int i = 0; i < k; i++) if (mask >> i & 1) p.push_back(mid[i]);
            p.push_back(n - 1);
            bool ok = true;
            for (size_t i = 0; i + 1 < p.size() && ok; i++) ok = edge[p[i]][p[i + 1]];
            if (ok) res.push_back(p);
        }
        return res;
    }
};

int main() {
    vector<vector<int>> g = {{4, 3, 1}, {3, 2, 4}, {3}, {4}, {}};
    for (auto& p : Solution().allPathsSourceTarget(g)) {
        for (int x : p) cout << x << " ";
        cout << "/ ";
    }
    cout << "\n";                                      // 0 4 / 0 1 4 / 0 3 4 / 0 1 3 4 / 0 1 2 3 4 /
}
\end{lstlisting}

\begin{complexitybox}
\textbf{Time:} $2^{n-2}$ subsets, \bigO{n} each: $8192 \cdot 15 \approx 1.2 \cdot 10^5$ steps for $n = 15$ --- fast, because $n$ is so small. But it pays $2^{n-2}$ even when the graph has a single path. \textbf{Space:} \bigO{n^2} plus the output. \textbf{Verdict:} \ac{}-fast for $n \le 15$ and an ideal reference for testing; useless as soon as $n$ grows.
\end{complexitybox}
\end{bruteforcebox}

\begin{algorithmbox}[Approach 2: backtracking with a path stack]
\textbf{Idea.} Template 1 of the pattern without the marks: push $u$, and if $u = n-1$ record a copy of the path, otherwise recurse into every successor; pop $u$ on exit.

\begin{lstlisting}[style=cpp]
#include <bits/stdc++.h>
using namespace std;

// All Paths From Source to Target (LeetCode 797): backtracking. In a DAG a path can never
// revisit a vertex, so no visited/onPath array is needed -- only the path stack.
class Solution {
    vector<vector<int>> res;
    vector<int> path;
    void dfs(const vector<vector<int>>& g, int u) {
        path.push_back(u);                             // enter u
        if (u == (int)g.size() - 1) res.push_back(path);
        else
            for (int v : g[u]) dfs(g, v);
        path.pop_back();                               // exit u
    }
public:
    vector<vector<int>> allPathsSourceTarget(vector<vector<int>>& graph) {
        res.clear();
        dfs(graph, 0);
        return res;
    }
};

int main() {
    vector<vector<int>> g1 = {{1, 2}, {3}, {3}, {}};
    vector<vector<int>> g2 = {{4, 3, 1}, {3, 2, 4}, {3}, {4}, {}};
    for (auto* g : {&g1, &g2}) {
        for (auto& p : Solution().allPathsSourceTarget(*g)) {
            for (int x : p) cout << x << " ";
            cout << "/ ";
        }
        cout << "\n";
    }
    // 0 1 3 / 0 2 3 /
    // 0 4 / 0 3 4 / 0 1 3 4 / 0 1 2 3 4 / 0 1 4 /
}
\end{lstlisting}

\begin{complexitybox}
\textbf{Time:} one call per path from 0 to \emph{any} vertex, plus a copy of every finished path. On the complete DAG with $n = 15$: $16\,384$ calls ($= 2^{14}$: there are $2^{v-1}$ paths from 0 to each $v \ge 1$, plus the root) and $8192$ copies of total length $69\,632$, in 3\,ms. \textbf{Space:} \bigO{n} for the stack and the path, plus the output. \textbf{Verdict:} \ac{}; within a factor 2 of the output size on the worst case.
\end{complexitybox}
\end{algorithmbox}

\begin{invariantbox}[Why it is correct]
\textbf{Invariant.} When \texttt{dfs(u)} starts, \texttt{path} holds a path $0 = x_0 \to x_1 \to \dots \to x_k = u$ of the graph; when it returns, \texttt{path} is $x_0, \dots, x_{k-1}$ again. (Push and pop are paired; induction over the calls, exactly as in the pattern's invariant.)

\emph{Soundness.} Every recorded \texttt{path} ends in $n-1$ and consists of edges, so it is a path $0 \leadsto n-1$; in a DAG it is automatically simple.

\emph{Completeness and no duplicates.} A path $0 = y_0 \to \dots \to y_m = n-1$ is produced by the chain of calls $\texttt{dfs}(y_0), \dots, \texttt{dfs}(y_m)$ --- each $y_{i+1}$ is in $\texttt{graph}[y_i]$, and no intermediate $y_i$ equals $n-1$ (vertex $n-1$ occurs only at the end of a simple path). Different paths are different call chains, so each is recorded once.

\emph{Termination without marks.} The recursion depth is the length of the current path, and in a DAG every path has at most $n$ vertices. In a graph with a cycle the same code would recurse forever --- the marks of template 1 are what Problem~\ref{ssec:ch2-max-gold} needs. \hfill$\blacksquare$
\end{invariantbox}

\begin{alternativebox}[Approach 3: memoised path lists --- DP on the DAG]
\textbf{Idea.} The paths from $u$ to $n-1$ depend only on $u$, not on how $u$ was reached. So compute $P(u)$ = the list of paths $u \leadsto n-1$ once per vertex: $P(n-1) = [[n-1]]$, and $P(u)$ = ``$u$ in front of every path of $P(v)$'' for every successor $v$. This is DP on a DAG (\dpbook{12}{DP on DAGs \& Graphs}) with \emph{lists} as values.

\begin{lstlisting}[style=cpp]
#include <bits/stdc++.h>
using namespace std;

// Alternative: P(u) = all paths u -> n-1, computed once per vertex (memoised DFS = DP on the DAG).
class Solution {
    vector<vector<vector<int>>> memo;
    vector<char> done;
    const vector<vector<int>>& paths(const vector<vector<int>>& g, int u) {
        if (done[u]) return memo[u];
        vector<vector<int>> r;
        if (u == (int)g.size() - 1) r.push_back(vector<int>(1, u));
        else
            for (int v : g[u])
                for (auto& p : paths(g, v)) {          // every path v ... n-1, with u in front
                    vector<int> q(1, u);
                    q.insert(q.end(), p.begin(), p.end());
                    r.push_back(q);
                }
        done[u] = 1;
        memo[u] = r;
        return memo[u];
    }
public:
    vector<vector<int>> allPathsSourceTarget(vector<vector<int>>& graph) {
        memo.assign(graph.size(), {});
        done.assign(graph.size(), 0);
        return paths(graph, 0);
    }
};

int main() {
    vector<vector<int>> g = {{4, 3, 1}, {3, 2, 4}, {3}, {4}, {}};
    for (auto& p : Solution().allPathsSourceTarget(g)) {
        for (int x : p) cout << x << " ";
        cout << "/ ";
    }
    cout << "\n";                                      // 0 4 / 0 3 4 / 0 1 3 4 / 0 1 2 3 4 / 0 1 4 /
}
\end{lstlisting}

\begin{complexitybox}
\textbf{Time:} each vertex is expanded once (15 expansions on the complete DAG instead of $16\,384$ calls), but the lists are copied: building $P(u)$ costs the total length of $P(u)$, so the work is the sum of the sizes of \emph{all} lists, about twice the output on the complete DAG (measured: 9\,ms against 3\,ms for Approach~2, the copying dominates). \textbf{Space:} all lists $P(v)$ are kept at the same time. \textbf{Verdict:} \ac{}. Not faster here --- the output dominates --- but the same memo with \emph{numbers} instead of lists (count the paths, find the longest) is \VE{} where backtracking would be exponential.
\end{complexitybox}
\end{alternativebox}

\begin{tricksbox}[Approach 4: Implementation Refinements]
\begin{itemize}[leftmargin=*]
    \item \textbf{Prune dead ends.} Mark the vertices that can reach $n-1$ (a memoised DFS, or one search from $n-1$ on the reversed edges, Pattern~5) and never enter the others. On the DAG ``$0 \to 14$, plus $0 \to$ a complete DAG on $1..13$ that never reaches 14'' the plain backtracking makes $8193$ calls for one path; with the pruning it makes 2.
\begin{lstlisting}[style=cpp]
#include <bits/stdc++.h>
using namespace std;

// Backtracking that never enters a vertex from which n-1 cannot be reached.
class Solution {
    vector<vector<int>> res;
    vector<int> path;
    vector<int> good;                                  // -1 unknown, 0 no, 1 yes: v reaches n-1
    long long calls = 0;
    bool reaches(const vector<vector<int>>& g, int u) {
        if (good[u] != -1) return good[u];
        int ok = (u == (int)g.size() - 1);
        for (int v : g[u]) if (reaches(g, v)) ok = 1;  // memoised: every vertex is decided once
        good[u] = ok;
        return ok;
    }
    void dfs(const vector<vector<int>>& g, int u) {
        calls++;
        path.push_back(u);
        if (u == (int)g.size() - 1) res.push_back(path);
        else
            for (int v : g[u]) if (good[v] == 1) dfs(g, v);
        path.pop_back();
    }
public:
    vector<vector<int>> allPathsSourceTarget(vector<vector<int>>& graph) {
        good.assign(graph.size(), -1);
        res.clear(); calls = 0;
        if (reaches(graph, 0)) dfs(graph, 0);
        return res;
    }
    long long lastCalls() const { return calls; }
};

int main() {
    int n = 15;
    vector<vector<int>> g(n);
    for (int i = 1; i < n - 1; i++) {
        g[0].push_back(i);
        for (int j = i + 1; j < n - 1; j++) g[i].push_back(j);   // 1..13: complete, no way to 14
    }
    g[0].push_back(14);
    Solution s;
    size_t found = s.allPathsSourceTarget(g).size();
    cout << found << " " << s.lastCalls() << "\n";      // 1 2
}
\end{lstlisting}
    \item \textbf{One shared path vector.} Pushing and popping one \texttt{vector<int>} costs \bigO{1} per step; passing the path \emph{by value} into every call copies it at every node of the search tree.
    \item \textbf{Iterative form.} Not needed: the depth is at most $n = 15$. For the counting variant on big DAGs the memo is iterative over a topological order anyway.
    \item \textbf{Graphs with cycles.} Add \texttt{onPath} (template 1); without it the recursion never ends.
\end{itemize}
\end{tricksbox}

\subsubsection*{5. Trace}

\begin{tracebox}[{Example 2: graph = [[4,3,1],[3,2,4],[3],[4],[]]}]
\begin{center}
\begin{tikzpicture}[x=0.8cm, y=0.8cm]
    \tikzset{gv/.append style={minimum size=6.2mm, font=\scriptsize\bfseries}}
    \foreach \k/\t/\sa/\sb/\sc/\sd/\se in {
        0/{call 2: path [0,4] --- output}/gvseen/gv/gv/gv/gvcur,
        1/{call 4: path [0,3,4] --- output}/gvseen/gv/gv/gvseen/gvcur,
        2/{call 10: path [0,1,2,3,4] --- output}/gvseen/gvseen/gvseen/gvseen/gvcur}{
        \begin{scope}[xshift=\k*5.5cm]
            \node[\sa] (t0) at (0,0) {0}; \node[\sb] (t1) at (1.4,1.2) {1}; \node[\sc] (t2) at (2.9,1.2) {2};
            \node[\sd] (t3) at (2.9,-0.6) {3}; \node[\se] (t4) at (4.3,0) {4};
            \draw[gde] (t0) -- (t1); \draw[gde] (t0) -- (t3); \draw[gde] (t0) to[bend right=50] (t4);
            \draw[gde] (t1) -- (t2); \draw[gde] (t1) -- (t3); \draw[gde] (t1) -- (t4);
            \draw[gde] (t2) -- (t3); \draw[gde] (t3) -- (t4);
            \ifnum\k=0 \draw[gpath, -{Stealth[length=3mm]}] (t0) to[bend right=50] (t4); \fi
            \ifnum\k=1 \draw[gpath, -{Stealth[length=3mm]}] (t0) -- (t3); \draw[gpath, -{Stealth[length=3mm]}] (t3) -- (t4); \fi
            \ifnum\k=2 \foreach \a/\b in {0/1, 1/2, 2/3, 3/4} \draw[gpath, -{Stealth[length=3mm]}] (t\a) -- (t\b); \fi
            \node[font=\scriptsize\bfseries] at (2.15,2.0) {\t};
        \end{scope}
    }
\end{tikzpicture}
\end{center}
\begin{center}
\begin{tikzpicture}[tn/.style={rectangle, rounded corners=2pt, draw=darkgray!70, fill=white, font=\scriptsize\ttfamily, inner sep=2pt},
                    to/.style={tn, draw=accentorange, fill=answerfill, line width=0.9pt}]
    \node[tn] (c1) at (5,0) {1: 0};
    \node[to] (c2) at (0.6,-0.8) {2: 0 4};
    \node[tn] (c3) at (3.2,-0.8) {3: 0 3};
    \node[tn] (c5) at (8.4,-0.8) {5: 0 1};
    \node[to] (c4) at (3.2,-1.6) {4: 0 3 4};
    \node[tn] (c6) at (6.2,-1.6) {6: 0 1 3};
    \node[tn] (c8) at (8.4,-1.6) {8: 0 1 2};
    \node[to] (c11) at (10.8,-1.6) {11: 0 1 4};
    \node[to] (c7) at (6.2,-2.4) {7: 0 1 3 4};
    \node[tn] (c9) at (8.4,-2.4) {9: 0 1 2 3};
    \node[to] (c10) at (8.4,-3.2) {10: 0 1 2 3 4};
    \foreach \a/\b in {c1/c2, c1/c3, c1/c5, c3/c4, c5/c6, c5/c8, c5/c11, c6/c7, c8/c9, c9/c10}
        \draw[darkgray!80, -{Stealth[length=1.6mm]}] (\a) -- (\b);
    \node[note, text width=4.2cm, anchor=west] at (-0.4,-2.8) {the search tree: 11 calls, numbered in the order they happen; orange = an output};
\end{tikzpicture}
\end{center}
\begin{center}\small
\begin{tabular}{@{}rllll@{}}
\toprule
\textbf{call} & \textbf{enter} & \textbf{path after push} & \textbf{action} & \textbf{res} \\
\midrule
1 & 0 & \texttt{[0]} & try 4, 3, 1 & --- \\
2 & 4 & \texttt{[0,4]} & target: copy, pop & 1 path \\
3 & 3 & \texttt{[0,3]} & try 4 & \\
4 & 4 & \texttt{[0,3,4]} & target: copy, pop & 2 \\
5 & 1 & \texttt{[0,1]} & try 3, 2, 4 & \\
6 & 3 & \texttt{[0,1,3]} & try 4 & \\
7 & 4 & \texttt{[0,1,3,4]} & target & 3 \\
8 & 2 & \texttt{[0,1,2]} & try 3 & \\
9 & 3 & \texttt{[0,1,2,3]} & try 4 & \\
10 & 4 & \texttt{[0,1,2,3,4]} & target & 4 \\
11 & 4 & \texttt{[0,1,4]} & target & 5 \\
\bottomrule
\end{tabular}
\end{center}
Vertex 3 is entered three times (calls 3, 6, 9) and expands the same subtree $3 \to 4$ each time --- exactly what Approach~3 avoids: its memo computes $P(3) = [[3,4]]$ once and reuses it.
\end{tracebox}

\begin{diagrambox}
\begin{tikzpicture}[x=1.15cm, y=1cm]
    \node[gvsrc] (m0) at (0,0) {0}; \node[gv] (m1) at (1.4,1.2) {1}; \node[gv] (m2) at (2.9,1.2) {2};
    \node[gv] (m3) at (2.9,-0.6) {3}; \node[gvdst] (m4) at (4.3,0) {4};
    \draw[gde] (m0) -- (m1); \draw[gde] (m0) -- (m3); \draw[gde] (m0) to[bend right=50] (m4);
    \draw[gde] (m1) -- (m2); \draw[gde] (m1) -- (m3); \draw[gde] (m1) -- (m4);
    \draw[gde] (m2) -- (m3); \draw[gde] (m3) -- (m4);
    \node[font=\scriptsize\ttfamily, anchor=west, text=deepblue] at (4.7,0) {P(4) = [4]};
    \node[font=\scriptsize\ttfamily, anchor=north, text=deepblue] at (2.9,-0.95) {P(3) = [3 4]};
    \node[font=\scriptsize\ttfamily, anchor=south, text=deepblue] at (2.9,1.5) {P(2) = [2 3 4]};
    \node[font=\scriptsize\ttfamily, anchor=south east, text=deepblue, align=right] at (1.3,1.5) {P(1) = [1 3 4] [1 2 3 4] [1 4]};
    \node[font=\scriptsize\ttfamily, anchor=north east, text=deepblue, align=right] at (-0.2,-0.4) {P(0) = [0 4] [0 3 4]\\ \ \ \ \ [0 1 3 4] [0 1 2 3 4] [0 1 4]};
\end{tikzpicture}
\figcap{Approach 3 on Example 2: each vertex is expanded once (5 expansions instead of 11 calls); $P(0)$ comes out in the same order as the backtracking output.}
\end{diagrambox}

\subsubsection*{6. Test suite}

\begin{testbox}
\centering\small
\begin{tabular}{@{}p{5.6cm}p{3.6cm}p{5.2cm}@{}}
\toprule
\textbf{Input} & \textbf{Our output} & \textbf{Why} \\
\midrule
\texttt{[[1],[]]} & \texttt{[[0,1]]} & smallest graph, $n = 2$ \\
\texttt{[[],[]]} & \texttt{[]} & 0 has no edges: no path at all \\
\texttt{[[1],[],[1]]} & \texttt{[]} & $n-1 = 2$ is a source: unreachable \\
\texttt{[[2],[0],[]]} & \texttt{[[0,2]]} & vertex 0 is not the first vertex of the topological order \\
complete DAG, $n = 15$ & 8192 paths & the worst case: $2^{n-2}$ paths, 16\,384 calls \\
dead-end DAG above, $n = 15$ & \texttt{[[0,14]]} & 8193 calls plain, 2 calls with pruning \\
\bottomrule
\end{tabular}
\end{testbox}
All outputs come from running the programs above. Approaches 1--3 (and the pruned version) were also compared on $2000$ random DAGs with $n \le 10$ (random topological labels, random edge density, shuffled adjacency lists): the three sets of paths, sorted, were identical every time.

\subsubsection*{7. Common mistakes}

\begin{warningbox}[Mistake 1 --- a visited array copied from plain DFS]
\begin{lstlisting}[style=snippet]
if (!visited[v]) { visited[v] = true; dfs(g, v); }       // WRONG: each vertex used by one path only
\end{lstlisting}
On Example 2 this returns only \texttt{[[0,4]]}: vertex 4 is marked by the first path and can never end another one, and vertex 3 is never re-entered from 1. In a DAG drop the array; in a graph with cycles un-mark on exit.
\end{warningbox}

\begin{warningbox}[Mistake 2 --- recording the path without copying]
Storing a pointer or a reference to \texttt{path} instead of a copy gives a list of identical (finally empty) paths, because \texttt{path} keeps changing. \texttt{res.push\_back(path)} copies, which is what we want.
\end{warningbox}

\begin{warningbox}[Mistake 3 --- forgetting the pop on the target branch]
Returning early when $u = n-1$ (``\texttt{if (u == n-1) \{ res.push\_back(path); return; \}}'') skips the \texttt{pop\_back}; the vertex $n-1$ stays in the path and every later path contains it in the middle. Pop before returning, or keep a single exit as in Approach~2.
\end{warningbox}

\begin{warningbox}[Mistake 4 --- assuming 0 is a source]
Nothing says that 0 comes first in the topological order. \texttt{[[2],[0],[]]} has the edge $1 \to 0$; the backtracking does not care, but a brute force that orders vertices by number would.
\end{warningbox}

\subsubsection*{8. Takeaways}

\begin{itemize}[leftmargin=*]
    \item Backtracking is DFS plus a path stack; in a DAG the marks are unnecessary.
    \item When the output is exponential, the right goal is ``proportional to the output'': at most $2^{n-2}$ paths, and the search makes at most about twice as many calls.
    \item The future of a path in a DAG depends only on its last vertex --- memoise. With lists it gives the same output; with numbers it turns the exponential count into \VE.
    \item Prune what cannot reach the target before searching: a reverse search costs \VE{} and can save an exponential amount.
\end{itemize}

\begin{variantsbox}
\variant{Only the number of paths}{\varnone}
  {return how many paths lead from 0 to $n-1$, with $n$ up to $10^5$.}
  {$\text{cnt}(u)$ = number of paths $u \leadsto n-1$, depending only on $u$.}
  {$\text{cnt}(n-1) = 1$, $\text{cnt}(u) = \sum_{u \to v} \text{cnt}(v)$, memoised; \VE{}. On the complete DAG with $n = 15$ it returns 8192, the length of the list of Approach~2.}
\variant{The graph may have cycles}{\varin{https://atcoder.jp/contests/abc284/tasks/abc284_e}}
  {count the simple paths from vertex 1 in an undirected graph, capped at $10^6$.}
  {paths can now return to used vertices; the past matters.}
  {un-mark on exit (template 1) and stop at the cap: every call is one path (Exercise E4.4).}
\variant{Visit every vertex}{\varin{https://atcoder.jp/contests/abc054/tasks/abc054_c}}
  {count the paths from 1 that visit all $n \le 8$ vertices of an undirected graph.}
  {a path is a goal only when its length is $n$.}
  {backtracking with marks, answer counted at depth $n$ (Exercise E4.1).}
\variant{The longest path}{\varnone}
  {return one path from 0 to $n-1$ with the most vertices, $n$ up to $10^5$.}
  {$\text{best}(u)$ = the most vertices on a path $u \leadsto n-1$, plus the successor that achieves it.}
  {DP on the DAG, $\text{best}(u) = 1 + \max_{u \to v} \text{best}(v)$, then follow the stored successors; \VE. Both DP variants were checked against the lists of Approach~2 on 2000 random DAGs.}
\end{variantsbox}

%============================================================
% CH2 · P4 · PROBLEM 2 — PATH WITH MAXIMUM GOLD (LeetCode 1219)
%============================================================
\subsection{Problem 2.4.2: Path with Maximum Gold}
\problemlabel{2.4.2}{ssec:ch2-max-gold}

\problemcard{2.4.2}{Path with Maximum Gold}{LeetCode 1219}{https://leetcode.com/problems/path-with-maximum-gold/}{\exmedium\ \lc{Medium}\ (a grid with cycles: the marks are needed now, and the cost must be measured)}{Backtracking from every gold cell; mark a cell by setting it to 0 and restore it on exit; prune with the gold that is still reachable}

\subsubsection*{1. Problem statement}

A gold mine is an $m \times n$ grid; cell $(r,c)$ holds $\texttt{grid}[r][c] \ge 0$ units of gold. Starting at any cell with gold, you may repeatedly step up, down, left or right, collecting all the gold of every cell you stand on. You may never visit a cell twice and never step onto a cell with $0$ gold; you may stop anywhere. Return the maximum total gold.

\textbf{Constraints:} $1 \le m, n \le 15$; $0 \le \texttt{grid}[r][c] \le 100$; \textbf{at most 25 cells contain gold}.

\begin{table}[H]
\centering\small
\begin{tabular}{@{}lcl@{}}
\toprule
\textbf{grid} & \textbf{Output} & \textbf{A best path} \\
\midrule
\texttt{[[0,6,0],[5,8,7],[0,9,0]]} & \texttt{24} & $9 \to 8 \to 7$ \\
\texttt{[[1,0,7],[2,0,6],[3,4,5],[0,3,0],[9,0,20]]} & \texttt{28} & $1 \to 2 \to 3 \to 4 \to 5 \to 6 \to 7$ \\
\bottomrule
\end{tabular}
\end{table}

\begin{diagrambox}
\begin{tikzpicture}[x=0.75cm, y=0.75cm]
    \tikzset{gc/.style={gcell, minimum size=7.5mm, fill=answerfill!60}, zc/.style={gcell, minimum size=7.5mm, fill=lightgray, text=midgray},
             pc/.style={gc, fill=currentfill, draw=accentorange, line width=1.1pt}}
    \begin{scope}
        \foreach \v/\s [count=\i from 0] in {0/zc,6/gc,0/zc} \node[\s] at (\i,0) {\v};
        \foreach \v/\s [count=\i from 0] in {5/gc,8/pc,7/pc} \node[\s] at (\i,-1) {\v};
        \foreach \v/\s [count=\i from 0] in {0/zc,9/pc,0/zc} \node[\s] at (\i,-2) {\v};
        \draw[gpath, -{Stealth[length=2.5mm]}] (1,-1.7) -- (1,-1.3);
        \draw[gpath, -{Stealth[length=2.5mm]}] (1.3,-1) -- (1.7,-1);
        \node[font=\scriptsize\bfseries] at (1,-3.0) {Example 1: $9 + 8 + 7 = 24$};
    \end{scope}
    \begin{scope}[xshift=4.2cm, yshift=1.1cm]
        \foreach \row/\y in {{1/pc,0/zc,7/pc}/0, {2/pc,0/zc,6/pc}/-1, {3/pc,4/pc,5/pc}/-2, {0/zc,3/gc,0/zc}/-3, {9/gc,0/zc,20/gc}/-4}
            \foreach \v/\s [count=\i from 0] in \row \node[\s] at (\i,\y) {\v};
        \draw[gpath, -{Stealth[length=2.5mm]}] (0,-0.3) -- (0,-0.7);
        \draw[gpath, -{Stealth[length=2.5mm]}] (0,-1.3) -- (0,-1.7);
        \draw[gpath, -{Stealth[length=2.5mm]}] (0.3,-2) -- (0.7,-2);
        \draw[gpath, -{Stealth[length=2.5mm]}] (1.3,-2) -- (1.7,-2);
        \draw[gpath, -{Stealth[length=2.5mm]}] (2,-1.7) -- (2,-1.3);
        \draw[gpath, -{Stealth[length=2.5mm]}] (2,-0.7) -- (2,-0.3);
        \node[font=\scriptsize\bfseries] at (1,-5.0) {Example 2: 28};
    \end{scope}
    \node[note, text width=5.2cm, align=left, anchor=west] at (9.6,-1.3) {Example 2: the 20 and the 9 have no gold neighbour at all --- each is a path of one cell, worth less than the chain. The 3 below the 4 is a dead end: a path can end there, but then it misses one side of the chain.};
\end{tikzpicture}
\figcap{The two examples; grey cells hold no gold. Orange: one optimal path.}
\end{diagrambox}

\subsubsection*{2. The graph model}

\begin{modelbox}[Graph Model --- Path with Maximum Gold]
\renewcommand{\arraystretch}{1.35}
\begin{tabularx}{\linewidth}{@{}>{\bfseries}l X@{}}
Vertices & the gold cells only, $k \le 25$ of them; vertex $(r,c)$ has the value $\texttt{grid}[r][c] \in [1, 100]$ \\
Edges & undirected, between gold cells that share a side; implicit (look at the 4 neighbours); a grid graph is \textbf{bipartite} (chessboard colours) and has degree $\le 4$ \\
Weights & on the vertices, all positive \\
Question & the maximum-value \textbf{simple path}, any start, any end --- template 2 of the pattern \\
Size & $k \le 25$, so the graph is tiny; the search tree is not: up to $3\,060\,417$ nodes (measured below) \\
\end{tabularx}
\end{modelbox}

\begin{diagrambox}
\begin{tikzpicture}[x=1.2cm, y=1.1cm]
    \tikzset{gv/.append style={minimum size=7mm}}
    \node[gv] (g01) at (1,0) {6};
    \node[gv] (g10) at (0,-1) {5}; \node[gv] (g11) at (1,-1) {8}; \node[gv] (g12) at (2,-1) {7};
    \node[gv] (g21) at (1,-2) {9};
    \foreach \a/\b in {01/11, 10/11, 11/12, 11/21} \draw[ge] (g\a) -- (g\b);
    \node[font=\scriptsize\bfseries] at (1,-2.8) {Example 1: a star};
    \begin{scope}[xshift=4.6cm, yshift=0.6cm]
        \foreach \v/\n/\x/\y in {1/a/0/0, 7/b/2/0, 2/c/0/-1, 6/d/2/-1, 3/e/0/-2, 4/f/1/-2, 5/g/2/-2, 3/h/1/-3, 9/i/0/-4, 20/j/2/-4}
            \node[gv, minimum size=6.5mm, font=\scriptsize\bfseries] (n\n) at (\x,\y) {\v};
        \foreach \a/\b in {a/c, c/e, e/f, f/g, g/d, d/b, f/h} \draw[ge] (n\a) -- (n\b);
        \node[font=\scriptsize\bfseries, align=center] at (1,-4.85) {Example 2: a tree plus\\two isolated vertices};
    \end{scope}
    \node[note, text width=4.6cm, align=left, anchor=west] at (8.4,-1) {Grey cells are not vertices at all. The real graph is often a forest of small pieces; only a big \emph{compact} block of gold makes the search tree large.};
\end{tikzpicture}
\figcap{The graphs of the examples: one vertex per gold cell, one edge per pair of side-adjacent gold cells.}
\end{diagrambox}

\subsubsection*{3. Key observations}

\begin{itemize}[leftmargin=*]
    \item \textbf{The path's history matters.} ``Never visit a cell twice'' depends on the whole path; the grid has cycles (any $2 \times 2$ block of gold). So the marks are needed --- and un-marked on exit (Tool~4), because another path may use the cell.
    \item \textbf{The grid is its own \texttt{onPath} array.} Set the cell to 0 when entering: the rule ``never step on 0'' then forbids it automatically. Restore the saved value on exit.
    \item \textbf{An optimal path cannot be extended.} All values are positive, so a best path has no free gold neighbour at either end. This makes the bound ``everything still reachable from here'' strong.
    \item \textbf{How big is the search?} With $\le 25$ gold cells and degree $\le 4$ the crude bound $25 \cdot 3^{24}$ is meaningless. Measured instead: a hill-climbing search over 25-cell shapes inside a $9 \times 9$ board (moving one cell at a time, keeping the change when the number of calls grows) ended at the compact $5 \times 5$ block, with $3\,060\,417$ calls of the backtracking from all 25 starts. That is the realistic worst case: fast.
\end{itemize}

\subsubsection*{4. Approaches}

\begin{bruteforcebox}
\textbf{Idea: every simple path is a prefix of an ordering of the gold cells.} Enumerate all $k!$ orderings with \texttt{next\_permutation}; walk each one while consecutive cells are adjacent and keep the best prefix sum.

\begin{lstlisting}[style=cpp]
#include <bits/stdc++.h>
using namespace std;

// Brute force: every ordering of the gold cells; the longest prefix of side-adjacent cells is a path.
class Solution {
public:
    int getMaximumGold(vector<vector<int>>& grid) {
        vector<pair<int, int>> cells;
        for (int r = 0; r < (int)grid.size(); r++)
            for (int c = 0; c < (int)grid[0].size(); c++)
                if (grid[r][c] > 0) cells.push_back({r, c});
        int best = 0;                                  // cells is sorted: the first permutation
        do {
            int sum = 0;
            for (size_t i = 0; i < cells.size(); i++) {
                if (i > 0 && abs(cells[i].first - cells[i - 1].first) +
                             abs(cells[i].second - cells[i - 1].second) != 1) break;
                sum += grid[cells[i].first][cells[i].second];
                best = max(best, sum);
            }
        } while (next_permutation(cells.begin(), cells.end()));
        return best;
    }
};

int main() {
    vector<vector<int>> g1 = {{0, 6, 0}, {5, 8, 7}, {0, 9, 0}};
    vector<vector<int>> g2 = {{1, 0, 7}, {2, 0, 6}, {3, 4, 5}, {0, 3, 0}, {9, 0, 20}};
    cout << Solution().getMaximumGold(g1) << " " << Solution().getMaximumGold(g2) << "\n";   // 24 28
}
\end{lstlisting}

\begin{complexitybox}
\textbf{Time:} \bigO{k! \cdot k}: $5! = 120$ orderings for Example 1, $10! = 3.6 \cdot 10^6$ for Example 2, $25! \approx 1.6 \cdot 10^{25}$ in the worst case. \textbf{Space:} \bigO{k}. \textbf{Verdict:} \tle{} beyond $k \approx 11$; a perfect reference for tiny random grids, because it shares no code with the search.
\end{complexitybox}
\end{bruteforcebox}

\begin{algorithmbox}[Approach 2: backtracking with temporary zeroing]
\textbf{Idea.} Template 2 of the pattern on the grid. \texttt{dfs(r, c)} returns the most gold collectable by a path that starts at $(r,c)$ and avoids the cells currently set to 0: take this cell's gold, set it to 0, try the four neighbours, restore the gold.

\begin{lstlisting}[style=cpp]
#include <bits/stdc++.h>
using namespace std;

// Path with Maximum Gold (LeetCode 1219): backtracking from every gold cell.
// A cell on the current path is set to 0 (the rule "never step on 0" then forbids it) and restored on exit.
class Solution {
    int R, C;
    int dfs(vector<vector<int>>& g, int r, int c) {
        static const int dr[4] = {-1, 1, 0, 0}, dc[4] = {0, 0, -1, 1};
        int gold = g[r][c];
        g[r][c] = 0;                                   // enter: (r, c) is on the path
        int best = 0;
        for (int d = 0; d < 4; d++) {
            int nr = r + dr[d], nc = c + dc[d];
            if (nr >= 0 && nr < R && nc >= 0 && nc < C && g[nr][nc] > 0)
                best = max(best, dfs(g, nr, nc));
        }
        g[r][c] = gold;                                // exit: restore the gold
        return gold + best;
    }
public:
    int getMaximumGold(vector<vector<int>>& grid) {
        R = grid.size(); C = grid[0].size();
        int ans = 0;
        for (int r = 0; r < R; r++)
            for (int c = 0; c < C; c++)
                if (grid[r][c] > 0) ans = max(ans, dfs(grid, r, c));
        return ans;
    }
};

int main() {
    vector<vector<int>> g1 = {{0, 6, 0}, {5, 8, 7}, {0, 9, 0}};
    vector<vector<int>> g2 = {{1, 0, 7}, {2, 0, 6}, {3, 4, 5}, {0, 3, 0}, {9, 0, 20}};
    cout << Solution().getMaximumGold(g1) << " " << Solution().getMaximumGold(g2) << "\n";   // 24 28
}
\end{lstlisting}

\begin{complexitybox}
\textbf{Time:} one call per simple path that starts at a gold cell: 25 calls for Example 1, 66 for Example 2. Worst case found: the $5 \times 5$ block of gold, $3\,060\,417$ calls, $0.12$--$0.19$\,s here (all-equal or random values make no difference: without pruning the values do not change the search). Other shapes measured: $4 \times 6$ block 1\,603\,536 calls, $3 \times 8$ 826\,020, $2 \times 12$ 145\,100, a $1 \times 25$ line only 625. \textbf{Space:} \bigO{k} recursion depth, the grid itself is the mark array. \textbf{Verdict:} \ac{}-fast.
\end{complexitybox}
\end{algorithmbox}

\begin{invariantbox}[Why it is correct]
\textbf{Invariant.} When \texttt{dfs(r, c)} is called, the cells set to 0 by the active calls are exactly the cells of the current path before $(r,c)$; when it returns, the grid is exactly as it was at the call (every \texttt{g[r][c] = 0} is undone by the matching \texttt{g[r][c] = gold}).

\textbf{Claim.} \texttt{dfs(r, c)} returns the maximum gold of a simple path that starts at $(r,c)$ and uses only gold cells that are not on the current path. By induction on the number of available cells: the path either stops at $(r,c)$ (value \texttt{gold}) or continues to an available gold neighbour $(r', c')$ and then is a path from $(r', c')$ avoiding the current path \emph{and} $(r,c)$ --- exactly what the recursive call computes, because $(r,c)$ is 0 during it. Taking the maximum over the neighbours covers all cases.

The outer loop tries every possible first cell, so the answer is the maximum over all simple paths. \hfill$\blacksquare$
\end{invariantbox}

\begin{alternativebox}[Approach 3: memoise on (set of used cells, current cell) --- bitmask DP]
\textbf{Idea.} The future of a path depends only on which cells it has used and where it stands, so $\text{best}(S, u)$ = ``most gold still collectable from $u$ when the cells of $S$ are used'' can be memoised (\dpbook{13}{Bitmask DP}). Number the gold cells $0..k-1$; the key is $S \cdot 32 + u$ in a hash map, since $2^{25} \cdot 25$ states would not fit in an array.

\begin{lstlisting}[style=cpp]
#include <bits/stdc++.h>
using namespace std;
typedef long long ll;

// Alternative: memoise best(S, u) = max gold still collectable from u when the cells in S are used.
class Solution {
    vector<int> val;
    vector<vector<int>> adj;                           // gold cells numbered 0..k-1
    unordered_map<ll, int> memo;
    int best(ll S, int u) {
        ll key = S * 32 + u;
        auto it = memo.find(key);
        if (it != memo.end()) return it->second;
        int res = 0;
        for (int v : adj[u])
            if (!(S >> v & 1)) res = max(res, val[v] + best(S | (1LL << v), v));
        memo[key] = res;
        return res;
    }
public:
    int getMaximumGold(vector<vector<int>>& grid) {
        int R = grid.size(), C = grid[0].size();
        vector<vector<int>> id(R, vector<int>(C, -1));
        val.clear();
        for (int r = 0; r < R; r++)
            for (int c = 0; c < C; c++)
                if (grid[r][c] > 0) { id[r][c] = val.size(); val.push_back(grid[r][c]); }
        int k = val.size();
        adj.assign(k, vector<int>());
        int dr[4] = {-1, 1, 0, 0}, dc[4] = {0, 0, -1, 1};
        for (int r = 0; r < R; r++)
            for (int c = 0; c < C; c++)
                if (id[r][c] >= 0)
                    for (int d = 0; d < 4; d++) {
                        int nr = r + dr[d], nc = c + dc[d];
                        if (nr >= 0 && nr < R && nc >= 0 && nc < C && id[nr][nc] >= 0)
                            adj[id[r][c]].push_back(id[nr][nc]);
                    }
        memo.clear();
        int ans = 0;
        for (int u = 0; u < k; u++) ans = max(ans, val[u] + best(1LL << u, u));
        return ans;
    }
};

int main() {
    vector<vector<int>> g1 = {{0, 6, 0}, {5, 8, 7}, {0, 9, 0}};
    vector<vector<int>> g2 = {{1, 0, 7}, {2, 0, 6}, {3, 4, 5}, {0, 3, 0}, {9, 0, 20}};
    cout << Solution().getMaximumGold(g1) << " " << Solution().getMaximumGold(g2) << "\n";   // 24 28
}
\end{lstlisting}

\begin{complexitybox}
\textbf{Time:} one evaluation per reachable state. On the $5 \times 5$ block: $959\,085$ distinct states and $1\,187\,569$ calls --- only $2.6$ times fewer than the $3\,060\,417$ calls of Approach~2, and \emph{slower} ($0.46$--$0.75$\,s), because every call pays for a hash lookup. On a grid most pairs (set, last cell) are reached by just one path, so there is little to share. \textbf{Space:} the hash map, about a million entries. \textbf{Verdict:} \ac{}, but worse than plain backtracking here. Bitmask DP wins when the graph is dense ($K_n$: $n!$ paths but only $2^n n$ states) --- the Hamiltonian problems of Chapter~23.
\end{complexitybox}
\end{alternativebox}

\begin{tricksbox}[Approach 4: Implementation Refinements --- pruning, measured]
\textbf{Bound by the reachable gold.} At a node, flood-fill the gold cells that are still reachable from the current cell without crossing the path. If $\text{sum} + \text{reachable} \le \text{best}$, no extension can improve the answer: do not recurse. The fill costs \bigO{k} per node and is a valid optimistic bound, because any extension uses only reachable cells.

\begin{lstlisting}[style=cpp]
#include <bits/stdc++.h>
using namespace std;

// Backtracking with a bound: gold still reachable from (r, c) through free cells. If even all of it
// cannot beat the best path found so far, the branch is cut.
class Solution {
    int R, C, best;
    vector<vector<int>> g, mark;
    int stamp = 0;
    const int dr[4] = {-1, 1, 0, 0}, dc[4] = {0, 0, -1, 1};
    bool ok(int r, int c) { return r >= 0 && r < R && c >= 0 && c < C && g[r][c] > 0; }
    int reachable(int r, int c) {                      // gold reachable from the neighbours of (r, c)
        stamp++;
        int s = 0;
        vector<pair<int, int>> st(1, make_pair(r, c));
        while (!st.empty()) {
            int x = st.back().first, y = st.back().second;
            st.pop_back();
            for (int d = 0; d < 4; d++) {
                int nx = x + dr[d], ny = y + dc[d];
                if (ok(nx, ny) && mark[nx][ny] != stamp) {
                    mark[nx][ny] = stamp; s += g[nx][ny]; st.push_back(make_pair(nx, ny));
                }
            }
        }
        return s;
    }
    void dfs(int r, int c, int sum) {
        calls++;
        int gold = g[r][c];
        g[r][c] = 0; sum += gold;                      // enter
        best = max(best, sum);
        if (sum + reachable(r, c) > best)              // the bound: otherwise nothing can improve
            for (int d = 0; d < 4; d++)
                if (ok(r + dr[d], c + dc[d])) dfs(r + dr[d], c + dc[d], sum);
        g[r][c] = gold;                                // exit
    }
public:
    long long calls = 0;
    int getMaximumGold(vector<vector<int>>& grid) {
        g = grid; R = g.size(); C = g[0].size(); best = 0;
        mark.assign(R, vector<int>(C, 0));
        for (int r = 0; r < R; r++)
            for (int c = 0; c < C; c++)
                if (g[r][c] > 0) dfs(r, c, 0);
        return best;
    }
};

int main() {
    vector<vector<int>> block(5, vector<int>(5, 1));   // the worst shape for Approach 2
    vector<vector<int>> holes(4, vector<int>(6, 1));
    holes[0][1] = holes[3][4] = 0;                     // 22 cells, the best path takes 21
    for (auto* g : {&block, &holes}) {
        Solution s;
        int ans = s.getMaximumGold(*g);
        cout << ans << " " << s.calls << "\n";         // 25 65  /  21 5227
    }
}
\end{lstlisting}

\begin{center}
\small
\renewcommand{\arraystretch}{1.2}
\begin{tabular}{@{}lrrr@{}}
\toprule
\textbf{grid} & \textbf{answer} & \textbf{calls, Approach 2} & \textbf{calls, with the bound} \\
\midrule
$5 \times 5$ block, all 1 & 25 & 3\,060\,417 & 65 \\
$5 \times 5$ block, random 1..100 & 1023 & 3\,060\,417 & 65 \\
$5 \times 5$ block without its centre & 24 & 664\,288 & 63 \\
$4 \times 6$ block minus two cells & 21 & 135\,422 & 5\,227 \\
\bottomrule
\end{tabular}
\end{center}
When a path through \emph{all} gold exists, the bound stops everything as soon as one is found. The last grid is the hard case for the bound: no path takes all 22 cells, so it cannot cut until the best 21 is known, and it still saves a factor 26. The weaker bound ``everything not yet on the path'' (template 2) only helps in the first three rows; on the last grid it makes the same $135\,422$ calls as Approach~2.

\textbf{A tempting start rule.} ``An optimal path can always be chosen to start at a cell with at most 2 gold neighbours.'' On the $5 \times 5$ block it reduces the calls from $3\,060\,417$ to $614\,980$. We compared it with the full search on $260\,000$ random grids (up to $7 \times 7$, all densities, values 1--3 and 1--100, plus all-equal values) and a targeted search for the shapes that a counter-example would need, and it never lost the optimum --- but we have no proof. It is therefore \emph{not} part of the solution: the bound above is provably correct and stronger.

\textbf{Small things.} The direction arrays are \texttt{static} (or members) so they are not rebuilt per call; the grid is passed by reference --- a copy per call would cost $225$ cells at each of the 3 million nodes.
\end{tricksbox}

\subsubsection*{5. Trace}

\begin{tracebox}[{Example 1: grid = [[0,6,0],[5,8,7],[0,9,0]], directions up, down, left, right}]
\begin{center}
\begin{tikzpicture}[x=0.62cm, y=0.62cm]
    \tikzset{g2p4c/.style={gcell, minimum size=6mm, font=\scriptsize}, g2p4z/.style={g2p4c, fill=lightgray, text=midgray},
             g2p4g/.style={g2p4c, fill=answerfill!60}, g2p4p/.style={g2p4c, fill=hitfill}, g2p4u/.style={g2p4c, fill=currentfill, draw=accentorange, line width=1pt}}
    \foreach \k/\t/\a/\b/\cc/\d/\e/\vb/\vc/\vd in {
        0/{start (0,1): 6}/g2p4u/g2p4g/g2p4g/g2p4g/g2p4g/6/8/7,
        1/{6, 8: grid shows 0 0}/g2p4p/g2p4g/g2p4u/g2p4g/g2p4g/0/0/7,
        2/{6, 8, 9: leaf}/g2p4p/g2p4g/g2p4p/g2p4g/g2p4u/0/0/7,
        3/{start (1,2): 7, 8, 9}/g2p4g/g2p4g/g2p4p/g2p4p/g2p4u/6/0/0}{
        \begin{scope}[xshift=\k*3.4cm]
            \node[g2p4z] at (0,0) {0}; \node[\a] at (1,0) {\vb}; \node[g2p4z] at (2,0) {0};
            \node[\b] at (0,-1) {5}; \node[\cc] at (1,-1) {\vc}; \node[\d] at (2,-1) {\vd};
            \node[g2p4z] at (0,-2) {0}; \node[\e] at (1,-2) {9}; \node[g2p4z] at (2,-2) {0};
            \node[font=\tiny\bfseries, align=center] at (1,1.0) {\t};
        \end{scope}
    }
\end{tikzpicture}
\end{center}
{\small Yellow: on the path (already zeroed in \texttt{g}); orange: the current cell. In panel 2 the grid stores 0 at $(0,1)$ and $(1,1)$ --- that is how the search knows not to step back. Panel 4 shows the path $7, 8, 9$, the first one to reach 24.}

\begin{center}
\begin{tikzpicture}[tn/.style={rectangle, rounded corners=2pt, draw=darkgray!70, fill=white, font=\scriptsize, inner sep=2pt}]
    \node[tn] (a) at (0,0) {7 \,$\to$ 24};
    \node[tn] (b) at (0,-0.8) {8 \,$\to$ 17};
    \node[tn] (c1) at (-2.2,-1.6) {6 \,$\to$ 6};
    \node[tn, fill=currentfill, draw=accentorange] (c2) at (0,-1.6) {9 \,$\to$ 9};
    \node[tn] (c3) at (2.2,-1.6) {5 \,$\to$ 5};
    \draw[->] (a) -- (b); \draw[->] (b) -- (c1); \draw[->] (b) -- (c2); \draw[->] (b) -- (c3);
    \node[note, text width=6.2cm, anchor=west] at (3.6,-0.8) {the calls from the start $(1,2)$: each node shows the cell's gold and the value it returns, $\text{gold} + \max(\text{children})$; $7 + (8 + 9) = 24$};
\end{tikzpicture}
\end{center}

\begin{center}\small
\begin{tabular}{@{}lllr@{}}
\toprule
\textbf{start} & \textbf{calls (in order)} & \textbf{returns} & \textbf{best} \\
\midrule
$(0,1)$ gold 6 & 6, 8, 9, 5, 7 & $8 + \max(9, 5, 7) = 17$, so $6 + 17 = 23$ & 23 \\
$(1,0)$ gold 5 & 5, 8, 6, 9, 7 & $8 + \max(6, 9, 7) = 17$, so $5 + 17 = 22$ & 23 \\
$(1,1)$ gold 8 & 8, 6, 9, 5, 7 & $8 + \max(6, 9, 5, 7) = 17$ & 23 \\
$(1,2)$ gold 7 & 7, 8, 6, 9, 5 & $8 + \max(6, 9, 5) = 17$, so $7 + 17 = 24$ & \textbf{24} \\
$(2,1)$ gold 9 & 9, 8, 6, 5, 7 & $8 + \max(6, 5, 7) = 15$, so $9 + 15 = 24$ & 24 \\
\bottomrule
\end{tabular}
\end{center}
25 calls in total, 5 per start: from a leaf the search goes to 8 and then to each of the other three leaves. The answer is \textbf{24}, found from $(1,2)$ and again from $(2,1)$ --- the same path $7, 8, 9$ read in both directions.
\end{tracebox}

\begin{diagrambox}
\begin{tikzpicture}[x=0.55cm, y=0.55cm]
    \tikzset{g2p4c/.style={gcell, minimum size=5.3mm, font=\tiny}, g2p4h/.style={g2p4c, fill=answerfill!60}, g2p4p/.style={g2p4c, fill=currentfill, draw=accentorange}}
    \begin{scope}
        \foreach \x in {0,...,4} \foreach \y in {0,...,4} \node[g2p4h] at (\x,-\y) {1};
        \node[font=\scriptsize\bfseries, align=center] at (2,1.1) {$5 \times 5$ block: 3\,060\,417 calls};
    \end{scope}
    \begin{scope}[xshift=4.6cm]
        \foreach \x in {0,...,4} \foreach \y in {0,...,4} \node[g2p4c] (q\x\y) at (\x,-\y) {};
        \foreach \x/\y [count=\i] in {0/0,0/1,0/2,0/3,0/4,1/4,1/3,1/2,1/1,1/0,2/0,2/1,2/2,2/3,2/4,3/4,3/3,3/2,3/1,3/0,4/0,4/1,4/2,4/3,4/4}
            \node[g2p4p] at (\x,-\y) {\i};
        \node[font=\scriptsize\bfseries, align=center] at (2,1.1) {with the bound: the first path};
    \end{scope}
    \node[note, text width=5.6cm, anchor=west, align=left] at (13.6,-2) {The first path the search builds from the corner snakes through all 25 cells; its value equals the total, so $\text{sum} + \text{reachable} \le \text{best}$ everywhere afterwards: 65 calls in all (25 of them build this path, the rest are cut at once).};
\end{tikzpicture}
\figcap{Why the bound is so strong on the worst shape of Approach~2.}
\end{diagrambox}

\subsubsection*{6. Test suite}

\begin{testbox}
\centering\small
\begin{tabular}{@{}p{5.8cm}lp{6.2cm}@{}}
\toprule
\textbf{grid} & \textbf{Output} & \textbf{Why} \\
\midrule
\texttt{[[5]]} & 5 & one cell \\
\texttt{[[0,0],[0,0]]} & 0 & no gold at all: the loop never starts a search \\
\texttt{[[1,0,1],[0,1,0]]} & 1 & three isolated cells \\
\texttt{[[1,1],[1,1]]} & 4 & a cycle of four: the marks are essential \\
\texttt{[[10,1],[1,10]]} & 22 & the two 10s are not adjacent; every path of 4 takes all \\
$1 \times 25$ line of 1s & 25 & 625 calls: a path graph is cheap \\
$5 \times 5$ block of 1s & 25 & 3\,060\,417 calls (Approach~2), 65 with the bound \\
\bottomrule
\end{tabular}
\end{testbox}
All outputs come from running the programs. Approaches 1--4 were compared on $3000$ random grids of size up to $4 \times 4$ with at most 9 gold cells (density about 65\,\%, values 1--100): all four answers agreed every time, and the grid was always restored after the call.

\subsubsection*{7. Common mistakes}

\begin{warningbox}[Mistake 1 --- restoring the wrong value]
\begin{lstlisting}[style=snippet]
g[r][c] = 0; ... ; g[r][c] = 1;        // WRONG: the cell held gold, not 1
int gold = g[r][c]; g[r][c] = 0; ... ; g[r][c] = gold;   // RIGHT
\end{lstlisting}
The first version silently changes the mine for every later start.
\end{warningbox}

\begin{warningbox}[Mistake 2 --- a global visited array that is never cleared]
Marking cells visited ``forever'' turns the search into a plain DFS: on \texttt{[[1,1],[1,1]]} it still finds 4, but on Example 1 the first start $(0,1)$ marks all five cells and every later start finds nothing new --- the answer becomes 23 instead of 24.
\end{warningbox}

\begin{warningbox}[Mistake 3 --- starting only at the largest cell]
The best path need not contain the largest value: in Example 2 the 20 is isolated, and the answer 28 avoids it. Every gold cell must be tried as a start (or the bound must be strong enough to make that cheap).
\end{warningbox}

\begin{warningbox}[Mistake 4 --- an unproved pruning rule]
``Start only at cells with $\le 2$ gold neighbours'' passed all our random tests; ``start only at cells with \emph{exactly one} gold neighbour'' fails at once: in the $2 \times 2$ block every cell has 2 neighbours and the program would return 0. Test every rule against the unpruned search before using it.
\end{warningbox}

\subsubsection*{8. Takeaways}

\begin{itemize}[leftmargin=*]
    \item A grid with ``each cell at most once'' is a graph with cycles: marks, un-marked on exit. The grid itself can be the mark array (temporary zeroing).
    \item Do not estimate the search with $3^{24}$: \textbf{measure} it on the worst shape. Here the worst is about $3 \cdot 10^6$ calls.
    \item The best bound is ``value so far $+$ everything still reachable''. It is provable, cheap with a flood fill, and turns millions of calls into dozens.
    \item Memoising on (set, last cell) is correct but pays off only when many paths share a state --- not on sparse grids.
\end{itemize}

\begin{variantsbox}
\variant{Spell a word instead of collecting gold}{\varin{https://leetcode.com/problems/word-search/}}
  {the path must spell a given word letter by letter.}
  {the same grid search, but a cell is only entered if its letter is the next one of the word.}
  {stop at the first success; prune by letter counts (Exercise E4.2).}
\variant{Weighted edges instead of cells}{\varin{https://atcoder.jp/contests/abc317/tasks/abc317_c}}
  {a general graph of $n \le 10$ towns, roads with lengths; the longest simple route.}
  {vertices are towns, every start allowed.}
  {template 2 with edge weights instead of vertex values (Exercise E4.3).}
\variant{A path that visits every cell}{\varin{https://cses.fi/problemset/task/1625}}
  {count the 48-step paths through all 49 squares of a $7 \times 7$ grid matching a pattern.}
  {the same search with a fixed start, a fixed end and an exact length.}
  {strong dead-end pruning is needed: $46.8$ million calls without the second rule (Exercise E4.7).}
\variant{Must end on a given cell}{\varnone}
  {the path must finish at a given gold cell $t$.}
  {only nodes whose current cell is $t$ are answers.}
  {record the sum only at $t$; bound with ``reachable gold'' only if $t$ is still reachable (otherwise cut). Checked against Approach~1 restricted to orderings ending at $t$ on 1000 random grids.}
\end{variantsbox}

%============================================================
% CH2 · PATTERN 4 — EXERCISES, HINTS, ANSWERS
%============================================================
\subsection{Pattern 4: Exercises and Going Further}
\label{ssec:ch2-p4-ex}

\begin{patternbox}{Pattern 4 Summary Card}
\renewcommand{\arraystretch}{1.25}
\begin{tabularx}{\linewidth}{@{}>{\bfseries}l X@{}}
Recognise & ``all paths'', ``the best path'', ``each vertex / cell at most once'' on a \emph{tiny} input ($n \le 10$--$15$, $\le 25$ cells, a word of $\le 15$ letters), or a capped count \\
Core & DFS whose marks are removed on exit: \texttt{onPath[]} $=$ the current path; one call per simple path \\
State & everything the answer needs, kept on the path: add on enter, undo on exit (path list, sum, colour counts, prefix sums) \\
Cost & (number of simple paths) $\times$ degree; count it on the worst input before coding \\
Cut & bounds (value so far $+$ reachable rest), dead ends (unreachable target, forced cells, split regions), parity, caps \\
Linear cases & DAG: no marks, memoise on the last vertex (DP); tree: one path per vertex, the path stack is a plain DFS \\
Top traps & forgetting to un-mark or to restore; going on past the target; plain DFS marks in a backtracking; a prune without proof; deep recursion on $10^5$ \\
\end{tabularx}
\end{patternbox}

\exlegend
\begin{enumerate}[leftmargin=*, label=\textbf{E4.\arabic*}, itemsep=6pt]
    \item \exeasy\ \pclassic \exprob{One-stroke Path}{AtCoder ABC054 C}{https://atcoder.jp/contests/abc054/tasks/abc054_c}. A simple undirected graph with $N \le 8$ vertices. Count the paths that start at vertex 1 and visit every vertex exactly once.
    \item \exmedium\ \pclassic \lc{Medium}\ \exprob{Word Search}{LeetCode 79}{https://leetcode.com/problems/word-search/}. An $m \times n$ board of letters ($m, n \le 6$) and a word of length $\le 15$. Can the word be read along side-adjacent cells, each cell used at most once?
    \item \exeasy\ \pconcept \exprob{Remembering the Days}{AtCoder ABC317 C}{https://atcoder.jp/contests/abc317/tasks/abc317_c}. $N \le 10$ towns, roads with lengths up to $10^8$. Start anywhere and travel without visiting a town twice; maximise the total length of the roads used.
    \item \exmedium\ \pconcept \exprob{Count Simple Paths}{AtCoder ABC284 E}{https://atcoder.jp/contests/abc284/tasks/abc284_e}. A simple undirected graph with $N, M \le 2\cdot10^5$ and every degree $\le 10$. Let $K$ be the number of simple paths starting at vertex 1 (the path ``1'' included). Print $\min(K, 10^6)$.
    \item \exmedium\ \pconcept \exprob{Unique Color}{AtCoder ABC198 E}{https://atcoder.jp/contests/abc198/tasks/abc198_e}. A tree with $N \le 10^5$ vertices; vertex $i$ has colour $C_i$. Vertex $x$ is good if no other vertex on the path from 1 to $x$ has the colour of $x$. Print all good vertices.
    \item \exhard\ \pconcept \cf{1700}\ \exprob{Path Prefixes}{Codeforces 1714G}{https://codeforces.com/problemset/problem/1714/G}. A tree rooted at 1 ($\sum n \le 2\cdot10^5$); every edge has two costs $a_j, b_j \ge 1$. For every vertex $i \ge 2$ let $A_i$ be the sum of $a$ on the root path; $r_i$ is the length of the longest prefix of that path whose sum of $b$ is at most $A_i$. Print $r_2, \dots, r_n$.
    \item \exhard\ \ptricky \exprob{Grid Paths}{CSES 1625}{https://cses.fi/problemset/task/1625}. Count the paths of 48 steps from the upper-left to the lower-left square of a $7 \times 7$ grid that visit all 49 squares, matching a 48-character description over \texttt{D}, \texttt{U}, \texttt{L}, \texttt{R} and \texttt{?} (any direction). There are $88\,418$ such paths in total; the time limit is 1 second.
\end{enumerate}

\subsubsection*{Hints}
\begin{enumerate}[leftmargin=*, label=\textbf{E4.\arabic*}]
    \item $7! = 5040$ orders at most. Template 1 without a target: count a path when its length reaches $N$.
    \item Backtrack letter by letter; mark a cell by overwriting it. Before searching at all, compare letter counts.
    \item Template 2 with edge weights; every node of the search tree is a candidate, and every town a start.
    \item Every call of the backtracking is one new simple path. What does the cap do to the number of calls? How deep can the recursion become?
    \item The path from 1 to $x$ in a tree is the DFS stack when $x$ is entered. Keep a count per colour.
    \item On the stack, keep the prefix sums of $b$ along the root path. They are increasing.
    \item Two prunes are enough: (1) moving straight into a blocked square while both side squares are free splits the free squares; (2) a free neighbour with at most one other free neighbour must be entered right now.
\end{enumerate}

\subsubsection*{Answers}

\begin{solutionbox}[E4.1 One-stroke Path]
Backtracking from vertex 1 with \texttt{onPath} marks; a path is complete when it has $N$ vertices. At most $(N-1)! = 5040$ complete paths, and at most $13\,700$ calls in total (the number of simple paths from one vertex of $K_8$, from the table of the pattern theory).
\begin{center}
\begin{tikzpicture}[x=1.1cm, y=0.9cm]
    \tikzset{gv/.append style={minimum size=6.5mm, font=\scriptsize\bfseries}}
    \node[gvsrc] (v1) at (0,0) {1}; \node[gv] (v3) at (1.2,0) {3}; \node[gv] (v4) at (2.4,0) {4};
    \node[gv] (v5) at (3.4,0.8) {5}; \node[gv] (v6) at (3.4,-0.8) {6}; \node[gv] (v7) at (4.8,-0.8) {7}; \node[gvdst] (v2) at (4.8,0.6) {2};
    \foreach \a/\b in {4/6} \draw[ge] (v\a) -- (v\b);
    \foreach \a/\b in {1/3, 3/4, 4/5, 5/6, 6/7, 7/2} \draw[gpath] (v\a) -- (v\b);
    \node[note, text width=4.4cm, anchor=west] at (5.4,0) {Sample 2: the only one-stroke path from 1 is $1, 3, 4, 5, 6, 7, 2$ (the edge $4{-}6$ is unused).};
\end{tikzpicture}
\end{center}
\begin{lstlisting}[style=cpp]
#include <bits/stdc++.h>
using namespace std;

// One-stroke Path (ABC054 C): count the paths that start at 1 and visit every vertex exactly once.
int n, m, answer = 0;
bool adj[9][9];
bool onPath[9];

void go(int u, int visited) {                // visited = number of vertices on the path
    if (visited == n) { answer++; return; }  // the path covers everything: one more Hamiltonian path
    for (int v = 1; v <= n; v++)
        if (adj[u][v] && !onPath[v]) {
            onPath[v] = true;                // enter v
            go(v, visited + 1);
            onPath[v] = false;               // leave v: other paths may use it
        }
}

int main() {
    cin >> n >> m;
    for (int i = 0; i < m; i++) {
        int a, b;
        cin >> a >> b;
        adj[a][b] = adj[b][a] = true;
    }
    onPath[1] = true;
    go(1, 1);
    cout << answer << "\n";
}
\end{lstlisting}
Samples: \textbf{2} (the triangle: $1,2,3$ and $1,3,2$) and \textbf{1}. Checked against a brute force over all $(N-1)!$ orders of $2..N$ (\texttt{next\_permutation}) on 500 random graphs with $N \le 8$: identical.
\end{solutionbox}

\begin{solutionbox}[E4.2 Word Search]
Backtracking over cells: \texttt{match(i, r, c)} checks that cell $(r,c)$ holds \texttt{word[i]}, overwrites it with \texttt{'\#'}, tries the four neighbours for \texttt{word[i+1]} and restores the letter. Two prunes, both provable: if the board has fewer copies of some letter than the word needs, the answer is \texttt{false} without any search; and a path read backwards is still a path, so the word may be searched from its rarer end.
\begin{lstlisting}[style=cpp]
#include <bits/stdc++.h>
using namespace std;

// Word Search (LeetCode 79): backtracking over cells; a used cell is marked '#' and restored on exit.
class Solution {
    int R, C;
    bool match(vector<vector<char>>& b, const string& w, int i, int r, int c) {
        if (r < 0 || r >= R || c < 0 || c >= C || b[r][c] != w[i]) return false;
        if (i + 1 == (int)w.size()) return true;
        char keep = b[r][c];
        b[r][c] = '#';                                       // enter: the cell is on the path
        bool ok = match(b, w, i + 1, r + 1, c) || match(b, w, i + 1, r - 1, c) ||
                  match(b, w, i + 1, r, c + 1) || match(b, w, i + 1, r, c - 1);
        b[r][c] = keep;                                      // exit: restore it
        return ok;
    }
public:
    bool exist(vector<vector<char>>& board, string word) {
        R = board.size(); C = board[0].size();
        map<char, int> cnt;
        for (auto& row : board) for (char ch : row) cnt[ch]++;
        map<char, int> need = cnt;
        for (char ch : word) if (--need[ch] < 0) return false;   // prune 1: not enough letters
        if (cnt[word.back()] < cnt[word[0]])                     // prune 2: start at the rarer end
            reverse(word.begin(), word.end());
        for (int r = 0; r < R; r++)
            for (int c = 0; c < C; c++)
                if (match(board, word, 0, r, c)) return true;
        return false;
    }
};

int main() {
    vector<vector<char>> b = {{'A', 'B', 'C', 'E'}, {'S', 'F', 'C', 'S'}, {'A', 'D', 'E', 'E'}};
    Solution s;
    for (string w : {"ABCCED", "SEE", "ABCB"}) cout << boolalpha << s.exist(b, w) << " ";
    cout << "\n";                                                // true true false
}
\end{lstlisting}
Examples: \textbf{true, true, false}. Measured on a $6 \times 6$ board of \texttt{A}s and the word \texttt{AAAAAAAAAAAAAAB}: without prune 1 the search makes $9\,030\,868$ calls before answering \texttt{false}; with it, none. Checked against a brute force over all ordered tuples of distinct cells (boards up to $3 \times 3$, words up to 5 letters over \texttt{AB} or \texttt{ABC}) on 3000 random cases, 1056 of them \texttt{true}: identical, and the board was restored every time.
\end{solutionbox}

\begin{solutionbox}[E4.3 Remembering the Days]
Template 2 with edge weights: \texttt{go(u, len)} updates the answer at every node and extends the path along every road to an unused town. All $N$ towns are tried as starts. $N \le 10$ is small enough even for the complete graph: $K_{10}$ has $986\,410$ simple paths from each town (the table of the pattern theory), so $10$ starts make about $10^7$ calls --- under half a second here. Lengths up to $10^8$ times 9 roads need \texttt{long long}.
\begin{lstlisting}[style=cpp]
#include <bits/stdc++.h>
using namespace std;
typedef long long ll;

// Remembering the Days (ABC317 C): maximum total length of a simple path, any start, N <= 10.
int n, m;
ll w[11][11];                    // 0 = no road
bool onPath[11];
ll best = 0;

void go(int u, ll len) {
    best = max(best, len);       // every simple path is a candidate, wherever it stops
    for (int v = 1; v <= n; v++)
        if (w[u][v] && !onPath[v]) {
            onPath[v] = true;
            go(v, len + w[u][v]);
            onPath[v] = false;
        }
}

int main() {
    cin >> n >> m;
    for (int i = 0; i < m; i++) {
        int a, b;
        ll c;
        cin >> a >> b >> c;
        w[a][b] = w[b][a] = c;
    }
    for (int s = 1; s <= n; s++) {
        onPath[s] = true;
        go(s, 0);
        onPath[s] = false;
    }
    cout << best << "\n";
}
\end{lstlisting}
Samples: \textbf{1110} ($4 \to 1 \to 3 \to 2$), \textbf{1}, \textbf{20}. Checked against a brute force over all $N!$ orders of the towns (best prefix that follows roads) on 500 random graphs with $N \le 8$: identical.
\end{solutionbox}

\begin{solutionbox}[E4.4 Count Simple Paths]
Every call of the backtracking enters a vertex along a new simple path, so ``count the calls and stop at $10^6$'' needs at most $10^6$ calls, each scanning $\le 10$ neighbours. The trap is that the obvious recursive code can still die: a path graph $1 - 2 - \dots - 2\cdot10^5$ has only $2\cdot10^5$ simple paths from 1 (below the cap) but makes the recursion $2\cdot10^5$ deep. The iterative version below keeps one iterator per vertex and \textbf{resets it when the vertex is left}, because the vertex may be entered again on another path.
\begin{lstlisting}[style=cpp]
#include <bits/stdc++.h>
using namespace std;

// Count Simple Paths (ABC284 E): number of simple paths from vertex 1, capped at 10^6.
// Every call of the backtracking IS one path, so the cap bounds the number of calls.
// Iterative: a path can be 2*10^5 vertices long.
int main() {
    ios::sync_with_stdio(false);
    cin.tie(nullptr);
    const long long CAP = 1000000;
    int n, m;
    cin >> n >> m;
    vector<vector<int>> adj(n + 1);
    for (int i = 0; i < m; i++) {
        int u, v;
        cin >> u >> v;
        adj[u].push_back(v);
        adj[v].push_back(u);
    }
    vector<char> onPath(n + 1, 0);
    vector<size_t> it(n + 1, 0);
    vector<int> st;
    long long count = 1;                       // the path "1" of length 0
    st.push_back(1); onPath[1] = 1;
    while (!st.empty() && count < CAP) {
        int u = st.back();
        if (it[u] == adj[u].size()) {          // u has no more extensions: leave it
            onPath[u] = 0; it[u] = 0;          // reset its iterator: u may be entered again later
            st.pop_back();
            continue;
        }
        int v = adj[u][it[u]++];
        if (onPath[v]) continue;
        onPath[v] = 1; st.push_back(v);        // enter v: a new simple path 1 ... u v
        count++;
    }
    cout << min(count, CAP) << "\n";
}
\end{lstlisting}
Samples: \textbf{3}, \textbf{16}, \textbf{2023}. Large tests: $K_{11}$ (9\,864\,101 paths) prints $1\,000\,000$ in 0.18\,s; the path on $2\cdot10^5$ vertices prints $200\,000$; a random graph with $N = M = 2\cdot10^5$ and all degrees $\le 10$ prints $1\,000\,000$ in 0.9\,s, of which 0.55\,s is reading the input. Checked against a bitmask DP ($\text{dp}[S][v]$ = number of simple paths from 1 with vertex set $S$ ending at $v$, summed and capped) on 500 random graphs with $N \le 11$: identical.
\end{solutionbox}

\begin{solutionbox}[E4.5 Unique Color]
The path from 1 to $x$ is exactly the DFS stack when $x$ is entered. Keep $\texttt{cnt}[c]$ = how often colour $c$ occurs on the stack: $x$ is good iff $\texttt{cnt}[C_x] = 0$ just before $x$ is pushed. Increment on enter, decrement on exit --- template 3 of the pattern. A tree has one path per vertex, so this is linear; the tree can be a path of $10^5$ vertices, hence the iterative DFS.
\begin{lstlisting}[style=cpp]
#include <bits/stdc++.h>
using namespace std;

// Unique Color (ABC198 E): x is good iff no vertex on the path 1..parent(x) has colour C[x].
// Path stack: cnt[c] = occurrences of colour c on the current root path. Iterative (n = 1e5).
int main() {
    ios::sync_with_stdio(false);
    cin.tie(nullptr);
    int n;
    cin >> n;
    vector<int> c(n + 1);
    for (int i = 1; i <= n; i++) cin >> c[i];
    vector<vector<int>> adj(n + 1);
    for (int i = 0; i < n - 1; i++) {
        int a, b;
        cin >> a >> b;
        adj[a].push_back(b);
        adj[b].push_back(a);
    }
    vector<int> cnt(100001, 0), par(n + 1, 0);
    vector<size_t> it(n + 1, 0);
    vector<char> good(n + 1, 0);
    vector<int> st(1, 1);
    good[1] = 1; cnt[c[1]]++;                       // enter the root
    while (!st.empty()) {
        int u = st.back();
        if (it[u] == adj[u].size()) {               // leave u: its colour leaves the path
            cnt[c[u]]--;
            st.pop_back();
            continue;
        }
        int v = adj[u][it[u]++];
        if (v == par[u]) continue;
        par[v] = u;
        good[v] = (cnt[c[v]] == 0);                 // the path 1..u is exactly the stack
        cnt[c[v]]++;                                // enter v
        st.push_back(v);
    }
    for (int i = 1; i <= n; i++)
        if (good[i]) cout << i << "\n";
}
\end{lstlisting}
Samples: \textbf{1 2 3 4 6} and \textbf{1 2 3 5 6 7 8}. A path of $10^5$ vertices runs in 0.5\,s (mostly input). Checked against a brute force (BFS parents, then walk up from every $x$ and compare colours, \bigO{n^2}) on 500 random trees with $N \le 12$ and 3 colours: identical.
\end{solutionbox}

\begin{solutionbox}[E4.6 Path Prefixes]
Keep on the stack the prefix sums of $b$ along the root path: $\texttt{pre}[0] = 0$, $\texttt{pre}[d]$ = sum of $b$ on the first $d$ edges. Entering $v$ pushes $\texttt{pre.back()} + b_v$, leaving $v$ pops it. Because all $b_j \ge 1$, \texttt{pre} is strictly increasing, so $r_v$ = (the largest $d$ with $\texttt{pre}[d] \le A_v$) is one \texttt{upper\_bound}. Total \bigO{n \log n}; the tree can be a chain of $2\cdot10^5$ vertices, so the DFS is iterative.
\begin{center}
\begin{tikzpicture}[x=1.05cm, y=0.8cm]
    \tikzset{gv/.append style={minimum size=6.2mm, font=\scriptsize\bfseries}}
    \node[gvseen] (p1) at (0,0) {1}; \node[gvseen] (p2) at (1.8,0) {2}; \node[gvseen] (p4) at (3.6,0) {4}; \node[gvcur] (p3) at (5.4,0) {3};
    \draw[gtree] (p1) -- node[gw, above]{5 / 6} (p2);
    \draw[gtree] (p2) -- node[gw, above]{9 / 10} (p4);
    \draw[gtree] (p4) -- node[gw, above]{5 / 1} (p3);
    \node[font=\scriptsize, anchor=west] at (6.0,0) {$A_3 = 5 + 9 + 5 = 19$};
    \foreach \v/\i in {0/0, 6/1, 16/2, 17/3} {
        \node[cell] (q\i) at (\i*1.8,-1.2) {\v};
        \node[idx] at (\i*1.8,-1.85) {\texttt{pre[\i]}};
    }
    \node[font=\scriptsize, anchor=west, align=left] at (6.0,-1.2) {largest $d$ with $\texttt{pre}[d] \le 19$:\\$d = 3$, so $r_3 = 3$};
\end{tikzpicture}
\end{center}
{\small Sample 1 at vertex 3 (edge labels $a / b$): the stack is the root path $1, 2, 4, 3$, and \texttt{pre} holds its $b$-prefix sums.}
\begin{lstlisting}[style=cpp]
#include <bits/stdc++.h>
using namespace std;
typedef long long ll;

// Path Prefixes (CF 1714G). Keep the prefix sums of b along the root path on a stack:
// pre[d] = b-sum of the first d edges of the path. Since b > 0, pre is increasing, so
// r_i = (largest d with pre[d] <= A_i) = upper_bound(pre, A_i) - 1. Iterative DFS.
int main() {
    ios::sync_with_stdio(false);
    cin.tie(nullptr);
    int t;
    cin >> t;
    while (t--) {
        int n;
        cin >> n;
        vector<vector<int>> child(n + 1);
        vector<ll> a(n + 1), b(n + 1);
        for (int j = 2; j <= n; j++) {
            int p;
            cin >> p >> a[j] >> b[j];
            child[p].push_back(j);
        }
        vector<ll> r(n + 1, 0), A(n + 1, 0);
        vector<ll> pre(1, 0);                    // the path stack: pre[0] = 0 for the empty prefix
        vector<size_t> it(n + 1, 0);
        vector<int> st(1, 1);
        while (!st.empty()) {
            int u = st.back();
            if (it[u] == child[u].size()) {      // leave u: its edge leaves the path
                st.pop_back();
                pre.pop_back();
                continue;
            }
            int v = child[u][it[u]++];
            A[v] = A[u] + a[v];
            pre.push_back(pre.back() + b[v]);    // enter v: push the b-sum of the path 1..v
            r[v] = (upper_bound(pre.begin(), pre.end(), A[v]) - pre.begin()) - 1;
            st.push_back(v);
        }
        for (int i = 2; i <= n; i++) cout << r[i] << " \n"[i == n];
    }
}
\end{lstlisting}
The sample (4 test cases) is reproduced: \texttt{0 3 1 2 1 1 2 3}, \texttt{0 0 3}, \texttt{1 2 2}, \texttt{0 1 2 1 1 2 2 1 1}. (When the root is popped, \texttt{pre} becomes empty --- harmless, the loop ends.) A chain of $2\cdot10^5$ vertices runs in 1.2\,s here, 0.74\,s of it reading the input. Checked against a brute force (walk every root path explicitly, \bigO{n^2}) on 500 random multi-test inputs with $n \le 10$: identical.
\end{solutionbox}

\begin{solutionbox}[E4.7 Grid Paths]
A plain backtracking that only respects the pattern is far too slow (the free squares have to be avoided wisely). Three cuts, all provable:
\begin{itemize}[leftmargin=*]
    \item \textbf{Target too early.} Standing on the lower-left square before step 48 is a failure: return.
    \item \textbf{Split.} If the square straight ahead (in the direction of the last move) is visited or outside, while the two side squares are free, return: the path so far together with the border is a closed curve through the current square that separates the two side squares, and both must still be visited.
    \item \textbf{Forced square.} A free neighbour $w$ (not the target) with at most one free neighbour of its own can only be entered from here, now --- later it would be a dead end. If two neighbours are forced, return; if one is, take it.
\end{itemize}
\begin{center}
\begin{tikzpicture}[x=0.55cm, y=0.55cm]
    \tikzset{c/.style={gcell, minimum size=5.3mm, font=\tiny}, v/.style={c, fill=hitfill}, cu/.style={c, fill=currentfill, draw=accentorange, line width=0.9pt},
             fr/.style={c, fill=basefill}}
    \begin{scope}
        \foreach \x in {0,...,6} \foreach \y in {0,...,6} \node[c] at (\x,-\y) {};
        \foreach \x in {0,...,3} \node[v] at (\x,0) {};
        \foreach \y in {1,...,5} \node[v] at (3,-\y) {};
        \node[cu] at (3,-6) {}; \node[fr] at (2,-6) {L}; \node[fr] at (4,-6) {R};
        \draw[warningred, line width=1.5pt] (2.5,-6.55) -- (3.5,-6.55);
        \node[font=\scriptsize\bfseries] at (3,1.0) {prune 1: moving down into the border};
        \node[font=\scriptsize, align=center, text width=4.4cm] at (3,-7.6) {L and R lie on different sides of the path: whichever way it turns, the other side is lost};
    \end{scope}
    \begin{scope}[xshift=5.6cm]
        \foreach \x in {0,...,6} \foreach \y in {0,...,6} \node[c] at (\x,-\y) {};
        \foreach \x in {0,...,6} \node[v] at (\x,0) {};
        \foreach \x in {2,...,6} \node[v] at (\x,-1) {};
        \node[cu] at (1,-1) {}; \node[fr] at (0,-1) {w};
        \node[font=\scriptsize\bfseries] at (3,1.0) {prune 2: a forced square};
        \node[font=\scriptsize, align=center, text width=4.4cm] at (3,-7.6) {$w$ has one free neighbour left (below it): enter it now, or it becomes a dead end};
    \end{scope}
\end{tikzpicture}
\end{center}
\begin{lstlisting}[style=cpp]
#include <bits/stdc++.h>
using namespace std;

// Grid Paths (CSES 1625): paths of 48 steps from the upper-left to the lower-left square of a
// 7x7 grid that visit every square once and match a pattern of D, U, L, R and '?'.
const int K = 7, STEPS = K * K - 1;
bool vis[K + 2][K + 2];          // a frame of visited squares around the grid: no bounds checks
string pat;
long long answer = 0;
const int dr[4] = {1, -1, 0, 0}, dc[4] = {0, 0, -1, 1};   // D U L R
const string DIRS = "DULR";

bool freeSq(int r, int c) { return !vis[r][c]; }

void go(int r, int c, int step, int lastDir) {
    if (r == K && c == 1) {                      // at the target (row K, column 1 in framed coordinates)
        if (step == STEPS) answer++;             // ... a full path; otherwise the target came too early
        return;
    }
    if (step == STEPS) return;
    // prune 1: the square straight ahead is blocked, but both side squares are free.
    // The free squares then fall into two parts, and the path can never visit both.
    if (lastDir >= 0) {
        int ar = r + dr[lastDir], ac = c + dc[lastDir];
        if (!freeSq(ar, ac)) {
            if (lastDir < 2) {                   // moving vertically: the sides are left and right
                if (freeSq(r, c - 1) && freeSq(r, c + 1)) return;
            } else {                             // moving horizontally: the sides are up and down
                if (freeSq(r - 1, c) && freeSq(r + 1, c)) return;
            }
        }
    }
    // prune 2: a free neighbour w (not the target) with at most one free neighbour of its own
    // must be entered now, from here; later it could only be a dead end. Two such squares: fail.
    int forced = -1;
    for (int d = 0; d < 4; d++) {
        int wr = r + dr[d], wc = c + dc[d];
        if (vis[wr][wc] || (wr == K && wc == 1)) continue;
        int freeNb = 0;
        for (int e = 0; e < 4; e++) freeNb += freeSq(wr + dr[e], wc + dc[e]);
        if (freeNb <= 1) {
            if (forced >= 0) return;
            forced = d;
        }
    }
    for (int d = 0; d < 4; d++) {
        if (forced >= 0 && d != forced) continue;
        if (pat[step] != '?' && pat[step] != DIRS[d]) continue;
        int nr = r + dr[d], nc = c + dc[d];
        if (vis[nr][nc]) continue;
        vis[nr][nc] = true;                      // enter
        go(nr, nc, step + 1, d);
        vis[nr][nc] = false;                     // leave
    }
}

int main() {
    cin >> pat;
    for (int i = 0; i < K + 2; i++) vis[0][i] = vis[K + 1][i] = vis[i][0] = vis[i][K + 1] = true;
    vis[1][1] = true;
    go(1, 1, 0, -1);
    cout << answer << "\n";
}
\end{lstlisting}
\begin{center}
\small
\renewcommand{\arraystretch}{1.2}
\begin{tabular}{@{}lrrr@{}}
\toprule
\textbf{input} & \textbf{answer} & \textbf{calls, prune 1 only} & \textbf{calls, prunes 1 + 2} \\
\midrule
the sample \texttt{??????R??????U???\dots LD????D?} & 201 & 2\,642\,437 & 117\,844 \\
48 times \texttt{?} & 88\,418 & 46\,816\,433 (1.0\,s) & 2\,866\,561 (0.22\,s) \\
\bottomrule
\end{tabular}
\end{center}
With prune 1 alone the all-\texttt{?} input takes a full second here --- exactly the time limit; prune 2 cuts the calls by a factor 16. Checked by running the same program with $K = 5$ and $K = 6$ (all-\texttt{?}: 86 and 1770 paths) against a version without any pruning, on 300 and 100 random patterns (each character \texttt{?} with probability 0.8): identical.
\end{solutionbox}


%============================================================
\clearpage
\section{Pattern 5: Reverse Search --- Start from the Target}
\label{sec:ch2-p5}
%============================================================
%============================================================
% CH2 · PATTERN 5 — REVERSE SEARCH: THEORY
%============================================================
\subsection{Pattern Theory}
\label{ssec:ch2-p5-theory}

% local styles of this pattern: a reversed edge, and the whole graph D drawn reversed
\tikzset{
    g2p5-re/.style={ge, {Stealth[length=2.4mm]}-},
    g2p5-rt/.style={gtree, {Stealth[length=2.4mm]}-},
    g2p5-Drev/.style={g2-e12/.style=g2p5-re, g2-e15/.style=g2p5-re, g2-e16/.style=g2p5-re, g2-e23/.style=g2p5-re,
        g2-e24/.style=g2p5-re, g2-e31/.style=g2p5-re, g2-e43/.style=g2p5-re, g2-e54/.style=g2p5-re,
        g2-e56/.style=g2p5-re, g2-e67/.style=g2p5-re, g2-e76/.style=g2p5-re, g2-e87/.style=g2p5-re},
    % grid cells, used as  \foreach \s [count=\c from 0] in {w,e,o,...} \node[g2p5-\s] at (\c,-\r) {};
    g2p5-e/.style={gcell, minimum size=6.5mm, font=\scriptsize\bfseries},
    g2p5-w/.style={g2p5-e, fill=midgray!70},
    g2p5-o/.style={g2p5-e, fill=computedfill},
    g2p5-i/.style={g2p5-e, fill=answerfill},
    g2p5-G/.style={g2p5-e, fill=basefill, label={center:G}},
    g2p5-B/.style={g2p5-e, fill=dupfill, label={center:B}},
    g2p5-Go/.style={g2p5-e, fill=computedfill, draw=successgreen, line width=1.1pt, label={center:G}},
    g2p5-x/.style={g2p5-e, fill=midgray!70, draw=recursionred, line width=1.1pt}
}

\begin{patternbox}{Pattern 5 at a Glance}
``Which vertices can reach the target $t$?'' is $n$ questions, one per start vertex. Turn every edge around: $u$ reaches $t$ in $G$ \textbf{iff} $t$ reaches $u$ in the reversed graph $G^R$. One search from $t$ in $G^R$ answers all $n$ questions at once --- and a search from a \emph{set} of targets answers ``who can reach any of them?'' just as cheaply.

\begin{center}
\begin{tikzpicture}[x=1.0cm, y=0.95cm]
    \begin{scope}
        \pic {g2-D};
        \node[font=\scriptsize\bfseries] at (2.4,2.25) {$D$: who can reach 4? ($8$ forward searches)};
    \end{scope}
    \begin{scope}[xshift=7.2cm, g2p5-Drev, g2-v4/.style=gvdst, g2-v1/.style=gvdone, g2-v2/.style=gvdone,
                  g2-v3/.style=gvdone, g2-v5/.style=gvdone, g2-v6/.style=gvdead, g2-v7/.style=gvdead, g2-v8/.style=gvdead,
                  g2-e24/.style=g2p5-rt, g2-e54/.style=g2p5-rt, g2-e15/.style=g2p5-rt, g2-e31/.style=g2p5-rt]
        \pic {g2-D};
        \node[font=\scriptsize\bfseries] at (2.4,2.25) {$D^R$: one search from 4};
    \end{scope}
    \draw[flow] (5.3,0.2) -- node[lbl, above]{reverse} (6.8,0.2);
\end{tikzpicture}
\end{center}
\begin{recurrencebox}
\centering\small
build \texttt{radj[v]} $= \{u : u \to v\}$\quad;\quad push every target, mark it\quad;\quad pop $v$, push every unmarked $u \in \texttt{radj}[v]$\\[3pt]
marked $=$ $\{u : u \leadsto T\}$ \qquad one search, \VE{} instead of $n$ searches, \bigO{n(n+m)}
\end{recurrencebox}
\end{patternbox}

%------------------------------------------------------------
\subsubsection*{From the lessons to the pattern}

Tool~5 of Lesson~\ref{ssec:ch2-toolbox} showed the trick in two lines. This pattern is about \emph{seeing} it in a statement, because the statement almost never says ``reversed graph''. It says ``every good person must be able to escape'', ``water flows to the ocean'', ``a region is captured unless it touches the border'', ``list all ancestors''. In each of these the question is asked from the point of view of many starting points, all heading to the same few places. The pattern is to change the point of view:
\begin{itemize}[leftmargin=*]
    \item \textbf{reverse the edges} when the graph is directed (an edge $u \to v$ becomes $v \to u$; in a grid with rules like ``water flows downhill'' the reversed rule is ``climb uphill'');
    \item \textbf{reverse the roles} when the graph is undirected: the search does not have to change at all, only its start --- from the exit, the border, the ocean, instead of from every cell;
    \item \textbf{start from all targets at once} (several exits, the whole border, every cell of a set $T$): one search with many sources costs the same as a search with one.
\end{itemize}

\begin{diagrambox}
\begin{tikzpicture}[x=1cm, y=1cm]
    \begin{scope}
        \node[font=\small\bfseries] at (2.0,1.55) {forward: one search per start};
        \foreach \i/\y in {1/1.0, 2/0.35, 3/-0.3, 4/-0.95}{
            \node[gvsrc, minimum size=6mm, font=\scriptsize\bfseries] (s\i) at (0,\y) {$s_\i$};
            \draw[gtree, -{Stealth[length=2.2mm]}, decorate, decoration={snake, amplitude=0.6mm, segment length=4mm, post length=2mm}] (s\i) -- (3.4,0.05);
        }
        \node[gvdst] (t) at (3.8,0.05) {$t$};
        \node[note, text width=4.2cm] at (2.0,-1.65) {$n$ searches, each may cost \bigO{n+m}: \bigO{n(n+m)}};
    \end{scope}
    \begin{scope}[xshift=7.6cm]
        \node[font=\small\bfseries] at (2.0,1.55) {reverse: one search from $t$};
        \foreach \i/\y in {1/1.0, 2/0.35, 3/-0.3, 4/-0.95}{
            \node[gvdone, minimum size=6mm, font=\scriptsize\bfseries] (r\i) at (0,\y) {$s_\i$};
        }
        \node[gvdst] (t2) at (3.8,0.05) {$t$};
        \foreach \i/\y in {1/1.0, 2/0.35, 3/-0.3, 4/-0.95}
            \draw[gtree, {Stealth[length=2.2mm]}-] (r\i) -- (t2);
        \node[note, text width=4.2cm] at (2.0,-1.65) {one search in $G^R$, every vertex and edge once: \VE};
    \end{scope}
\end{tikzpicture}
\figcap{The same question asked from both ends. The answer is the same set of vertices; only the cost differs.}
\end{diagrambox}

%------------------------------------------------------------
\subsubsection*{When to recognise it (and when not)}

\begin{itemize}[leftmargin=*]
    \item ``\textbf{from which} vertices / cells can you reach the exit, the border, the ocean, a vertex of $T$?'' --- one search from the targets in $G^R$;
    \item ``\textbf{for every} vertex, all its ancestors / all vertices that can reach it'' --- one reverse search \emph{per vertex}, \bigO{n(n+m)}, fine for $n \le 1000$--$2000$ (Problem~\ref{ssec:ch2-all-ancestors});
    \item ``every good one must get out, no bad one may'' --- first make the bad ones harmless, then one search from the exit checks all good ones (Problem~\ref{ssec:ch2-solve-maze});
    \item ``regions that do \emph{not} touch the border'', ``lakes'', ``captured'', ``surrounded'' --- search from the border; whatever stays unmarked is enclosed;
    \item ``can the walk go on forever?'', ``which states are safe / doomed?'' --- peel dead ends backwards (reverse peeling, below);
    \item ``does everybody reach 1 \emph{and} does 1 reach everybody?'' --- one search from 1 in $G$, one in $G^R$.
\end{itemize}
\textbf{Not this pattern:} \emph{shortest distances} to a target from all vertices need BFS in $G^R$ (and from many sources a multi-source BFS, Chapter~3) or Dijkstra with weights; DFS gives the right \emph{set} but not the distances. Full \emph{strong connectivity} --- which pairs reach each other --- is the strongly connected components of Chapter~19; the two-search test above only answers ``is the whole graph one component?''. Reachability for all $n^2$ pairs on large graphs needs a DAG of components and bitsets, not $n$ searches.

%------------------------------------------------------------
\subsubsection*{The algorithm and its invariant}

\begin{invariantbox}[Why one reverse search is enough]
\textbf{Lemma.} $u \leadsto v$ in $G$ iff $v \leadsto u$ in $G^R$. A path $u = x_0 \to x_1 \to \dots \to x_k = v$ in $G$ is, read backwards, the path $v = x_k \to x_{k-1} \to \dots \to x_0 = u$ in $G^R$, since every edge was turned around; the converse is the same argument applied to $G^R$, whose reverse is $G$.

\textbf{Multi-source.} Add a new vertex $\star$ with edges $\star \to t$ for every $t \in T$ in $G^R$. A search from $\star$ marks exactly the vertices $v$ with $t \leadsto v$ in $G^R$ for some $t$, i.e.\ with $v \leadsto T$ in $G$. Pushing all of $T$ on the stack before the loop is exactly ``entering $\star$''.

\textbf{Invariant of the loop.} Every marked vertex can reach $T$ in $G$ (it was marked along a reversed edge from a marked vertex), and when the stack is empty no reversed edge leaves the marked set --- so, by the closure argument of Chapter~1, every vertex that can reach $T$ is marked.

\textbf{Cost.} Building \texttt{radj} is \bigO{m}; every vertex is pushed once, every reversed edge looked at once: \VE, the same as one forward search.
\end{invariantbox}

\begin{diagrambox}
\begin{tikzpicture}[x=1.0cm, y=0.95cm]
    \begin{scope}[g2p5-Drev, g2-v2/.style=gvdst, g2-v8/.style=gvdst, g2-v1/.style=gvdone, g2-v3/.style=gvdone,
                  g2-v4/.style=gvdone, g2-v5/.style=gvdone, g2-v6/.style=gvdead, g2-v7/.style=gvdead,
                  g2-e12/.style=g2p5-rt, g2-e31/.style=g2p5-rt, g2-e43/.style=g2p5-rt, g2-e54/.style=g2p5-rt]
        \pic {g2-D};
        \node[font=\scriptsize\bfseries] at (2.4,2.25) {$D^R$ from $T = \{2, 8\}$ at once};
    \end{scope}
    \node[note, text width=6.2cm, align=left, anchor=west] at (6.0,0.3) {Stack at the start: \texttt{2 8}. Pop 8: nothing points into 8. Pop 2: mark 1 (edge $1 \to 2$). Pop 1: mark 3 ($3 \to 1$). Pop 3: 2 is marked, mark 4 ($4 \to 3$). Pop 4: 2 marked, mark 5 ($5 \to 4$). Pop 5: 1 is marked.\\[3pt]
    Marked (template output): \texttt{1 2 3 4 5 8}. The targets are marked too --- a vertex of $T$ ``reaches'' $T$ in zero steps. Vertices 6 and 7 cannot reach $T$: their only exits lead to each other.};
\end{tikzpicture}
\figcap{A multi-target reverse search on $D$. Thick blue: the reversed edges along which a new vertex was marked.}
\end{diagrambox}

%------------------------------------------------------------
\subsubsection*{Reversing roles in grids}

In an undirected grid nothing has to be reversed, yet the same idea saves the day: ``which empty cells are \emph{not} connected to the border?'' asked cell by cell is one search per cell, \bigO{(RC)^2}. Asked from the other end, it is one search that starts from \emph{all} empty border cells at once; the unmarked empty cells are the enclosed ones (template 2).

\begin{diagrambox}
\begin{tikzpicture}[x=6.5mm, y=6.5mm]
    \begin{scope}
        \foreach \r/\row in {0/{w,w,w,w,w}, 1/{w,e,e,w,e}, 2/{w,e,w,e,e}, 3/{w,w,w,w,w}}
            \foreach \s [count=\c from 0] in \row \node[g2p5-\s] at (\c,-\r) {};
        \node[font=\scriptsize\bfseries] at (2,1.1) {input (grey $=$ wall)};
    \end{scope}
    \begin{scope}[xshift=5.2cm]
        \foreach \r/\row in {0/{w,w,w,w,w}, 1/{w,i,i,w,o}, 2/{w,i,w,o,o}, 3/{w,w,w,w,w}}
            \foreach \s [count=\c from 0] in \row \node[g2p5-\s] at (\c,-\r) {};
        \draw[flow] (5.1,-1) -- (4.45,-1);
        \draw[flow] (5.1,-2) -- (4.45,-2);
        \node[font=\scriptsize\bfseries] at (2,1.1) {search from the border};
        \node[note, text width=3.6cm, align=left, anchor=west] at (5.9,-1.5) {start: the two empty border cells (arrows); blue $=$ connected to the border, orange $=$ enclosed};
    \end{scope}
\end{tikzpicture}
\figcap{Template 2 on the grid \texttt{\#\#\#\#\# / \#..\#. / \#.\#.. / \#\#\#\#\#}: output \texttt{\#\#\#\#\# \#..\#o \#.\#oo \#\#\#\#\#}. The three orange cells are enclosed. This is Surrounded Regions (E5.1) and the lakes of E5.3.}
\end{diagrambox}

The two-ocean version (E5.2) runs this twice: once from the Pacific border, once from the Atlantic border, with the rule reversed (``flows down'' becomes ``climbs up''), and intersects the two marked sets.

%------------------------------------------------------------
\subsubsection*{Block the bad, then check the good}

Some problems mix two kinds of starts: \emph{good} ones that must reach the target and \emph{bad} ones that must not. Asking ``can good $g$ get out?'' for every $g$ separately is again $n$ searches --- and the answer depends on which cells we decide to block. The reverse view untangles both:
\begin{enumerate}[leftmargin=*]
    \item make every bad start \textbf{harmless} by the cheapest local change that is forced anyway (in Problem~\ref{ssec:ch2-solve-maze}: wall in the four neighbours of every bad person);
    \item then run \textbf{one} search from the target and check that it marks every good start and no bad one.
\end{enumerate}
Step 1 is safe because any solution must already cut every bad start from the target, and a good start whose path passes \emph{next to} a bad one is lost in every solution. The proof is in the problem.

\begin{diagrambox}
\begin{tikzpicture}[x=6.5mm, y=6.5mm]
    \begin{scope}
        \foreach \r/\row in {0/{G,e,e,e}, 1/{e,e,B,e}, 2/{e,w,e,e}}
            \foreach \s [count=\c from 0] in \row \node[g2p5-\s] at (\c,-\r) {};
        \node[font=\scriptsize\bfseries] at (1.5,1.1) {exit $=$ bottom right};
    \end{scope}
    \begin{scope}[xshift=3.3cm]
        \foreach \r/\row in {0/{G,e,x,e}, 1/{e,x,B,x}, 2/{e,w,x,e}}
            \foreach \s [count=\c from 0] in \row \node[g2p5-\s] at (\c,-\r) {};
        \node[font=\scriptsize\bfseries] at (1.5,1.1) {1. wall in B (red)};
    \end{scope}
    \begin{scope}[xshift=6.6cm]
        \foreach \r/\row in {0/{G,e,x,e}, 1/{e,x,B,x}, 2/{e,w,x,o}}
            \foreach \s [count=\c from 0] in \row \node[g2p5-\s] at (\c,-\r) {};
        \node[font=\scriptsize\bfseries] at (1.5,1.1) {2. search from the exit};
    \end{scope}
    \node[note, text width=4.3cm, align=left, anchor=west] at (15.0,-1) {Both neighbours of the exit were walled, so the search from the exit (blue) marks only the exit: G is cut off and the answer is \emph{no} --- every route from G to the exit passes next to B, so no choice of walls does better.};
\end{tikzpicture}
\figcap{The two steps of ``block the bad, check the good''. The blue cells are reached from the exit; G is not among them.}
\end{diagrambox}

%------------------------------------------------------------
\subsubsection*{All ancestors: one reverse search per vertex}

``For every $v$, list all vertices that can reach $v$'' has an output of size up to $n^2/2$, so no algorithm beats $\Theta(n^2)$ in the worst case. One reverse search from every $v$ costs \bigO{n(n+m)} and produces each list already \emph{sorted} if we collect it by scanning $u = 0, 1, \dots, n-1$ (template 3). The mirror image works as well: a forward search from every $u$ in increasing order appends $u$ to every vertex it reaches --- also sorted for free.

\begin{diagrambox}
\begin{tikzpicture}[x=1.15cm, y=1cm]
    \begin{scope}
        \node[gv] (a0) at (0,0.7) {0}; \node[gv] (a1) at (0,-0.7) {1};
        \node[gv] (a2) at (1.4,0.7) {2}; \node[gv] (a3) at (2.8,0.7) {3}; \node[gv] (a4) at (1.4,-0.7) {4};
        \draw[gde] (a0) -- (a2); \draw[gde] (a1) -- (a2); \draw[gde] (a2) -- (a3); \draw[gde] (a1) -- (a4);
        \node[font=\scriptsize\bfseries] at (1.4,1.5) {a DAG};
    \end{scope}
    \begin{scope}[xshift=4.6cm]
        \node[gvdone] (b0) at (0,0.7) {0}; \node[gvdone] (b1) at (0,-0.7) {1};
        \node[gvdone] (b2) at (1.4,0.7) {2}; \node[gvdst] (b3) at (2.8,0.7) {3}; \node[gvdead] (b4) at (1.4,-0.7) {4};
        \draw[g2p5-rt] (b0) -- (b2); \draw[g2p5-rt] (b1) -- (b2); \draw[g2p5-rt] (b2) -- (b3); \draw[g2p5-re] (b1) -- (b4);
        \node[font=\scriptsize\bfseries] at (1.4,1.5) {reverse search from 3};
    \end{scope}
    \node[note, text width=3.6cm, align=left, anchor=west] at (8.2,0.1) {template 3 prints\\\texttt{0: | 1: | 2: 0 1 |}\\\texttt{3: 0 1 2 | 4: 1 |}\\[2pt]one search per vertex; 4 is not an ancestor of 3 although both are below 1.};
\end{tikzpicture}
\figcap{Ancestors of 3 = the vertices marked by a search from 3 along reversed edges.}
\end{diagrambox}

In a DAG there is a faster-looking alternative: process the vertices in topological order (Chapter~5) and set $\mathrm{anc}(v) = \bigcup_{u \to v} \big(\mathrm{anc}(u) \cup \{u\}\big)$ with bitsets, \bigO{m \cdot n / 64}. For $n = 1000$, $m = 2000$ it ran in 7\,ms against 13\,ms for the $n$ reverse searches --- both instantaneous. The bitsets need $n^2/8$ bytes, though: fine for $n = 10^4$ (12.5\,MB), impossible for $n = 10^5$. And the per-vertex searches also work when the graph has cycles.

%------------------------------------------------------------
\subsubsection*{Reverse peeling: where can a walk go on forever?}

A vertex of out-degree 0 is a dead end: every walk that enters it stops. Delete it. Its \emph{predecessors} lose one exit each, and some of them may become dead ends in turn --- the reversed edges tell us exactly whose out-degree drops. When nothing more can be deleted, every remaining vertex has an exit to another remaining vertex, so a walk can go on forever; every deleted vertex leads only into dead ends. This is Kahn's algorithm (Chapter~5) run on $G^R$: peeling in-degree-0 vertices of $G^R$ $=$ peeling out-degree-0 vertices of $G$ (template 4, Exercise E5.4).

\begin{diagrambox}
\begin{tikzpicture}[x=1.25cm, y=1cm]
    \node[gvsel] (p1) at (0,0.8) {1}; \node[gvsel] (p2) at (1.2,1.4) {2}; \node[gvsel] (p3) at (1.2,0.1) {3};
    \node[gvdead] (p4) at (2.6,0.1) {4}; \node[gvdead] (p5) at (4.0,0.1) {5};
    \node[gvdead] (p6) at (2.6,-1.3) {6}; \node[gvdead] (p7) at (4.0,-1.3) {7};
    \draw[gde] (p1) -- (p2); \draw[gde] (p2) -- (p3); \draw[gde] (p3) -- (p1);
    \draw[gfaded, -{Stealth[length=2.2mm]}] (p3) -- (p4); \draw[gfaded, -{Stealth[length=2.2mm]}] (p4) -- (p5);
    \draw[gfaded, -{Stealth[length=2.2mm]}] (p6) -- (p4); \draw[gfaded, -{Stealth[length=2.2mm]}] (p6) -- (p7);
    \draw[gfaded, -{Stealth[length=2.2mm]}] (p7) -- (p5);
    \node[gdist, text=recursionred, above=1pt of p5] {round 1};
    \node[gdist, text=recursionred, above=1pt of p4] {round 2};
    \node[gdist, text=recursionred, right=1pt of p7] {round 2};
    \node[gdist, text=recursionred, left=1pt of p6] {round 3};
    \node[note, text width=5.4cm, align=left, anchor=west] at (5.2,0.1) {5 has out-degree 0. Removing it drops 4 and 7 to out-degree 0; removing those drops 6 to 0. Vertex 3 keeps its edge $3 \to 1$.\\[2pt]Template 4 deletes $5, 7, 4, 6$ and prints \texttt{1 2 3}: only the cycle and what still leads into it survive.};
\end{tikzpicture}
\figcap{Reverse peeling. Grey: deleted vertices; purple: vertices from which an endless walk exists.}
\end{diagrambox}

%------------------------------------------------------------
\subsubsection*{Variations map}

\begin{diagrambox}
\begin{tikzpicture}[
    var/.style={rectangle, rounded corners=2pt, draw=patternborder, fill=patternbg,
                font=\scriptsize, align=center, text width=2.95cm, minimum height=1.3cm, inner sep=3pt}]
    \node[var, fill=answerfill, draw=accentorange, text width=5.4cm, font=\scriptsize\bfseries] (root) at (0,0)
        {who can reach the target? search FROM the target in $G^R$};
    \node[var] (v1) at (-5.1,-2) {\textbf{all ancestors}\\one reverse search each\\(2.5.1)};
    \node[var] (v2) at (-1.7,-2) {\textbf{good in, bad out}\\wall the bad, search from the exit (2.5.2)};
    \node[var] (v3) at (1.7,-2)  {\textbf{from the border}\\enclosed $=$ unmarked\\(E5.1, E5.3)};
    \node[var] (v4) at (5.1,-2)  {\textbf{two targets}\\two searches, intersect\\(E5.2)};
    \node[var] (v5) at (-5.1,-3.8) {\textbf{endless walks}\\peel dead ends backwards\\(E5.4)};
    \node[var] (v6) at (-1.7,-3.8) {\textbf{a target set $T$}\\all of $T$ on the stack\\(E5.5)};
    \node[var] (v7) at (1.7,-3.8)  {\textbf{only cheap edges}\\search inside a binary search (E5.6)};
    \node[var] (v8) at (5.1,-3.8)  {\textbf{strongly connected?}\\one search in $G$, one in $G^R$ (template 5)};
    \node[draw=patternborder, dashed, rounded corners=4pt, fit=(v1)(v4)(v5)(v8), inner sep=5pt] (grp) {};
    \draw[flow] (root.south) -- (grp.north);
    \node[note, below=2mm of grp] {distances to the target: BFS / Dijkstra on $G^R$ (Chapters 3, 7); all pairs: SCCs + bitsets (Chapter 19)};
\end{tikzpicture}
\end{diagrambox}

\textbf{Twists and how to handle them}
\begin{itemize}[leftmargin=*]
    \item \textbf{The rule, not the graph, is directed} (``flow to a cell that is not higher''): reverse the rule (``climb to a cell that is not lower''), never build $G^R$ explicitly.
    \item \textbf{Targets are a whole region} (the border, all cells of one colour): push them all before the loop.
    \item \textbf{Two different targets} (two oceans, ``reaches $a$ and reaches $b$''): two reverse searches and an AND of the marks.
    \item \textbf{Bad starts must not escape}: neutralise them first (wall their neighbourhood), then search once from the target.
    \item \textbf{Output per vertex is a list}: one reverse search per vertex, \bigO{n(n+m)}; collect by scanning $u$ in order to get sorted lists without sorting.
    \item \textbf{``Forever'', ``safe'', ``doomed''}: peel out-degree-0 vertices using the reversed edges.
    \item \textbf{Strongly connected?} ``Everyone reaches 1 and 1 reaches everyone'' is one search from 1 in $G$ and one in $G^R$; both must mark all $n$ vertices (on $D$ the template marks 7 and 5: $D$ is not strongly connected).
\end{itemize}

%------------------------------------------------------------
\subsubsection*{Code templates}

\begin{lstlisting}[style=cpp]
#include <bits/stdc++.h>
using namespace std;

// (1) WHO CAN REACH T?  One iterative DFS in the reversed graph, from ALL targets at once.
//     radj[v] lists the u with an edge u -> v. Returns can[u] = 1 iff u reaches some target.
vector<char> canReach(const vector<vector<int>>& radj, const vector<int>& targets) {
    vector<char> can(radj.size(), 0);
    vector<int> st;
    for (int t : targets) if (!can[t]) { can[t] = 1; st.push_back(t); }
    while (!st.empty()) {
        int v = st.back(); st.pop_back();
        for (int u : radj[v])                    // reversed edge v -> u = original u -> v
            if (!can[u]) { can[u] = 1; st.push_back(u); }
    }
    return can;
}

// (2) ROLES REVERSED IN A GRID: mark every '.' cell connected to the border, from the border.
vector<string> markFromBorder(vector<string> g) {
    int R = g.size(), C = g[0].size();
    int dr[] = {1, -1, 0, 0}, dc[] = {0, 0, 1, -1};
    vector<pair<int, int>> st;
    for (int r = 0; r < R; r++)
        for (int c = 0; c < C; c++) {
            bool border = (r == 0 || c == 0 || r == R - 1 || c == C - 1);
            if (border && g[r][c] == '.') { g[r][c] = 'o'; st.push_back({r, c}); }
        }
    while (!st.empty()) {
        int r = st.back().first, c = st.back().second; st.pop_back();
        for (int d = 0; d < 4; d++) {
            int nr = r + dr[d], nc = c + dc[d];
            if (nr < 0 || nc < 0 || nr >= R || nc >= C || g[nr][nc] != '.') continue;
            g[nr][nc] = 'o'; st.push_back({nr, nc});
        }
    }
    return g;                                    // 'o' = reaches the border, '.' left = enclosed
}

// (3) ALL ANCESTORS: one reverse search per vertex, O(n (n + m)). Lists come out sorted.
vector<vector<int>> allAncestors(const vector<vector<int>>& radj) {
    int n = radj.size();
    vector<vector<int>> anc(n);
    for (int v = 0; v < n; v++) {
        vector<char> can = canReach(radj, vector<int>(1, v));
        for (int u = 0; u < n; u++) if (u != v && can[u]) anc[v].push_back(u);
    }
    return anc;
}

// (4) REVERSE PEELING: vertices from which an infinite walk exists.
//     Delete vertices of out-degree 0; the reversed edges tell whose out-degree drops.
vector<char> endless(const vector<vector<int>>& adj, const vector<vector<int>>& radj) {
    int n = adj.size();
    vector<int> out(n), q;
    vector<char> alive(n, 1);
    for (int v = 0; v < n; v++) { out[v] = adj[v].size(); if (!out[v]) q.push_back(v); }
    while (!q.empty()) {
        int v = q.back(); q.pop_back();
        alive[v] = 0;                            // v is a dead end: every walk from v stops
        for (int u : radj[v]) if (--out[u] == 0) q.push_back(u);
    }
    return alive;
}

int main() {
    // the graph D of Part A, 0-indexed (vertex k is k-1)
    vector<vector<int>> D = {{1, 4, 5}, {2, 3}, {0}, {2}, {3, 5}, {6}, {5}, {6}};
    int n = D.size();
    vector<vector<int>> R(n);
    for (int u = 0; u < n; u++) for (int v : D[u]) R[v].push_back(u);
    vector<char> a = canReach(R, {3});
    for (int u = 0; u < n; u++) if (a[u]) cout << u + 1 << " ";      // 1 2 3 4 5
    cout << "\n";
    vector<char> b = canReach(R, {1, 7});                              // targets 2 and 8
    for (int u = 0; u < n; u++) if (b[u]) cout << u + 1 << " ";      // 1 2 3 4 5 8
    cout << "\n";
    vector<string> g = markFromBorder({"#####", "#..#.", "#.#..", "#####"});
    for (auto& s : g) cout << s << " ";                               // ##### #..#o #.#oo #####
    cout << "\n";
    // a DAG: 0->2, 1->2, 2->3, 1->4
    vector<vector<int>> Rd = {{}, {}, {0, 1}, {2}, {1}};
    vector<vector<int>> an = allAncestors(Rd);
    for (int v = 0; v < 5; v++) {                // 0: | 1: | 2: 0 1 | 3: 0 1 2 | 4: 1 |
        cout << v << ":";
        for (int u : an[v]) cout << " " << u;
        cout << " | ";
    }
    cout << "\n";
    // P: a cycle 1 -> 2 -> 3 -> 1 with a tail 3 -> 4 -> 5 and 6 -> 4, 6 -> 7 -> 5
    vector<vector<int>> P = {{1}, {2}, {0, 3}, {4}, {}, {3, 6}, {4}}, PR(7);
    for (int u = 0; u < 7; u++) for (int v : P[u]) PR[v].push_back(u);
    vector<char> e = endless(P, PR);
    for (int u = 0; u < 7; u++) if (e[u]) cout << u + 1 << " ";      // 1 2 3
    cout << "\n";
    // (5) "1 reaches everyone and everyone reaches 1": one search in D, one in D^R
    vector<char> f = canReach(D, {0}), r = canReach(R, {0});
    cout << count(f.begin(), f.end(), 1) << " " << count(r.begin(), r.end(), 1) << "\n";  // 7 5
}
\end{lstlisting}

Note the last two lines: \texttt{canReach} is just a search, so called with \texttt{D} instead of \texttt{R} it is the ordinary forward search. From 1, $D$ reaches 7 vertices (not 8) and only 5 vertices reach 1 --- two reasons why $D$ is not strongly connected.

\begin{insightbox}[Which template when?]
\renewcommand{\arraystretch}{1.25}
\begin{tabularx}{\linewidth}{@{}>{\bfseries}l X@{}}
(1) \texttt{canReach} & ``who can reach $t$ / any vertex of $T$?''; with \texttt{adj} instead of \texttt{radj} it is the forward search (template 5 uses both) \\
(2) \texttt{markFromBorder} & grids: enclosed regions, lakes, ``surrounded''; with a height rule it becomes the ocean search of E5.2 \\
(3) \texttt{allAncestors} & a list per vertex, $n \le$ a few thousand; for DAGs with $n \le 10^4$ bitsets in topological order are an option \\
(4) \texttt{endless} & ``walk forever'', ``safe states'', ``which positions are doomed'' --- Kahn on $G^R$ \\
\end{tabularx}
\end{insightbox}

%------------------------------------------------------------
\subsubsection*{Common mistakes}

\begin{warningbox}[Mistake 1 --- building the reversed graph the wrong way round]
\begin{lstlisting}[style=snippet]
radj[a].push_back(b);     // WRONG: this is the original edge a -> b again
radj[b].push_back(a);     // RIGHT: edge a -> b is stored at b, pointing back to a
\end{lstlisting}
On $D$ the wrong version answers ``who can reach 4?'' with $\{3, 4, 1, 2, 5, 6, 7\}$ --- the vertices reachable \emph{from} 4.
\end{warningbox}

\begin{warningbox}[Mistake 2 --- reversing the rule only halfway in a grid]
``Water flows from $a$ to a neighbour $b$ with $h_b \le h_a$.'' The reverse search goes from $b$ to $a$ when $h_a \ge h_b$. Writing \texttt{h[nr][nc] <= h[r][c]} in the reverse search runs the \emph{forward} rule from the ocean: it marks the cells that water from the coast could flow \emph{to}, which is a different set.
\end{warningbox}

\begin{warningbox}[Mistake 3 --- one search per start ``because it is simpler'']
For $n = 10^5$, $m = 2 \cdot 10^5$ random edges and three targets, a forward search from every vertex (stopping as soon as a target is found) needed $3.8 \cdot 10^8$ stack pops and 63--95 seconds (three runs); the reverse search took 8--12\,ms (Exercise E5.5).
\end{warningbox}

\begin{warningbox}[Mistake 4 --- forgetting that the targets are marked too]
Template 1 marks the targets themselves. If the question is ``which \emph{other} vertices'' (ancestors, E5.6 ``every node except 0''), exclude them when collecting the answer, or count them correctly.
\end{warningbox}

\begin{warningbox}[Mistake 5 --- peeling with in-degrees of $G$]
For ``endless walks'' the dead ends are vertices of out-degree 0, and when one is deleted the out-degree of its \emph{predecessors} drops: iterate over \texttt{radj}. Peeling in-degree-0 vertices of $G$ answers a different question: it keeps the vertices that can be reached \emph{from} a cycle, not those that can reach one (on the graph of the figure it would keep $1, \dots, 5$ and delete 6 and 7).
\end{warningbox}

\begin{checklistbox}[``I'm stuck'' checklist for Pattern 5]
\begin{itemize}[leftmargin=*]
    \item[$\square$] Is the question asked from many starts towards few targets? Then search from the targets.
    \item[$\square$] Directed graph: did I store every edge $a \to b$ in \texttt{radj[b]}? Directed rule in a grid: did I flip the comparison?
    \item[$\square$] Are there several targets (a border, an ocean, a set)? Push them all before the loop.
    \item[$\square$] Two conditions (``reaches $a$ and $b$'')? Two reverse searches and an AND.
    \item[$\square$] Bad starts that must not escape? Neutralise them first, then search once.
    \item[$\square$] ``Forever'' or ``safe''? Peel out-degree-0 vertices along reversed edges.
    \item[$\square$] Do I need distances rather than reachability? Then it is BFS (Chapter~3), not DFS.
\end{itemize}
\end{checklistbox}

%============================================================
% CH2 · P5 · PROBLEM 1 — ALL ANCESTORS OF A NODE IN A DAG (LeetCode 2192)
%============================================================
\subsection{Problem 2.5.1: All Ancestors of a Node in a Directed Acyclic Graph}
\problemlabel{2.5.1}{ssec:ch2-all-ancestors}

% the sample DAG of this problem as a pic; vertex styles g2p5-a<k>, edge styles g2p5-f<uv>
\tikzset{
    g2p5-a0/.style=gv, g2p5-a1/.style=gv, g2p5-a2/.style=gv, g2p5-a3/.style=gv,
    g2p5-a4/.style=gv, g2p5-a5/.style=gv, g2p5-a6/.style=gv, g2p5-a7/.style=gv,
    g2p5-f03/.style=gde, g2p5-f04/.style=gde, g2p5-f13/.style=gde, g2p5-f24/.style=gde, g2p5-f27/.style=gde,
    g2p5-f35/.style=gde, g2p5-f36/.style=gde, g2p5-f37/.style=gde, g2p5-f46/.style=gde,
    g2p5-allrev/.style={g2p5-f03/.style=g2p5-re, g2p5-f04/.style=g2p5-re, g2p5-f13/.style=g2p5-re,
        g2p5-f24/.style=g2p5-re, g2p5-f27/.style=g2p5-re, g2p5-f35/.style=g2p5-re, g2p5-f36/.style=g2p5-re,
        g2p5-f37/.style=g2p5-re, g2p5-f46/.style=g2p5-re},
    g2p5-S1/.pic={
        \node[g2p5-a1] (s1) at (0,2) {1}; \node[g2p5-a0] (s0) at (1.6,2) {0}; \node[g2p5-a2] (s2) at (3.2,2) {2};
        \node[g2p5-a3] (s3) at (0.8,1) {3}; \node[g2p5-a4] (s4) at (2.4,1) {4};
        \node[g2p5-a5] (s5) at (0,0) {5}; \node[g2p5-a6] (s6) at (1.6,0) {6}; \node[g2p5-a7] (s7) at (3.2,0) {7};
        \draw[g2p5-f03] (s0) -- (s3); \draw[g2p5-f04] (s0) -- (s4); \draw[g2p5-f13] (s1) -- (s3);
        \draw[g2p5-f24] (s2) -- (s4); \draw[g2p5-f27] (s2) -- (s7); \draw[g2p5-f35] (s3) -- (s5);
        \draw[g2p5-f36] (s3) -- (s6); \draw[g2p5-f37] (s3) -- (s7); \draw[g2p5-f46] (s4) -- (s6);
    }
}

\problemcard{2.5.1}{All Ancestors of a Node in a Directed Acyclic Graph}{LeetCode 2192}{https://leetcode.com/problems/all-ancestors-of-a-node-in-a-directed-acyclic-graph/}{\exmedium\ \lc{Medium}}{For every vertex, one search in the reversed graph collects its ancestors; scanning the marks in index order gives sorted lists for free}

\subsubsection*{1. Problem statement}

A directed acyclic graph has $n$ vertices $0, \dots, n-1$ and a list of edges $[\mathit{from}, \mathit{to}]$. A vertex $u$ is an \textbf{ancestor} of $v$ if $u$ can reach $v$ along the edges. Return, for every $v$, the list of its ancestors in increasing order.

\textbf{Constraints:} $1 \le n \le 1000$, $0 \le m \le \min(2000, n(n-1)/2)$; no self-loops, no duplicate edges; the graph is acyclic.

\begin{table}[H]
\centering\small
\begin{tabular}{@{}p{5.6cm}p{5.6cm}l@{}}
\toprule
\textbf{Input} & \textbf{Official output} & \textbf{Ours} \\
\midrule
$n = 8$, \texttt{[[0,3],[0,4],[1,3],[2,4],}\newline\texttt{[2,7],[3,5],[3,6],[3,7],[4,6]]} & \texttt{[[],[],[],[0,1],[0,2],[0,1,3],}\newline\texttt{[0,1,2,3,4],[0,1,2,3]]} & identical \\
$n = 5$, all 10 edges $i \to j$ with $i < j$ & \texttt{[[],[0],[0,1],[0,1,2],}\newline\texttt{[0,1,2,3]]} & identical \\
\bottomrule
\end{tabular}
\end{table}

\begin{diagrambox}
\begin{tikzpicture}[x=1.15cm, y=1.1cm]
    \begin{scope}
        \pic {g2p5-S1};
        \node[font=\scriptsize\bfseries] at (1.6,2.75) {Sample 1};
    \end{scope}
    \begin{scope}[xshift=5.6cm, g2p5-a6/.style=gvdst, g2p5-a0/.style=gvsel, g2p5-a1/.style=gvsel,
                  g2p5-a2/.style=gvsel, g2p5-a3/.style=gvsel, g2p5-a4/.style=gvsel,
                  g2p5-f03/.style=gpath, g2p5-f13/.style=gpath, g2p5-f36/.style=gpath,
                  g2p5-f04/.style=gpath, g2p5-f24/.style=gpath, g2p5-f46/.style=gpath]
        \pic {g2p5-S1};
        \node[font=\scriptsize\bfseries] at (1.6,2.75) {the ancestors of 6};
    \end{scope}
    \begin{scope}[xshift=11.2cm, x=0.95cm, y=0.7cm]
        \foreach \v/\a [count=\i from 0] in {0/{---}, 1/{---}, 2/{---}, 3/{0 1}, 4/{0 2}, 5/{0 1 3}, 6/{0 1 2 3 4}, 7/{0 1 2 3}}
            \node[font=\scriptsize, anchor=west] at (0,2.4-0.59*\i) {\textbf{\v:}\ \a};
        \node[font=\scriptsize\bfseries] at (1.1,3.05) {answer};
    \end{scope}
\end{tikzpicture}
\figcap{Sample 1. Vertex 6 is reached from 3 and from 4, hence from all of 0, 1, 2, 3, 4 (purple, orange paths). Vertices 0, 1, 2 have in-degree 0 and no ancestors.}
\end{diagrambox}

\subsubsection*{2. The graph model}

\begin{modelbox}[Graph Model --- All Ancestors]
\renewcommand{\arraystretch}{1.35}
\begin{tabularx}{\linewidth}{@{}>{\bfseries}l X@{}}
Vertices & the $n$ given vertices \\
Edges & directed, as given; the graph is a DAG (we will not need that) \\
Weights & none \\
Question & for every $v$: the set $\{u \ne v : u \leadsto v\}$ $=$ the vertices reachable from $v$ in the \textbf{reversed} graph $G^R$, sorted \\
Size & $n \le 1000$, $m \le 2000$; the output alone can have $\binom{n}{2} \approx 5 \cdot 10^5$ numbers, so \bigO{n(n+m)} $\approx 3 \cdot 10^6$ is the natural target \\
\end{tabularx}
\end{modelbox}

\begin{diagrambox}
\begin{tikzpicture}[x=1.15cm, y=1.1cm]
    \begin{scope}
        \pic {g2p5-S1};
        \node[font=\scriptsize\bfseries] at (1.6,2.75) {$G$: ``who reaches $v$?''};
    \end{scope}
    \draw[flow] (4.0,1) -- node[lbl, above]{turn every edge} (5.6,1);
    \node[note, text width=2.2cm, align=center] at (4.8,0.25) {\texttt{radj[to]}\\\texttt{.push\_back(from)}};
    \begin{scope}[xshift=7.0cm, g2p5-allrev]
        \pic {g2p5-S1};
        \node[font=\scriptsize\bfseries] at (1.6,2.75) {$G^R$: ``whom does $v$ reach?''};
    \end{scope}
\end{tikzpicture}
\figcap{The model: ancestors in $G$ are descendants in $G^R$. A plain reachability search from $v$ in $G^R$ is all we need.}
\end{diagrambox}

\subsubsection*{3. Key observations}

\begin{itemize}[leftmargin=*]
    \item \textbf{Ancestors of $v$ $=$ vertices reachable from $v$ in $G^R$} (the lemma of the pattern). One search per vertex: $n$ searches of \bigO{n + m} each.
    \item \textbf{The output is quadratic anyway.} A chain $0 \to 1 \to \dots \to 999$ has $\sum_v v = 499\,500$ ancestor entries; nothing can be faster than writing them. So \bigO{n(n+m)} is not ``slow'', it is the size of the answer plus a factor $\le 3$.
    \item \textbf{Sorted for free.} After the search from $v$, scan $u = 0, 1, \dots, n-1$ and collect the marked ones: the list is sorted without \texttt{sort}. Collecting in \emph{visiting} order would not be: from 6 the search visits $3, 4, 0, 2, 1$.
    \item \textbf{One marker array, $n$ searches.} Instead of clearing \texttt{seen} before every search, store \texttt{mark[u] = v} (``seen in the search from $v$''): no \bigO{n} reset, no reallocation.
    \item \textbf{Acyclicity is not needed} for the search approaches. It only matters for the topological-order alternative.
\end{itemize}

\subsubsection*{4. Approaches}

\begin{bruteforcebox}
\textbf{Idea: ask every pair.} For every $v$ and every $u \ne v$, run a fresh search from $u$ and check whether it reaches $v$.

\begin{lstlisting}[style=cpp]
#include <bits/stdc++.h>
using namespace std;

class Solution {
    vector<vector<int>> adj;
    bool reach(int s, int t) {                    // a fresh iterative DFS from s
        vector<char> seen(adj.size(), 0);
        vector<int> st(1, s);
        seen[s] = 1;
        while (!st.empty()) {
            int u = st.back(); st.pop_back();
            if (u == t) return true;
            for (int v : adj[u]) if (!seen[v]) { seen[v] = 1; st.push_back(v); }
        }
        return false;
    }
public:
    vector<vector<int>> getAncestors(int n, vector<vector<int>>& edges) {
        adj.assign(n, vector<int>());
        for (auto& e : edges) adj[e[0]].push_back(e[1]);
        vector<vector<int>> ans(n);
        for (int v = 0; v < n; v++)
            for (int u = 0; u < n; u++)                 // increasing u: already sorted
                if (u != v && reach(u, v)) ans[v].push_back(u);
        return ans;
    }
};
\end{lstlisting}

\begin{complexitybox}
\textbf{Time:} $n^2$ searches of \bigO{n+m}: \bigO{n^2(n+m)}. On a chain of $1000$ vertices plus $1001$ random forward edges ($m = 2000$) it took \textbf{5.0\,s} on our machine. \textbf{Space:} \bigO{n+m}. \textbf{Verdict:} \tle; kept as the reference for testing.
\end{complexitybox}
\end{bruteforcebox}

\begin{algorithmbox}[Approach 2: one reverse search per vertex]
\textbf{Idea.} Build \texttt{radj}. For every $v$, search from $v$ in $G^R$, marking with \texttt{mark[u] = v}; then scan $u$ in increasing order and collect the marked $u \ne v$.

\begin{lstlisting}[style=cpp]
#include <bits/stdc++.h>
using namespace std;

class Solution {
public:
    vector<vector<int>> getAncestors(int n, vector<vector<int>>& edges) {
        vector<vector<int>> radj(n);
        for (auto& e : edges) radj[e[1]].push_back(e[0]);   // from -> to, kept at "to"
        vector<vector<int>> ans(n);
        vector<int> mark(n, -1), st;                      // mark[u] = v: seen from v
        for (int v = 0; v < n; v++) {
            mark[v] = v;
            st.assign(1, v);
            while (!st.empty()) {                               // iterative DFS in G^R
                int x = st.back(); st.pop_back();
                for (int u : radj[x])
                    if (mark[u] != v) { mark[u] = v; st.push_back(u); }
            }
            for (int u = 0; u < n; u++)                   // index order = sorted
                if (u != v && mark[u] == v) ans[v].push_back(u);
        }
        return ans;
    }
};
\end{lstlisting}

\begin{complexitybox}
\textbf{Time:} $n$ searches of \bigO{n+m} plus $n$ scans of \bigO{n}: \bigO{n(n+m)} $\approx 3 \cdot 10^6$. The same chain input: \textbf{13\,ms}. \textbf{Space:} \bigO{n+m} besides the output. \textbf{Verdict:} fast enough by a factor of several hundred. The search is iterative, so even a chain of $10^5$ vertices would not overflow the stack.
\end{complexitybox}
\end{algorithmbox}

\begin{invariantbox}[Why it is correct]
\textbf{Claim.} After the search from $v$, $\texttt{mark}[u] = v$ iff $u \leadsto v$ in $G$.

\emph{Only ancestors are marked.} $v$ is marked first. A vertex $u$ is marked only from a popped $x$ with $u \in \texttt{radj}[x]$, i.e.\ an edge $u \to x$; by induction $x \leadsto v$, so $u \to x \leadsto v$.

\emph{Every ancestor is marked.} Suppose $u \leadsto v$ along $u = y_0 \to y_1 \to \dots \to y_k = v$ but $u$ is unmarked. Take the largest $i$ with $y_i$ unmarked; $y_{i+1}$ is marked, was pushed and later popped, and then it scanned $\texttt{radj}[y_{i+1}] \ni y_i$ and would have marked it. Contradiction.

\emph{Old marks do not interfere.} Marks from earlier searches hold values $< v$, so ``\texttt{mark[u] != v}'' is exactly ``not yet seen in this search''. The final scan in index order outputs the set sorted. \hfill$\blacksquare$
\end{invariantbox}

\begin{alternativebox}[Approach 3: forward searches that hand out their name]
\textbf{Idea.} Turn it around once more: run a forward search from every $u = 0, 1, \dots, n-1$ \emph{in this order}, and append $u$ to the list of every vertex it reaches. Since $u$ increases, every list receives its entries in increasing order --- sorted for free again, no reversed graph needed. Recursion depth is at most $n \le 1000$, so the recursive form is safe here.

\begin{lstlisting}[style=cpp]
#include <bits/stdc++.h>
using namespace std;

class Solution {
    vector<vector<int>> adj, ans;
    vector<int> mark;                                  // mark[x] = u: x reached from u
    void dfs(int x, int u) {
        for (int v : adj[x])
            if (mark[v] != u) { mark[v] = u; ans[v].push_back(u); dfs(v, u); }
    }
public:
    vector<vector<int>> getAncestors(int n, vector<vector<int>>& edges) {
        adj.assign(n, vector<int>());
        ans.assign(n, vector<int>());
        mark.assign(n, -1);
        for (auto& e : edges) adj[e[0]].push_back(e[1]);
        for (int u = 0; u < n; u++) { mark[u] = u; dfs(u, u); }   // u increasing
        return ans;
    }
};
\end{lstlisting}

\begin{center}
\begin{tikzpicture}[x=1.15cm, y=1.1cm]
    \begin{scope}[g2p5-a0/.style=gvsrc, g2p5-a3/.style=gvdone, g2p5-a5/.style=gvdone, g2p5-a6/.style=gvdone,
                  g2p5-a7/.style=gvdone, g2p5-a4/.style=gvdone,
                  g2p5-f03/.style={gtree, -{Stealth[length=2.4mm]}}, g2p5-f35/.style={gtree, -{Stealth[length=2.4mm]}},
                  g2p5-f36/.style={gtree, -{Stealth[length=2.4mm]}}, g2p5-f37/.style={gtree, -{Stealth[length=2.4mm]}},
                  g2p5-f04/.style={gtree, -{Stealth[length=2.4mm]}}, g2p5-f46/.style=gfaded]
        \pic {g2p5-S1};
        \node[font=\scriptsize\bfseries] at (1.6,2.75) {forward search from $u = 0$};
    \end{scope}
    \node[note, text width=7.2cm, align=left, anchor=west] at (4.3,1.0) {visits $3, 5, 6, 7, 4$ and appends 0 to each of their lists; the edge $4 \to 6$ (grey) finds 6 already marked with 0, so 0 is \emph{not} appended twice.\\[3pt]
    Visiting orders of the next searches: $u=1$: $3,5,6,7$; $u=2$: $4,6,7$; $u=3$: $5,6,7$; $u=4$: $6$; $u = 5, 6, 7$: nothing.\\[3pt]
    List of 6 grows as $[0] \to [0,1] \to [0,1,2] \to [0,1,2,3] \to [0,1,2,3,4]$.};
\end{tikzpicture}
\end{center}

\begin{complexitybox}
\textbf{Time:} \bigO{n(n+m)}, 15\,ms on the chain input. \textbf{Space:} \bigO{n+m} plus recursion depth $\le n$. \textbf{Verdict:} as good as Approach~2. Which one to write is a matter of taste: Approach~2 needs \texttt{radj} and a scan, Approach~3 needs the increasing order of $u$.
\end{complexitybox}
\end{alternativebox}

\begin{tricksbox}[Approach 4: Topological order and bitsets]
In a DAG, ancestors can be \emph{merged} instead of searched: if $u \to v$, then every ancestor of $u$ and $u$ itself are ancestors of $v$. Processing the vertices in topological order (Kahn's algorithm, Chapter~5) guarantees that $\mathrm{anc}(u)$ is complete before it is merged into its successors. With \texttt{bitset<1000>} a merge is $1000/64 \approx 16$ word operations.

\begin{lstlisting}[style=cpp]
#include <bits/stdc++.h>
using namespace std;

class Solution {
public:
    vector<vector<int>> getAncestors(int n, vector<vector<int>>& edges) {
        vector<vector<int>> adj(n);
        vector<int> indeg(n, 0), order;
        for (auto& e : edges) { adj[e[0]].push_back(e[1]); indeg[e[1]]++; }
        for (int v = 0; v < n; v++) if (!indeg[v]) order.push_back(v);
        for (size_t i = 0; i < order.size(); i++)         // Kahn: topological order
            for (int v : adj[order[i]]) if (--indeg[v] == 0) order.push_back(v);
        vector<bitset<1000>> A(n);                        // A[v][u]: u is an ancestor
        for (int u : order)
            for (int v : adj[u]) { A[v] |= A[u]; A[v][u] = 1; }   // anc(u) + u into v
        vector<vector<int>> ans(n);
        for (int v = 0; v < n; v++)
            for (int u = 0; u < n; u++) if (A[v][u]) ans[v].push_back(u);
        return ans;
    }
};
\end{lstlisting}
\textbf{Time:} \bigO{n + m \cdot n/64 + n^2} --- the final $n^2$ scan dominates; 7\,ms on the chain input. \textbf{Space:} $n^2$ bits $= 125$\,KB. Unlike the searches it needs the graph to be acyclic, and a bitset size fixed at compile time; for $n = 10^5$ the $n^2$ bits (1.25\,GB) would not fit.
\end{tricksbox}

\subsubsection*{5. Trace}

\begin{tracebox}[{Sample 1, Approach 2: the searches from $v = 6$ and $v = 7$}]
\texttt{radj} in input order: $\texttt{radj}[3] = 0, 1$;\ $\texttt{radj}[4] = 0, 2$;\ $\texttt{radj}[5] = 3$;\ $\texttt{radj}[6] = 3, 4$;\ $\texttt{radj}[7] = 2, 3$; the lists of $0, 1, 2$ are empty.
\begin{center}
\begin{tikzpicture}[x=0.82cm, y=0.82cm]
    \tikzset{gv/.append style={minimum size=6.2mm, font=\scriptsize\bfseries}}
    \begin{scope}[g2p5-allrev, g2p5-a6/.style=gvcur, g2p5-a3/.style=gvseen, g2p5-a4/.style=gvseen,
                  g2p5-f36/.style=g2p5-rt, g2p5-f46/.style=g2p5-rt]
        \pic {g2p5-S1};
        \node[font=\scriptsize\bfseries] at (1.6,2.8) {1. pop 6: mark 3, 4};
    \end{scope}
    \begin{scope}[xshift=4.0cm, g2p5-allrev, g2p5-a6/.style=gvdst, g2p5-a3/.style=gvseen, g2p5-a4/.style=gvcur,
                  g2p5-a0/.style=gvseen, g2p5-a2/.style=gvseen,
                  g2p5-f36/.style=g2p5-rt, g2p5-f46/.style=g2p5-rt, g2p5-f04/.style=g2p5-rt, g2p5-f24/.style=g2p5-rt]
        \pic {g2p5-S1};
        \node[font=\scriptsize\bfseries] at (1.6,2.8) {2. pop 4: mark 0, 2};
    \end{scope}
    \begin{scope}[xshift=8.0cm, g2p5-allrev, g2p5-a6/.style=gvdst, g2p5-a3/.style=gvcur, g2p5-a4/.style=gvdone,
                  g2p5-a0/.style=gvdone, g2p5-a2/.style=gvdone, g2p5-a1/.style=gvseen,
                  g2p5-f36/.style=g2p5-rt, g2p5-f46/.style=g2p5-rt, g2p5-f04/.style=g2p5-rt, g2p5-f24/.style=g2p5-rt,
                  g2p5-f13/.style=g2p5-rt]
        \pic {g2p5-S1};
        \node[font=\scriptsize\bfseries] at (1.6,2.8) {3. pops 2, 0; pop 3: mark 1};
    \end{scope}
    \begin{scope}[xshift=12.0cm, g2p5-allrev, g2p5-a7/.style=gvdst, g2p5-a3/.style=gvdone, g2p5-a2/.style=gvdone,
                  g2p5-a0/.style=gvdone, g2p5-a1/.style=gvdone,
                  g2p5-f37/.style=g2p5-rt, g2p5-f27/.style=g2p5-rt, g2p5-f03/.style=g2p5-rt, g2p5-f13/.style=g2p5-rt]
        \pic {g2p5-S1};
        \node[font=\scriptsize\bfseries] at (1.6,2.8) {4. the search from 7};
    \end{scope}
\end{tikzpicture}
\end{center}
\begin{center}\small
\begin{tabular}{@{}cllll@{}}
\toprule
\textbf{search} & \textbf{pop} & \textbf{newly marked} & \textbf{stack after} & \textbf{note} \\
\midrule
$v = 6$ & 6 & 3, 4 & \texttt{3 4} & \\
        & 4 & 0, 2 & \texttt{3 0 2} & \\
        & 2 & --- & \texttt{3 0} & in-degree 0 \\
        & 0 & --- & \texttt{3} & \\
        & 3 & 1 & \texttt{1} & 0 is already marked \\
        & 1 & --- & empty & scan: \textbf{0 1 2 3 4} \\
\midrule
$v = 7$ & 7 & 2, 3 & \texttt{2 3} & \\
        & 3 & 0, 1 & \texttt{2 0 1} & \\
        & 1, 0, 2 & --- & empty & scan: \textbf{0 1 2 3} \\
\midrule
$v = 5$ & 5; 3; 1, 0 & 3; 0, 1; --- & & scan: \textbf{0 1 3} \\
\bottomrule
\end{tabular}
\end{center}
The visiting order from 6 is $3, 4, 0, 2, 1$; the index-order scan turns it into the sorted list. Over all eight searches $0+0+0+2+2+3+5+4 = 16$ ancestors are written --- exactly the size of the official output.
\end{tracebox}

\subsubsection*{6. Test suite}

\begin{testbox}
\centering\small
\begin{tabular}{@{}p{5.2cm}p{4.2cm}p{4.4cm}@{}}
\toprule
\textbf{Input} & \textbf{Our output} & \textbf{Why} \\
\midrule
$n = 1$, no edges & \texttt{[[]]} & smallest input \\
$n = 3$, no edges & \texttt{[[],[],[]]} & every list empty \\
chain $0 \to 1 \to 2 \to 3$ & \texttt{[[],[0],[0,1],[0,1,2]]} & long reverse paths \\
chain $3 \to 2 \to 1 \to 0$ & \texttt{[[1,2,3],[2,3],[3],[]]} & visiting order $\ne$ sorted order \\
diamond $0{\to}1, 0{\to}2, 1{\to}3, 2{\to}3$ & \texttt{[[],[0],[0],[0,1,2]]} & 0 reaches 3 twice, listed once \\
$n = 6$: $0{\to}1$, $2{\to}3$, $4{\to}5$ & \texttt{[[],[0],[],[2],[],[4]]} & disconnected \\
chain of 1000 $+$ 1001 random edges & 499\,500 entries, 13\,ms & the largest output \\
\bottomrule
\end{tabular}
\end{testbox}
All outputs come from running Approaches 2 and 3 (they agree). In addition, Approaches 1--4 were run on 3000 random DAGs with $n \le 9$ (a random vertex order, each forward pair an edge with probability $1/3$, edges shuffled) and produced identical lists every time.

\subsubsection*{7. Common mistakes}

\begin{warningbox}[Mistake 1 --- forward search without per-source marks]
Appending $u$ on every \emph{path} instead of every \emph{vertex} gives duplicates: without \texttt{mark} the forward version prints \texttt{[0,0,1,2,3,4]} for vertex 6 of Sample~1 (0 reaches 6 via 3 and via 4). Worse, the running time becomes the number of paths: a ladder of $k$ layers of two vertices, each layer joined completely to the next, has $2^{k-1}$ paths from a top vertex to a bottom vertex, and $m \le 2000$ allows $k = 500$.
\end{warningbox}

\begin{warningbox}[Mistake 2 --- collecting in visiting order]
Pushing $u$ into the answer when it is \emph{marked} produces \texttt{[3,4,0,2,1]} for vertex 6. Either scan the marks in index order (Approach~2), or let the outer loop run in increasing order (Approach~3), or sort.
\end{warningbox}

\begin{warningbox}[Mistake 3 --- a fresh \texttt{vector<bool> seen(n)} per search, inside a recursive helper that copies it]
A new array per search is fine (\bigO{n} each, \bigO{n^2} total). Passing it \emph{by value} into a recursive DFS copies it at every call: \bigO{n} per visited vertex, up to \bigO{n^3} $= 10^9$ in total, and \tle. Pass by reference, or use the \texttt{mark[u] = v} trick.
\end{warningbox}

\begin{warningbox}[Mistake 4 --- including $v$ itself]
The search from $v$ marks $v$. ``Ancestor'' requires $u \ne v$: skip it in the scan.
\end{warningbox}

\subsubsection*{8. Takeaways}

\begin{itemize}[leftmargin=*]
    \item ``Everything that reaches $v$'' is ``everything reachable from $v$ in $G^R$''.
    \item When the output is quadratic, $n$ linear searches are optimal up to a constant.
    \item Reuse one marker array across searches with \texttt{mark[u] = v}.
    \item Sorted output comes for free from index-order scans or an increasing outer loop.
\end{itemize}

\begin{variantsbox}
\variant{Only the counts}{\varnone}
  {return, for every $v$, the number of ancestors instead of the list.}
  {the same reachability sets; only their sizes are needed.}
  {count in the same $n$ reverse searches. (Summing ancestor counts of a DAG by DP double-counts shared ancestors --- on the diamond it gives 4 for vertex 3 instead of 3.)}
\variant{Descendants instead}{\varnone}
  {for every $v$, the vertices $v$ can reach.}
  {no reversal: forward searches in $G$.}
  {Approach~2 with \texttt{adj} in place of \texttt{radj}; or Approach~3 with the roles of the two lists swapped.}
\variant{Ancestors of one vertex}{\varnone}
  {only one query vertex $q$.}
  {a single reverse search.}
  {template 1 of the pattern with $T = \{q\}$: \VE.}
\variant{Cycles allowed}{\varnone}
  {the graph is a general directed graph.}
  {``ancestor'' still means ``can reach''; vertices on a cycle through $v$ become ancestors of $v$.}
  {Approaches 1--3 are unchanged (no acyclicity was used); Approach~4 breaks, since Kahn's order misses the cycle vertices --- it would first need the strongly connected components (Chapter~19).}
\variant{Common ancestors}{\varnone}
  {for a pair $(a, b)$, the vertices that reach both.}
  {two target sets.}
  {two reverse searches, from $a$ and from $b$, and an AND of the marks (like E5.2).}
\end{variantsbox}

%============================================================
% CH2 · P5 · PROBLEM 2 — SOLVE THE MAZE (Codeforces 1365D)
%============================================================
\subsection{Problem 2.5.2: Solve The Maze}
\problemlabel{2.5.2}{ssec:ch2-solve-maze}

% cells with a text: green G reached / not reached, visit numbers
\tikzset{
    g2p5-Gn/.style={g2p5-e, fill=basefill, draw=successgreen, line width=1.1pt},
    g2p5-Bn/.style={g2p5-e, fill=dupfill, draw=recursionred, line width=1.1pt},
    g2p5-on/.style={g2p5-e, fill=computedfill}
}
% \gptworow{y}{style/text, ...}: one row of cells at height y
\newcommand{\gptworow}[2]{\foreach \s/\t [count=\c from 0] in {#2} \node[g2p5-\s] at (\c,-#1) {\t};}

\problemcard{2.5.2}{Solve The Maze}{Codeforces 1365D}{https://codeforces.com/problemset/problem/1365/D}{\exhard\ \cf{1700}}{Wall in every bad person, then one search from the exit must reach every good person and no bad one}

\subsubsection*{1. Problem statement}

A maze is an $n \times m$ grid of cells: empty (\texttt{.}), wall (\texttt{\#}), good person (\texttt{G}) or bad person (\texttt{B}). People move between side-adjacent non-wall cells; the only exit is the cell $(n, m)$, which is guaranteed to be empty. You may turn any set of \emph{empty} cells (the exit included) into walls; \texttt{G} and \texttt{B} cells cannot be blocked and can be walked through. Can you do it so that every good person can reach the exit and no bad person can? Print \texttt{Yes} or \texttt{No}.

\textbf{Constraints:} $1 \le t \le 100$ test cases, $1 \le n, m \le 50$.

\begin{table}[H]
\centering\small
\begin{tabular}{@{}clll@{}}
\toprule
\textbf{\#} & \textbf{Maze} & \textbf{Official / our output} & \textbf{Why} \\
\midrule
1 & \texttt{.} & \texttt{Yes} / \texttt{Yes} & nobody there \\
2 & \texttt{G.} & \texttt{Yes} / \texttt{Yes} & G is next to the exit \\
3 & \texttt{\#B / G.} & \texttt{No} / \texttt{No} & walling B in walls the exit, G is stuck \\
4 & \texttt{G.\# / B\#.} & \texttt{No} / \texttt{No} & G stands next to B \\
5 & \texttt{\#B. / \#.. / GG.} & \texttt{Yes} / \texttt{Yes} & wall the two cells next to B \\
6 & \texttt{\#B / B.} & \texttt{Yes} / \texttt{Yes} & no G: wall the exit itself \\
\bottomrule
\end{tabular}
\end{table}

\begin{diagrambox}
\begin{tikzpicture}[x=6.5mm, y=6.5mm]
    \begin{scope}
        \gptworow{0}{w/, B/}\gptworow{1}{G/, e/}
        \node[font=\scriptsize\bfseries] at (0.5,1.1) {sample 3};
        \node[note] at (0.5,-2.1) {No};
    \end{scope}
    \begin{scope}[xshift=2.3cm]
        \gptworow{0}{G/, e/, w/}\gptworow{1}{B/, w/, e/}
        \node[font=\scriptsize\bfseries] at (1,1.1) {sample 4};
        \node[note] at (1,-2.1) {No};
    \end{scope}
    \begin{scope}[xshift=5.2cm]
        \gptworow{0}{w/, B/, e/}\gptworow{1}{w/, e/, e/}\gptworow{2}{G/, G/, e/}
        \node[font=\scriptsize\bfseries] at (1,1.1) {sample 5};
        \node[note] at (1,-2.8) {Yes};
    \end{scope}
    \begin{scope}[xshift=8.1cm]
        \gptworow{0}{w/, B/, x/}\gptworow{1}{w/, x/, o/}\gptworow{2}{G/, G/, o/}
        \node[font=\scriptsize\bfseries] at (1,1.1) {5: our walls};
        \node[note] at (1,-2.8) {G, G reached};
    \end{scope}
    \begin{scope}[xshift=11.0cm]
        \gptworow{0}{w/, B/}\gptworow{1}{B/, x/}
        \node[font=\scriptsize\bfseries] at (0.5,1.1) {sample 6};
        \node[note] at (0.5,-2.1) {Yes (exit walled)};
    \end{scope}
\end{tikzpicture}
\figcap{Some samples (exit $=$ bottom-right cell; grey $=$ wall, red frame $=$ a wall we add, blue $=$ reached from the exit). In sample 5 the cells right of B and below B are walled; the search from the exit then reaches both G.}
\end{diagrambox}

\subsubsection*{2. The graph model}

\begin{modelbox}[Graph Model --- Solve The Maze]
\renewcommand{\arraystretch}{1.35}
\begin{tabularx}{\linewidth}{@{}>{\bfseries}l X@{}}
Vertices & the non-wall cells, up to $50 \cdot 50 = 2500$ per test \\
Edges & undirected, between side-adjacent non-wall cells; implicit (computed from the grid) \\
Weights & none \\
Question & choose a set $W$ of empty cells to delete so that in the remaining graph every G is in the component of the exit and no B is --- equivalently: one search \textbf{from the exit} marks all G and no B \\
Size & $2500$ vertices, $\le 5000$ edges per test, 100 tests: any linear algorithm is instant; trying sets $W$ is not ($2^{2500}$) \\
\end{tabularx}
\end{modelbox}

\begin{diagrambox}
\begin{tikzpicture}[x=9mm, y=9mm]
    \begin{scope}
        \gptworow{0}{w/, B/, e/}\gptworow{1}{w/, e/, e/}\gptworow{2}{G/, G/, e/}
        \node[font=\scriptsize\bfseries] at (1,0.95) {sample 5 as a grid};
    \end{scope}
    \begin{scope}[xshift=4.2cm]
        \foreach \r/\c/\l in {0/1/B, 0/2/{}, 1/1/{}, 1/2/{}, 2/0/G, 2/1/G, 2/2/{}}
            \node[gv, minimum size=6mm, font=\scriptsize\bfseries] (c\r\c) at (\c,-\r) {\l};
        \foreach \a/\b in {01/02, 01/11, 02/12, 11/12, 11/21, 12/22, 20/21, 21/22} \draw[ge] (c\a) -- (c\b);
        \node[gvdst, minimum size=6mm, font=\scriptsize\bfseries] at (2,-2) {ex};
        \node[font=\scriptsize\bfseries] at (1,0.95) {the cell graph};
    \end{scope}
    \node[note, text width=5.2cm, align=left, anchor=west] at (8.3,-1) {7 open cells, 8 edges. B is connected to the exit through $(1,3)$ and through $(2,2)$: both must be deleted --- they are exactly the empty neighbours of B. What remains of the exit's component is $\{(2,3), (3,3), (3,2), (3,1)\}$, containing both G.};
\end{tikzpicture}
\figcap{The model: cells are vertices, sides are edges, walls are missing vertices. Adding a wall deletes a vertex.}
\end{diagrambox}

\subsubsection*{3. Key observations}

\begin{itemize}[leftmargin=*]
    \item \textbf{Reverse the roles.} ``Every G can reach the exit'' is one question per G; ``the exit reaches every G'' is one search. The graph is undirected, so nothing has to be reversed except the starting point.
    \item \textbf{A G next to a B is fatal.} If G escapes, B can step onto G's cell (it cannot be walled) and follow G's route. Answer \texttt{No}.
    \item \textbf{A G route that passes next to a B is useless.} If a route from G to the exit visits a cell $c$ adjacent to B, then B can step to $c$ and follow the rest of the route out. So in \emph{every} valid solution, every escaping route avoids all neighbours of all B.
    \item \textbf{Therefore: wall all empty neighbours of every B.} It costs nothing (no valid route uses them) and it traps every B (its four sides are now walls, borders, or other B --- never G, see above). After that, no more walls are needed: just check with one search from the exit.
    \item \textbf{No G at all:} answer \texttt{Yes} --- wall the exit (allowed) or simply rely on the B being walled in. Do not run into the trap of answering \texttt{No} because the exit itself became a wall.
\end{itemize}

\begin{diagrambox}
\begin{tikzpicture}[x=6.5mm, y=6.5mm]
    \begin{scope}
        \gptworow{0}{G/, e/, B/, e/}
        \node[font=\scriptsize\bfseries] at (1.5,1.0) {\texttt{G.B.}: the exit is next to B};
        \gptworow{2}{G/, e/, B/, x/}
        \node[note, text width=3.6cm, anchor=north] at (1.5,-2.7) {walling B in walls the exit: \texttt{No}};
    \end{scope}
    \begin{scope}[xshift=4.6cm]
        \gptworow{0}{G/, e/, e/, w/}\gptworow{1}{e/, w/, e/, e/}\gptworow{2}{e/, B/, e/, e/}
        \node[font=\scriptsize\bfseries] at (1.5,1.0) {two routes of G};
        \draw[gpath, -{Stealth[length=2mm]}] (0.15,0.15) -- (2.15,0.15) -- (2.15,-0.85) -- (3.15,-0.85) -- (3.15,-1.8);
        \draw[gback] (-0.15,-0.15) -- (-0.15,-2.15) -- (2.9,-2.15);
        \node[note, text width=3.8cm, anchor=north] at (1.5,-2.7) {the bottom route (red) passes next to B, so B could follow it out: useless. The top route (orange) avoids B's neighbours};
    \end{scope}
    \begin{scope}[xshift=9.6cm]
        \gptworow{0}{G/, e/, e/, w/}\gptworow{1}{e/, w/, e/, e/}\gptworow{2}{x/, B/, x/, e/}
        \node[font=\scriptsize\bfseries] at (1.5,1.0) {wall in B};
        \node[note, text width=3.6cm, anchor=north] at (1.5,-2.7) {G still escapes along the top: \texttt{Yes}};
    \end{scope}
\end{tikzpicture}
\figcap{Why walling the neighbours of B is never a mistake: a cell next to B is useless for G, because B could use it too.}
\end{diagrambox}

\subsubsection*{4. Approaches}

\begin{bruteforcebox}
\textbf{Idea: try every set of walls.} For each subset of the empty cells, wall it, run one search from the exit (unless it is walled) and check the condition. Only for tiny grids ($\le 15$ empty cells), but a perfect reference: it uses no insight at all.

\begin{lstlisting}[style=cpp]
#include <bits/stdc++.h>
using namespace std;

int n, m;
const int dr[] = {1, -1, 0, 0}, dc[] = {0, 0, 1, -1};

bool ok(const vector<string>& h) {               // search from the exit, check G and B
    vector<vector<char>> seen(n, vector<char>(m, 0));
    if (h[n - 1][m - 1] != '#') {
        vector<pair<int, int>> st(1, make_pair(n - 1, m - 1));
        seen[n - 1][m - 1] = 1;
        while (!st.empty()) {
            int r = st.back().first, c = st.back().second; st.pop_back();
            for (int d = 0; d < 4; d++) {
                int a = r + dr[d], b = c + dc[d];
                if (a < 0 || b < 0 || a >= n || b >= m) continue;
                if (seen[a][b] || h[a][b] == '#') continue;
                seen[a][b] = 1; st.push_back(make_pair(a, b));
            }
        }
    }
    for (int r = 0; r < n; r++)
        for (int c = 0; c < m; c++) {
            if (h[r][c] == 'G' && !seen[r][c]) return false;
            if (h[r][c] == 'B' && seen[r][c]) return false;
        }
    return true;
}

int main() {
    int t;
    cin >> t;
    while (t--) {
        cin >> n >> m;
        vector<string> g(n);
        for (auto& s : g) cin >> s;
        vector<pair<int, int>> e;                // the empty cells
        for (int r = 0; r < n; r++)
            for (int c = 0; c < m; c++)
                if (g[r][c] == '.') e.push_back(make_pair(r, c));
        bool any = false;
        for (long long mask = 0; mask < (1LL << e.size()) && !any; mask++) {
            vector<string> h = g;
            for (size_t i = 0; i < e.size(); i++)
                if (mask >> i & 1) h[e[i].first][e[i].second] = '#';
            any = ok(h);
        }
        cout << (any ? "Yes" : "No") << "\n";
    }
}
\end{lstlisting}

\begin{complexitybox}
\textbf{Time:} $2^{E} \cdot \bigO{nm}$ for $E$ empty cells. Fine for the samples ($E \le 4$); a $50 \times 50$ maze can have $E = 2500$. \textbf{Space:} \bigO{nm}. \textbf{Verdict:} \tle{} beyond $E \approx 20$; our reference for random tests.
\end{complexitybox}
\end{bruteforcebox}

\begin{algorithmbox}[Approach 2: wall the bad, one search from the exit]
\textbf{Idea.} (1) For every B, look at its four neighbours: a G means \texttt{No}; an empty cell becomes a wall. (2) No G: \texttt{Yes}. (3) The exit is now a wall: \texttt{No}. (4) Search from the exit; \texttt{Yes} iff every G and no B is reached.

\begin{lstlisting}[style=cpp]
#include <bits/stdc++.h>
using namespace std;

// Solve The Maze (CF 1365D): wall in every B, then one DFS from the exit
int n, m;
vector<string> g;
const int dr[] = {1, -1, 0, 0}, dc[] = {0, 0, 1, -1};

bool inside(int r, int c) { return r >= 0 && c >= 0 && r < n && c < m; }

bool solve() {
    int goods = 0;
    for (int r = 0; r < n; r++)
        for (int c = 0; c < m; c++) {
            if (g[r][c] == 'G') goods++;
            if (g[r][c] != 'B') continue;
            for (int d = 0; d < 4; d++) {
                int nr = r + dr[d], nc = c + dc[d];
                if (!inside(nr, nc)) continue;
                if (g[nr][nc] == 'G') return false;        // B would follow G out
                if (g[nr][nc] == '.') g[nr][nc] = '#';     // block B in
            }
        }
    if (goods == 0) return true;                           // nobody must escape
    if (g[n - 1][m - 1] == '#') return false;              // the exit itself was walled
    vector<vector<char>> seen(n, vector<char>(m, 0));
    vector<pair<int, int>> st(1, make_pair(n - 1, m - 1));
    seen[n - 1][m - 1] = 1;
    int reachedG = 0, reachedB = 0;
    while (!st.empty()) {                                  // search from the exit
        int r = st.back().first, c = st.back().second; st.pop_back();
        if (g[r][c] == 'G') reachedG++;
        if (g[r][c] == 'B') reachedB++;
        for (int d = 0; d < 4; d++) {
            int nr = r + dr[d], nc = c + dc[d];
            if (!inside(nr, nc) || seen[nr][nc] || g[nr][nc] == '#') continue;
            seen[nr][nc] = 1;
            st.push_back(make_pair(nr, nc));
        }
    }
    return reachedG == goods && reachedB == 0;
}

int main() {
    int t;
    cin >> t;
    while (t--) {
        cin >> n >> m;
        g.assign(n, "");
        for (auto& s : g) cin >> s;
        cout << (solve() ? "Yes" : "No") << "\n";
    }
}
\end{lstlisting}

\begin{complexitybox}
\textbf{Time:} \bigO{nm} per test: the B scan looks at $4nm$ neighbours, the search visits each cell once. 100 tests of $50 \times 50$ cells (all G) ran in 0.2\,s including I/O. \textbf{Space:} \bigO{nm}. \textbf{Verdict:} fast; the stack holds at most $nm = 2500$ cells, and a recursive DFS (depth $\le 2500$) would be safe too.
\end{complexitybox}
\end{algorithmbox}

\begin{invariantbox}[Why it is correct]
Let $W_0$ be the set of empty neighbours of all B, and suppose no G is adjacent to a B.

\emph{If any valid wall set $W$ exists, then $W_0$ is valid.} With $W_0$ every B is trapped: its neighbours are walls, borders or B cells, so its component contains only B cells and never the (empty) exit. Now take a good person $g$. In the solution $W$ it has a route $P$ to the exit. No cell of $P$ is adjacent to a B --- otherwise that B steps onto $P$ and follows it out, contradicting the validity of $W$. So no cell of $P$ is in $W_0$, and $P$ is still open after walling $W_0$: $g$ escapes. Hence $W_0$ is valid.

\emph{If $W_0$ is not valid, nothing is.} By the contrapositive of the above. And if some G is adjacent to a B, no $W$ is valid (B follows G). If there is no G, $W_0$ is valid (every B is trapped) --- even when it walls the exit.

\emph{The search decides $W_0$.} With $W_0$ in place, the component of the exit is exactly the set marked by the search (Chapter~1), so ``every G marked, no B marked'' is the definition of valid. \hfill$\blacksquare$
\end{invariantbox}

\begin{alternativebox}[Approach 3: components with a DSU]
After walling $W_0$, the question is only ``which component is each person in?'', so any connectivity method works. A DSU over the open cells unites every pair of adjacent open cells; then every G must share the exit's root and no B may.

\begin{lstlisting}[style=cpp]
#include <bits/stdc++.h>
using namespace std;

int n, m;
vector<string> g;
vector<int> par;
int fnd(int x) { while (par[x] != x) x = par[x] = par[par[x]]; return x; }

bool solve() {
    const int dr[] = {1, -1, 0, 0}, dc[] = {0, 0, 1, -1};
    int goods = 0;
    for (int r = 0; r < n; r++)
        for (int c = 0; c < m; c++) {
            if (g[r][c] == 'G') goods++;
            if (g[r][c] != 'B') continue;
            for (int d = 0; d < 4; d++) {
                int a = r + dr[d], b = c + dc[d];
                if (a < 0 || b < 0 || a >= n || b >= m) continue;
                if (g[a][b] == 'G') return false;
                if (g[a][b] == '.') g[a][b] = '#';
            }
        }
    if (goods == 0) return true;
    if (g[n - 1][m - 1] == '#') return false;
    par.resize(n * m);
    iota(par.begin(), par.end(), 0);
    for (int r = 0; r < n; r++)                            // unite right and down
        for (int c = 0; c < m; c++) {
            if (g[r][c] == '#') continue;
            int x = r * m + c;
            if (r + 1 < n && g[r + 1][c] != '#') par[fnd(x)] = fnd(x + m);
            if (c + 1 < m && g[r][c + 1] != '#') par[fnd(x)] = fnd(x + 1);
        }
    int ex = fnd(n * m - 1);
    for (int r = 0; r < n; r++)
        for (int c = 0; c < m; c++) {
            if (g[r][c] == 'G' && fnd(r * m + c) != ex) return false;
            if (g[r][c] == 'B' && fnd(r * m + c) == ex) return false;
        }
    return true;
}

int main() {
    int t;
    cin >> t;
    while (t--) {
        cin >> n >> m;
        g.assign(n, "");
        for (auto& s : g) cin >> s;
        cout << (solve() ? "Yes" : "No") << "\n";
    }
}
\end{lstlisting}

\begin{center}
\begin{tikzpicture}[x=6.5mm, y=6.5mm]
    \gptworow{0}{Gn/, on/, x/, on/, on/}
    \gptworow{1}{on/, w/, Bn/, x/, on/}
    \gptworow{2}{on/, on/, w/, on/, Gn/}
    \gptworow{3}{on/, on/, on/, on/, on/}
    \draw[gtree] (0,-3) -- (4,-3); \draw[gtree] (0,0) -- (0,-3); \draw[gtree] (4,0) -- (4,-3);
    \draw[gtree] (0,0) -- (1,0); \draw[gtree] (3,0) -- (4,0); \draw[gtree] (0,-2) -- (1,-2);
    \draw[gtree] (1,-2) -- (1,-3); \draw[gtree] (3,-2) -- (4,-2); \draw[gtree] (3,-2) -- (3,-3);
    \node[note, text width=6.4cm, align=left, anchor=west] at (5.4,-1.5) {The trace maze of Section~5 after walling: the DSU joins all 15 open cells (blue lines: the adjacencies it unites) into the exit's set, both G included. B is alone in its own set.};
\end{tikzpicture}
\end{center}

\begin{complexitybox}
\textbf{Time:} \bigO{nm\,\alpha(nm)} per test. \textbf{Space:} \bigO{nm}. \textbf{Verdict:} same answers as Approach~2 on all tests below; more code for no gain here, but the DSU form is the one to use when walls are \emph{removed} one by one (Chapter~6).
\end{complexitybox}
\end{alternativebox}

\begin{tricksbox}[Approach 4: Implementation Refinements]
\begin{itemize}[leftmargin=*]
    \item \textbf{Check G before walling the exit.} The order ``no G $\Rightarrow$ \texttt{Yes}'' before ``exit walled $\Rightarrow$ \texttt{No}'' matters: sample 6 has a walled exit and the answer \texttt{Yes}.
    \item \textbf{Counting instead of a second scan.} Count G while walling; the search counts the G it reaches; compare the two numbers.
    \item \textbf{Modifying the grid in place} is fine (each test reads a new grid). If the grid must survive, copy it first.
    \item \textbf{The ``no B reached'' test is a safety net.} After the walling a B can only be connected to other B, never to the exit, so \texttt{reachedB} is always 0 in Approach~2 --- but the check costs nothing and protects against a bug in the walling loop.
\end{itemize}
\end{tricksbox}

\subsubsection*{5. Trace}

\begin{tracebox}[{A $4 \times 5$ maze: \texttt{G.... / .\#B.. / ..\#.G / .....}}]
Step 1 walls the empty neighbours of B $= (2,3)$: the cells $(1,3)$ and $(2,4)$ (the cells $(2,2)$ and $(3,3)$ are walls already). Two G, no G next to B, the exit is open. Step 2 is the search from the exit, neighbours tried in the order down, up, right, left. The numbers are the order in which cells are popped:
\begin{center}
\begin{tikzpicture}[x=6.5mm, y=6.5mm]
    \begin{scope}
        \gptworow{0}{G/, e/, e/, e/, e/}\gptworow{1}{e/, w/, B/, e/, e/}\gptworow{2}{e/, e/, w/, e/, G/}\gptworow{3}{e/, e/, e/, e/, e/}
        \node[font=\scriptsize\bfseries] at (2,1.0) {input};
    \end{scope}
    \begin{scope}[xshift=4.2cm]
        \gptworow{0}{G/, e/, x/, e/, e/}\gptworow{1}{e/, w/, B/, x/, e/}\gptworow{2}{e/, e/, w/, e/, G/}\gptworow{3}{e/, e/, e/, e/, e/}
        \node[font=\scriptsize\bfseries] at (2,1.0) {1. wall in B};
    \end{scope}
    \begin{scope}[xshift=8.4cm]
        \gptworow{0}{Gn/8, on/9, x/, on/15, on/14}
        \gptworow{1}{on/7, w/, Bn/B, x/, on/13}
        \gptworow{2}{on/6, on/10, w/, on/11, Gn/12}
        \gptworow{3}{on/5, on/4, on/3, on/2, on/1}
        \node[font=\scriptsize\bfseries] at (2,1.0) {2. search from the exit};
    \end{scope}
\end{tikzpicture}
\end{center}
\begin{center}\small
\begin{tabular}{@{}cll@{}}
\toprule
\textbf{pop} & \textbf{cell} & \textbf{stack after (top on the right)} \\
\midrule
1 & $(4,5)$ exit & $(3,5)\ (4,4)$ \\
2 & $(4,4)$ & $(3,5)\ (3,4)\ (4,3)$ \\
3 & $(4,3)$ & $(3,5)\ (3,4)\ (4,2)$\quad ($(3,3)$ is a wall) \\
4--5 & $(4,2)$, $(4,1)$ & $(3,5)\ (3,4)\ (3,2)\ (3,1)$ \\
6--7 & $(3,1)$, $(2,1)$ & $(3,5)\ (3,4)\ (3,2)\ (1,1)$ \\
8 & $(1,1)$ \textbf{G} & $(3,5)\ (3,4)\ (3,2)\ (1,2)$ \\
9--11 & $(1,2)$, $(3,2)$, $(3,4)$ & $(3,5)$\quad ($(1,3)$ is our wall) \\
12 & $(3,5)$ \textbf{G} & $(2,5)$ \\
13--15 & $(2,5)$, $(1,5)$, $(1,4)$ & empty \\
\bottomrule
\end{tabular}
\end{center}
Both G are reached (pops 8 and 12), B is not: \texttt{Yes}. With one more B in the corner, \texttt{G.... / .\#B.. / ..\#.G / B....}, the cells $(3,1)$ and $(4,2)$ are walled as well, the G in $(1,1)$ is cut off and the program prints \texttt{No}:
\begin{center}
\begin{tikzpicture}[x=6.5mm, y=6.5mm]
    \gptworow{0}{G/, e/, x/, on/8, on/7}
    \gptworow{1}{e/, w/, Bn/B, x/, on/6}
    \gptworow{2}{x/, e/, w/, on/4, Gn/5}
    \gptworow{3}{Bn/B, x/, on/3, on/2, on/1}
    \node[note, text width=6.8cm, align=left, anchor=west] at (5.4,-1.5) {The search from the exit pops only 8 cells. The region $\{(1,1), (1,2), (2,1)\}$ with the G is closed by our walls $(1,3)$ and $(3,1)$ --- and both are neighbours of a B, so no solution could keep them open. The cell $(3,2)$ is isolated.};
\end{tikzpicture}
\end{center}
\end{tracebox}

\subsubsection*{6. Test suite}

\begin{testbox}
\centering\small
\begin{tabular}{@{}p{4.6cm}lp{6.6cm}@{}}
\toprule
\textbf{Maze} & \textbf{Ours} & \textbf{Why} \\
\midrule
\texttt{.} & \texttt{Yes} & $1 \times 1$, nobody \\
\texttt{GB.} & \texttt{No} & G next to B \\
\texttt{B. / ..} & \texttt{Yes} & no G; B walled in \\
\texttt{G.B / ...} & \texttt{No} & B next to the exit: the exit is walled, G exists \\
\texttt{G\#. / .\#.} & \texttt{No} & G separated by walls, no B at all \\
\texttt{G.B.} & \texttt{No} & B stands between G and the exit \\
the $4 \times 5$ trace maze & \texttt{Yes} & G escapes around B \\
same with \texttt{B} at $(4,1)$ & \texttt{No} & the second B cuts G off \\
100 tests $50 \times 50$, all G (half with one B) & 50 \texttt{Yes}, 50 \texttt{No} & 0.2\,s \\
\bottomrule
\end{tabular}
\end{testbox}
All outputs come from running Approach~2. Approaches 2 and 3 were also compared with the brute force on $40\,000$ random grids of size up to $4 \times 4$ (cells \texttt{.}/\texttt{\#}/\texttt{G}/\texttt{B} with probabilities $0.5/0.2/0.2/0.1$, the exit forced empty): no difference.

\subsubsection*{7. Common mistakes}

\begin{warningbox}[Mistake 1 --- answering \texttt{No} when the exit gets walled but there is no G]
Sample 6 (\texttt{\#B / B.}): the exit is a neighbour of B and becomes a wall, but nobody needs to escape. Test ``no G'' first.
\end{warningbox}

\begin{warningbox}[Mistake 2 --- searching from every G]
Correct, but $(nm)^2 = 6.25 \cdot 10^6$ cell visits per test and up to $6 \cdot 10^8$ over 100 tests. The one search from the exit is $nm$.
\end{warningbox}

\begin{warningbox}[Mistake 3 --- walling only the cells that lie between B and the exit]
``Block the shortest path of B'' or ``block where B's search meets G's search'' leaves other routes open or wastes cells G needs. The neighbours of B are the only cells whose walling is provably free.
\end{warningbox}

\begin{warningbox}[Mistake 4 --- treating G and B cells as obstacles]
People can be walked through. A search that stops at G cells misses a G that is only reachable through another G (sample 5: the left G is reached through the right one).
\end{warningbox}

\subsubsection*{8. Takeaways}

\begin{itemize}[leftmargin=*]
    \item ``Everyone good must reach the exit'' $=$ one search from the exit.
    \item Make the bad starts harmless with the change every solution needs anyway; then only check.
    \item The proof idea --- ``a cell next to B is useless for G'' --- is the whole problem.
    \item Edge cases live at the exit: it can be walled, and with no G that is fine.
\end{itemize}

\begin{variantsbox}
\variant{Several exits}{\varnone}
  {every border cell marked \texttt{E} is an exit.}
  {same graph; the target is a set.}
  {wall in B, then one search from \emph{all} open exits at once (template 1 of the pattern).}
\variant{Minimum number of walls}{\varnone}
  {as few walls as possible.}
  {now walling all of B's neighbours may be wasteful: a vertex cut between B and the exit.}
  {minimum vertex cut, i.e.\ maximum flow on a split-vertex graph (Chapter~16); the search alone no longer suffices.}
\variant{Directed corridors}{\varnone}
  {some cells are one-way (arrows).}
  {directed cell graph.}
  {the escape check becomes a search from the exit in the \emph{reversed} graph; the ``B follows G'' argument still holds for cells B can reach.}
\variant{Enclosed regions}{\varin{https://leetcode.com/problems/surrounded-regions/}}
  {no people; find the cells that cannot reach the border.}
  {the targets are the whole border.}
  {one search from all border cells (Exercise E5.1).}
\end{variantsbox}

%============================================================
% CH2 · PATTERN 5 — EXERCISES, HINTS, ANSWERS
%============================================================
\subsection{Pattern 5: Exercises and Going Further}
\label{ssec:ch2-p5-ex}

\begin{patternbox}{Pattern 5 Summary Card}
\renewcommand{\arraystretch}{1.25}
\begin{tabularx}{\linewidth}{@{}>{\bfseries}l X@{}}
Recognise & many starts, few targets: ``from which cells / vertices can you reach \dots?'', ``every good one must escape'', ``touches the border'', ``ancestors'', ``forever'' \\
Reverse & directed: store $a \to b$ in \texttt{radj[b]}; directed rule in a grid: flip the comparison; undirected: only the start changes \\
Many targets & push all of them before the loop --- one search, \VE \\
Per vertex & one reverse search per vertex, \bigO{n(n+m)}, sorted lists by an index-order scan \\
Good vs.\ bad & neutralise the bad starts with a change every solution needs, then one search from the target \\
Forever & peel out-degree-0 vertices along reversed edges (Kahn on $G^R$) \\
Top traps & \texttt{radj[a]} instead of \texttt{radj[b]}; half-reversed grid rule; targets counted as answers; a walled exit with nobody to save \\
\end{tabularx}
\end{patternbox}

\exlegend
\begin{enumerate}[leftmargin=*, label=\textbf{E5.\arabic*}, itemsep=6pt]
    \item \exeasy\ \pclassic \lc{Medium}\ \exprob{Surrounded Regions}{LeetCode 130}{https://leetcode.com/problems/surrounded-regions/}. A board of \texttt{X} and \texttt{O}. Every region of \texttt{O} cells (4-connected) that does not touch the border is captured: flip it to \texttt{X}. Modify the board in place. ($1 \le R, C \le 200$.)
    \item \exmedium\ \pclassic \lc{Medium}\ \exprob{Pacific Atlantic Water Flow}{LeetCode 417}{https://leetcode.com/problems/pacific-atlantic-water-flow/}. An $R \times C$ height map; the Pacific touches the top and left edges, the Atlantic the bottom and right edges. Rain flows to a side-neighbour of height $\le$ the current one, and from a border cell into the adjacent ocean. List all cells from which water can reach both oceans. ($R, C \le 200$.)
    \item \exmedium\ \ptricky \cf{1600}\ \exprob{Lakes in Berland}{Codeforces 723D}{https://codeforces.com/problemset/problem/723/D}. A map of water (\texttt{.}) and land (\texttt{*}) surrounded by the ocean. A lake is a maximal 4-connected water region with no cell on the border. Fill as few water cells as possible so that exactly $k$ lakes remain; print the number and the new map. ($n, m \le 50$.)
    \item \exmedium\ \pconcept \exprob{Endless Walk}{AtCoder ABC245 F}{https://atcoder.jp/contests/abc245/tasks/abc245_f}. A simple directed graph with $N$ vertices and $M$ edges. Count the vertices from which one can walk along edges forever (choosing the edges cleverly). ($N, M \le 2 \cdot 10^5$.)
    \item \exmedium\ \pconcept \exown{Who can reach $T$?} Given a directed graph and a set $T$ of target vertices, mark every vertex that can reach some vertex of $T$. Implement the multi-target reverse search and the naive ``one forward search per vertex'', check that they agree on random graphs, and measure both for $n = 10^3, 10^4, 10^5$ with $m = 2n$ random edges and $|T| = 3$.
    \item \exhard\ \precent \lc{Medium}\ \exprob{Minimize the Maximum Edge Weight of Graph}{LeetCode 3419}{https://leetcode.com/problems/minimize-the-maximum-edge-weight-of-graph/}. A directed weighted graph on $n$ vertices and an integer \texttt{threshold} $\ge 1$. Remove edges so that vertex 0 is reachable from every vertex and every vertex keeps at most \texttt{threshold} outgoing edges, minimising the largest remaining weight. Return it, or $-1$. ($n, m \le 10^5$.)
\end{enumerate}

\subsubsection*{Hints}
\begin{enumerate}[leftmargin=*, label=\textbf{E5.\arabic*}]
    \item Which \texttt{O} cells are \emph{safe}? Find those first, starting from the border.
    \item Asking ``where does water from this cell go?'' for every cell is $R \cdot C$ searches. Ask instead ``which cells can send water to the Pacific?'' --- once, from the coast, uphill.
    \item The water regions that touch the border are not lakes; find them from the border (or flag them during a flood fill). Which lakes are cheapest to fill?
    \item A vertex with no outgoing edge is hopeless. What happens to a vertex all of whose edges lead to hopeless vertices?
    \item Template 1 of the pattern. For the naive version, stop each search as soon as it finds a target --- it is still slow.
    \item Binary search on the answer $W$. With only the edges of weight $\le W$: does everybody reach 0? And is \texttt{threshold} ever the real obstacle?
\end{enumerate}

\subsubsection*{Answers}

\begin{solutionbox}[E5.1 Surrounded Regions]
\textbf{Reverse the roles.} ``Is this region captured?'' per region becomes one search from all \texttt{O} cells on the border: everything it reaches is safe (temporarily marked \texttt{S}); every \texttt{O} left unmarked is captured. Finally \texttt{S} $\to$ \texttt{O}, \texttt{O} $\to$ \texttt{X}.
\begin{lstlisting}[style=cpp]
#include <bits/stdc++.h>
using namespace std;

class Solution {
public:
    void solve(vector<vector<char>>& board) {
        int R = board.size(), C = board[0].size();
        int dr[] = {1, -1, 0, 0}, dc[] = {0, 0, 1, -1};
        vector<pair<int, int>> st;
        for (int r = 0; r < R; r++)
            for (int c = 0; c < C; c++) {
                bool border = (r == 0 || c == 0 || r == R - 1 || c == C - 1);
                if (border && board[r][c] == 'O') {
                    board[r][c] = 'S'; st.push_back({r, c});
                }
            }
        while (!st.empty()) {                     // S = safe: connected to the border
            int r = st.back().first, c = st.back().second; st.pop_back();
            for (int d = 0; d < 4; d++) {
                int nr = r + dr[d], nc = c + dc[d];
                if (nr < 0 || nc < 0 || nr >= R || nc >= C) continue;
                if (board[nr][nc] != 'O') continue;
                board[nr][nc] = 'S'; st.push_back({nr, nc});
            }
        }
        for (auto& row : board)
            for (char& ch : row) ch = (ch == 'S') ? 'O' : 'X';   // captured O -> X
    }
};
\end{lstlisting}
Example 1 \texttt{XXXX / XOOX / XXOX / XOXX} $\to$ \texttt{XXXX / XXXX / XXXX / XOXX} (the bottom \texttt{O} touches the border and survives); example 2 \texttt{X} $\to$ \texttt{X}. Checked against a brute force (one search per \texttt{O} cell) on 5000 random boards up to $6 \times 6$. \bigO{RC}; iterative, since a $200 \times 200$ snake of \texttt{O} would make a recursive DFS $\approx 20\,000$ levels deep.
\end{solutionbox}

\begin{solutionbox}[E5.2 Pacific Atlantic Water Flow]
Water flows from $a$ to a neighbour $b$ if $h_b \le h_a$. Reversed: from $b$ we may \textbf{climb} to $a$ if $h_a \ge h_b$. Run this climb from all Pacific border cells and, separately, from all Atlantic border cells; the answer is the intersection.
\begin{center}
\begin{tikzpicture}[x=6.5mm, y=6.5mm]
    \foreach \r/\row in {0/{1,2,2,3,5}, 1/{3,2,3,4,4}, 2/{2,4,5,3,1}, 3/{6,7,1,4,5}, 4/{5,1,1,2,4}}
        \foreach \h [count=\c from 0] in \row \node[g2p5-e] at (\c,-\r) {\h};
    \foreach \r/\c in {0/4, 1/3, 1/4, 2/2, 3/0, 3/1, 4/0}
        \node[g2p5-e, fill=answerfill, draw=accentorange, line width=1.1pt] at (\c,-\r) {};
    \foreach \r/\row in {0/{1,2,2,3,5}, 1/{3,2,3,4,4}, 2/{2,4,5,3,1}, 3/{6,7,1,4,5}, 4/{5,1,1,2,4}}
        \foreach \h [count=\c from 0] in \row \node[font=\scriptsize\bfseries] at (\c,-\r) {\h};
    \draw[primaryblue, line width=2pt] (-0.55,-4.5) -- (-0.55,0.55) -- (4.5,0.55) node[right, font=\scriptsize] {Pacific};
    \draw[successgreen, line width=2pt] (-0.5,-4.55) -- (4.55,-4.55) -- (4.55,0.5);
    \node[font=\scriptsize, anchor=west] at (4.7,-4.4) {Atlantic};
\end{tikzpicture}
\end{center}
\begin{lstlisting}[style=cpp]
#include <bits/stdc++.h>
using namespace std;

class Solution {
    int R, C;
    // reverse search from an ocean's border: climb to neighbours that are NOT LOWER
    vector<vector<char>> climb(const vector<vector<int>>& h, vector<pair<int, int>> st) {
        vector<vector<char>> ok(R, vector<char>(C, 0));
        for (auto& p : st) ok[p.first][p.second] = 1;
        int dr[] = {1, -1, 0, 0}, dc[] = {0, 0, 1, -1};
        while (!st.empty()) {
            int r = st.back().first, c = st.back().second; st.pop_back();
            for (int d = 0; d < 4; d++) {
                int nr = r + dr[d], nc = c + dc[d];
                if (nr < 0 || nc < 0 || nr >= R || nc >= C || ok[nr][nc]) continue;
                if (h[nr][nc] < h[r][c]) continue;     // no flow from nr up to r
                ok[nr][nc] = 1; st.push_back({nr, nc});
            }
        }
        return ok;
    }
public:
    vector<vector<int>> pacificAtlantic(vector<vector<int>>& heights) {
        R = heights.size(); C = heights[0].size();
        vector<pair<int, int>> pac, atl;
        for (int r = 0; r < R; r++) { pac.push_back({r, 0}); atl.push_back({r, C - 1}); }
        for (int c = 0; c < C; c++) { pac.push_back({0, c}); atl.push_back({R - 1, c}); }
        sort(pac.begin(), pac.end());              // corners appear twice: deduplicate
        pac.erase(unique(pac.begin(), pac.end()), pac.end());
        sort(atl.begin(), atl.end());
        atl.erase(unique(atl.begin(), atl.end()), atl.end());
        vector<vector<char>> P = climb(heights, pac), A = climb(heights, atl);
        vector<vector<int>> res;
        for (int r = 0; r < R; r++)
            for (int c = 0; c < C; c++)
                if (P[r][c] && A[r][c]) res.push_back({r, c});
        return res;
    }
};
\end{lstlisting}
Example 1 gives \texttt{[0,4] [1,3] [1,4] [2,2] [3,0] [3,1] [4,0]} (orange above), example 2 \texttt{[[1]]} gives \texttt{[0,0]}; both equal the official outputs. A brute force that floods downhill from every cell agreed on 5000 random maps up to $6 \times 6$ with heights $0..4$ (many ties, which is where the $\le$ vs.\ $<$ bug lives). \bigO{RC}.
\end{solutionbox}

\begin{solutionbox}[E5.3 Lakes in Berland]
\textbf{The trap:} water connected to the ocean is not a lake, even if it is large, and filling it is never needed. Flood-fill every water component and flag those with a border cell (equivalently: first search from all border water cells, then the unmarked components are the lakes). With $L$ lakes, fill the $L - k$ \emph{smallest} ones: each lake must be filled completely to disappear (filling part of it leaves a smaller lake, or splits it into more lakes), so the cost is the sum of the sizes of the lakes we remove, minimised by removing the smallest.
\begin{lstlisting}[style=cpp]
#include <bits/stdc++.h>
using namespace std;

int n, m, k;
vector<string> g;
const int dr[] = {1, -1, 0, 0}, dc[] = {0, 0, 1, -1};

// iterative flood fill of '.' cells from (r, c); border = some cell is on the edge
vector<pair<int, int>> fill(int r, int c, vector<vector<char>>& seen, bool& border) {
    vector<pair<int, int>> cells, st(1, make_pair(r, c));
    seen[r][c] = 1; border = false;
    while (!st.empty()) {
        int x = st.back().first, y = st.back().second; st.pop_back();
        cells.push_back(make_pair(x, y));
        if (x == 0 || y == 0 || x == n - 1 || y == m - 1) border = true;
        for (int d = 0; d < 4; d++) {
            int a = x + dr[d], b = y + dc[d];
            if (a < 0 || b < 0 || a >= n || b >= m) continue;
            if (seen[a][b] || g[a][b] != '.') continue;
            seen[a][b] = 1; st.push_back(make_pair(a, b));
        }
    }
    return cells;
}

int main() {
    cin >> n >> m >> k;
    g.resize(n);
    for (auto& s : g) cin >> s;
    vector<vector<char>> seen(n, vector<char>(m, 0));
    vector<vector<pair<int, int>>> lakes;
    for (int r = 0; r < n; r++)
        for (int c = 0; c < m; c++)
            if (g[r][c] == '.' && !seen[r][c]) {
                bool border;
                vector<pair<int, int>> cells = fill(r, c, seen, border);
                if (!border) lakes.push_back(cells);   // touches the ocean: not a lake
            }
    sort(lakes.begin(), lakes.end(),
         [](const vector<pair<int, int>>& a, const vector<pair<int, int>>& b) {
             return a.size() < b.size();
         });
    int filled = 0;
    for (size_t i = 0; i + k < lakes.size(); i++)      // fill all but the k largest
        for (auto& p : lakes[i]) { g[p.first][p.second] = '*'; filled++; }
    cout << filled << "\n";
    for (auto& s : g) cout << s << "\n";
}
\end{lstlisting}
Both samples are reproduced exactly (\texttt{1} and the maps). In sample 1 the bottom-left water touches the border and is ignored, although it is as large as the lake of size 2. A brute force that tries all sets of water cells to fill (grids up to $5 \times 5$ with $\le 12$ water cells, random $k$) agreed on 1500 tests: same minimum, a map with exactly $k$ lakes, only water cells changed.
\end{solutionbox}

\begin{solutionbox}[E5.4 Endless Walk]
A vertex with out-degree 0 cannot start an endless walk; delete it, and decrement the out-degree of its predecessors --- found through the reversed edges. Repeat. The survivors all have an edge to another survivor, so a walk can go on forever from each of them; every deleted vertex leads only into deleted vertices. This is template 4 of the pattern (Kahn's algorithm on $G^R$).
\begin{center}
\begin{tikzpicture}[x=1.2cm]
    \node[gvsel] (w1) at (0,0) {1}; \node[gvsel] (w2) at (1.3,0) {2}; \node[gvsel] (w3) at (2.6,0.6) {3};
    \node[gvsel] (w4) at (2.6,-0.6) {4}; \node[gvdead] (w5) at (3.9,-0.6) {5};
    \draw[gde] (w1) -- (w2); \draw[gde] (w2) -- (w3); \draw[gde] (w3) -- (w4); \draw[gde] (w4) -- (w2);
    \draw[gfaded, -{Stealth[length=2.2mm]}] (w4) -- (w5);
    \node[note, text width=5.6cm, align=left, anchor=west] at (4.6,0) {Sample 1: only 5 has out-degree 0. Deleting it leaves 4 with out-degree 1 (edge $4 \to 2$), so the peeling stops: answer \textbf{4}.};
\end{tikzpicture}
\end{center}
\begin{lstlisting}[style=cpp]
#include <bits/stdc++.h>
using namespace std;

int main() {
    ios::sync_with_stdio(false);
    cin.tie(nullptr);
    int n, m;
    cin >> n >> m;
    vector<vector<int>> radj(n + 1);
    vector<int> out(n + 1, 0);
    for (int i = 0; i < m; i++) {
        int u, v;
        cin >> u >> v;
        radj[v].push_back(u);              // remember who points INTO v
        out[u]++;
    }
    vector<int> q;
    for (int v = 1; v <= n; v++) if (out[v] == 0) q.push_back(v);
    int removed = 0;
    while (!q.empty()) {
        int v = q.back(); q.pop_back();
        removed++;                         // every walk from v is finite
        for (int u : radj[v])
            if (--out[u] == 0) q.push_back(u);   // u's last exit led to a dead end
    }
    cout << n - removed << "\n";
}
\end{lstlisting}
Samples: \texttt{4} and \texttt{2} (sample 2: $1 \rightleftarrows 2$ and an isolated 3). A brute force with transitive closure (a vertex counts iff it reaches a vertex that lies on a cycle) agreed on 2000 random graphs with $N \le 8$. \bigO{N + M}, no recursion. A DFS alternative: colour the graph and count the vertices that reach a back edge --- correct, but it needs the iterative DFS and a second pass; the peeling is simpler.
\end{solutionbox}

\begin{solutionbox}[E5.5 Who can reach $T$?]
\begin{lstlisting}[style=cpp]
#include <bits/stdc++.h>
using namespace std;

// (a) one multi-target search in G^R
vector<char> reverseSearch(const vector<vector<int>>& radj, const vector<int>& T) {
    vector<char> can(radj.size(), 0);
    vector<int> st;
    for (int t : T) if (!can[t]) { can[t] = 1; st.push_back(t); }
    while (!st.empty()) {
        int v = st.back(); st.pop_back();
        for (int u : radj[v]) if (!can[u]) { can[u] = 1; st.push_back(u); }
    }
    return can;
}

// (b) a forward search from every vertex, stopped at the first target
vector<char> forwardEach(const vector<vector<int>>& adj, const vector<char>& inT,
                         long long& work) {
    int n = adj.size();
    vector<char> can(n, 0);
    vector<int> mark(n, -1), st;
    for (int s = 0; s < n; s++) {
        mark[s] = s; st.assign(1, s);
        while (!st.empty() && !can[s]) {
            int u = st.back(); st.pop_back(); work++;
            if (inT[u]) can[s] = 1;
            for (int v : adj[u]) if (mark[v] != s) { mark[v] = s; st.push_back(v); }
        }
    }
    return can;
}

int main() {
    mt19937 rng(11);
    int checks = 0, bad = 0;
    for (int t = 0; t < 2000; t++) {                 // small graphs: both must agree
        int n = rng() % 9 + 1, m = rng() % 20;
        vector<vector<int>> adj(n), radj(n);
        for (int i = 0; i < m; i++) {
            int a = rng() % n, b = rng() % n;
            adj[a].push_back(b); radj[b].push_back(a);
        }
        vector<int> T; vector<char> inT(n, 0);
        for (int v = 0; v < n; v++) if (rng() % 4 == 0) { T.push_back(v); inT[v] = 1; }
        long long w = 0;
        bad += reverseSearch(radj, T) != forwardEach(adj, inT, w);
        checks++;
    }
    cout << checks << " small graphs, " << bad << " differences\n";
    for (int n : {1000, 10000, 100000}) {             // timing: m = 2n, |T| = 3
        vector<vector<int>> adj(n), radj(n);
        for (int i = 0; i < 2 * n; i++) {
            int a = rng() % n, b = rng() % n;
            adj[a].push_back(b); radj[b].push_back(a);
        }
        vector<int> T = {0, 1, 2};
        vector<char> inT(n, 0);
        inT[0] = inT[1] = inT[2] = 1;
        long long work = 0;
        clock_t t0 = clock();
        vector<char> a = reverseSearch(radj, T);
        clock_t t1 = clock();
        vector<char> b = forwardEach(adj, inT, work);
        clock_t t2 = clock();
        cout << "n=" << n << ": reach T " << count(a.begin(), a.end(), 1)
             << ", equal " << (a == b)
             << ", reverse " << (t1 - t0) * 1000 / CLOCKS_PER_SEC << " ms, forward "
             << (t2 - t1) * 1000 / CLOCKS_PER_SEC << " ms (" << work << " pops)\n";
    }
}
\end{lstlisting}
Output on our machine:
\begin{lstlisting}[style=trace]
2000 small graphs, 0 differences
n=1000: reach T 805, equal 1, reverse 0 ms, forward 1 ms (59355 pops)
n=10000: reach T 7993, equal 1, reverse 0 ms, forward 724 ms (14748395 pops)
n=100000: reach T 79599, equal 1, reverse 8 ms, forward 63193 ms (384723015 pops)
\end{lstlisting}
About 80\% of the vertices reach $T$ (with $m = 2n$ most vertices lie in or before the giant strongly connected part). The forward searches that \emph{fail} --- the remaining 20\% --- are the expensive ones: each explores everything its start can reach before giving up, and that is a large part of the graph. Early stopping does not save the naive method; the reverse search is about $8000$ times faster at $n = 10^5$ (earlier runs of the same experiment: 10--12\,ms against 64--95\,s) and its cost grows linearly.
\end{solutionbox}

\begin{solutionbox}[E5.6 Minimize the Maximum Edge Weight of Graph]
\textbf{Two observations.} (1) ``0 is reachable from every vertex'' is one search from 0 in $G^R$. (2) The \texttt{threshold} is a decoy: if everybody reaches 0, the reverse search tree gives every vertex $\ne 0$ exactly one outgoing edge (towards its parent, i.e.\ towards 0) and vertex 0 none; keep only those $n - 1$ edges, and every out-degree is $\le 1 \le$ \texttt{threshold}. Removing edges never hurts the maximum weight. So the answer is the smallest $W$ for which the edges of weight $\le W$ let everybody reach 0 --- monotone in $W$, hence a binary search over the distinct weights, each step one reverse search.
\begin{center}
\begin{tikzpicture}[x=1.25cm]
    \node[gvdst] (z0) at (0,0) {0}; \node[gv] (z4) at (1.4,0) {4}; \node[gv] (z3) at (2.8,0) {3};
    \node[gv] (z2) at (2.8,-1.2) {2}; \node[gv] (z1) at (1.4,-1.2) {1};
    \draw[g2p5-rt] (z0) -- node[gw, above]{1} (z4);
    \draw[g2p5-rt] (z4) -- node[gw, above]{2} (z3);
    \draw[g2p5-rt] (z3) -- node[gw, right]{2} (z2);
    \draw[g2p5-rt] (z2) -- node[gw, below]{1} (z1);
    \draw[gfaded, -{Stealth[length=2.2mm]}] (z1) -- node[gw, left, pos=0.4]{3} (z3);
    \draw[gfaded, -{Stealth[length=2.2mm]}] (z1) -- node[gw, left]{5} (z4);
    \node[note, text width=5.0cm, align=left, anchor=west] at (3.6,-0.6) {Example 3 with $W = 2$: the reverse search from 0 reaches $4, 3, 2, 1$ (blue); the edges $1 \to 3$ (3) and $1 \to 4$ (5) are dropped and every vertex keeps one outgoing edge. With $W = 1$ vertex 3 is cut off. Answer \textbf{2}.};
\end{tikzpicture}
\end{center}
\begin{lstlisting}[style=cpp]
#include <bits/stdc++.h>
using namespace std;

class Solution {
    // only edges of weight <= W: does every vertex reach 0?  = one search from 0 in G^R
    bool allReachZero(int n, const vector<vector<pair<int, int>>>& radj, int W) {
        vector<char> seen(n, 0);
        vector<int> st(1, 0);
        seen[0] = 1;
        int cnt = 1;
        while (!st.empty()) {
            int v = st.back(); st.pop_back();
            for (auto& e : radj[v])                   // original edge e.first -> v
                if (e.second <= W && !seen[e.first]) {
                    seen[e.first] = 1; cnt++; st.push_back(e.first);
                }
        }
        return cnt == n;
    }
public:
    int minMaxWeight(int n, vector<vector<int>>& edges, int threshold) {
        vector<vector<pair<int, int>>> radj(n);
        vector<int> ws;
        for (auto& e : edges) {
            radj[e[1]].push_back({e[0], e[2]});
            ws.push_back(e[2]);
        }
        sort(ws.begin(), ws.end());
        ws.erase(unique(ws.begin(), ws.end()), ws.end());
        if (!allReachZero(n, radj, ws.back())) return -1;
        int lo = 0, hi = ws.size() - 1;               // smallest index that works
        while (lo < hi) {
            int mid = (lo + hi) / 2;
            if (allReachZero(n, radj, ws[mid])) hi = mid; else lo = mid + 1;
        }
        return ws[lo];                                // threshold never matters
    }
};
\end{lstlisting}
The four examples give \texttt{1 -1 2 -1}, as required. A brute force over all subsets of at most 10 edges (checking out-degrees and reachability directly) agreed on 3000 random graphs with $n \le 6$ and random thresholds. \bigO{(n+m)\log m}, iterative (a path of $10^5$ vertices is a legal input). Alternative without binary search: a Prim/Dijkstra-like ``minimax'' search from 0 in $G^R$ that always takes the cheapest edge leaving the reached set; the answer is the largest edge it takes.
\end{solutionbox}


%============================================================
\clearpage
\section{Chapter Summary \& What Comes Next}
\label{sec:ch2-summary}
%============================================================
%============================================================
% CH2 — CHAPTER SUMMARY & WHAT COMES NEXT
%   cheat sheet, the five patterns in pictures, recognition table, complexity table,
%   mind map, growth map, self-check quiz (answers at the end of the chapter),
%   readiness checklist, final words
%============================================================

\subsection{The Chapter on One Page}
\label{ssec:ch2-summary-main}

\begin{chaptertheorybox}{Cheat Sheet --- Chapter 2}
\begin{center}
\small
\renewcommand{\arraystretch}{1.45}
\begin{tabularx}{\linewidth}{@{}>{\bfseries}l >{\raggedright\arraybackslash}p{3.4cm} >{\raggedright\arraybackslash}X l@{}}
\toprule
Pattern & What DFS records & Core move & Cost \\
\midrule
P1 Cycles & colours $0/1/2$, \texttt{parent} & edge to a \emph{gray} vertex (undirected: not the parent edge) $\Rightarrow$ cycle $=$ stack from $v$ to $u$ & \VE \\
P2 Subtrees & parent, values per vertex & combine children when \emph{leaving}; pass values down when \emph{entering}; order array for big trees & \bigO{n} \\
P3 Timestamps & $\tin{}$, $\tout{}$, pre-order, sizes & $u$ ancestor of $v \iff \tin{u} \le \tin{v} \wedge \tout{v} \le \tout{u}$; subtree $=$ block of the pre-order & \bigO{n} $+ \bigO{1}$/query \\
P4 Backtracking & the current path & un-mark on exit $\Rightarrow$ every simple path; prune, cap, keep state on the path & exponential \\
P5 Reverse search & reversed edges & one search from the target(s) in $G^R$ instead of one per start & \VE \\
\bottomrule
\end{tabularx}
\end{center}

\textbf{The rhythm:} DFS enters a vertex, disappears into each white neighbour, and leaves; the gray vertices are always the path from the root, and the call stack is that path.

\textbf{The theorems:} parenthesis (intervals nest or are disjoint), white path (descendants $=$ vertices reachable by a white path at entry time), cycles (a directed graph has a cycle iff DFS finds a back edge), undirected DFS has only tree and back edges.

\textbf{The depth rule:} recursion depth $=$ longest DFS path, up to $n$; on a 2\,MB stack it dies near $6.5\cdot10^4$. Iterative DFS with \texttt{it[u]}, or the order array, when $n$ is large.

\textbf{The golden rule:} decide \emph{which hook} the answer lives in --- entering, examining an edge, or leaving --- before writing a line of code.
\end{chaptertheorybox}

\subsection{The Five Patterns in Pictures}

\begin{diagrambox}
\begin{tikzpicture}[
    sv/.style={minimum size=6mm, font=\scriptsize\bfseries},
    pnl/.style={draw=patternborder, fill=patternbg!45, rounded corners=3pt, line width=0.8pt},
    ptitle/.style={font=\scriptsize\bfseries\color{patternborder!80!black}},
    pcap/.style={font=\scriptsize, text width=4.5cm, align=center, text=darkgray}]
    % ---------------- P1 cycles ----------------
    \begin{scope}[shift={(0,4.6)}]
        \draw[pnl] (0,0) rectangle (4.8,4.2);
        \node[ptitle] at (2.4,3.9) {P1 Cycle Detection};
        \node[gvseen, sv] (a) at (0.7,3.0) {1};
        \node[gvsel, sv] (b) at (1.9,3.0) {2};
        \node[gvsel, sv] (c) at (3.1,3.0) {3};
        \node[gvcur, sv] (d) at (4.1,2.1) {4};
        \node[gvdone, sv] (e) at (1.9,1.7) {5};
        \draw[gtree, -{Stealth[length=2mm]}] (a) -- (b); \draw[gtree, -{Stealth[length=2mm]}] (b) -- (c);
        \draw[gtree, -{Stealth[length=2mm]}] (c) -- (d); \draw[gtree, -{Stealth[length=2mm]}] (a) -- (e);
        \draw[gback] (d) to[bend left=25] (b);
        \draw[gcross] (d) to[bend left=10] (e);
        \node[pcap] at (2.4,0.55) {Edge to a gray vertex closes a cycle\\($2,3,4$); to a black one: harmless.};
    \end{scope}
    % ---------------- P2 subtrees ----------------
    \begin{scope}[shift={(5.1,4.6)}]
        \draw[pnl] (0,0) rectangle (4.8,4.2);
        \node[ptitle] at (2.4,3.9) {P2 Subtree Aggregation};
        \node[gvsrc, sv] (r) at (2.4,3.2) {7};
        \node[gvdone, sv] (x) at (1.3,2.2) {3};
        \node[gvdone, sv] (y) at (3.5,2.2) {3};
        \node[gvdone, sv] (x1) at (0.8,1.25) {1}; \node[gvdone, sv] (x2) at (1.8,1.25) {1};
        \node[gvdone, sv] (y1) at (3.0,1.25) {1}; \node[gvdone, sv] (y2) at (4.0,1.25) {1};
        \foreach \a/\b in {r/x, r/y, x/x1, x/x2, y/y1, y/y2} \draw[gtree] (\a) -- (\b);
        \draw[->, thick, primaryblue] (4.5,1.2) -- (4.5,3.2);
        \node[pcap] at (2.4,0.55) {Sizes flow up when a vertex is left:\\$\mathrm{sz}[u] = 1 + \sum \mathrm{sz}[\text{child}]$.};
    \end{scope}
    % ---------------- P3 timestamps ----------------
    \begin{scope}[shift={(10.2,4.6)}]
        \draw[pnl] (0,0) rectangle (4.8,4.2);
        \node[ptitle] at (2.4,3.9) {P3 Entry \& Exit Times};
        \foreach \l/\a/\b/\c in {3/0.4/4.4/primaryblue, 2/0.7/2.3/successgreen, 2/2.6/4.1/successgreen, 1/0.95/1.5/accentorange, 1/2.85/3.4/accentorange}
            \draw[\c, line width=2pt] (\a,\l*0.65+0.6) -- (\b,\l*0.65+0.6);
        \node[font=\tiny\ttfamily] at (2.4,1.15) {( ( ( ) ) ( ( ) ) )};
        \node[pcap] at (2.4,0.55) {Intervals $[\tin{}, \tout{}]$ nest:\\ancestor test in two comparisons.};
    \end{scope}
    % ---------------- P4 backtracking ----------------
    \begin{scope}[shift={(0,0)}]
        \draw[pnl] (0,0) rectangle (4.8,4.2);
        \node[ptitle] at (2.4,3.9) {P4 Backtracking over Paths};
        \node[gvsrc, sv] (s) at (0.6,2.4) {s};
        \node[gv, sv] (p) at (1.8,3.1) {}; \node[gv, sv] (q) at (1.8,1.7) {};
        \node[gv, sv] (w) at (3.0,2.4) {};
        \node[gvdst, sv] (t) at (4.2,2.4) {t};
        \draw[gpath] (s) -- (p); \draw[gpath] (p) -- (w); \draw[gpath] (w) -- (t);
        \draw[ge] (s) -- (q); \draw[ge] (q) -- (w); \draw[ge] (p) -- (q);
        \node[pcap] at (2.4,0.55) {Mark on enter, \textbf{un-mark on exit}:\\every simple path $s \leadsto t$ is tried.};
    \end{scope}
    % ---------------- P5 reverse ----------------
    \begin{scope}[shift={(5.1,0)}]
        \draw[pnl] (0,0) rectangle (4.8,4.2);
        \node[ptitle] at (2.4,3.9) {P5 Reverse Search};
        \node[gvdone, sv] (a1) at (0.7,3.0) {}; \node[gvdone, sv] (a2) at (0.7,1.8) {};
        \node[gvdone, sv] (a3) at (2.0,2.4) {}; \node[gvdead, sv] (a4) at (2.0,1.3) {};
        \node[gvsrc, sv] (tt) at (3.6,2.4) {t};
        \draw[gde] (a1) -- (a3); \draw[gde] (a2) -- (a3); \draw[gde] (a3) -- (tt); \draw[gde] (tt) -- (a4);
        \draw[flow, recursionred] (3.6,3.2) to[bend right=30] node[above, font=\tiny, text=recursionred] {search on $G^R$} (0.9,3.5);
        \node[pcap] at (2.4,0.55) {One search from $t$ along reversed\\edges finds everyone who reaches $t$.};
    \end{scope}
    % ---------------- the three hooks ----------------
    \begin{scope}[shift={(10.2,0)}]
        \draw[pnl, draw=definitionborder, fill=definitionbg] (0,0) rectangle (4.8,4.2);
        \node[ptitle, text=definitionborder] at (2.4,3.9) {Before Any Code: Which Hook?};
        \node[font=\scriptsize, text width=4.3cm, align=left, anchor=north west] at (0.25,3.55) {%
            \textbf{enter $u$:} depth, root-path values, $\tin{u}$, mark on path\\[3pt]
            \textbf{examine $u \to v$:} colour of $v$ $\Rightarrow$ tree / back / forward / cross\\[3pt]
            \textbf{leave $u$:} subtree values, $\tout{u}$, finishing order, un-mark};
        \node[pcap] at (2.4,0.55) {Every pattern fills one or two\\of these three slots.};
    \end{scope}
\end{tikzpicture}
\figcap{The five patterns of the chapter and the one habit that connects them. Orange $=$ current vertex, light orange $=$ gray (on the stack), purple $=$ on the cycle, blue $=$ finished, thick orange $=$ the path being tried.}
\end{diagrambox}

\subsection{Which Pattern? A Recognition Table}

\begin{center}
\small
\renewcommand{\arraystretch}{1.35}
\begin{tabularx}{\linewidth}{@{}>{\raggedright\arraybackslash}p{4.7cm} >{\raggedright\arraybackslash}X c >{\raggedright\arraybackslash}p{2.6cm}@{}}
\toprule
\textbf{The statement says \dots} & \textbf{The graph in disguise} & \textbf{Pattern} & \textbf{Seen in} \\
\midrule
``a round trip'', ``a route back to the start'' & any cycle; undirected: skip the parent, directed: gray target & P1 & Round Trip, Round Trip II \\
``is the order consistent?'', ``contradiction'', ``deadlock'' & constraints as arrows; consistent iff acyclic & P1 & E1.5 \\
``is it a tree / exactly one cycle?'' & counting: $m$ versus $n - c$ & P1 & E1.2, E1.4 \\
``for every vertex, in its subtree \dots'' & post-order aggregation & P2 & White-Black, Highest Score \\
``remove vertex / edge $v$: sizes of the parts'' & $\mathrm{sz}[\text{child}]$ and $n - \mathrm{sz}[v]$ & P2 & Highest Score, E2.2 \\
``add $x$ to every vertex of the subtree of $p$'' & push-down from the root & P2 & Ki \\
``the $k$-th vertex in DFS order of $u$'s subtree'' & pre-order position $+$ size & P3 & Military Problem \\
``descendants of $u$ at depth $d$'' & $\tin{}$ interval $\times$ sorted list per depth & P3 & Count Descendants \\
``remove $k$ cells, keep the rest connected'' & finishing order / a connected prefix & P3 & Maze \\
``all paths'', ``best path'', $n \le 15$ or $\le 25$ cells & backtracking, un-mark on exit & P4 & All Paths, Max Gold \\
``which vertices can reach / are reachable from \dots{} the target'' & reversed graph, one search & P5 & All Ancestors \\
``good ones must escape, bad ones must not'' & block around the bad, search from the exit & P5 & Solve The Maze \\
\bottomrule
\end{tabularx}
\end{center}

\begin{warningbox}[When the table points to the wrong row]
``Shortest'' in the statement means BFS (Chapter~3), not DFS. ``All paths'' with $n = 10^5$ is never backtracking --- it is counting by DP on a DAG (\dpbook{12}{DP on DAGs \& Graphs}). And ``subtree'' questions with \emph{updates} between the queries need the data structures of Chapter~15. Always check the constraints before trusting a row.
\end{warningbox}

\subsection{Complexity at a Glance}

\begin{center}
\small
\renewcommand{\arraystretch}{1.35}
\begin{tabularx}{\linewidth}{@{}l >{\raggedright\arraybackslash}X c c@{}}
\toprule
\textbf{Pattern} & \textbf{Algorithm} & \textbf{Time} & \textbf{Memory} \\
\midrule
P1 Cycles & colour DFS (iterative), parent walk for printing & \VE & \bigO{n+m} \\
P1, yes/no undirected & count components: cycle iff $m > n - c$ & \VE & \bigO{n} \\
P2 Subtrees & one DFS, or BFS order + reverse loop & \bigO{n} & \bigO{n} \\
P3 Timestamps & one DFS; queries by comparisons or binary search & \bigO{n + q \log n} & \bigO{n} \\
P4 Backtracking & all simple paths & up to $\bigO{n!}$ & \bigO{n} \\
P5 Reverse search & one search on $G^R$ (per target, or multi-source) & \VE & \bigO{n+m} \\
\bottomrule
\end{tabularx}
\end{center}

\noindent Rule of thumb for P4: count the paths before coding. A $6 \times 6$ grid graph already has $1\,262\,816$ simple corner-to-corner paths (Lesson~\ref{ssec:ch2-toolbox}); $15$ vertices of a DAG have at most $2^{13} = 8192$ paths from the first to the last.

\subsection{Mind Map}

\begin{diagrambox}
\begin{tikzpicture}[
    mindmap, grow cyclic, every node/.style={concept, execute at begin node=\hskip0pt},
    concept color=primaryblue!70, text=white, font=\scriptsize\bfseries,
    level 1/.append style={level distance=4.1cm, sibling angle=72, font=\scriptsize\bfseries},
    level 2/.append style={level distance=2.3cm, sibling angle=40, font=\tiny}]
\node[font=\small\bfseries, minimum size=2.6cm] {DFS}
    child[concept color=successgreen!80!black] { node {Mechanics}
        child { node {call stack} }
        child { node {iterative it[u]} }
        child { node {depth limits} } }
    child[concept color=accentorange] { node {Timestamps}
        child { node {colours} }
        child { node {tin / tout} }
        child { node {parenthesis} } }
    child[concept color=softpurple] { node {DFS tree}
        child { node {tree, back} }
        child { node {forward, cross} }
        child { node {white path} } }
    child[concept color=patternborder] { node {5 Patterns}
        child { node {cycles, subtrees} }
        child { node {times, paths} }
        child { node {reverse} } }
    child[concept color=warningred!80!black] { node {Traps}
        child { node {gray vs visited} }
        child { node {parent edge} }
        child { node {stack overflow} } };
\end{tikzpicture}
\end{diagrambox}

\subsection{How Each Pattern Grows Through the Book}

\begin{diagrambox}
\begin{tikzpicture}[x=3.15cm, y=1.25cm,
    pat/.style={rectangle, rounded corners=3pt, draw=patternborder, fill=patternbg, line width=1pt,
                font=\scriptsize\bfseries, align=center, text width=2.2cm, minimum height=9mm},
    lvl/.style={rectangle, rounded corners=2pt, draw=#1, fill=#1!10, font=\tiny, align=center,
                text width=2.45cm, minimum height=8mm}]
    \node[font=\scriptsize\bfseries, color=deepblue] at (0,0.8) {Chapter 2};
    \node[font=\scriptsize\bfseries, color=successgreen] at (1,0.8) {Part I (Ch.~3--6)};
    \node[font=\scriptsize\bfseries, color=softpurple] at (2,0.8) {Part II--III (Ch.~7--17)};
    \node[font=\scriptsize\bfseries, color=warningred] at (3,0.8) {Part IV--VI (Ch.~18--35)};
    \node[pat] (p1) at (0,0)  {P1 Cycles};
    \node[pat] (p2) at (0,-1) {P2 Subtrees};
    \node[pat] (p3) at (0,-2) {P3 Timestamps};
    \node[pat] (p4) at (0,-3) {P4 Backtracking};
    \node[pat] (p5) at (0,-4) {P5 Reverse};
    \node[lvl=successgreen] (a1) at (1,0)  {topological order, cycle reporting (Ch.~5)};
    \node[lvl=successgreen] (a2) at (1,-1) {components as units (Ch.~4)};
    \node[lvl=successgreen] (a3) at (1,-2) {finishing order = topological order (Ch.~5)};
    \node[lvl=successgreen] (a4) at (1,-3) {BFS over states instead (Ch.~3)};
    \node[lvl=successgreen] (a5) at (1,-4) {multi-source BFS (Ch.~3)};
    \node[lvl=softpurple] (b1) at (2,0)  {negative cycles (Ch.~8)};
    \node[lvl=softpurple] (b2) at (2,-1) {tree fundamentals, rerooting (Ch.~13)};
    \node[lvl=softpurple] (b3) at (2,-2) {Euler tour + range queries (Ch.~13, 15)};
    \node[lvl=softpurple] (b4) at (2,-3) {DP on DAGs, bitmask DP (DP book)};
    \node[lvl=softpurple] (b5) at (2,-4) {reversed Dijkstra (Ch.~7)};
    \node[lvl=warningred] (c1) at (3,0)  {low-links: bridges, cut vertices (Ch.~18)};
    \node[lvl=warningred] (c2) at (3,-1) {small-to-large, centroids (Ch.~15, 17)};
    \node[lvl=warningred] (c3) at (3,-2) {SCC by finishing times (Ch.~19)};
    \node[lvl=warningred] (c4) at (3,-3) {Hamiltonian paths on small graphs (Ch.~23)};
    \node[lvl=warningred] (c5) at (3,-4) {Kosaraju's second pass on $G^R$ (Ch.~19)};
    \foreach \i in {1,...,5}{\draw[flow] (p\i) -- (a\i); \draw[flow] (a\i) -- (b\i); \draw[flow] (b\i) -- (c\i);}
\end{tikzpicture}
\figcap{Every arrow means ``the same DFS with one more piece of bookkeeping''. Bridges and strongly connected components, the two most famous DFS algorithms, are built entirely from this chapter's timestamps and back edges.}
\end{diagrambox}

\subsection{Self-Check Quiz}

Answer without looking back.
\begin{enumerate}[leftmargin=*]
    \item What is on the call stack of a recursive DFS at any moment, in graph terms?
    \item Why does the ``push all neighbours'' iterative DFS give no exit times?
    \item State the parenthesis theorem. What does $\tin{u} < \tin{v} < \tout{v} < \tout{u}$ tell you?
    \item Name the four edge types of a directed DFS and how each one is recognised when the edge is examined.
    \item Why does an undirected DFS have neither forward nor cross edges?
    \item In a directed graph, why is ``the target is already visited'' not a cycle test? Give a 3-vertex example.
    \item In an undirected multigraph, why must the parent \emph{edge} be skipped rather than the parent vertex?
    \item An undirected graph has $n$ vertices, $m$ edges and $c$ components. When does it contain a cycle?
    \item Give the formula for the sizes of the parts left when a vertex $v$ is deleted from a rooted tree.
    \item With pre-order numbers $\mathrm{pos}$ and sizes $\mathrm{sz}$, which positions does the subtree of $u$ occupy?
    \item What exactly changes between DFS and backtracking, and what does it cost?
    \item How do you find all vertices that can reach $t$ with a single search?
\end{enumerate}

\noindent\textit{Answers are at the end of the chapter, after the solution sketches (page~\pageref{ssec:ch2-quiz-answers}).}

\subsection{Ready for Chapter 3?}

\begin{checklistbox}[Move on when every box is ticked]
\begin{itemize}[leftmargin=*]
    \item[$\square$] I can write recursive DFS, iterative DFS with \texttt{it[u]} and the order-array loop from memory, and I know when recursion is too deep.
    \item[$\square$] I can compute $\tin{}$, $\tout{}$, pre-order and subtree sizes, and use the ancestor test.
    \item[$\square$] I can classify any edge of a DFS as tree, back, forward or cross, by hand and in code.
    \item[$\square$] I can find and print a cycle in directed and undirected graphs, including multigraphs and self-loops.
    \item[$\square$] I can solve a ``for every subtree'' problem on $2\cdot10^5$ vertices without recursion.
    \item[$\square$] I can recognise when backtracking is affordable, and when a search should be run backwards.
    \item[$\square$] I have solved at least 20 of the practice problems of Tiers 1--3, including at least 4 from Tier 3.
\end{itemize}
\end{checklistbox}

\subsection{What Comes Next}

DFS answers ``what is the structure?'': cycles, subtrees, ancestors, all the paths. It does \emph{not} answer ``how far?'' --- the path DFS finds to a vertex can be absurdly long (Lesson~\ref{ssec:ch2-toolbox}). Chapter~3 (BFS) explores in layers instead, and every vertex is reached by a shortest path; most of the ``minimum number of moves'' problems of contests live there. Chapters~4--6 then combine both searches into components, bipartite graphs, topological orders and disjoint-set unions. Later, Chapter~18 (bridges and articulation points) and Chapter~19 (strongly connected components) return to the DFS tree of this chapter and squeeze out much more of its structure with one extra number per vertex.


%============================================================
\clearpage
\section{Practice Problem Set}
\label{sec:ch2-problems}
%============================================================
%============================================================
% CH2 — PRACTICE PROBLEM SET (30 problems: 3 tiers + Stretch)   (tag g2s)
% Every problem carries a level badge and exactly one kind tag (AUTHORING.md 6b).
% None of these is a solved problem or an exercise anywhere else in the book, and
% no two of them share the same model + algorithm. Every sketch in
% problem_set_solutions.tex was implemented, reproduces the official samples and
% was compared with a brute force on random small inputs.
%============================================================
\subsection{Practice Problems}
\label{ssec:ch2-problem-set}

\begin{insightbox}[How to use this set]
\begin{itemize}[leftmargin=*]
    \item Work tier by tier: \exeasy{} problems should feel routine before you start the \exmedium{} ones; \exhard{} problems (up to \cf{1900}) need one non-obvious observation on top of a pattern. The \exstretch{} group is \emph{above} this chapter's level: optional, never needed to continue.
    \item Every tier mixes well-known classics, problems that isolate one idea of the chapter, recent contest problems and problems with a trap --- an edge case, a tempting wrong idea or a twist in the model.
    \item Before coding, write down the Graph Model \emph{and} the DFS story: what is computed when a vertex is entered, what is returned when it is left, and which edges matter (tree, back, reversed). Many of these problems never mention a graph; several hide a tree inside an array or a recursion.
    \item Watch the depth: most inputs allow $10^5$ or more vertices on one path. Use the iterative DFS or the order-array trick of Lesson~\ref{ssec:ch2-call-stack} unless the recursion is provably shallow.
    \item The hint names the pattern and the key idea. Cover it first; peek only after 20 minutes. Stuck after an hour? Draw the samples, write the slowest correct solution, and compare it with your fast one on small random inputs.
    \item Design sketches for all 30 problems follow the set (page~\pageref{ssec:ch2-sketches}) --- compare them with your model \emph{after} a real attempt.
\end{itemize}
\end{insightbox}

\subsection*{Tier 1 --- Warm-up \exeasy}

\problembox{P2.01}{\exeasy\ \pclassic Gold Rush --- Codeforces 1829D \cf{1000}}{Recursion over an implicit tree: a pile $x$ has the two children $x/3$ and $2x/3$ (only when $3 \mid x$). Search it from $n$; the depth is only $\log_{3/2} n$.}{https://codeforces.com/problemset/problem/1829/D}
\problembox{P2.02}{\exeasy\ \precent Huge Pile --- Codeforces 2184C \cf{1100}}{The split tree of $n$ is huge, but at depth $d$ every pile is $\lfloor n/2^d \rfloor$ or $\lceil n/2^d \rceil$: merge equal piles and walk down level by level.}{https://codeforces.com/problemset/problem/2184/C}
\problembox{P2.03}{\exeasy\ \pconcept Permutation Transformation --- Codeforces 1490D \cf{1200}}{The tree is defined by a recursion on segments: let the recursive call carry the depth down (P2 push-down, without ever building the tree).}{https://codeforces.com/problemset/problem/1490/D}
\problembox{P2.04}{\exeasy\ \ptricky Minimum Score of a Path Between Two Cities --- LeetCode 2492 \lc{Medium}}{A path may repeat roads. What does that do to the answer? One DFS over the component of city $1$ suffices.}{https://leetcode.com/problems/minimum-score-of-a-path-between-two-cities/}
\problembox{P2.05}{\exeasy\ \precent Count Simple Paths --- AtCoder ABC378 D}{P4 on a grid: from every empty cell, backtrack over paths of exactly $K \le 11$ moves; mark on the way in, un-mark on the way out.}{https://atcoder.jp/contests/abc378/tasks/abc378_d}
\problembox{P2.06}{\exeasy\ \precent Disappearing Permutation --- Codeforces 2086C \cf{1300}}{P1 on a permutation: once one position of a cycle of $i \to p_i$ is zero, the whole cycle must become the identity. Walk each cycle at most once.}{https://codeforces.com/problemset/problem/2086/C}
\problembox{P2.07}{\exeasy\ \pconcept Prerequisites --- AtCoder ABC315 E}{P3 orders: DFS from book $1$ along the prerequisite arrows and print the books in \emph{post-order} --- a book is finished only after everything it needs.}{https://atcoder.jp/contests/abc315/tasks/abc315_e}

\subsection*{Tier 2 --- Core \exmedium}

\problembox{P2.08}{\exmedium\ \ptricky Phase Shift --- Codeforces 1735C \cf{1400}}{P1 while building a functional graph: greedily choose the smallest predecessor letter, but never close a cycle shorter than $26$.}{https://codeforces.com/problemset/problem/1735/C}
\problembox{P2.09}{\exmedium\ \pclassic Nastya and Potions --- Codeforces 1851E \cf{1500}}{Recipes form a DAG. Memoised DFS: the cost of a potion is known when its DFS call returns --- $\min(c_i, \sum \text{cost of the ingredients})$.}{https://codeforces.com/problemset/problem/1851/E}
\problembox{P2.10}{\exmedium\ \precent Add Zeros --- Codeforces 2027C \cf{1500}}{Vertices are array \emph{lengths}: position $i$ turns length $a_i + i - 1$ into $a_i + 2i - 2$. DFS from $n$ in a \texttt{map}-based graph; answer the largest length reached.}{https://codeforces.com/problemset/problem/2027/C}
\problembox{P2.11}{\exmedium\ \precent Idiot First Search --- Codeforces 2195E \cf{1500}}{P2 up and down: $f(v)$, the time to leave a fresh subtree, is computed bottom-up; the answer for $k$ is the sum of $f$ along the path from $k$ to the root.}{https://codeforces.com/problemset/problem/2195/E}
\problembox{P2.12}{\exmedium\ \pclassic Valera and Elections --- Codeforces 369C \cf{1600}}{P2 bottom-up: elect $v$ only if the road above $v$ needs repair and nobody below $v$ is elected yet.}{https://codeforces.com/problemset/problem/369/C}
\problembox{P2.13}{\exmedium\ \pconcept Subtree K-th Max --- AtCoder ABC239 E}{P2 with a capped summary: $K \le 20$, so each subtree returns only its $20$ largest values.}{https://atcoder.jp/contests/abc239/tasks/abc239_e}
\problembox{P2.14}{\exmedium\ \ptricky Sum in the tree --- Codeforces 1098A \cf{1600}}{An erased $s_v$ should be as large as possible: the minimum $s$ of its children (or its parent's $s$ for a leaf). Then check that no $a_v$ is negative.}{https://codeforces.com/problemset/problem/1098/A}
\problembox{P2.15}{\exmedium\ \pconcept Height of Binary Tree After Subtree Removal Queries --- LeetCode 2458 \lc{Hard}}{P3: a subtree is an interval $[\mathit{tin}, \mathit{tout}]$ of the pre-order. The height after removing it is the largest depth \emph{outside} the interval: prefix and suffix maxima.}{https://leetcode.com/problems/height-of-binary-tree-after-subtree-removal-queries/}
\problembox{P2.16}{\exmedium\ \pconcept Game Master --- Codeforces 1608C \cf{1700}}{P5: draw ``beats'' arrows between neighbours in both sorted orders. A player can win iff he reaches the strongest player of map~1; search backwards from that one player.}{https://codeforces.com/problemset/problem/1608/C}

\subsection*{Tier 3 --- Challenge \exhard}

\problembox{P2.17}{\exhard\ \precent Tree Pruning --- Codeforces 2018C \cf{1700}}{Fix the final leaf depth $D$: vertex $v$ survives iff $\text{depth}(v) \le D \le$ the deepest depth in its subtree. Add $+1$ on that interval of $D$ for every vertex.}{https://codeforces.com/problemset/problem/2018/C}
\problembox{P2.18}{\exhard\ \precent Iris and Game on the Tree --- Codeforces 2006A \cf{1700}}{A leaf's weight is non-zero iff its value differs from the root's. Count the leaves; the game is about who is forced to colour the root.}{https://codeforces.com/problemset/problem/2006/A}
\problembox{P2.19}{\exhard\ \precent Reachability and Tree --- Codeforces 2112D \cf{1700}}{Every edge gives one good pair, so exactly one directed path of length $2$ may exist. Centre it at a vertex of degree $2$ and alternate sources and sinks everywhere else.}{https://codeforces.com/problemset/problem/2112/D}
\problembox{P2.20}{\exhard\ \precent Trapmigiano Reggiano --- Codeforces 2071C \cf{1700}}{Root the tree at the trap and put the cheese on the vertices in order of decreasing depth. Is $-1$ ever the answer?}{https://codeforces.com/problemset/problem/2071/C}
\problembox{P2.21}{\exhard\ \ptricky Mad City --- Codeforces 1873H \cf{1700}}{P1: $n$ edges, connected --- exactly one cycle. Valeriu is safe forever iff he reaches the cycle (at his entry point) strictly before Marcel can get there.}{https://codeforces.com/problemset/problem/1873/H}
\problembox{P2.22}{\exhard\ \ptricky Bakry and Partitioning --- Codeforces 1592C \cf{1700}}{Total xor $0$: one cut suffices. Otherwise you need three parts with xor $X$: cut greedily in post-order whenever a remaining subtree has xor $X$ --- and $k \ge 3$.}{https://codeforces.com/problemset/problem/1592/C}
\problembox{P2.23}{\exhard\ \pconcept Fake Plastic Trees --- Codeforces 1693B \cf{1700}}{P2 greedy bottom-up: the children can lift $v$ up to the sum of their values for free; start a new operation only when that sum is below $l_v$.}{https://codeforces.com/problemset/problem/1693/B}
\problembox{P2.24}{\exhard\ \pclassic Uncle Bogdan and Country Happiness --- Codeforces 1388C \cf{1800}}{P2: subtree sums give the number of people through each city; $h_v$ then fixes the number of happy ones, which must not grow from a city to its children.}{https://codeforces.com/problemset/problem/1388/C}
\problembox{P2.25}{\exhard\ \pclassic Edge Weight Assignment --- Codeforces 1338B \cf{1800}}{Minimum: depth parities of the leaves seen from one leaf. Maximum: only leaves hanging on the same vertex are forced to share a weight.}{https://codeforces.com/problemset/problem/1338/B}
\problembox{P2.26}{\exhard\ \precent Root was Built by Love, Broken by Destiny --- Codeforces 2127D \cf{1800}}{Any cycle forces a crossing, so the graph must be a tree, and its non-leaf vertices must form a path (a caterpillar). Then count orders of the leaves hanging on each spine vertex.}{https://codeforces.com/problemset/problem/2127/D}

\subsection*{Stretch --- Above This Chapter \exstretch}

\problembox{P2.27}{\exstretch\ \pclassic Reachability from the Capital --- Codeforces 999E \cf{2000}}{Delete what the capital reaches. Among the rest, a vertex that reaches the most others is a safe first target of a new road; repeat greedily.}{https://codeforces.com/problemset/problem/999/E}
\problembox{P2.28}{\exstretch\ \pconcept Tree Shuffling --- Codeforces 1363E \cf{2000}}{Push down $\min(a)$ along root paths, then fix wrong $0$s and $1$s in pairs as deep as possible (P2 up and down).}{https://codeforces.com/problemset/problem/1363/E}
\problembox{P2.29}{\exstretch\ \ptricky Ice Cave --- Codeforces 540C \cf{2000}}{A flood fill whose target must be entered twice when it is intact: count its intact neighbours. Mind the case start $=$ target.}{https://codeforces.com/problemset/problem/540/C}
\problembox{P2.30}{\exstretch\ \ptricky Coloring Edges --- Codeforces 1217D \cf{2100}}{P1 + timestamps: colour back edges $2$ and all others $1$. Why can neither colour class contain a cycle?}{https://codeforces.com/problemset/problem/1217/D}

\begin{tcolorbox}[enhanced, breakable, colback=white, colframe=accentorange, boxrule=1pt, arc=2mm,
    title={Topic index of the practice set --- one idea per problem}, fonttitle=\bfseries, colbacktitle=accentorange]
\small
\renewcommand{\arraystretch}{1.2}
\begin{tabularx}{\linewidth}{@{}>{\bfseries}l X@{}}
P1 Cycles & 06 a zero fixes its whole permutation cycle $\cdot$ 08 never close a short cycle $\cdot$ 21 race to the unique cycle $\cdot$ 26 a cycle makes the answer $0$ $\cdot$ 30 back edges get their own colour \\
P2 Subtrees & 03 depth carried down a recursion $\cdot$ 09 memoised cost on a DAG $\cdot$ 11 bottom-up $f$, top-down prefix sum $\cdot$ 12 elect only if nobody below $\cdot$ 13 top-$20$ summary $\cdot$ 14 erased sums $:=$ min of the children $\cdot$ 17 deepest vertex below $v$ $\cdot$ 18 leaves vs.\ root, who moves first $\cdot$ 22 greedy xor cuts $\cdot$ 23 free lift from the children $\cdot$ 24 happy people never increase downwards $\cdot$ 25 leaf parities and sibling leaves $\cdot$ 28 cheapest ancestor fixes pairs \\
P3 Times \& orders & 07 post-order is a valid reading order $\cdot$ 15 subtree $=$ pre-order interval $\cdot$ 17 survival interval of depths $\cdot$ 20 decreasing depth from the trap $\cdot$ 30 exit times decrease along non-back edges \\
P4 Search trees & 01 split tree of a pile $\cdot$ 02 merge equal piles level by level $\cdot$ 05 backtracking over simple paths of length $K$ $\cdot$ 10 implicit graph of array lengths \\
P5 Reverse \& reach & 04 the whole component counts, not one path $\cdot$ 16 reverse search from the strongest player $\cdot$ 19 one length-$2$ path at a degree-$2$ vertex $\cdot$ 27 greedy cover of what the capital misses $\cdot$ 29 enter the target twice \\
\end{tabularx}
\end{tcolorbox}

\clearpage
%============================================================
% CH2 — SOLUTION SKETCHES FOR THE PRACTICE SET   (tag g2s)
% Graph model, algorithm and complexity for each problem.
% Sketches, not full code: the point is to check your MODEL.
% Every sketch was implemented in C++, reproduces the official samples
% (LeetCode: the examples) and was compared with a brute force or a
% checker on random small inputs. Order follows problem_set.tex.
%============================================================
\subsection{Solution Sketches}
\label{ssec:ch2-sketches}

\begin{warningbox}[Read a sketch only after a real attempt]
Each sketch gives the \emph{design} --- the graph model, what the DFS computes on the way in and on the way out, the complexity and the one trap of the problem --- in a few lines, so you can compare it with your own model. Writing and debugging the code is left to you. If your model differs but you can argue why it answers the question, it is probably also right. ``Brute force'' below means an independent slow program (exhaustive search, simulation or a closed formula) that the fast solution was compared with on random small inputs.
\end{warningbox}

\begingroup
\providecommand{\sk}{}
\renewcommand{\sk}[2]{\par\smallskip\noindent\textbf{\textcolor{deepblue}{#1}}\quad #2\par}

\begin{solutionbox}[Tier 1 --- Warm-up]
\sk{P2.01 Gold Rush (CF 1829D).}{Implicit tree: the children of a pile $x$ with $3 \mid x$ are $x/3$ and $2x/3$; a pile never grows. \texttt{can(x)}: true if $x = m$, false if $x < m$ or $3 \nmid x$, else \texttt{can(x/3) || can(2x/3)}. Every call divides by $3$, and $3^{15} > 10^7$, so the depth is at most $14$ and there are fewer than $2^{15}$ calls: plain recursion is fine. Samples: \texttt{YES YES NO NO YES YES NO YES YES NO NO}; checked against the closed form ``$n = m \cdot 3^a / 2^b$ with $b \le a$'' on $10\,000$ random cases.}
\sk{P2.02 Huge Pile (CF 2184C).}{The split tree of $n$ has up to $n$ leaves, but by induction every pile at depth $d$ is $\lfloor n/2^d \rfloor$ or $\lceil n/2^d \rceil$. Keep the set of distinct pile sizes per level (at most $2$), go down until $k$ appears (answer $d$) or all piles are $1$ (answer $-1$). \bigO{\log n} per test. Samples: $2, 1, -1, 29$; checked against a DFS over the full, unmerged split tree on $8000$ random cases.}
\sk{P2.03 Permutation Transformation (CF 1490D).}{The tree is the recursion itself: \texttt{build(l, r, dep)} finds the maximum position $m$ of $a[l..r]$, records $d_m = \textit{dep}$ and recurses on $[l, m-1]$ and $[m+1, r]$ with $\textit{dep}+1$ --- the depth is pushed \emph{down} as a parameter, no tree is stored. \bigO{n^2} worst case, fine for $n \le 100$. Samples: \texttt{1 0 2 3 1}, \texttt{0}, \texttt{0 1 3 2}; checked against ``$a_i$ is an ancestor of $a_j$ iff $a_i$ is the maximum between them'' on $4000$ random permutations.}
\sk{P2.04 Minimum Score of a Path Between Two Cities.}{A path may reuse roads and pass through $1$ and $n$ several times, so it can make a detour over \emph{any} road of the component of city $1$ (which contains $n$). Answer: the smallest weight of an edge in that component --- one DFS from $1$, taking the minimum over every edge scanned. \emph{Trap:} a shortest-path or ``best simple path'' search answers a different question (example~2 needs the detour $1 \to 2 \to 1$). \bigO{n + m}. Examples: $5$, $2$; checked against a search over states (city, smallest road so far) on $1886$ random graphs.}
\sk{P2.05 Count Simple Paths (ABC378 D).}{Grid graph of the empty cells. From every empty start cell, backtrack: mark the cell, try the $4$ neighbours that are empty and unmarked with one move less, un-mark on return; count when $K$ moves are done. At most $HW \cdot 4 \cdot 3^{K-1}$ calls, about $2 \cdot 10^7$ for $K = 11$; depth $K$. Samples: $2$, $0$, $218070$; checked against enumerating all sequences of $K+1$ cells on $200$ random grids up to $3 \times 3$.}
\sk{P2.06 Disappearing Permutation (CF 2086C).}{Arrows $i \to p_i$: disjoint cycles. If $p_d$ becomes $0$, position $d$ must be set to $d$; then value $d$ is no longer free for the position $j$ with $p_j = d$, which must be fixed too, and so on around the cycle --- the whole cycle of $d$ becomes the identity, while untouched cycles stay valid. Keep a \texttt{fixed} flag; each query walks the cycle of $d$ only if it is not fixed yet and adds its length. Every vertex is walked once: \bigO{n}. Samples: \texttt{1 2 3}, \texttt{2 4 4 5 5}, \texttt{4 4 4 4 7 7 7}; checked against trying all subsets of positions to reset on $1000$ random cases ($n \le 8$).}
\sk{P2.07 Prerequisites (ABC315 E).}{Arrow $i \to P_{i,j}$ (``$i$ needs $P_{i,j}$''). The books to read are exactly those reachable from $1$; a book's DFS call ends after the calls of all its prerequisites, so the \emph{post-order} (without book $1$) is a valid reading order. $N$ up to $2 \cdot 10^5$ on a chain: iterative DFS with a per-vertex iterator. \bigO{N + \sum C_i}. Samples: we print \texttt{3 5 2 4} (the sample's \texttt{5 3 4 2} is another valid order), \texttt{6 5 4 3 2}, \texttt{5}; a checker (same set as a reachability search, every prerequisite earlier) accepted all samples and $200$ random DAGs.}
\end{solutionbox}

\begin{solutionbox}[Tier 2 --- Core]
\sk{P2.08 Phase Shift (CF 1735C).}{Functional graph on letters: \texttt{nxt[y] = x} means $y$ is encrypted as $x$; the final graph must be one cycle through all $26$ letters. Scan $t$: if $x = t_i$ has no predecessor yet, try $y = \texttt{a}, \texttt{b}, \dots$ with no successor yet and $y \ne x$; reject $y$ if walking \texttt{nxt} from $x$ ends at $y$ after fewer than $26$ letters (the arrow $y \to x$ would close a short cycle). \emph{Trap:} the natural greedy ``smallest unused letter'' fails exactly there --- sample~2 must answer \texttt{ac}, not \texttt{ab}. \bigO{26^2} per new letter plus \bigO{n}. Samples: \texttt{b}, \texttt{ac}, \texttt{abcdebfadg}, and the two $26$-letter answers; checked against all circular orders of a $5$-letter alphabet (same code with $26$ replaced by $5$) on $3000$ random strings.}
\sk{P2.09 Nastya and Potions (CF 1851E).}{Arrow from a potion to each ingredient; no potion is obtainable from itself, so the graph is a DAG. $\textit{cost}(i) = 0$ for owned potions, else $\min(c_i, \sum_j \textit{cost}(e_j))$ (just $c_i$ without a recipe), computed when the DFS call of $i$ returns, memoised. Iterative post-order (chains of length $2 \cdot 10^5$); sums in \texttt{long long}. \bigO{n + \sum m_i}. Samples: \texttt{23 8 0 5 10}, \texttt{0 143 0}, \texttt{5 0 1 3 4}, \texttt{0 0 0 0}, \texttt{0 0 0 0 0 2}, \texttt{0 0 0 0 0 0}, \texttt{0 0}; checked against relaxing all recipes until nothing changes on $1500$ random DAGs.}
\sk{P2.10 Add Zeros (CF 2027C).}{The appended zeros never satisfy the condition ($a_i \ge 1$), so only the original positions matter. Position $i \ge 2$ is usable exactly when the length is $u_i = a_i + i - 1$, and it moves the length to $u_i + i - 1$: a directed edge $u_i \to u_i + i - 1$ between lengths up to $\approx 10^{12}$. Store the edges in a \texttt{map}, DFS from length $n$, answer the largest length visited. \bigO{n \log n}. Samples: $10, 11, 10, 1$; checked against a search over the actual arrays on $1500$ random cases.}
\sk{P2.11 Idiot First Search (CF 2195E).}{Let $f(v)$ be the number of moves from entering $v$ with its whole subtree blank until Bob steps to the parent: $f = 1$ for a leaf, $f(v) = f(l) + f(r) + 3$ otherwise (down left, the left tour, down right, the right tour, up). Whenever Bob climbs into an ancestor, everything there is blank again, so $\tau(k) = f(k) + f(\textit{par}(k)) + \dots + f(1)$: compute $f$ bottom-up and push the prefix sum down, all modulo $10^9+7$. \bigO{n} with an explicit stack (the tree can be a path of depth $1.5 \cdot 10^5$). Samples: \texttt{1}, \texttt{9 10 14 15 15}, \texttt{13 22 14 27 23 28 28}; checked against a literal simulation of the walk on $1200$ random trees.}
\sk{P2.12 Valera and Elections (CF 369C).}{Root at district $1$. A problem road $(\textit{par}(v), v)$ is repaired iff some elected candidate lives in the subtree of $v$. Post-order: if the road above $v$ is a problem road and nothing in $v$'s subtree is elected, elect $v$ (the deepest choice covers the most roads); propagate ``something below is elected'' to the parent. \bigO{n}. Samples: \texttt{1 / 5}, \texttt{1 / 3}, \texttt{4 / 2 3 4 5} (any order is accepted); a checker (all problem roads repaired, size equal to the minimum over all subsets) accepted $300$ random trees.}
\sk{P2.13 Subtree K-th Max (ABC239 E).}{Since $K \le 20$, each vertex returns the $20$ largest values of its subtree: merge its own value with its children's lists, sort, truncate. Answer queries from the stored lists. \bigO{20\,n \log n}; iterative post-order (a path of $10^5$ vertices). Samples: \texttt{4 5}, \texttt{9 10}, \texttt{1 10 100 1000}; checked against sorting each queried subtree on $300$ random trees with $30$ queries each.}
\sk{P2.14 Sum in the tree (CF 1098A).}{Erased sums sit at even depth. Raising $s_v$ for such a $v$ lowers the sum of all children's $a$ by one each, so take $s_v$ as large as allowed: $\min$ of the children's $s$, or $s_{\textit{par}(v)}$ for a leaf (then $a_v = 0$). Since $p_i < i$, one pass in index order fills the values, a second computes $a_v = s_v - s_{\textit{par}(v)}$; any negative $a_v$ means $-1$. \emph{Trap:} choosing $s_v = s_{\textit{par}(v)}$ for every erased vertex looks minimal but piles the cost on the children. \bigO{n}. Samples: $1, 2, -1$; checked against trying all values of the erased sums on $300$ random trees (including corrupted inputs).}
\sk{P2.15 Height of Binary Tree After Subtree Removal Queries.}{Pre-order DFS with timestamps: the subtree of $v$ is the interval $[\mathit{tin}_v, \mathit{tout}_v]$ of the pre-order, and removing it leaves exactly the vertices outside. Store the depth at each pre-order position, build prefix and suffix maxima, and answer $\max(\textit{pre}[\mathit{tin}_q], \textit{suf}[\mathit{tout}_q + 1])$. \bigO{n + m}. Examples: $[2]$, $[3, 2, 3, 2]$; checked against recomputing the height for every possible query on $3000$ random trees.}
\sk{P2.16 Game Master (CF 1608C).}{Arrow $u \to v$ if $u$ beats $v$ on some map. Player $x$ wins iff he reaches every player (eliminate them in reverse DFS order, each by the one who discovered him). Only arrows between neighbours of the two sorted orders are needed (beating is transitive along a sorted order), so the graph has $2(n-1)$ arrows. The strongest player $s$ on map~1 reaches everybody, so $x$ wins iff $x$ reaches $s$: one reverse search from $s$ (arrows from each player to the next stronger one on either map) marks all winners. \bigO{n \log n} for sorting. Samples: \texttt{0001}, \texttt{1111}, \texttt{1}; checked against a bitmask search over tournament orders on $1500$ random cases ($n \le 7$).}
\end{solutionbox}

\begin{solutionbox}[Tier 3 --- Challenge]
\sk{P2.17 Tree Pruning (CF 2018C).}{Fix the final common leaf depth $D$. A vertex $v$ can stay iff $\text{depth}(v) \le D$ and some vertex of depth $\ge D$ lies below it, i.e.\ $D \le \textit{mx}(v)$ (deepest depth in the subtree, computed bottom-up); keeping exactly these is optimal and reachable by leaf removals. Add $+1$ on $[\text{depth}(v), \textit{mx}(v)]$ with a difference array; the answer is $n - \max_D$. \bigO{n}, iterative DFS. Samples: $2, 2, 5$; checked against all root-containing subtrees on $1500$ random trees ($n \le 11$).}
\sk{P2.18 Iris and Game on the Tree (CF 2006A).}{The changes $\mathtt{10}$/$\mathtt{01}$ along a root--leaf string cancel except for the ends: the weight is non-zero iff the leaf's value differs from the root's. Count leaves $c_0, c_1$ (known) and $q$ (\texttt{?}), and $u$ = \texttt{?} vertices that are neither root nor leaf (harmless moves). Root known: $(\text{leaves of the other value}) + \lceil q/2 \rceil$. Root \texttt{?} and $c_0 \ne c_1$: Iris colours the root at once, $\max(c_0, c_1) + \lfloor q/2 \rfloor$. Root \texttt{?} and $c_0 = c_1$: nobody wants to colour the root first; the parity of $u$ decides who must, giving $c_0 + \lceil q/2 \rceil$ for odd $u$ and $c_0 + \lfloor q/2 \rfloor$ for even $u$. \bigO{n}. Samples: $2, 1, 1, 2, 1, 0$; checked against a full minimax over the game on $1200$ random trees ($n \le 8$).}
\sk{P2.19 Reachability and Tree (CF 2112D).}{Every directed edge is a good pair, so $n - 1$ pairs are free and exactly one directed path of length $2$ must exist (two would give $\ge n+1$ pairs). Pick a vertex $c$ of degree $2$ with neighbours $a, b$; orient $a \to c \to b$, make $a$ a source and $b$ a sink, and alternate sources and sinks by DFS on both sides --- then $(a, b)$ is the only extra pair. Without a degree-$2$ vertex the answer is \texttt{NO} (a length-$2$ path through a vertex of degree $\ge 3$ creates at least two extra pairs). \bigO{n}. Samples: \texttt{YES YES NO YES} with our own orientations; a checker (same edges, exactly $n$ good pairs, \texttt{NO} only if no orientation of all $2^{n-1}$ works) accepted $1200$ random trees.}
\sk{P2.20 Trapmigiano Reggiano (CF 2071C).}{Root the tree at $\textrm{en}$ and output the vertices by decreasing depth. Invariant: after the cheese at depth $d$ the mouse is at depth $\le d$ --- if it is deeper, its step towards a vertex of depth $d$ goes up, and if it is at depth $\le d$, one step changes the depth by $1$ and ends at most at $d$. After the last cheese (depth $0$) the mouse is at $\textrm{en}$. \emph{Trap:} the statement offers $-1$, but it is never needed. \bigO{n \log n}. Samples: our orders differ from the samples; a simulator of the mouse accepted the samples and $1500$ random trees ($n \le 12$, never $-1$).}
\sk{P2.21 Mad City (CF 1873H).}{$n$ vertices, $n$ edges, connected: exactly one cycle (peel degree-$1$ vertices; what remains is the cycle). On the cycle Valeriu can always run away from Marcel. Let $c$ be the cycle vertex nearest to $b$ (where he enters it); he escapes iff $\textit{dist}(b, c) < \textit{dist}(a, c)$, two BFS runs. \emph{Trap:} with equality Marcel waits at $c$; and $b$ already on the cycle gives distance $0$, so the answer is \texttt{YES} unless $a = b$. \bigO{n}. Samples: \texttt{YES NO YES NO NO YES}; checked against a retrograde solution of the full pursuit game on $1500$ random unicyclic graphs ($n \le 8$).}
\sk{P2.22 Bakry and Partitioning (CF 1592C).}{Let $X$ be the xor of all values. $X = 0$: cut any edge, both parts have equal xor. $X \ne 0$: an even number of parts would xor to $0$, so we need $3$ parts of xor $X$ (more can be merged in threes). Post-order: when the remaining xor of a subtree equals $X$, cut it off and pass $0$ upwards; we need at least two cuts and $k \ge 3$. \emph{Trap:} forgetting $k \ge 3$, or counting the root's remainder as a cut. \bigO{n}. Samples: \texttt{NO YES NO YES NO}; checked against all sets of deleted edges on $1800$ random trees ($n \le 9$).}
\sk{P2.23 Fake Plastic Trees (CF 1693B).}{Post-order greedy. Each child $c$ ends with a value $\textit{val}(c)$, and operations ending inside the children can lift $v$ by up to $S = \sum \textit{val}(c)$ for free (the arrays are non-decreasing). If $S < l_v$, start one new operation at $v$ and set $\textit{val}(v) = r_v$; else $\textit{val}(v) = \min(r_v, S)$. Since $p_i < i$, a loop from $n$ down to $1$ is the post-order. \bigO{n}. Samples: $1, 2, 2, 5$; checked against a BFS over all value vectors on $900$ random trees ($n \le 4$, $r \le 3$).}
\sk{P2.24 Uncle Bogdan and Country Happiness (CF 1388C).}{Subtree sums give $\textit{sub}_v$, the number of people passing city $v$. With $g_v$ happy ones, $h_v = 2 g_v - \textit{sub}_v$, so $g_v = (\textit{sub}_v + h_v)/2$ must be an integer in $[0, \textit{sub}_v]$; and since a mood never improves, the happy people in the children cannot exceed those in $v$: $\sum_c g_c \le g_v$. These local checks are also sufficient. Two post-order passes. \bigO{n}. Samples: \texttt{YES YES}, \texttt{NO NO}; checked against enumerating where every person's mood breaks on $1200$ random cases.}
\sk{P2.25 Edge Weight Assignment (CF 1338B).}{Root at a leaf. Minimum: if every leaf is at even depth, all weights equal work ($f = 1$); otherwise $f = 2$ is impossible and $f = 3$ always suffices (with weights $1, 2, 3$, whose xor is $0$). Maximum: start from $n - 1$ distinct weights; leaves hanging on the same vertex must have equal edge weights, so subtract $(\text{leaf neighbours} - 1)$ for every vertex with at least two of them --- everything else can be made distinct. \bigO{n}. Samples: \texttt{1 4}, \texttt{3 3}, \texttt{1 6}; checked against all assignments of weights $1..15$ on $60$ random trees ($n \le 6$).}
\sk{P2.26 Root was Built by Love, Broken by Destiny (CF 2127D).}{Every bridge crosses the river, so the graph is bipartite with the two sides as parts, and every cycle drawn this way has crossing bridges: if $m \ne n - 1$, answer $0$. In a tree, the non-leaf vertices must form a path (a caterpillar): a vertex with three non-leaf neighbours forces a crossing, giving $0$. Otherwise the spine alternates sides, the leaves of each spine vertex sit in a block between its spine neighbours in any order: $\textit{ans} = 2 \cdot (2 \text{ if the spine has} \ge 2 \text{ vertices}) \cdot \prod_v (\text{leaves at } v)!$, with $2$ for $n = 2$. \bigO{n + m}. Samples: $2, 0, 8, 12$; checked against all side assignments and orders on $1200$ random graphs ($n \le 6$).}
\end{solutionbox}

\begin{solutionbox}[Stretch --- above this chapter]
\sk{P2.27 Reachability from the Capital (CF 999E).}{Mark everything $s$ reaches. A new road may always start at $s$. For each unmarked $v$, count the unmarked vertices it reaches (one DFS each, \bigO{n(n+m)} for $n, m \le 5000$); process them by decreasing count, and whenever $v$ is still unmarked add a road $s \to v$ and mark what $v$ reaches. A vertex with the largest count is a source of the condensation, so every road added this way is necessary. Samples: $3$, $1$; checked against trying every set of new roads $s \to v$ on $300$ random graphs ($n \le 9$).}
\sk{P2.28 Tree Shuffling (CF 1363E).}{An operation at $u$ may be replaced by one at the cheapest ancestor, so push down $c_v = \min(a_v, c_{\textit{par}(v)})$. Post-order: count the wrong digits in the subtree, $x$ of type $0 \to 1$ and $y$ of type $1 \to 0$; at $v$ fix $\min(x, y)$ pairs at cost $2 c_v$ each (deeper means cheaper) and pass the rest up. Leftovers at the root mean $-1$ (different numbers of $1$s). \bigO{n}, iterative DFS. Samples: $4$, $24000$, $-1$; checked against Dijkstra over all $2^n$ digit patterns with every possible shuffle on $300$ random trees ($n \le 6$).}
\sk{P2.29 Ice Cave (CF 540C).}{Flood fill from the start over intact cells, treating the target as enterable but never walking through it. Cracked target: answer = reached. Intact target: we arrive, crack it, must step out to an intact neighbour and come back --- reached and at least two intact neighbours, counting the start cell as one if it is adjacent (we can arrive from the start directly). \emph{Trap:} start $=$ target needs just one intact neighbour; and an intact target is not enough to reach once. \bigO{nm}. Samples: \texttt{YES NO YES}; checked against a search over (position, set of cracked cells) on $300$ random grids up to $3 \times 4$.}
\sk{P2.30 Coloring Edges (CF 1217D).}{If the digraph has no cycle, one colour suffices. Otherwise two colours always suffice: colour the back edges of a DFS (edges into a grey vertex, Lesson~\ref{ssec:ch2-cycles}) with $2$ and all other edges with $1$. Tree, forward and cross edges go from a vertex with a larger exit time to one with a smaller exit time, back edges the other way, so neither class can close a cycle. \emph{Trap:} trying to minimise the number of back edges is unnecessary --- only $k$ is minimised. \bigO{n + m}. Samples: \texttt{1 / 1 1 1 1 1}, \texttt{2 / 1 1 2}; a checker ($k$ equal to $1$ iff the graph is acyclic, every colour class acyclic) accepted $300$ random digraphs.}
\end{solutionbox}
\endgroup

%------------------------------------------------------------
\subsection{Answers to the Self-Check Quiz}
\label{ssec:ch2-quiz-answers}

\begin{solutionbox}[Answers]
\begin{enumerate}[leftmargin=*]
    \item The path from the start vertex to the current vertex, along the tree edges through which DFS went deeper --- exactly the gray vertices (Lesson~\ref{ssec:ch2-call-stack}).
    \item It never knows when a vertex is finished: a vertex is popped once and all its neighbours are pushed at once, so there is no moment ``all of $u$'s neighbours are done''. Its stack is also not a path and may hold a vertex several times.
    \item Two DFS intervals $[\tin{}, \tout{}]$ are either disjoint or nested, never partially overlapping; nested means ancestor and descendant. $\tin{u} < \tin{v} < \tout{v} < \tout{u}$: $v$ is a descendant of $u$.
    \item Tree (target white), back (target gray), forward (target black, $\tin{u} < \tin{v}$), cross (target black, $\tin{v} < \tin{u}$).
    \item If $\{u, v\}$ is an edge and $u$ is entered first, $u$ scans $v$ while gray: either $v$ is still white (tree edge) or it was entered meanwhile and is a descendant, so the edge is first seen from $v$'s side as a back edge to the gray $u$.
    \item A black target was finished earlier; it only means two paths lead there. Example: $1 \to 2 \to 3$ and $1 \to 3$ --- vertex 3 is black when $1 \to 3$ is examined, and there is no cycle.
    \item Two parallel edges between $p$ and $v$ form a cycle of length 2. Skipping only the edge used to enter $v$ lets the second one be found as a back edge; skipping the vertex $p$ hides it.
    \item A forest has $n - c$ edges, so a cycle exists iff $m > n - c$.
    \item The children's sizes $\mathrm{sz}[c]$ for every child $c$ of $v$, and the part above, $n - \mathrm{sz}[v]$ (if $v$ is not the root).
    \item The block $\mathrm{pos}[u], \mathrm{pos}[u] + 1, \dots, \mathrm{pos}[u] + \mathrm{sz}[u] - 1$.
    \item The mark of a vertex is removed when it is left. Then the marks mean ``on the current path'' and the search tries every simple path --- exponentially many in general, instead of \VE.
    \item Reverse every edge and run one search from $t$ in $G^R$; the marked vertices are exactly those that can reach $t$ in $G$.
\end{enumerate}
\end{solutionbox}


